\input{preamble}

% OK, start here.
%
\begin{document}

\title{Catégories dérivées des schémas}


\maketitle

\phantomsection
\label{section-phantom}

\tableofcontents

\section{Introduction}
\label{section-introduction}

\noindent
Dans ce chapitre, nous étudions les catégories dérivées de modules sur les schémas.
La plupart des résultats présentés ici se trouvent dans
\cite{TT}, \cite{Bokstedt-Neeman}, \cite{BvdB} et \cite{LN}.
Il existe bien entendu de nombreuses autres références.


\section{Conventions}
\label{section-conventions}

\noindent
Si $\mathcal{A}$ est une catégorie abélienne et $M$ un objet
de $\mathcal{A}$, nous notons aussi $M$ l'objet de $K(\mathcal{A})$
et/ou de $D(\mathcal{A})$ correspondant au complexe qui vaut
$M$ en degré $0$ et qui est nul en tout autre degré.

\medskip\noindent
Pour un anneau $A$, on note $K(A)$ la catégorie homotopique
des complexes de $A$-modules et $D(A)$ la catégorie dérivée associée.
De même, pour un espace annelé $(X, \mathcal{O}_X)$, le symbole
$K(\mathcal{O}_X)$ désigne la catégorie homotopique des complexes de
$\mathcal{O}_X$-modules et $D(\mathcal{O}_X)$ la catégorie dérivée
associée.










\section{Catégorie dérivée des modules quasi-cohérents}
\label{section-derived-quasi-coherent}

\noindent
Dans cette section, nous étudions les rapports entre les modules quasi-cohérents
et les modules quelconques sur un schéma $X$. On pourra consulter
\cite[appendice B]{TT}. D'après la discussion de
Schémas, section \ref{schemes-section-quasi-coherent},
le plongement
$\QCoh(\mathcal{O}_X) \subset \textit{Mod}(\mathcal{O}_X)$
fait de $\QCoh(\mathcal{O}_X)$ une sous-catégorie de Serre faible de
la catégorie des $\mathcal{O}_X$-modules. Notons
$$
D_\QCoh(\mathcal{O}_X) \subset D(\mathcal{O}_X)
$$
la sous-catégorie des complexes dont les faisceaux de cohomologie sont quasi-cohérents ; voir
Catégories dérivées, section \ref{derived-section-triangulated-sub}.
Nous obtenons ainsi un foncteur canonique
\begin{equation}
\label{equation-compare}
D(\QCoh(\mathcal{O}_X))
\longrightarrow
D_\QCoh(\mathcal{O}_X)
\end{equation}
voir Catégories dérivées, équation (\ref{derived-equation-compare}).

\begin{lemma}
\label{lemma-quasi-coherence-direct-sums}
Soit $X$ un schéma. Alors $D_\QCoh(\mathcal{O}_X)$
admet des sommes directes.
\end{lemma}

\begin{proof}
D'après Injectifs, lemme \ref{injectives-lemma-derived-products},
la catégorie dérivée $D(\mathcal{O}_X)$ admet des sommes directes,
qui se calculent en prenant les sommes directes terme à terme de représentants quelconques.
Il est donc clair que le faisceau de cohomologie d'une somme directe est la
somme directe des faisceaux de cohomologie, puisque le foncteur somme directe
est exact (dans toute catégorie abélienne de Grothendieck). Le lemme
en résulte, car une somme directe de faisceaux quasi-cohérents est quasi-cohérente ; voir
Schémas, section \ref{schemes-section-quasi-coherent}.
\end{proof}

\noindent
Nous aurons besoin de quelques résultats sur les limites dérivées. Signalons
que, dans le lemme ci-dessous, la limite dérivée ne sera généralement pas
un objet de $D_\QCoh$.

\begin{lemma}
\label{lemma-Rlim-quasi-coherent}
Soit $X$ un schéma. Soit $(K_n)$ un système projectif dans
$D_\QCoh(\mathcal{O}_X)$, de limite dérivée
$K = R\lim K_n$ dans $D(\mathcal{O}_X)$. Supposons que $H^q(K_{n + 1}) \to H^q(K_n)$
soit surjectif pour tous $q \in \mathbf{Z}$ et $n \geq 1$.
Alors
\begin{enumerate}
\item $H^q(K) = \lim H^q(K_n)$,
\item $R\lim H^q(K_n) = \lim H^q(K_n)$, et
\item pour tout ouvert affine $U \subset X$, on a
$H^p(U, \lim H^q(K_n)) = 0$ pour $p > 0$.
\end{enumerate}
\end{lemma}

\begin{proof}
Soit $\mathcal{B}$ l'ensemble des ouverts affines de $X$.
Puisque $H^q(K_n)$ est quasi-cohérent, on a $H^p(U, H^q(K_n)) = 0$
pour $U \in \mathcal{B}$ d'après Cohomologie des schémas, lemme
\ref{coherent-lemma-quasi-coherent-affine-cohomology-zero}.
De plus, les applications $H^0(U, H^q(K_{n + 1})) \to H^0(U, H^q(K_n))$
sont surjectives pour $U \in \mathcal{B}$ d'après
Schémas, lemme \ref{schemes-lemma-equivalence-quasi-coherent}.
L'assertion (1) résulte de Cohomologie, lemme
\ref{cohomology-lemma-derived-limit-suitable-system},
dont nous venons de vérifier les hypothèses.
Les assertions (2) et (3) résultent de
Cohomologie, lemme \ref{cohomology-lemma-inverse-limit-is-derived-limit}.
\end{proof}

\noindent
Le lemme suivant nous aidera à « calculer » un foncteur dérivé à droite
sur un objet de $D_\QCoh(\mathcal{O}_X)$.

\begin{lemma}
\label{lemma-nice-K-injective}
Soit $X$ un schéma. Soit $E$ un objet de
$D_\QCoh(\mathcal{O}_X)$. Alors le morphisme $E \to R\lim \tau_{\geq -n}E$ de
Catégories dérivées, remarque
\ref{derived-remark-map-into-derived-limit-truncations},
est un isomorphisme\footnote{En particulier,
$E$ admet un représentant K-injectif d'après
Catégories dérivées, lemme \ref{derived-lemma-difficulty-K-injectives}.}.
\end{lemma}

\begin{proof}
Notons $\mathcal{H}^i = H^i(E)$ le $i$-ième faisceau de cohomologie de $E$.
Soit $\mathcal{B}$ l'ensemble des ouverts affines de $X$. Alors
$H^p(U, \mathcal{H}^i) = 0$ pour tous $p > 0$, tous $i \in \mathbf{Z}$
et tous $U \in \mathcal{B}$ ; voir
Cohomologie des schémas, lemme
\ref{coherent-lemma-quasi-coherent-affine-cohomology-zero}.
Le lemme résulte donc de
Cohomologie, lemme \ref{cohomology-lemma-is-limit-dimension}.
\end{proof}

\begin{lemma}
\label{lemma-application-nice-K-injective}
Soit $X$ un schéma. Soient $F : \textit{Mod}(\mathcal{O}_X) \to \textit{Ab}$
un foncteur additif et $N \geq 0$ un entier. Supposons que
\begin{enumerate}
\item $F$ commute aux produits directs dénombrables,
\item $R^pF(\mathcal{F}) = 0$ pour tout $p \geq N$ et tout $\mathcal{F}$
quasi-cohérent.
\end{enumerate}
Alors, pour $E \in D_\QCoh(\mathcal{O}_X)$,
\begin{enumerate}
\item $H^i(RF(\tau_{\leq a}E)) \to H^i(RF(E))$ est un isomorphisme
pour $i \leq a$,
\item $H^i(RF(E)) \to H^i(RF(\tau_{\geq b - N + 1}E))$ est un isomorphisme
pour $i \geq b$,
\item si $H^i(E) = 0$ pour $i \not \in [a, b]$, où
$-\infty \leq a \leq b \leq \infty$, alors $H^i(RF(E)) = 0$
pour $i \not \in [a, b + N - 1]$.
\end{enumerate}
\end{lemma}

\begin{proof}
L'assertion (1) est
Catégories dérivées, lemme \ref{derived-lemma-negative-vanishing}.

\medskip\noindent
Démontrons l'assertion (2). Posons $E_n = \tau_{\geq -n}E$. On a
$E = R\lim E_n$ ; voir le lemme \ref{lemma-nice-K-injective}. Par suite,
$RF(E) = R\lim RF(E_n)$ dans $D(\textit{Ab})$ d'après Injectifs, lemme
\ref{injectives-lemma-RF-commutes-with-Rlim}. Ainsi, pour tout $i \in \mathbf{Z}$,
on a une suite exacte courte
$$
0 \to R^1\lim H^{i - 1}(RF(E_n)) \to H^i(RF(E)) \to \lim H^i(RF(E_n)) \to 0
$$
voir Compléments sur l'algèbre, remarque
\ref{more-algebra-remark-compare-derived-limit}.
Pour démontrer (2), nous montrerons que le terme de gauche est nul
et que le terme de droite est égal à $H^i(RF(E_{-b + N - 1}))$
pour tout $b$ tel que $i \geq b$.

\medskip\noindent
Pour tout $n$, on a un triangle distingué
$$
H^{-n}(E)[n] \to E_n \to E_{n - 1} \to H^{-n}(E)[n + 1]
$$
(Catégories dérivées, remarque
\ref{derived-remark-truncation-distinguished-triangle})
dans $D(\mathcal{O}_X)$. Puisque $H^{-n}(E)$ est quasi-cohérent, on a
$$
H^i(RF(H^{-n}(E)[n])) = R^{i + n}F(H^{-n}(E)) = 0
$$
pour $i + n \geq N$ et
$$
H^i(RF(H^{-n}(E)[n + 1])) = R^{i + n + 1}F(H^{-n}(E)) = 0
$$
pour $i + n + 1 \geq N$. On en déduit que
$$
H^i(RF(E_n)) \to H^i(RF(E_{n - 1}))
$$
est un isomorphisme pour $n \geq N - i$. Les systèmes $H^i(RF(E_n))$ vérifient donc tous
la condition ML, et le terme en $R^1\lim$ de notre suite exacte
courte est nul (voir la discussion de
Compléments sur l'algèbre, section \ref{more-algebra-section-Rlim}).
De plus, le système $H^i(RF(E_n))$ est constant à partir
de $n = N - i - 1$, comme voulu.

\medskip\noindent
Démontrons (3). Sous l'hypothèse faite sur $E$, on a
$\tau_{\leq a - 1}E = 0$, et l'annulation
de $H^i(RF(E))$ pour $i \leq a - 1$ résulte de (1).
De même, on a $\tau_{\geq b + 1}E = 0$, d'où
l'annulation de $H^i(RF(E))$ pour $i \geq b + N$ d'après
l'assertion (2).
\end{proof}

\noindent
Le lemme suivant joue un rôle essentiel dans nombre des
résultats de ce chapitre.

\begin{lemma}
\label{lemma-affine-compare-bounded}
Soit $X = \Spec(A)$ un schéma affine. Tous les foncteurs du diagramme
$$
\xymatrix{
D(\QCoh(\mathcal{O}_X)) \ar[rr]_{(\ref{equation-compare})}
& &
D_\QCoh(\mathcal{O}_X) \ar[ld]^{R\Gamma(X, -)} \\
& D(A) \ar[lu]^{\widetilde{\ \ }}
}
$$
sont des équivalences de catégories triangulées. De plus, pour $E$ dans
$D_\QCoh(\mathcal{O}_X)$, on a $H^0(X, E) = H^0(X, H^0(E))$.
\end{lemma}

\begin{proof}
Le foncteur $R\Gamma(X, -)$ donne un foncteur
$D(\mathcal{O}_X) \to D(A)$, donc, par restriction, un foncteur
\begin{equation}
\label{equation-back}
R\Gamma(X, -) : D_\QCoh(\mathcal{O}_X) \longrightarrow D(A).
\end{equation}
Nous montrerons que ce foncteur est quasi-inverse de (\ref{equation-compare})
moyennant l'équivalence entre les modules quasi-cohérents sur $X$ et
la catégorie des $A$-modules.

\medskip\noindent
Précisons ce point. Notons $(Y, \mathcal{O}_Y)$ l'espace réduit à un point, muni du faisceau
d'anneaux défini par $A$. Notons
$\pi : (X, \mathcal{O}_X) \to (Y, \mathcal{O}_Y)$
le morphisme évident d'espaces annelés.
Alors $R\Gamma(X, -)$ s'identifie à $R\pi_*$ et le foncteur
(\ref{equation-compare}), moyennant l'équivalence
$\textit{Mod}(\mathcal{O}_Y) = \text{Mod}_A = \QCoh(\mathcal{O}_X)$,
s'identifie à $L\pi^* = \pi^* = \widetilde{\ }$ (voir
Modules, lemme \ref{modules-lemma-construct-quasi-coherent-sheaves} et
Schémas, lemmes \ref{schemes-lemma-compare-constructions} et
\ref{schemes-lemma-equivalence-quasi-coherent}). Ainsi les foncteurs
$$
\xymatrix{
D(A) \ar@<1ex>[r] & D(\mathcal{O}_X) \ar@<1ex>[l]
}
$$
sont adjoints (d'après Cohomologie, lemme \ref{cohomology-lemma-adjoint}). En
particulier, on obtient des morphismes d'adjonction canoniques
$$
a : \widetilde{R\Gamma(X, E)} \longrightarrow E
$$
pour $E$ dans $D(\mathcal{O}_X)$, et
$$
b : M^\bullet \longrightarrow R\Gamma(X, \widetilde{M^\bullet})
$$
pour $M^\bullet$ un complexe de $A$-modules.

\medskip\noindent
Soit $E$ un objet de $D_\QCoh(\mathcal{O}_X)$. On peut appliquer
le lemme \ref{lemma-application-nice-K-injective}
au foncteur $F(-) = \Gamma(X, -)$
avec $N = 1$, d'après Cohomologie des schémas, lemme
\ref{coherent-lemma-quasi-coherent-affine-cohomology-zero}.
Par suite,
$$
H^0(R\Gamma(X, E)) = H^0(R\Gamma(X, \tau_{\geq 0}E)) = \Gamma(X, H^0(E))
$$
(la dernière égalité résulte de la définition de la troncature canonique).
Nous en déduirons que les morphismes d'adjonction $a$ et $b$
induisent des isomorphismes $H^0(a)$ et $H^0(b)$. Ainsi $a$ et $b$
sont des quasi-isomorphismes (puisque l'énoncé est invariant par décalage),
ce qui démontre le lemme.

\medskip\noindent
Dans les deux cas, on utilise l'exactitude du foncteur $\widetilde{\ }$
(Schémas, lemme \ref{schemes-lemma-spec-sheaves}). Elle
entraîne en effet que
$$
H^0\left(\widetilde{R\Gamma(X, E)}\right) =
\widetilde{H^0(R\Gamma(X, E))} =
\widetilde{\Gamma(X, H^0(E))}
$$
qui est égal à $H^0(E)$, car $H^0(E)$ est quasi-cohérent. Ainsi
$H^0(a)$ est un isomorphisme. Dans l'autre sens, on a
$$
H^0(R\Gamma(X, \widetilde{M^\bullet})) =
\Gamma(X, H^0(\widetilde{M^\bullet})) =
\Gamma(X, \widetilde{H^0(M^\bullet)}) =
H^0(M^\bullet)
$$
ce qui prouve que $H^0(b)$ est un isomorphisme.
\end{proof}

\begin{lemma}
\label{lemma-affine-K-flat}
Soit $X = \Spec(A)$ un schéma affine. Si $K^\bullet$ est un complexe K-plat
de $A$-modules, alors $\widetilde{K^\bullet}$ est un complexe K-plat
de $\mathcal{O}_X$-modules.
\end{lemma}

\begin{proof}
D'après Compléments sur l'algèbre, lemme \ref{more-algebra-lemma-base-change-K-flat},
$K^\bullet \otimes_A A_\mathfrak p$ est un complexe K-plat
de $A_\mathfrak p$-modules pour tout $\mathfrak p \in \Spec(A)$.
On déduit donc de
Cohomologie, lemme \ref{cohomology-lemma-check-K-flat-stalks}
(et de
Schémas, lemme \ref{schemes-lemma-spec-sheaves})
que $\widetilde{K^\bullet}$ est K-plat.
\end{proof}

\begin{lemma}
\label{lemma-quasi-coherence-pushforward}
Si $f : X \to Y$ est un morphisme de schémas affines défini par le morphisme d'anneaux
$A \to B$, le diagramme
$$
\xymatrix{
D(B) \ar[d] \ar[r] & D_\QCoh(\mathcal{O}_X) \ar[d]^{Rf_*} \\
D(A) \ar[r] & D_\QCoh(\mathcal{O}_Y)
}
$$
est commutatif.
\end{lemma}

\begin{proof}
Cela résulte du lemme \ref{lemma-affine-compare-bounded},
compte tenu de l'égalité $R\Gamma(Y, Rf_*K) = R\Gamma(X, K)$ donnée par
Cohomologie, lemme \ref{cohomology-lemma-Leray-unbounded}.
\end{proof}

\begin{lemma}
\label{lemma-quasi-coherence-pullback}
Soit $f : Y \to X$ un morphisme de schémas.
\begin{enumerate}
\item Le foncteur $Lf^*$ envoie $D_\QCoh(\mathcal{O}_X)$
dans $D_\QCoh(\mathcal{O}_Y)$.
\item Si $X$ et $Y$ sont affines et si $f$ est défini par le morphisme d'anneaux
$A \to B$, le diagramme
$$
\xymatrix{
D(B) \ar[r] & D_\QCoh(\mathcal{O}_Y) \\
D(A) \ar[r] \ar[u]^{- \otimes_A^\mathbf{L} B} &
D_\QCoh(\mathcal{O}_X) \ar[u]_{Lf^*}
}
$$
est commutatif.
\end{enumerate}
\end{lemma}

\begin{proof}
Montrons d'abord que le diagramme
$$
\xymatrix{
D(B) \ar[r] & D(\mathcal{O}_Y) \\
D(A) \ar[r] \ar[u]^{- \otimes_A^\mathbf{L} B} &
D(\mathcal{O}_X) \ar[u]_{Lf^*}
}
$$
est commutatif. Cela résulte immédiatement du lemme \ref{lemma-affine-K-flat} et
de la construction des foncteurs en question. Pour démontrer (1), soit
$E$ un objet de $D_\QCoh(\mathcal{O}_X)$. Pour vérifier que
$Lf^*E$ a des faisceaux de cohomologie quasi-cohérents, on peut travailler localement sur $X$.
Notons que $Lf^*$ est compatible à la restriction aux sous-schémas ouverts.
On peut donc supposer que $f$ est un morphisme de schémas affines comme dans (2).
Le lemme \ref{lemma-affine-compare-bounded} montre alors que
$E$ provient d'un complexe de $A$-modules. La commutativité du premier
diagramme de la démonstration donne l'énoncé analogue pour $Lf^*E$, d'où (1).
\end{proof}

\begin{lemma}
\label{lemma-quasi-coherence-tensor-product}
Soit $X$ un schéma.
\begin{enumerate}
\item Pour des objets $K, L$ de $D_\QCoh(\mathcal{O}_X)$,
le produit tensoriel dérivé $K \otimes^\mathbf{L}_{\mathcal{O}_X} L$ appartient à
$D_\QCoh(\mathcal{O}_X)$.
\item Si $X = \Spec(A)$ est affine, on a
$$
\widetilde{M^\bullet} \otimes_{\mathcal{O}_X}^\mathbf{L} \widetilde{K^\bullet}
=
\widetilde{M^\bullet \otimes_A^\mathbf{L} K^\bullet}
$$
pour tout couple de complexes de $A$-modules $K^\bullet$, $M^\bullet$.
\end{enumerate}
\end{lemma}

\begin{proof}
L'égalité (2) résulte immédiatement du lemme \ref{lemma-affine-K-flat}
et de la construction du produit tensoriel dérivé.
Pour démontrer (1), soient $K, L$ des objets de $D_\QCoh(\mathcal{O}_X)$.
Pour vérifier que $K \otimes^\mathbf{L} L$ appartient à
$D_\QCoh(\mathcal{O}_X)$, on peut travailler localement sur $X$, donc
supposer $X = \Spec(A)$ affine. D'après
le lemme \ref{lemma-affine-compare-bounded}, on peut représenter
$K$ et $L$ par des complexes de $A$-modules. L'assertion (2) donne alors
le résultat.
\end{proof}



\section{Image directe totale}
\label{section-total-direct-image}

\noindent
Le lemme suivant est l'analogue de
Cohomologie des schémas, lemme
\ref{coherent-lemma-quasi-coherence-higher-direct-images}.

\begin{lemma}
\label{lemma-quasi-coherence-direct-image}
Soit $f : X \to S$ un morphisme de schémas.
Supposons que $f$ soit quasi-séparé et quasi-compact.
\begin{enumerate}
\item Le foncteur $Rf_*$ envoie $D_\QCoh(\mathcal{O}_X)$
dans $D_\QCoh(\mathcal{O}_S)$.
\item Si $S$ est quasi-compact, il existe un entier $N = N(X, S, f)$
tel que, pour tout objet $E$ de $D_\QCoh(\mathcal{O}_X)$
vérifiant $H^m(E) = 0$ pour $m > 0$, on ait
$H^m(Rf_*E) = 0$ pour $m \geq N$.
\item Plus précisément, si $S$ est quasi-compact, on peut trouver $N = N(X, S, f)$
tel que, pour tout morphisme de schémas $S' \to S$,
la même conclusion vaille pour le foncteur $R(f')_*$,
où $f' : X' \to S'$ est obtenu à partir de $f$ par changement de base.
\end{enumerate}
\end{lemma}

\begin{proof}
Soit $E$ un objet de $D_\QCoh(\mathcal{O}_X)$. Pour démontrer (1), il faut
montrer que $Rf_*E$ a des faisceaux de cohomologie quasi-cohérents. La question est locale
sur $S$ ; on peut donc supposer $S$ quasi-compact. Choisissons $N = N(X, S, f)$ comme dans
Cohomologie des schémas, lemme
\ref{coherent-lemma-quasi-coherence-higher-direct-images}.
Alors $R^pf_*\mathcal{F} = 0$ pour tout $\mathcal{O}_X$-module quasi-cohérent
$\mathcal{F}$ et tout $p \geq N$, et cela reste vrai après changement de base.

\medskip\noindent
Supposons d'abord $E$ borné inférieurement. Montrons que (1), (2) et (3) sont vérifiées
pour un tel $E$, avec le choix de $N$ fait ci-dessus. Dans ce cas, on peut par exemple utiliser la
suite spectrale
$$
R^pf_*H^q(E) \Rightarrow R^{p + q}f_*E
$$
(Catégories dérivées, lemme \ref{derived-lemma-two-ss-complex-functor}),
la quasi-cohérence de $R^pf_*H^q(E)$ et l'annulation de $R^pf_*H^q(E)$
pour $p \geq N$, pour obtenir (1), (2) et (3).

\medskip\noindent
Démontrons ensuite (2) et (3). Supposons $H^m(E) = 0$ pour $m > 0$.
Soit $U \subset S$ un ouvert affine. D'après Cohomologie des schémas, lemme
\ref{coherent-lemma-quasi-coherence-higher-direct-images-application},
et le choix de $N$,
on a $H^p(f^{-1}(U), \mathcal{F}) = 0$ pour $p \geq N$
et tout $\mathcal{O}_X$-module quasi-cohérent $\mathcal{F}$.
On peut donc appliquer le lemme \ref{lemma-application-nice-K-injective}
au foncteur $\Gamma(f^{-1}(U), -)$ pour voir que
$$
R\Gamma(U, Rf_*E) = R\Gamma(f^{-1}(U), E)
$$
a une cohomologie nulle en degrés $\geq N$. Puisque cela vaut pour
tout ouvert affine $U \subset S$, on en déduit que $H^m(Rf_*E) = 0$
pour $m \geq N$.

\medskip\noindent
Démontrons maintenant (1) dans le cas général. Rappelons qu'il existe un
triangle distingué
$$
\tau_{\leq -n - 1}E \to E \to \tau_{\geq -n}E \to
(\tau_{\leq -n - 1}E)[1]
$$
dans $D(\mathcal{O}_X)$ ; voir Catégories dérivées, remarque
\ref{derived-remark-truncation-distinguished-triangle}.
D'après (2), les faisceaux de cohomologie de $Rf_*\tau_{\leq -n - 1}E$
sont nuls en degrés $\geq -n + N$.
Ainsi, étant donné un entier $q$, on voit que $R^qf_*E$ est égal
à $R^qf_*\tau_{\geq -n}E$ pour un certain $n$, et le résultat
ci-dessus s'applique.
\end{proof}

\begin{lemma}
\label{lemma-acyclicity-lemma}
Soit $f : X \to S$ un morphisme quasi-séparé et quasi-compact
de schémas. Soit $\mathcal{F}^\bullet$ un complexe de
$\mathcal{O}_X$-modules quasi-cohérents dont chaque terme est acyclique à droite pour $f_*$.
Alors $f_*\mathcal{F}^\bullet$ représente $Rf_*\mathcal{F}^\bullet$
dans $D(\mathcal{O}_S)$.
\end{lemma}

\begin{proof}
On dispose toujours d'un morphisme canonique
$f_*\mathcal{F}^\bullet \to Rf_*\mathcal{F}^\bullet$.
Il s'agit de montrer qu'il induit un isomorphisme sur les faisceaux de cohomologie.
L'énoncé étant invariant par décalage, il suffit de montrer que
$H^0(f_*(\mathcal{F}^\bullet)) \to R^0f_*\mathcal{F}^\bullet$
est un isomorphisme. L'énoncé est local sur $S$ ; on peut donc
supposer $S$ affine. D'après
le lemme \ref{lemma-quasi-coherence-direct-image},
on a $R^0f_*\mathcal{F}^\bullet = R^0f_*\tau_{\geq -n}\mathcal{F}^\bullet$
pour tout $n$ assez grand. On peut donc supposer $\mathcal{F}^\bullet$
borné inférieurement. Puisque chaque $\mathcal{F}^n$ est acyclique à droite pour $f_*$ par
hypothèse, le morphisme $f_*\mathcal{F}^\bullet \to Rf_*\mathcal{F}^\bullet$
est un quasi-isomorphisme d'après le lemme d'acyclicité de Leray (Catégories dérivées, lemme
\ref{derived-lemma-leray-acyclicity}).
\end{proof}

\begin{lemma}
\label{lemma-acyclicity-lemma-global}
Soit $X$ un schéma quasi-séparé et quasi-compact.
Soit $\mathcal{F}^\bullet$ un complexe de
$\mathcal{O}_X$-modules quasi-cohérents dont chaque terme est acyclique à droite pour $\Gamma(X, -)$.
Alors $\Gamma(X, \mathcal{F}^\bullet)$ représente
$R\Gamma(X, \mathcal{F}^\bullet)$ dans $D(\Gamma(X, \mathcal{O}_X)$.
\end{lemma}

\begin{proof}
Appliquer le lemme \ref{lemma-acyclicity-lemma} au morphisme canonique
$X \to \Spec(\Gamma(X, \mathcal{O}_X))$. Nous omettons quelques détails.
\end{proof}

\begin{lemma}
\label{lemma-spectral-sequence}
Soit $X$ un schéma quasi-séparé et quasi-compact. Pour tout objet
$K$ de $D_\QCoh(\mathcal{O}_X)$, la suite spectrale
$$
E_2^{i, j} = H^i(X, H^j(K)) \Rightarrow H^{i + j}(X, K)
$$
de Cohomologie, exemple \ref{cohomology-example-spectral-sequence},
est bornée et converge.
\end{lemma}

\begin{proof}
La construction de la suite spectrale donnée dans
Cohomologie, lemme \ref{cohomology-lemma-spectral-sequence-filtered-object},
à l'aide de la filtration définie par $\tau_{\leq -p}K$, montre qu'il
suffit de vérifier que, pour tout $n \in \mathbf{Z}$, on a
$$
H^n(X, \tau_{\leq -p}K) = 0 \text{ pour } p \gg 0
$$
et
$$
H^n(X, K) = H^n(X, \tau_{\leq -p}K) \text{ pour } p \ll 0
$$
La première assertion résulte du lemme \ref{lemma-application-nice-K-injective}
appliqué à $F = \Gamma(X, -)$ et de la borne donnée dans
Cohomologie des schémas, lemme
\ref{coherent-lemma-quasi-coherence-higher-direct-images}.
La seconde est vérifiée dès que
$-p \leq n$, pour tout espace annelé $(X, \mathcal{O}_X)$ et tout
$K \in D(\mathcal{O}_X)$.
\end{proof}

\begin{lemma}
\label{lemma-quasi-coherence-pushforward-direct-sums}
Soit $f : X \to S$ un morphisme quasi-séparé et quasi-compact de
schémas. Alors
$Rf_* : D_\QCoh(\mathcal{O}_X) \to D_\QCoh(\mathcal{O}_S)$
commute aux sommes directes.
\end{lemma}

\begin{proof}
Soit $E_i$ une famille d'objets de $D_\QCoh(\mathcal{O}_X)$,
et posons $E = \bigoplus E_i$. Nous voulons montrer que le morphisme
$$
\bigoplus Rf_*E_i \longrightarrow Rf_*E
$$
est un isomorphisme. Nous montrerons qu'il induit un isomorphisme sur les
faisceaux de cohomologie en degré $0$, ce qui entraînera le lemme.
Pour cela, on peut travailler localement sur $S$, ce qui permet de
supposer $S$ quasi-compact.
Choisissons un entier $N$ comme dans le lemme \ref{lemma-quasi-coherence-direct-image}.
Alors $R^0f_*E = R^0f_*\tau_{\geq -N}E$ et
$R^0f_*E_i = R^0f_*\tau_{\geq -N}E_i$ d'après le lemme cité. Notons que
$\tau_{\geq -N}E = \bigoplus \tau_{\geq -N}E_i$.
On peut donc supposer que tous les $E_i$ ont des faisceaux de cohomologie
nuls en degrés $< -N$. Utilisons ensuite les suites spectrales
$$
R^pf_*H^q(E) \Rightarrow R^{p + q}f_*E
\quad\text{et}\quad
R^pf_*H^q(E_i) \Rightarrow R^{p + q}f_*E_i
$$
(Catégories dérivées, lemme \ref{derived-lemma-two-ss-complex-functor})
pour nous ramener au cas d'une somme directe de faisceaux quasi-cohérents.
Ce cas est traité dans
Cohomologie des schémas, lemme \ref{coherent-lemma-colimit-cohomology}.
\end{proof}








\section{Morphismes affines}
\label{section-affine-morphisms}

\noindent
Dans cette section, nous réunissons quelques résultats sur l'image directe
par un morphisme affine de schémas.

\begin{lemma}
\label{lemma-pushforward-affine-morphism}
Soit $f : X \to S$ un morphisme affine de schémas. Soit $\mathcal{F}^\bullet$
un complexe de $\mathcal{O}_X$-modules quasi-cohérents. Alors
$f_*\mathcal{F}^\bullet = Rf_*\mathcal{F}^\bullet$.
\end{lemma}

\begin{proof}
Combiner le lemme \ref{lemma-acyclicity-lemma} et
Cohomologie des schémas, lemme \ref{coherent-lemma-relative-affine-vanishing}.
On peut aussi travailler sur les ouverts affines de $S$
et utiliser le lemme \ref{lemma-quasi-coherence-pushforward}.
\end{proof}

\begin{lemma}
\label{lemma-affine-morphism}
Soit $f : X \to S$ un morphisme affine de schémas.
Alors
$Rf_* : D_\QCoh(\mathcal{O}_X) \to D_\QCoh(\mathcal{O}_S)$
est conservatif.
\end{lemma}

\begin{proof}
L'énoncé signifie qu'un morphisme $\alpha : E \to F$ de
$D_\QCoh(\mathcal{O}_X)$ est un isomorphisme si
$Rf_*\alpha$ est un isomorphisme. On peut le vérifier sur les faisceaux de cohomologie.
En particulier, la question est locale sur $S$. On peut donc supposer $S$
affine, et il en est alors de même de $X$. Dans ce cas, l'énoncé résulte immédiatement de
la description des catégories dérivées
$D_\QCoh(\mathcal{O}_X)$ et
$D_\QCoh(\mathcal{O}_S)$ donnée dans
le lemme \ref{lemma-affine-compare-bounded}.
Nous omettons quelques détails.
\end{proof}

\begin{lemma}
\label{lemma-affine-morphism-pull-push}
Soit $f : X \to S$ un morphisme affine de schémas.
Pour $E$ dans $D_\QCoh(\mathcal{O}_S)$, on a
$Rf_* Lf^* E = E \otimes^\mathbf{L}_{\mathcal{O}_S} f_*\mathcal{O}_X$.
\end{lemma}

\begin{proof}
Puisque $f$ est affine, le morphisme $f_*\mathcal{O}_X \to Rf_*\mathcal{O}_X$
est un isomorphisme
(Cohomologie des schémas, lemme \ref{coherent-lemma-relative-affine-vanishing}).
On dispose d'un morphisme canonique $E \otimes^\mathbf{L} f_*\mathcal{O}_X =
E \otimes^\mathbf{L} Rf_*\mathcal{O}_X \to Rf_* Lf^* E$
adjoint au morphisme
$$
Lf^*(E \otimes^\mathbf{L} Rf_*\mathcal{O}_X) =
Lf^*E \otimes^\mathbf{L} Lf^*Rf_*\mathcal{O}_X \longrightarrow
Lf^* E \otimes^\mathbf{L} \mathcal{O}_X = Lf^* E
$$
déduit de $1 : Lf^*E \to Lf^*E$ et du morphisme canonique
$Lf^*Rf_*\mathcal{O}_X \to \mathcal{O}_X$. Pour vérifier que le morphisme ainsi construit
est un isomorphisme, on peut travailler localement sur $S$. On peut donc supposer
$S$ affine, et il en est alors de même de $X$. Dans ce cas, l'énoncé résulte immédiatement de
la description des catégories dérivées
$D_\QCoh(\mathcal{O}_X)$ et
$D_\QCoh(\mathcal{O}_S)$ et du foncteur $Lf^*$ donnée dans
les lemmes \ref{lemma-affine-compare-bounded} et
\ref{lemma-quasi-coherence-pullback}.
Nous omettons quelques détails.
\end{proof}

\noindent
Soit $Y$ un schéma. Soit $\mathcal{A}$ un faisceau de $\mathcal{O}_Y$-algèbres.
Nous noterons $D_\QCoh(\mathcal{A})$ l'image réciproque de
$D_\QCoh(\mathcal{O}_X)$ par le foncteur de restriction
$D(\mathcal{A}) \to D(\mathcal{O}_X)$. Autrement dit, $K \in D(\mathcal{A})$
appartient à $D_\QCoh(\mathcal{A})$ si et seulement si ses faisceaux de cohomologie sont
quasi-cohérents en tant que $\mathcal{O}_X$-modules. Si $\mathcal{A}$ est lui-même quasi-cohérent,
cela revient à demander que les faisceaux de cohomologie soient quasi-cohérents
en tant que $\mathcal{A}$-modules ; voir
Morphismes, lemme \ref{morphisms-lemma-affine-equivalence-modules}.

\begin{lemma}
\label{lemma-affine-morphism-equivalence}
Soit $f : X \to Y$ un morphisme affine de schémas. Alors $f_*$ induit
une équivalence
$$
\Phi : D_\QCoh(\mathcal{O}_X) \longrightarrow D_\QCoh(f_*\mathcal{O}_X)
$$
dont le composé avec $D_\QCoh(f_*\mathcal{O}_X) \to D_\QCoh(\mathcal{O}_Y)$
est $Rf_* : D_\QCoh(\mathcal{O}_X) \to D_\QCoh(\mathcal{O}_Y)$.
\end{lemma}

\begin{proof}
Rappelons que, pour calculer $Rf_*$ sur un objet $K \in D_\QCoh(\mathcal{O}_X)$,
on choisit un complexe K-injectif $\mathcal{I}^\bullet$ de
$\mathcal{O}_X$-modules représentant $K$, et l'on prend $f_*\mathcal{I}^\bullet$.
On définit donc $\Phi(K)$ comme le complexe $f_*\mathcal{I}^\bullet$
considéré comme un complexe de $f_*\mathcal{O}_X$-modules.
Notons $g : (X, \mathcal{O}_X) \to (Y, f_*\mathcal{O}_X)$ le
morphisme évident d'espaces annelés. Alors $g$ est un morphisme plat
d'espaces annelés (voir ci-dessous la description des fibres), et
$\Phi$ est la restriction de $Rg_*$ à $D_\QCoh(\mathcal{O}_X)$.
Nous affirmons que $Lg^*$ en est un quasi-inverse. Observons d'abord que
$Lg^*$ envoie $D_\QCoh(f_*\mathcal{O}_X)$ dans $D_\QCoh(\mathcal{O}_X)$,
car $g^*$ transforme les modules quasi-cohérents en modules quasi-cohérents
(Modules, lemme \ref{modules-lemma-pullback-quasi-coherent}).
Pour achever la démonstration, il suffit de montrer que
les morphismes d'adjonction
$$
Lg^*\Phi(K) = Lg^*Rg_*K \to K
\quad\text{et}\quad
M \to Rg_*Lg^*M = \Phi(Lg^*M)
$$
sont des isomorphismes pour $K \in D_\QCoh(\mathcal{O}_X)$ et
$M \in D_\QCoh(f_*\mathcal{O}_X)$. La question étant locale,
on peut supposer $Y$ affine, et donc aussi $X$.

\medskip\noindent
Supposons $Y = \Spec(B)$ et $X = \Spec(A)$. Soit
$\mathfrak p = x \in \Spec(A) = X$ un point dont l'image est
$\mathfrak q = y \in \Spec(B) = Y$. Alors
$(f_*\mathcal{O}_X)_y = A_\mathfrak q$ et $\mathcal{O}_{X, x} = A_\mathfrak p$ ;
le morphisme $g$ est donc plat. Par suite, $g^*$ est exact et $H^i(Lg^*M) = g^*H^i(M)$
pour tout $M$ dans $D(f_*\mathcal{O}_X)$.
Pour $K \in D_\QCoh(\mathcal{O}_X)$, on a
$$
H^i(\Phi(K)) = H^i(Rf_*K) = f_*H^i(K)
$$
d'après l'annulation des images directes supérieures
(Cohomologie des schémas, lemme \ref{coherent-lemma-relative-affine-vanishing})
et le lemme \ref{lemma-application-nice-K-injective} (un petit détail est omis).
Il suffit donc de montrer que
$$
g^*g_*\mathcal{F} \to \mathcal{F}
\quad\text{et}\quad
\mathcal{G} \to g_*g^*\mathcal{G}
$$
sont des isomorphismes, où $\mathcal{F}$ est
un $\mathcal{O}_X$-module quasi-cohérent et $\mathcal{G}$
un $f_*\mathcal{O}_X$-module quasi-cohérent. Cela résulte de
Morphismes, lemme \ref{morphisms-lemma-affine-equivalence-modules}.
\end{proof}




\section{Cohomologie à support dans un sous-ensemble fermé}
\label{section-cohomology-support}

\noindent
Nous développons, pour les schémas et les modules quasi-cohérents,
les résultats de Cohomologie, sections
\ref{cohomology-section-cohomology-support} et
\ref{cohomology-section-cohomology-support-bis}.

\begin{definition}
\label{definition-supported-on}
Soit $X$ un schéma. Soit $E$ un objet de $D(\mathcal{O}_X)$.
Soit $T \subset X$ une partie fermée.
On dit que $E$ est {\it à support dans $T$} si les
faisceaux de cohomologie $H^i(E)$ sont à support dans $T$.
\end{definition}

\noindent
Reprenons une partie de la discussion de
Cohomologie, section \ref{cohomology-section-cohomology-support-bis},
dans la situation de la définition.
Soit $X$ un schéma. Soit $T \subset X$ une partie fermée.
La catégorie des $\mathcal{O}_X$-modules dont le
support est contenu dans $T$ est une sous-catégorie de Serre de la
catégorie de tous les $\mathcal{O}_X$-modules ; voir
Homologie, définition \ref{homology-definition-serre-subcategory}
et
Modules, lemme \ref{modules-lemma-support-section-closed}.
On peut donc appliquer la discussion de
Catégories dérivées, section \ref{derived-section-triangulated-sub},
et l'on obtient une sous-catégorie triangulée strictement pleine et saturée de
$D(\mathcal{O}_X)$, formée des objets à support dans $T$ ; nous
noterons cette catégorie $D_T(\mathcal{O}_X)$ dans la suite.

\medskip\noindent
Dans la situation de la définition \ref{definition-supported-on},
notons $i : T \to X$ l'inclusion. Rappelons que, d'après
Cohomologie, section \ref{cohomology-section-cohomology-support-bis},
on dispose dans cette situation d'un foncteur
$R\mathcal{H}_T : D(\mathcal{O}_X) \to D(i^{-1}\mathcal{O}_X)$
adjoint à droite de $i_* : D(i^{-1}\mathcal{O}_X) \to D(\mathcal{O}_X)$.

\begin{lemma}
\label{lemma-support-quasi-coherent}
Soit $X$ un schéma. Soit $T \subset X$ une partie fermée telle que
$X \setminus T$ soit un ouvert rétrocompact de $X$. Soit $i : T \to X$
l'inclusion.
\begin{enumerate}
\item Pour $E$ dans $D_\QCoh(\mathcal{O}_X)$, on a
$i_*R\mathcal{H}_T(E)$ dans $D_{\QCoh, T}(\mathcal{O}_X)$.
\item Le foncteur
$i_* \circ R\mathcal{H}_T : D_\QCoh(\mathcal{O}_X) \to
D_{\QCoh, T}(\mathcal{O}_X)$ est adjoint à droite du foncteur d'inclusion
$D_{\QCoh, T}(\mathcal{O}_X) \to D_\QCoh(\mathcal{O}_X)$.
\end{enumerate}
\end{lemma}

\begin{proof}
Posons $U = X \setminus T$ et notons $j : U \to X$ l'inclusion. D'après
Cohomologie, lemme \ref{cohomology-lemma-triangle-sections-with-support-sheaves},
il existe un triangle distingué
$$
i_*R\mathcal{H}_T(E) \to E \to Rj_*(E|_U) \to i_*R\mathcal{H}_Z(E)[1]
$$
dans $D(\mathcal{O}_X)$. D'après le lemme \ref{lemma-quasi-coherence-direct-image},
le complexe $Rj_*(E|_U)$ a des faisceaux de cohomologie quasi-cohérents
(c'est ici que l'on utilise la rétrocompacité de $U$ dans $X$).
On obtient ainsi (1). L'assertion (2) en résulte, compte tenu de
l'adjonction des foncteurs donnée dans
Cohomologie, lemme \ref{cohomology-lemma-complexes-with-support-on-closed}.
\end{proof}

\begin{lemma}
\label{lemma-support-direct-sums}
Soit $X$ un schéma. Soit $T \subset X$ une partie fermée telle que
$X \setminus T$ soit un ouvert rétrocompact de $X$. Alors, pour une famille
d'objets $E_i$, $i \in I$ de $D_\QCoh(\mathcal{O}_X)$, on a
$R\mathcal{H}_T(\bigoplus E_i) = \bigoplus R\mathcal{H}_T(E_i)$.
\end{lemma}

\begin{proof}
Posons $U = X \setminus T$ et notons $j : U \to X$ l'inclusion. D'après
Cohomologie, lemme \ref{cohomology-lemma-triangle-sections-with-support-sheaves},
il existe un triangle distingué
$$
i_*R\mathcal{H}_T(E) \to E \to Rj_*(E|_U) \to i_*R\mathcal{H}_Z(E)[1]
$$
dans $D(\mathcal{O}_X)$ pour tout $E$ dans $D(\mathcal{O}_X)$. Le foncteur
$E \mapsto Rj_*(E|_U)$ commute aux sommes directes sur $D_\QCoh(\mathcal{O}_X)$
d'après le lemme \ref{lemma-quasi-coherence-pushforward-direct-sums}.
Il en résulte qu'il en est de même du foncteur $i_* \circ R\mathcal{H}_T$
(nous omettons les détails). Puisque $i_* : D(i^{-1}\mathcal{O}_X) \to D_T(\mathcal{O}_X)$
est une équivalence
(Cohomologie, lemme \ref{cohomology-lemma-complexes-with-support-on-closed}),
on conclut.
\end{proof}

\begin{remark}
\label{remark-support-c-equations}
Soit $X$ un schéma. Soient $f_1, \ldots, f_c \in \Gamma(X, \mathcal{O}_X)$.
Notons $Z \subset X$ le sous-schéma fermé défini par $f_1, \ldots, f_c$.
Pour $0 \leq p < c$ et $1 \leq i_0 < \ldots < i_p \leq c$, on note
$U_{i_0 \ldots i_p} \subset X$ le sous-schéma ouvert sur lequel
$f_{i_0} \ldots f_{i_p}$ est inversible. Pour tout $\mathcal{O}_X$-module
$\mathcal{F}$, on pose
$$
\mathcal{F}_{i_0 \ldots i_p} =
(U_{i_0 \ldots i_p} \to X)_*(\mathcal{F}|_{U_{i_0 \ldots i_p}})
$$
Dans cette situation, le {\it complexe de {\v C}ech alterné augmenté}
est le complexe de $\mathcal{O}_X$-modules
\begin{equation}
\label{equation-extended-alternating}
0 \to \mathcal{F} \to
\bigoplus\nolimits_{i_0} \mathcal{F}_{i_0} \to
\ldots \to
\bigoplus\nolimits_{i_0 < \ldots < i_p} \mathcal{F}_{i_0 \ldots i_p} \to
\ldots \to \mathcal{F}_{1 \ldots c} \to 0
\end{equation}
où $\mathcal{F}$ est placé en degré $0$. Les morphismes se construisent
comme suit. Étant donnés
$1 \leq i_0 < \ldots < i_{p + 1} \leq c$ et $0 \leq j \leq p + 1$, on
dispose du morphisme canonique
$$
\mathcal{F}_{i_0 \ldots \hat i_j \ldots i_{p + 1}} \to
\mathcal{F}_{i_0 \ldots i_p}
$$
provenant de l'inclusion
$U_{i_0 \ldots i_p} \subset U_{i_0 \ldots \hat i_j \ldots i_{p + 1}}$.
Les différentielles du complexe alterné augmenté sont définies à l'aide de ces
morphismes canoniques, affectés du signe $(-1)^j$.
\end{remark}

\begin{lemma}
\label{lemma-extended-alternating-zero}
Avec $X$, $f_1, \ldots, f_c \in \Gamma(X, \mathcal{O}_X)$ et
$\mathcal{F}$ comme dans la remarque \ref{remark-support-c-equations},
la restriction du complexe (\ref{equation-extended-alternating}) à
$X \setminus Z$ est acyclique.
\end{lemma}

\noindent
Signalons que ce lemme vaut plus généralement pour tout
complexe de {\v C}ech alterné augmenté défini comme dans la
remarque \ref{remark-support-c-equations} à partir d'un recouvrement
ouvert fini $X \setminus Z = U_1 \cup \ldots \cup U_c$.

\begin{proof}
Soit $W \subset X \setminus Z$ un ouvert. En évaluant le complexe
de faisceaux (\ref{equation-extended-alternating}) sur $W$, on obtient le
complexe
$$
\mathcal{F}(W) \to \bigoplus\nolimits_{i_0} \mathcal{F}(U_{i_0} \cap W) \to
\bigoplus\nolimits_{i_0 < i_1} \mathcal{F}(U_{i_0i_1} \cap W) \to \ldots
$$
Autrement dit, on obtient le complexe de {\v C}ech ordonné augmenté
associé au recouvrement $W = \bigcup U_i \cap W$ et à l'ordre
usuel sur $\{1, \ldots, c\}$ ; voir
Cohomologie, section \ref{cohomology-section-alternating-cech}.
D'après Cohomologie, lemme \ref{cohomology-lemma-alternating-cech-trivial},
ce complexe est homotope à zéro dès que $W$ est contenu dans
$V(f_i)$ pour un certain $1 \leq i \leq c$. Cela achève la démonstration.
\end{proof}

\begin{remark}
\label{remark-extended-alternating-map-to-support}
Soient $X$, $f_1, \ldots, f_c \in \Gamma(X, \mathcal{O}_X)$ et
$\mathcal{F}$ comme dans la remarque \ref{remark-support-c-equations}.
Notons $\mathcal{F}^\bullet$ le complexe
(\ref{equation-extended-alternating}). D'après
le lemme \ref{lemma-extended-alternating-zero},
les faisceaux de cohomologie de $\mathcal{F}^\bullet$
sont à support dans $Z$ ; par suite, $\mathcal{F}^\bullet$ est un objet de
$D_Z(\mathcal{O}_X)$. D'autre part, l'égalité
$\mathcal{F}^0 = \mathcal{F}$ définit un morphisme canonique
$\mathcal{F}^\bullet \to \mathcal{F}$ dans $D(\mathcal{O}_X)$.
Puisque $i_* \circ R\mathcal{H}_Z$ est adjoint à droite du
foncteur d'inclusion $D_Z(\mathcal{O}_X) \to D(\mathcal{O}_X)$, voir
Cohomologie, lemme \ref{cohomology-lemma-complexes-with-support-on-closed},
on obtient un diagramme commutatif canonique
$$
\xymatrix{
\mathcal{F}^\bullet \ar[rd] \ar[rr] & & \mathcal{F} \\
& i_*R\mathcal{H}_Z(\mathcal{F}) \ar[ru]
}
$$
dans $D(\mathcal{O}_X)$, fonctoriel en le $\mathcal{O}_X$-module $\mathcal{F}$.
\end{remark}

\begin{lemma}
\label{lemma-extended-alternating-represented}
Soient $X$, $f_1, \ldots, f_c \in \Gamma(X, \mathcal{O}_X)$ et
$\mathcal{F}$ comme dans la remarque \ref{remark-support-c-equations}.
Si $\mathcal{F}$ est quasi-cohérent, le complexe
(\ref{equation-extended-alternating}) représente
$i_* R\mathcal{H}_Z(\mathcal{F})$ dans $D_Z(\mathcal{O}_X)$.
\end{lemma}

\begin{proof}
Notons $\mathcal{F}^\bullet$ le complexe
(\ref{equation-extended-alternating}).
L'énoncé du lemme signifie que le morphisme
$\mathcal{F}^\bullet \to i_*R\mathcal{H}_Z(\mathcal{F})$
de la remarque \ref{remark-extended-alternating-map-to-support}
est un isomorphisme. Puisque $\mathcal{F}^\bullet$ appartient à
$D_Z(\mathcal{O}_X)$ (voir la remarque citée), on a
$i_*R\mathcal{H}_Z(\mathcal{F}^\bullet) = \mathcal{F}^\bullet$
d'après Cohomologie, lemme \ref{cohomology-lemma-complexes-with-support-on-closed}.
Le morphisme $U_{i_0 \ldots i_p} \to X$ est affine,
car il est donné, au-dessus des ouverts affines de $X$, par inversion de la fonction
$f_{i_0} \ldots f_{i_p}$. On a donc
$$
\mathcal{F}_{i_0 \ldots i_p} =
(U_{i_0 \ldots i_p} \to X)_*\mathcal{F}|_{U_{i_0 \ldots i_p}} =
R(U_{i_0 \ldots i_p} \to X)_*\mathcal{F}|_{U_{i_0 \ldots i_p}}
$$
d'après Cohomologie des schémas, lemme \ref{coherent-lemma-relative-affine-vanishing},
et l'hypothèse de quasi-cohérence de $\mathcal{F}$. On en déduit que
$R\mathcal{H}_Z(\mathcal{F}_{i_0 \ldots i_p}) = 0$ d'après Cohomologie, lemme
\ref{cohomology-lemma-sections-support-in-closed-disjoint-open}.
Ainsi $i_*R\mathcal{H}_Z(\mathcal{F}^p) = 0$ pour $p > 0$.
En réunissant ces résultats, on obtient
$$
\mathcal{F}^\bullet = i_*R\mathcal{H}_Z(\mathcal{F}^\bullet) =
i_*R\mathcal{H}_Z(\mathcal{F})
$$
comme voulu.
\end{proof}

\begin{lemma}
\label{lemma-supported-trivial-vanishing}
Soit $X$ un schéma. Soit $T \subset X$ une partie fermée qui peut
être définie localement par au plus $c$ éléments du faisceau structural.
Alors $\mathcal{H}^i_Z(\mathcal{F}) = 0$ pour $i > c$ et tout
$\mathcal{O}_X$-module quasi-cohérent $\mathcal{F}$.
\end{lemma}

\begin{proof}
Cela résulte immédiatement de la description locale de
$R\mathcal{H}_T(\mathcal{F})$ donnée dans
le lemme \ref{lemma-extended-alternating-represented}.
\end{proof}

\begin{lemma}
\label{lemma-supported-vanishing}
Soit $X$ un schéma. Soit $Z \subset X$ une partie fermée qui peut
être définie localement par une suite régulière au sens de Koszul ayant $c$ éléments.
Alors $\mathcal{H}^i_Z(\mathcal{F}) = 0$ pour $i \not = c$ et tout
$\mathcal{O}_X$-module plat et quasi-cohérent $\mathcal{F}$.
\end{lemma}

\begin{proof}
D'après la description de $R\mathcal{H}_Z(\mathcal{F})$ donnée dans
le lemme \ref{lemma-extended-alternating-represented}, on se ramène
à l'énoncé algébrique suivant : étant donnés un anneau $R$,
une suite régulière au sens de Koszul $f_1, \ldots, f_c \in R$ et un
$R$-module plat $M$, le complexe de {\v C}ech alterné augmenté
$M \to \bigoplus\nolimits_{i_0} M_{f_{i_0}} \to
\bigoplus\nolimits_{i_0 < i_1} M_{f_{i_0}f_{i_1}} \to
\ldots \to M_{f_1 \ldots f_c}$
de Compléments sur l'algèbre, section \ref{more-algebra-section-alternating-cech},
n'a de cohomologie qu'en degré $c$. D'après Compléments sur l'algèbre, lemme
\ref{more-algebra-lemma-vanishing-extended-alternating-koszul},
on obtient l'annulation voulue pour le complexe de
{\v C}ech alterné augmenté de $R$. Puisque le complexe associé à $M$ s'obtient
en tensorisant celui-ci par le $R$-module plat $M$
(Compléments sur l'algèbre, lemme
\ref{more-algebra-lemma-extended-alternating-form-module}),
on conclut.
\end{proof}

\begin{remark}
\label{remark-supported-map-c-equations}
Soient $X$, $f_1, \ldots, f_c \in \Gamma(X, \mathcal{O}_X)$ et
$\mathcal{F}$ comme dans la remarque \ref{remark-support-c-equations}.
Il existe un morphisme $\mathcal{O}_X|_Z$-linéaire canonique
$$
c_{f_1, \ldots, f_c} :
i^*\mathcal{F}
\longrightarrow
\mathcal{H}^c_Z(\mathcal{F})
$$
fonctoriel en $\mathcal{F}$. En effet, si l'on note $\mathcal{F}^\bullet$ le
complexe de {\v C}ech alterné augmenté (\ref{equation-extended-alternating}),
on dispose du morphisme canonique
$\mathcal{F}^\bullet \to i_*R\mathcal{H}_Z(\mathcal{F})$
de la remarque \ref{remark-extended-alternating-map-to-support}.
Il définit un morphisme canonique
$$
\Coker\left(\bigoplus \mathcal{F}_{1 \ldots \hat i \ldots c} \to
\mathcal{F}_{1 \ldots c}\right)
\longrightarrow
i_*\mathcal{H}^c_Z(\mathcal{F})
$$
sur les faisceaux de cohomologie en degré $c$.
Étant donnée une section locale $s$ de $\mathcal{F}$, on peut considérer la
section locale
$$
\frac{s}{f_1 \ldots f_c}
$$
de $\mathcal{F}_{1 \ldots c}$. La classe de cette section dans le conoyau
ci-dessus ne dépend que de $s$ modulo l'image de
$(f_1, \ldots, f_c) : \mathcal{F}^{\oplus c} \to \mathcal{F}$.
Puisque $i_*i^*\mathcal{F}$ est égal au conoyau de
$(f_1, \ldots, f_c) : \mathcal{F}^{\oplus c} \to \mathcal{F}$,
on obtient un morphisme de $\mathcal{O}_X$-modules
$i_*i^*\mathcal{F} \to i_*\mathcal{H}_Z^c(\mathcal{F})$.
La pleine fidélité de $i_*$ fournit alors le morphisme $c_{f_1, \ldots, f_c}$.
\end{remark}

\begin{example}
\label{example-affine-supported-map-c-equations}
Soient $X = \Spec(A)$ affine, $f_1, \ldots, f_c \in A$ et
$\mathcal{F} = \widetilde{M}$ pour un $A$-module $M$. Le morphisme
$c_{f_1, \ldots, f_c}$ de la remarque \ref{remark-supported-map-c-equations}
se décrit comme le morphisme
$$
M/(f_1, \ldots, f_c)M
\longrightarrow
\Coker\left(
\bigoplus M_{f_1 \ldots \hat f_i \ldots f_c} \to M_{f_1 \ldots f_c}
\right)
$$
qui envoie la classe de $s \in M$ sur la classe de $s/f_1 \ldots f_c$
dans le conoyau.
\end{example}

\begin{lemma}
\label{lemma-supported-map-determinant}
Soient $X$, $f_1, \ldots, f_c \in \Gamma(X, \mathcal{O}_X)$ et
$\mathcal{F}$ comme dans la remarque \ref{remark-support-c-equations}.
Soient $a_{ji} \in \Gamma(X, \mathcal{O}_X)$ pour $1 \leq i, j \leq c$,
et posons $g_j = \sum_{i = 1, \ldots, c} a_{ji}f_i$. Supposons que $g_1, \ldots, g_c$
définissent $Z$ comme sous-schéma. Si $\mathcal{F}$ est quasi-cohérent, on a
$$
c_{f_1, \ldots, f_c} = \det(a_{ji}) c_{g_1, \ldots, g_c}
$$
où $c_{f_1, \ldots, f_c}$ et $c_{g_1, \ldots, g_c}$ sont définis dans la
remarque \ref{remark-supported-map-c-equations}.
\end{lemma}

\begin{proof}
Nous montrerons que $c_{f_1, \ldots, f_c}(s) =
\det(a_{ij}) c_{g_1, \ldots, g_c}(s)$ en tant que sections globales de
$\mathcal{H}_Z(\mathcal{F})$, pour tout $s \in \mathcal{F}(X)$.
Cela suffit, car on obtient alors le même résultat pour les
sections sur tout sous-schéma ouvert de $X$.
Pour cela, pour $1 \leq i_0 < \ldots < i_p \leq c$
et $1 \leq j_0 < \ldots < j_q \leq c$, notons
$U_{i_0 \ldots i_p} \subset X$,
$V_{j_0 \ldots j_q} \subset X$ et
$W_{i_0 \ldots i_p, j_0 \ldots j_q} \subset X$
les sous-schémas ouverts où, respectivement, $f_{i_0} \ldots f_{i_p}$ est inversible,
$g_{j_0} \ldots g_{j_q}$ est inversible, et
$f_{i_0} \ldots f_{i_p}g_{j_0} \ldots g_{j_q}$ est inversible.
On note $\mathcal{F}_{i_0 \ldots i_p}$,
resp.\ $\mathcal{F}'_{j_0 \ldots j_q}$,
$\mathcal{F}''_{i_0 \ldots i_p, j_0 \ldots j_q}$
l'image directe sur $X$ de la restriction de $\mathcal{F}$ à
$U_{i_0 \ldots i_p}$, resp.\ $V_{j_0 \ldots j_q}$,
resp.\ $W_{i_0 \ldots i_p, j_0 \ldots j_q}$.
On obtient alors trois complexes de {\v C}ech alternés augmentés
$$
\mathcal{F}^\bullet :
\mathcal{F} \to \bigoplus\nolimits_{i_0} \mathcal{F}_{i_0} \to
\bigoplus\nolimits_{i_0 < i_1} \mathcal{F}_{i_0i_1} \to \ldots
$$
et
$$
(\mathcal{F}')^\bullet :
\mathcal{F} \to \bigoplus\nolimits_{j_0} \mathcal{F}'_{j_0} \to
\bigoplus\nolimits_{j_0 < j_1} \mathcal{F}'_{j_0j_1} \to \ldots
$$
et
$$
\begin{gathered}
(\mathcal{F}'')^\bullet :
\mathcal{F} \to
\bigoplus\nolimits_{i_0} \mathcal{F}_{i_0} \oplus
\bigoplus\nolimits_{j_0} \mathcal{F}'_{j_0} \\
{}\to
\bigoplus\nolimits_{i_0 < i_1} \mathcal{F}_{i_0i_1} \oplus
\bigoplus\nolimits_{i_0, j_0} \mathcal{F}''_{i_0, j_0} \oplus
\bigoplus\nolimits_{j_0 < j_1} \mathcal{F}'_{j_0j_1} \to
\ldots
\end{gathered}
$$
dont les différentielles sont celles qui interviennent dans la définition de
(\ref{equation-extended-alternating}).
On dispose de morphismes de complexes
$$
(\mathcal{F}'')^\bullet \to \mathcal{F}^\bullet
\quad\text{et}\quad
(\mathcal{F}'')^\bullet \to (\mathcal{F}')^\bullet
$$
donnés terme à terme par les projections (et induisant donc
l'identité en degré $0$). Notons que, d'après
le lemme \ref{lemma-extended-alternating-represented},
chacun de ces complexes représente
$i_*R\mathcal{H}_Z(\mathcal{F})$, et ces morphismes
représentent l'identité de cet objet. Il suffit donc
de trouver un élément
$$
\sigma \in H^c((\mathcal{F}'')^\bullet(X))
$$
dont les images par ces deux morphismes soient $c_{f_1, \ldots, f_c}(s)$ et $\det(a_{ji})c_{g_1, \ldots, g_c}(s)$.
On peut en fait donner explicitement un cocycle
représentant $\sigma$. On prend
$$
\sigma_{1 \ldots c} = \frac{s}{f_1 \ldots f_c} \in \mathcal{F}_{1 \ldots c}(X)
\quad\text{et}\quad
\sigma'_{1 \ldots c} = \frac{\det(a_{ji})s}{g_1 \ldots g_c} \in
\mathcal{F}'_{1 \ldots c}(X)
$$
et l'on prend
$$
\sigma_{i_0 \ldots i_p, j_0 \ldots j_{c - p - 2}} =
\frac{\lambda(i_0 \ldots i_p, j_0 \ldots j_{c - p - 2})s}%
{f_{i_0} \ldots f_{i_p}g_{j_0} \ldots g_{j_{c - p - 2}}}
\in \mathcal{F}''_{i_0 \ldots i_p, j_0 \ldots j_{c - p - 2}}(X)
$$
où $\lambda(i_0 \ldots i_p, j_0 \ldots j_{c - p - 2})$
est le coefficient de $e_1 \wedge \ldots \wedge e_c$ dans
l'expression formelle
$$
e_{i_0} \wedge \ldots \wedge e_{i_p} \wedge
(a_{j_01} e_1 + \ldots + a_{j_0c}e_c) \wedge \ldots \wedge
(a_{j_{c - p - 2}1} e_1 + \ldots + a_{j_{c - p - 2}c}e_c)
$$
Pour vérifier que $\sigma$ est un cocycle, il faut montrer que, pour
$1 \leq i_0 < \ldots < i_p \leq c$ et
$1 \leq j_0 < \ldots < j_{c - p - 1} \leq c$,
on a
\begin{align*}
0 & =
\sum\nolimits_{a = 0, \ldots, p} (-1)^a
f_{i_a} \lambda(i_0 \ldots \hat i_a \ldots i_p, j_0 \ldots j_{c - p - 1}) \\
& +
\sum\nolimits_{b = 0, \ldots, c - p - 1} (-1)^{p + b + 1}g_{j_b}
\lambda(i_0 \ldots i_p, j_0 \ldots \hat j_b \ldots j_{c - p - 1})
\end{align*}
Le plus simple est peut-être d'observer que l'expression formelle
$$
\xi = e_{i_0} \wedge \ldots \wedge e_{i_p} \wedge
(a_{j_01} e_1 + \ldots + a_{j_0c}e_c) \wedge \ldots \wedge
(a_{j_{c - p - 1}1} e_1 + \ldots + a_{j_{c - p - 1}c}e_c)
$$
est égale à $0$, puisqu'elle appartient à la $(c + 1)$-ième puissance extérieure du module libre
de base $e_1, \ldots, e_c$, et que l'expression ci-dessus est l'image de
$\xi$ par la différentielle de Koszul qui envoie $e_i \to f_i$. Nous omettons quelques
détails.
\end{proof}

\begin{lemma}
\label{lemma-supported-map-global}
Soit $X$ un schéma. Soit $Z \to X$ une immersion fermée de présentation finie
dont le faisceau conormal $\mathcal{C}_{Z/X}$ est localement libre de rang $c$.
Il existe alors un morphisme canonique
$$
c :
\wedge^c(\mathcal{C}_{Z/X})^\vee \otimes_{\mathcal{O}_Z} i^*\mathcal{F}
\longrightarrow
\mathcal{H}_Z^c(\mathcal{F})
$$
fonctoriel en le module quasi-cohérent $\mathcal{F}$.
\end{lemma}

\begin{proof}
Cela résulte de la construction de la
remarque \ref{remark-supported-map-c-equations}
et de l'indépendance du choix
des générateurs du faisceau d'idéaux établie dans le
lemme \ref{lemma-supported-map-determinant}.
Nous omettons quelques détails.
\end{proof}

\begin{remark}
\label{remark-supported-functorial}
Soit $g : X' \to X$ un morphisme de schémas. Soient
$f_1, \ldots, f_c \in \Gamma(X, \mathcal{O}_X)$.
Posons $f'_i = g^\sharp(f_i) \in \Gamma(X', \mathcal{O}_{X'})$.
Notons $Z \subset X$, resp.\ $Z' \subset X'$, le sous-schéma
fermé défini par $f_1, \ldots, f_c$, resp.\ $f'_1, \ldots, f'_c$.
Alors $Z' = Z \times_X X'$. Notons $h : Z' \to Z$ le morphisme
de schémas induit.
Soit $\mathcal{F}$ un $\mathcal{O}_X$-module.
Posons $\mathcal{F}' = g^*\mathcal{F}$.
Dans cette situation, si $\mathcal{F}$ est quasi-cohérent, le diagramme
$$
\xymatrix{
(i')^{-1}\mathcal{O}_{X'} \otimes_{h^{-1}i^{-1}\mathcal{O}_X}
h^{-1}\mathcal{H}^c_Z(\mathcal{F}) \ar[r] &
\mathcal{H}_{Z'}^c(\mathcal{F}') \\
h^*i^*\mathcal{F} \ar[r] \ar[u]_-{c_{f_1, \ldots, f_c}} &
(i')^*\mathcal{F}' \ar[u]^-{c_{f'_1, \ldots, f'_c}}
}
$$
est commutatif ; la flèche horizontale supérieure est
le morphisme de Cohomologie, remarque \ref{cohomology-remark-support-functorial},
sur les faisceaux de cohomologie en degré $c$. En effet,
notons $\mathcal{F}^\bullet$, resp.\ $(\mathcal{F}')^\bullet$,
le complexe de {\v C}ech alterné augmenté construit dans la
remarque \ref{remark-support-c-equations}
à l'aide de $\mathcal{F}, f_1, \ldots, f_c$,
resp.\ $\mathcal{F}', f'_1, \ldots, f'_c$.
Notons que $(\mathcal{F}')^\bullet = g^*\mathcal{F}^\bullet$.
Alors, sans supposer $\mathcal{F}$ quasi-cohérent, le diagramme
$$
\xymatrix{
i'_* L(g|_{Z'})^* R\mathcal{H}_Z(\mathcal{F}) \ar[r] \ar@{=}[d] &
i'_*R\mathcal{H}_{Z'}(Lg^*\mathcal{F}) \ar[d] \\
Lg^*i_*R\mathcal{H}_Z(\mathcal{F}) &
i'_*R\mathcal{H}_{Z'}(\mathcal{F}') \\
Lg^*(\mathcal{F}^\bullet) \ar[u] \ar[r] &
(\mathcal{F}')^\bullet \ar[u]
}
$$
est commutatif, où $g|_{Z'} : (Z', (i')^{-1}\mathcal{O}_{X'}) \to
(Z, i^{-1}\mathcal{O}_X)$ est le morphisme induit d'espaces annelés.
La flèche horizontale supérieure est donnée dans
Cohomologie, remarque \ref{cohomology-remark-support-functorial},
qui explique aussi le signe d'égalité.
Les flèches ascendantes proviennent de la
remarque \ref{remark-extended-alternating-map-to-support}.
La flèche horizontale inférieure est le morphisme $Lg^*\mathcal{F}^\bullet
\to g^*\mathcal{F}^\bullet = (\mathcal{F}')^\bullet$, et la flèche
descendante est induite par
$Lg^*\mathcal{F} \to g^*\mathcal{F} = \mathcal{F}'$.
Le diagramme est commutatif, car les deux parcours donnent
deux flèches $Lg^*\mathcal{F}^\bullet \to
i'_*R\mathcal{H}_{Z'}(\mathcal{F}')$ dont le composé avec
$i'_*R\mathcal{H}_{Z'}(\mathcal{F}') \to \mathcal{F}'$
est le morphisme canonique $Lg^*\mathcal{F}^\bullet \to \mathcal{F}'$.
Nous omettons quelques détails. La commutativité du premier diagramme
s'obtient alors en passant aux faisceaux de cohomologie en degré
$c$ dans ce diagramme et en utilisant la compatibilité à l'image réciproque de la construction du morphisme
$i^*\mathcal{F} \to
\Coker(\bigoplus \mathcal{F}_{1 \ldots \hat i \ldots c} \to
\mathcal{F}_{1 \ldots c})$
employée dans la remarque \ref{remark-supported-map-c-equations}.
\end{remark}







\section{Le cohérateur}
\label{section-coherator}

\noindent
Soit $X$ un schéma. Le {\it cohérateur} est un foncteur
$$
Q_X :
\textit{Mod}(\mathcal{O}_X)
\longrightarrow
\QCoh(\mathcal{O}_X)
$$
adjoint à droite du foncteur d'inclusion
$\QCoh(\mathcal{O}_X) \to \textit{Mod}(\mathcal{O}_X)$.
Il existe pour tout schéma $X$ ; de plus, le morphisme d'adjonction
$Q_X(\mathcal{F}) \to \mathcal{F}$ est un isomorphisme pour tout
module quasi-cohérent $\mathcal{F}$ ; voir
Propriétés, proposition \ref{properties-proposition-coherator}.
Puisque $Q_X$ est exact à gauche (comme adjoint à droite), on peut considérer son
foncteur dérivé à droite
$$
RQ_X :
D(\mathcal{O}_X)
\longrightarrow
D(\QCoh(\mathcal{O}_X)).
$$
Puisque $Q_X$ est adjoint à droite du foncteur d'inclusion
$\QCoh(\mathcal{O}_X) \to \textit{Mod}(\mathcal{O}_X)$,
le foncteur $RQ_X$ est adjoint à droite du foncteur canonique
$D(\QCoh(\mathcal{O}_X)) \to D(\mathcal{O}_X)$ d'après
Catégories dérivées, lemme \ref{derived-lemma-derived-adjoint-functors}.

\medskip\noindent
Dans cette section, nous étudierons le foncteur $RQ_X$. Dans la
section \ref{section-better-coherator},
nous étudierons l'adjoint à droite, étroitement lié au précédent, du foncteur d'inclusion
$D_\QCoh(\mathcal{O}_X) \to D(\mathcal{O}_X)$ (lorsqu'il existe).

\begin{lemma}
\label{lemma-affine-pushforward}
Soit $f : X \to Y$ un morphisme affine de schémas.
Alors $f_*$ définit un foncteur dérivé
$f_* : D(\QCoh(\mathcal{O}_X)) \to D(\QCoh(\mathcal{O}_Y))$.
Ce foncteur rend le diagramme
$$
\xymatrix{
D(\QCoh(\mathcal{O}_X)) \ar[d]_{f_*} \ar[r] &
D_\QCoh(\mathcal{O}_X) \ar[d]^{Rf_*} \\
D(\QCoh(\mathcal{O}_Y)) \ar[r] &
D_\QCoh(\mathcal{O}_Y)
}
$$
commutatif.
\end{lemma}

\begin{proof}
Le foncteur
$f_* : \QCoh(\mathcal{O}_X) \to \QCoh(\mathcal{O}_Y)$
est exact ; voir
Cohomologie des schémas, lemme \ref{coherent-lemma-relative-affine-vanishing}.
Par suite, $f_*$ définit un foncteur dérivé
$f_* : D(\QCoh(\mathcal{O}_X)) \to D(\QCoh(\mathcal{O}_Y))$
en appliquant simplement $f_*$ à un complexe représentant quelconque ; voir
Catégories dérivées, lemme \ref{derived-lemma-right-derived-exact-functor}.
Le diagramme est commutatif d'après le lemme \ref{lemma-pushforward-affine-morphism}.
\end{proof}

\begin{lemma}
\label{lemma-flat-pushforward-coherator}
Soit $f : X \to Y$ un morphisme de schémas. Supposons $f$
quasi-compact, quasi-séparé et plat. Alors, en notant
$$
\Phi : D(\QCoh(\mathcal{O}_X)) \to D(\QCoh(\mathcal{O}_Y))
$$
le foncteur dérivé à droite de
$f_* : \QCoh(\mathcal{O}_X) \to \QCoh(\mathcal{O}_Y)$,
on a $RQ_Y \circ Rf_* = \Phi \circ RQ_X$.
\end{lemma}

\begin{proof}
Nous le démontrerons en montrant que $RQ_Y \circ Rf_*$ et $\Phi \circ RQ_X$
sont adjoints à droite du même foncteur
$D(\QCoh(\mathcal{O}_Y)) \to D(\mathcal{O}_X)$.

\medskip\noindent
Puisque $f$ est quasi-compact et quasi-séparé,
$f_*$ préserve la quasi-cohérence ; voir
Schémas, lemme \ref{schemes-lemma-push-forward-quasi-coherent}.
Rappelons que $\QCoh(\mathcal{O}_X)$ est une catégorie abélienne de Grothendieck
(Propriétés, proposition \ref{properties-proposition-coherator}).
Par suite, tout $K$ dans $D(\QCoh(\mathcal{O}_X))$
peut être représenté par un complexe K-injectif $\mathcal{I}^\bullet$
de $\QCoh(\mathcal{O}_X)$ ; voir
Injectifs, théorème
\ref{injectives-theorem-K-injective-embedding-grothendieck}.
On peut alors définir $\Phi(K) = f_*\mathcal{I}^\bullet$.

\medskip\noindent
Puisque $f$ est plat, le foncteur $f^*$ est exact. Par suite, $f^*$ définit
$f^* : D(\mathcal{O}_Y) \to D(\mathcal{O}_X)$ ainsi que
$f^* : D(\QCoh(\mathcal{O}_Y)) \to D(\QCoh(\mathcal{O}_X))$.
Le foncteur $f^* = Lf^* : D(\mathcal{O}_Y) \to D(\mathcal{O}_X)$
est adjoint à gauche de
$Rf_* : D(\mathcal{O}_X) \to D(\mathcal{O}_Y)$ ;
voir Cohomologie, lemme \ref{cohomology-lemma-adjoint}.
De même, le foncteur
$f^* : D(\QCoh(\mathcal{O}_Y)) \to D(\QCoh(\mathcal{O}_X))$
est adjoint à gauche de
$\Phi : D(\QCoh(\mathcal{O}_X)) \to D(\QCoh(\mathcal{O}_Y))$
d'après Catégories dérivées, lemme \ref{derived-lemma-derived-adjoint-functors}.

\medskip\noindent
Soient $A$ un objet de $D(\QCoh(\mathcal{O}_Y))$ et
$E$ un objet de $D(\mathcal{O}_X)$. Alors
\begin{align*}
\Hom_{D(\QCoh(\mathcal{O}_Y))}(A, RQ_Y(Rf_*E))
& =
\Hom_{D(\mathcal{O}_Y)}(A, Rf_*E) \\
& =
\Hom_{D(\mathcal{O}_X)}(f^*A, E) \\
& =
\Hom_{D(\QCoh(\mathcal{O}_X))}(f^*A, RQ_X(E)) \\
& =
\Hom_{D(\QCoh(\mathcal{O}_Y))}(A, \Phi(RQ_X(E)))
\end{align*}
On en déduit l'assertion voulue.
\end{proof}

\begin{lemma}
\label{lemma-affine-coherator}
Soit $X = \Spec(A)$ un schéma affine. Alors
\begin{enumerate}
\item $Q_X : \textit{Mod}(\mathcal{O}_X) \to \QCoh(\mathcal{O}_X)$
est le foncteur
qui envoie $\mathcal{F}$ sur le $\mathcal{O}_X$-module quasi-cohérent
associé au $A$-module $\Gamma(X, \mathcal{F})$,
\item $RQ_X : D(\mathcal{O}_X) \to D(\QCoh(\mathcal{O}_X))$
est le foncteur qui envoie $E$ sur le complexe de
$\mathcal{O}_X$-modules quasi-cohérents associé à l'objet $R\Gamma(X, E)$ de $D(A)$,
\item restreint à $D_\QCoh(\mathcal{O}_X)$, le foncteur
$RQ_X$ définit un quasi-inverse de (\ref{equation-compare}).
\end{enumerate}
\end{lemma}

\begin{proof}
Le foncteur $Q_X$ est le foncteur
$$
\mathcal{F} \mapsto \widetilde{\Gamma(X, \mathcal{F})}
$$
d'après Schémas, lemme \ref{schemes-lemma-compare-constructions}.
Cela entraîne immédiatement (1) et (2). La troisième assertion
résulte de la démonstration du
lemme \ref{lemma-affine-compare-bounded}.
\end{proof}

\noindent
Nous pouvons maintenant établir un critère
pour que le foncteur $D(\QCoh(\mathcal{O}_X)) \to D_\QCoh(\mathcal{O}_X)$
soit une équivalence.

\begin{lemma}
\label{lemma-argument-proves}
Soit $X$ un schéma quasi-compact et quasi-séparé. Supposons que,
pour tout ouvert affine $U \subset X$, le foncteur dérivé à droite
$$
\Phi : D(\QCoh(\mathcal{O}_U)) \to D(\QCoh(\mathcal{O}_X))
$$
du foncteur exact à gauche
$j_* : \QCoh(\mathcal{O}_U) \to \QCoh(\mathcal{O}_X)$
s'inscrive dans un diagramme commutatif
$$
\xymatrix{
D(\QCoh(\mathcal{O}_U)) \ar[d]_\Phi \ar[r]_{i_U} &
D_\QCoh(\mathcal{O}_U) \ar[d]^{Rj_*} \\
D(\QCoh(\mathcal{O}_X)) \ar[r]^{i_X} &
D_\QCoh(\mathcal{O}_X)
}
$$
Alors le foncteur (\ref{equation-compare})
$$
D(\QCoh(\mathcal{O}_X))
\longrightarrow
D_\QCoh(\mathcal{O}_X)
$$
est une équivalence dont un quasi-inverse est donné par $RQ_X$.
\end{lemma}

\begin{proof}
Soit $E$ un objet de $D_\QCoh(\mathcal{O}_X)$ et
soit $A$ un objet de $D(\QCoh(\mathcal{O}_X))$.
Il faut montrer que les morphismes d'adjonction
$$
A \to RQ_X(i_X(A))
\quad\text{et}\quad
i_X(RQ_X(E)) \to E
$$
sont des isomorphismes. Considérons l'hypothèse
$H_n$ : les morphismes d'adjonction ci-dessus sont des isomorphismes
dès que $E$ et $i_X(A)$ sont à support
(définition \ref{definition-supported-on})
dans une partie fermée de $X$
contenue dans la réunion de $n$ ouverts affines de $X$.
Nous démontrerons $H_n$ par récurrence sur $n$.

\medskip\noindent
Cas initial : $n = 0$. Dans ce cas, $E = 0$, donc le morphisme
$i_X(RQ_X(E)) \to E$ est un isomorphisme. De même, $i_X(A) = 0$.
Les faisceaux de cohomologie de $i_X(A)$ sont donc nuls. Puisque le foncteur d'inclusion
$\QCoh(\mathcal{O}_X) \to \textit{Mod}(\mathcal{O}_X)$
est pleinement fidèle et exact, on en déduit que les objets de cohomologie
de $A$ sont nuls, c'est-à-dire que $A = 0$ ; le morphisme
$A \to RQ_X(i_X(A))$ est donc lui aussi un isomorphisme.

\medskip\noindent
Étape de récurrence. Supposons que $E$ et $i_X(A)$ soient à support dans une
partie fermée $T$ de $X$ contenue dans $U_1 \cup \ldots \cup U_n$,
où les $U_i \subset X$ sont des ouverts affines. Posons $U = U_n$.
Considérons les triangles distingués
$$
A \to \Phi(A|_U) \to A' \to A[1]
\quad\text{et}\quad
E \to Rj_*(E|_U) \to E' \to E[1]
$$
où $\Phi$ est comme dans l'énoncé du lemme.
Notons que $E \to Rj_*(E|_U)$ est un quasi-isomorphisme au-dessus de $U = U_n$.
Puisque $i_X \circ \Phi = Rj_* \circ i_U$ par hypothèse
et que $i_X(A)|_U = i_U(A|_U)$,
le morphisme $i_X(A) \to i_X(\Phi(A|_U))$ est un quasi-isomorphisme au-dessus de $U$.
Ainsi $i_X(A')$ et $E'$ sont à support dans la partie
fermée $T \setminus U$ de $X$, contenue dans
$U_1 \cup \ldots \cup U_{n - 1}$.
Par hypothèse de récurrence, l'énoncé est vrai pour $A'$ et $E'$. D'après
Catégories dérivées, lemme \ref{derived-lemma-third-isomorphism-triangle},
il suffit de montrer que les morphismes
$$
\Phi(A|_U) \to RQ_X(i_X(\Phi(A|_U)))
\quad\text{et}\quad
i_X(RQ_X(Rj_*E|_U)) \to Rj_*(E|_U)
$$
sont des isomorphismes. Par hypothèse et d'après
le lemme \ref{lemma-flat-pushforward-coherator}
(le morphisme d'inclusion $j : U \to X$ est plat, quasi-compact et
quasi-séparé), on a
$$
RQ_X(i_X(\Phi(A|_U))) = RQ_X(Rj_*(i_U(A|_U))) = \Phi(RQ_U(i_U(A|_U)))
$$
et
$$
i_X(RQ_X(Rj_*(E|_U))) = i_X(\Phi(RQ_U(E|_U))) = Rj_*(i_U(RQ_U(E|_U)))
$$
Enfin, les morphismes
$$
A|_U \to RQ_U(i_U(A|_U))
\quad\text{et}\quad
i_U(RQ_U(E|_U)) \to E|_U
$$
sont des isomorphismes d'après le lemme \ref{lemma-affine-coherator}. Le résultat en découle.
\end{proof}

\begin{proposition}
\label{proposition-quasi-compact-affine-diagonal}
Soit $X$ un schéma quasi-compact à diagonale affine.
Alors le foncteur (\ref{equation-compare})
$$
D(\QCoh(\mathcal{O}_X))
\longrightarrow
D_\QCoh(\mathcal{O}_X)
$$
est une équivalence dont un quasi-inverse est donné par $RQ_X$.
\end{proposition}

\begin{proof}
Soit $U \subset X$ un ouvert affine. Alors le morphisme
$U \to X$ est affine d'après
Morphismes, lemme \ref{morphisms-lemma-affine-permanence}.
L'hypothèse du lemme \ref{lemma-argument-proves}
est donc vérifiée d'après le lemme \ref{lemma-affine-pushforward}, ce qui conclut.
\end{proof}

\begin{lemma}
\label{lemma-direct-image-coherator}
Soit $f : X \to Y$ un morphisme de schémas.
Supposons $X$ et $Y$ quasi-compacts et à diagonale affine.
Alors, en notant
$$
\Phi : D(\QCoh(\mathcal{O}_X)) \to D(\QCoh(\mathcal{O}_Y))
$$
le foncteur dérivé à droite de
$f_* : \QCoh(\mathcal{O}_X) \to \QCoh(\mathcal{O}_Y)$,
le diagramme
$$
\xymatrix{
D(\QCoh(\mathcal{O}_X)) \ar[d]_\Phi \ar[r] &
D_\QCoh(\mathcal{O}_X) \ar[d]^{Rf_*} \\
D(\QCoh(\mathcal{O}_Y)) \ar[r] &
D_\QCoh(\mathcal{O}_Y)
}
$$
est commutatif.
\end{lemma}

\begin{proof}
Observons que les flèches horizontales du diagramme sont
des équivalences de catégories d'après la
proposition \ref{proposition-quasi-compact-affine-diagonal}.
On peut donc identifier ces catégories (et de même pour
les autres schémas quasi-compacts à diagonale affine).
L'énoncé du lemme affirme que le morphisme canonique
$\Phi(K) \to Rf_*(K)$ est un isomorphisme pour tout $K$ dans
$D(\QCoh(\mathcal{O}_X))$. Notons que, si $K_1 \to K_2 \to K_3 \to K_1[1]$
est un triangle distingué dans $D(\QCoh(\mathcal{O}_X))$ et
si l'énoncé est vrai pour deux des trois objets, il l'est
aussi pour le troisième.

\medskip\noindent
Soit $U \subset X$ un ouvert affine. Puisque la diagonale de $X$ est affine,
le morphisme d'inclusion $j : U \to X$
est affine (Morphismes, lemme \ref{morphisms-lemma-affine-permanence}).
De même, le composé $g = f \circ j : U \to Y$ est affine.
Soit $\mathcal{I}^\bullet$ un complexe K-injectif dans $\QCoh(\mathcal{O}_U)$.
Puisque $j_* : \QCoh(\mathcal{O}_U) \to \QCoh(\mathcal{O}_X)$
admet un adjoint à gauche exact
$j^* : \QCoh(\mathcal{O}_X) \to \QCoh(\mathcal{O}_U)$,
le complexe $j_*\mathcal{I}^\bullet$ est K-injectif
dans $\QCoh(\mathcal{O}_X)$ ; voir
Catégories dérivées, lemme \ref{derived-lemma-adjoint-preserve-K-injectives}.
Il en résulte que
$$
\Phi(j_*\mathcal{I}^\bullet) =
f_*j_*\mathcal{I}^\bullet =
g_*\mathcal{I}^\bullet
$$
D'après le lemme \ref{lemma-affine-pushforward},
$j_*\mathcal{I}^\bullet$ représente $Rj_*\mathcal{I}^\bullet$ et
$g_*\mathcal{I}^\bullet$ représente $Rg_*\mathcal{I}^\bullet$.
D'autre part, on a $Rf_* \circ Rj_* = Rg_*$.
Ainsi $f_*j_*\mathcal{I}^\bullet$ représente $Rf_*(j_*\mathcal{I}^\bullet)$.
On en déduit que le lemme est vrai pour tout complexe
de la forme $j_*\mathcal{G}^\bullet$, où $\mathcal{G}^\bullet$
est un complexe de modules quasi-cohérents sur $U$. (Notons que, si
$\mathcal{G}^\bullet \to \mathcal{I}^\bullet$ est un quasi-isomorphisme,
alors $j_*\mathcal{G}^\bullet \to j_*\mathcal{I}^\bullet$ est aussi un
quasi-isomorphisme, puisque $j_*$ est un foncteur exact
sur les modules quasi-cohérents.)

\medskip\noindent
Soit $\mathcal{F}^\bullet$ un complexe de
$\mathcal{O}_X$-modules quasi-cohérents. Soit $T \subset X$ une partie fermée
telle que le support de $\mathcal{F}^p$ soit contenu dans $T$
pour tout $p$. Nous raisonnons par récurrence sur le nombre minimal $n$
d'ouverts affines $U_1, \ldots, U_n$ tels que
$T \subset U_1 \cup \ldots \cup U_n$. Le cas initial $n = 0$ est immédiat.
Si $n \geq 1$, posons $U = U_1$ et notons $j : U \to X$
l'immersion ouverte, comme ci-dessus. Considérons le morphisme de complexes
$c : \mathcal{F}^\bullet \to j_*j^*\mathcal{F}^\bullet$.
On obtient deux suites exactes courtes de complexes :
$$
0 \to \Ker(c) \to \mathcal{F}^\bullet \to \Im(c) \to 0
$$
et
$$
0 \to \Im(c) \to j_*j^*\mathcal{F}^\bullet \to \Coker(c) \to 0
$$
Les complexes $\Ker(c)$ et $\Coker(c)$ sont à support
dans $T \setminus U \subset U_2 \cup \ldots \cup U_n$, et le résultat
est vrai pour eux par récurrence. Il est vrai pour
$j_*j^*\mathcal{F}^\bullet$ d'après la discussion du
paragraphe précédent. On conclut en considérant les triangles distingués
associés aux suites exactes courtes et en utilisant l'observation
faite au début de la démonstration.
\end{proof}

\begin{remark}[Avertissement]
\label{remark-warning-coherator}
Soit $X$ un schéma quasi-compact à diagonale affine. Bien que l'on sache
que $D(\QCoh(\mathcal{O}_X)) = D_\QCoh(\mathcal{O}_X)$ d'après la
proposition \ref{proposition-quasi-compact-affine-diagonal},
des phénomènes inattendus peuvent
se produire, et il est facile de se tromper dans ces questions. On se gardera notamment
de déduire sans précaution de cette égalité que les foncteurs dérivés coïncident.
Par exemple, soit $U \subset X$ un ouvert quasi-compact. On peut alors
considérer les foncteurs dérivés à droite supérieurs
$$
R^i(\QCoh)\Gamma(U, -) : \QCoh(\mathcal{O}_X) \to \textit{Ab}
$$
du foncteur exact à gauche $\Gamma(U, -)$. Puisqu'ils forment un
$\delta$-foncteur universel, et puisque les foncteurs $H^i(U, -)$ (définis pour tous les
faisceaux abéliens sur $X$), restreints à $\QCoh(\mathcal{O}_X)$, forment
un $\delta$-foncteur, on obtient des transformations canoniques
$$
t^i : R^i(\QCoh)\Gamma(U, -) \to H^i(U, -).
$$
Ces transformations ne sont pas en général des isomorphismes, même si $X = \Spec(A)$
est affine ! En effet, on a $R^1(\QCoh)\Gamma(U, \widetilde{I}) = 0$
si $I$ est un $A$-module injectif, par construction des foncteurs dérivés à droite
et par l'équivalence entre $\QCoh(\mathcal{O}_X)$ et $\text{Mod}_A$.
Mais Exemples, lemme \ref{examples-lemma-nonvanishing},
montre qu'il existe $A$, $I$ et $U$ tels que
$H^1(U, \widetilde{I}) \not = 0$.
\end{remark}




\section{Le cohérateur pour les schémas noethériens}
\label{section-coherator-Noetherian}

\noindent
Dans le cas des schémas noethériens, on peut utiliser le lemme suivant.

\begin{lemma}
\label{lemma-injective-quasi-coherent-sheaf-Noetherian}
Soit $X$ un schéma noethérien. Soit $\mathcal{J}$ un objet injectif
de $\QCoh(\mathcal{O}_X)$. Alors $\mathcal{J}$
est un faisceau flasque de $\mathcal{O}_X$-modules.
\end{lemma}

\begin{proof}
Soient $U \subset X$ un ouvert et $s \in \mathcal{J}(U)$ une section.
Soit $\mathcal{I} \subset \mathcal{O}_X$ le faisceau quasi-cohérent
d'idéaux définissant la structure de schéma réduit induite sur $X \setminus U$
(voir Schémas, définition \ref{schemes-definition-reduced-induced-scheme}).
D'après Cohomologie des schémas, lemme \ref{coherent-lemma-homs-over-open},
la section $s$ correspond à un morphisme $\sigma : \mathcal{I}^n \to \mathcal{J}$
pour un certain $n$. Puisque $\mathcal{J}$ est un objet injectif de
$\QCoh(\mathcal{O}_X)$, on peut prolonger $\sigma$ en un morphisme
$\tilde s : \mathcal{O}_X \to \mathcal{J}$. Alors $\tilde s$ correspond
à une section globale de $\mathcal{J}$ dont la restriction est $s$.
\end{proof}

\begin{lemma}
\label{lemma-Noetherian-pushforward}
Soit $f : X \to Y$ un morphisme de schémas noethériens.
Alors le foncteur $f_*$ sur les faisceaux quasi-cohérents admet un foncteur
dérivé à droite
$\Phi : D(\QCoh(\mathcal{O}_X)) \to D(\QCoh(\mathcal{O}_Y))$
tel que le diagramme
$$
\xymatrix{
D(\QCoh(\mathcal{O}_X)) \ar[d]_{\Phi} \ar[r] &
D_\QCoh(\mathcal{O}_X) \ar[d]^{Rf_*} \\
D(\QCoh(\mathcal{O}_Y)) \ar[r] &
D_\QCoh(\mathcal{O}_Y)
}
$$
soit commutatif.
\end{lemma}

\begin{proof}
Puisque $X$ et $Y$ sont des schémas noethériens, le morphisme est quasi-compact
et quasi-séparé (voir
Propriétés, lemme \ref{properties-lemma-locally-Noetherian-quasi-separated}
et
Schémas, remarque \ref{schemes-remark-quasi-compact-and-quasi-separated}).
Ainsi $f_*$ préserve la quasi-cohérence ; voir
Schémas, lemme \ref{schemes-lemma-push-forward-quasi-coherent}.
Soit maintenant $K$ un objet de $D(\QCoh(\mathcal{O}_X))$.
Puisque $\QCoh(\mathcal{O}_X)$ est une catégorie abélienne de Grothendieck
(Propriétés, proposition \ref{properties-proposition-coherator}), on peut
représenter $K$ par un complexe K-injectif $\mathcal{I}^\bullet$
tel que chaque $\mathcal{I}^n$ soit un objet injectif de
$\QCoh(\mathcal{O}_X)$ ; voir
Injectifs, théorème
\ref{injectives-theorem-K-injective-embedding-grothendieck}.
Le foncteur $\Phi$ est donc défini en posant
$$
\Phi(K) = f_*\mathcal{I}^\bullet
$$
où le membre de droite est considéré comme un objet de
$D(\QCoh(\mathcal{O}_Y))$. Pour achever la démonstration du lemme,
il suffit de montrer que le morphisme canonique
$$
f_*\mathcal{I}^\bullet \longrightarrow Rf_*\mathcal{I}^\bullet
$$
est un isomorphisme dans $D(\mathcal{O}_Y)$. Pour cela, d'après
le lemme \ref{lemma-acyclicity-lemma},
il suffit de montrer que $\mathcal{I}^n$ est acyclique à droite pour
$f_*$, pour tout $n \in \mathbf{Z}$.
Or $\mathcal{I}^n$ est flasque d'après
le lemme \ref{lemma-injective-quasi-coherent-sheaf-Noetherian},
et les modules flasques sont acycliques à droite pour $f_*$ d'après
Cohomologie, lemme \ref{cohomology-lemma-flasque-acyclic-pushforward}.
\end{proof}

\begin{proposition}
\label{proposition-Noetherian}
Soit $X$ un schéma noethérien. Alors le foncteur (\ref{equation-compare})
$$
D(\QCoh(\mathcal{O}_X))
\longrightarrow
D_\QCoh(\mathcal{O}_X)
$$
est une équivalence dont un quasi-inverse est donné par $RQ_X$.
\end{proposition}

\begin{proof}
Cela résulte du lemme \ref{lemma-argument-proves} et du
lemme \ref{lemma-Noetherian-pushforward}.
\end{proof}






\section{Complexes de Koszul}
\label{section-koszul}

\noindent
Soit $A$ un anneau et soit $f_1, \ldots, f_r$ une suite d'éléments
de $A$. Nous avons défini le complexe de Koszul
$K_\bullet(f_1, \ldots, f_r)$ dans
Compléments sur l'algèbre, définition \ref{more-algebra-definition-koszul-complex}.
C'est un complexe de chaînes concentré en degrés $r, \ldots, 0$.
Nous en faisons un complexe de cochaînes $K^\bullet(f_1, \ldots, f_r)$
en posant $K^{-n}(f_1, \ldots, f_r) = K_n(f_1, \ldots, f_r)$
et en conservant les mêmes différentielles. Dans la suite de cette section, tous
les complexes seront des complexes de cochaînes.

\medskip\noindent
Nous définissons un complexe $I^\bullet(f_1, \ldots, f_r)$
de sorte que l'on ait un triangle distingué
$$
I^\bullet(f_1, \ldots, f_r) \to
A \to
K^\bullet(f_1, \ldots, f_r) \to
I^\bullet(f_1, \ldots, f_r)[1]
$$
dans $K(A)$.
Autrement dit, on pose
$$
I^i(f_1, \ldots, f_r) =
\left\{
\begin{matrix}
K^{i - 1}(f_1, \ldots, f_r) & \text{si } i \leq 0 \\
0 & \text{sinon}
\end{matrix}
\right.
$$
et l'on prend l'opposée de la différentielle de $K^\bullet(f_1, \ldots, f_r)$.
Les morphismes du triangle distingué sont les morphismes évidents. Notons que
$I^0(f_1, \ldots, f_r) = A^{\oplus r} \to A$ est donné par
la multiplication par $f_i$ sur le $i$-ième facteur.
Ainsi $I^\bullet(f_1, \ldots, f_r) \to A$ se factorise en
$$
I^\bullet(f_1, \ldots, f_r) \to I \to A
$$
où $I = (f_1, \ldots, f_r)$. On a en fait une suite exacte courte
$$
0 \to H^{-1}(K^\bullet(f_1, \ldots, f_s)) \to
H^0(I^\bullet(f_1, \ldots, f_s)) \to I \to 0
$$
et, pour tout $i < 0$, on a
$H^i(I^\bullet(f_1, \ldots, f_r)) = H^{i - 1}(K^\bullet(f_1, \ldots, f_r))$.
Notons que, pour une seconde suite $g_1, \ldots, g_r$ d'éléments de $A$,
on dispose de morphismes canoniques
$$
I^\bullet(f_1g_1, \ldots, f_rg_r) \to I^\bullet(f_1, \ldots, f_r)
\quad\text{et}\quad
K^\bullet(f_1g_1, \ldots, f_rg_r) \to K^\bullet(f_1, \ldots, f_r)
$$
compatibles aux morphismes décrits ci-dessus. Le premier de ces morphismes est
donné par la multiplication par $g_i$ sur le $i$-ième facteur de
$I^0(f_1g_1, \ldots, f_rg_r) = A^{\oplus r}$. En particulier, étant donnés
$f_1, \ldots, f_r$, on obtient un système projectif de complexes
\begin{equation}
\label{equation-system}
I^\bullet(f_1, \ldots, f_r) \leftarrow
I^\bullet(f_1^2, \ldots, f_r^2) \leftarrow
I^\bullet(f_1^3, \ldots, f_r^3) \leftarrow \ldots
\end{equation}
qui jouera un rôle important dans la suite.
Fixons quelques notations pour énoncer commodément les lemmes suivants.

\begin{situation}
\label{situation-complex}
Ici $A$ est un anneau et $f_1, \ldots, f_r$ une suite d'éléments de $A$.
On pose $X = \Spec(A)$ et $U = D(f_1) \cup \ldots \cup D(f_r) \subset X$.
On note $\mathcal{U} : U = \bigcup_{i = 1, \ldots, r} D(f_i)$ le
recouvrement ouvert ainsi défini de $U$.
\end{situation}

\noindent
Notre premier lemme montre que les complexes ci-dessus permettent de calculer
la cohomologie des faisceaux quasi-cohérents sur $U$.

\begin{lemma}
\label{lemma-alternating-cech-complex}
Dans la situation \ref{situation-complex}, soit $M$ un $A$-module, et
notons $\mathcal{F}$ le $\mathcal{O}_X$-module associé. Alors
il existe un isomorphisme canonique de complexes
$$
\Psi : \colim_e \Hom_A(I^\bullet(f_1^e, \ldots, f_r^e), M)
\longrightarrow
\check{\mathcal{C}}_{alt}^\bullet(\mathcal{U}, \mathcal{F})
$$
fonctoriel en $M$, où les différentielles du complexe $\Hom$
sont les transposées des différentielles de
$I^\bullet(f_1^e, \ldots, f_r^e)$.
\end{lemma}

\begin{proof}
Rappelons que le complexe de {\v C}ech alterné est le sous-complexe
du complexe de {\v C}ech usuel formé des cochaînes alternées ; voir
Cohomologie, section \ref{cohomology-section-alternating-cech}.
Comme d'habitude, on considère une $p$-cochaîne de
$\check{\mathcal{C}}_{alt}^\bullet(\mathcal{U}, \mathcal{F})$
comme une fonction alternée $s$ sur $\{1, \ldots, r\}^{p + 1}$
dont la valeur $s_{i_0\ldots i_p}$ en $(i_0, \ldots, i_p)$ appartient à
$M_{f_{i_0}\ldots f_{i_p}} = \mathcal{F}(U_{i_0\ldots i_p})$.
D'autre part, une $p$-cochaîne $t$ de
$\Hom^\bullet(I^\bullet(f_1^e, \ldots, f_r^e), M)$
est une application $t : \wedge^{p + 1}(A^{\oplus r}) \to M$.
Notons $[i] \in A^{\oplus r}$ le $i$-ième vecteur de base et
posons
$$
[i_0, \ldots, i_p] = [i_0] \wedge \ldots \wedge [i_p]
\in \wedge^{p + 1}(A^{\oplus r})
$$
Pour $t$ comme ci-dessus, on pose
$$
\Psi(t)_{i_0 \ldots i_p} =
(-1)^p \frac{t([i_0, \ldots, i_p])}{f_{i_0}^e\ldots f_{i_p}^e}
$$
Il est clair que $\Psi(t)$ est une cochaîne alternée.
Cette règle est compatible aux morphismes de transition
du système, car le morphisme de transition
$$
I^\bullet(f_1^e, \ldots, f_r^e) \leftarrow
I^\bullet(f_1^{e + 1}, \ldots, f_r^{e + 1}),
$$
de (\ref{equation-system}) envoie $[i_0, \ldots, i_p]$
sur $f_{i_0}\ldots f_{i_p}[i_0, \ldots, i_p]$.
La description des localisés
$M_{f_{i_0} \ldots f_{i_p}}$ donnée dans
Algèbre, lemme \ref{algebra-lemma-localization-colimit},
montre que la règle $\Psi$ définit, après passage à la limite inductive, un isomorphisme des modules de cochaînes
en degré $p$. Pour achever la démonstration, il faut montrer que ce morphisme
est compatible aux différentielles. Pour $t$ comme ci-dessus, on calcule
\begin{align*}
d(\Psi(t))_{i_0 \ldots i_{p + 1}}
& =
\sum\nolimits_{j = 0}^{p + 1} (-1)^j
\Psi(t)_{i_0\ldots \hat i_j \ldots i_{p + 1}} \\
& =
(-1)^p
\sum\nolimits_{j = 0}^{p + 1} (-1)^j t([i_0 \ldots \hat i_j \ldots i_{p + 1}])
(f_{i_0} \ldots \hat f_{i_j} \ldots f_{i_p})^{-e}
\end{align*}
Rappelons que les différentielles de $I^\bullet(f_1^e, \ldots, f_r^e)$
sont les opposées de celles de $K^\bullet(f_1, \ldots, f_r)$.
Ainsi
\begin{align*}
\Psi(d(t))_{i_0 \ldots i_{p + 1}} & =
(-1)^{p + 1}
d(t)([i_0, \ldots, i_{p + 1}]) (f_{i_0} \ldots f_{i_{p + 1}})^{-e} \\
& =
(-1)^{p + 1}
t(d([i_0, \ldots, i_{p + 1}])) (f_{i_0} \ldots f_{i_{p + 1}})^{-e} \\
& =
(-1)^{p + 1}
t(-\sum\nolimits_{j = 0}^{p + 1}
(-1)^j f_{i_j}^e [i_0, \ldots, \hat i_j, \ldots i_{p + 1}])
(f_{i_0} \ldots f_{i_{p + 1}})^{-e} \\
& =
-(-1)^{p + 1}
\sum\nolimits_{j = 0}^{p + 1}
(-1)^j t([i_0, \ldots, \hat i_j, \ldots i_{p + 1}])
(f_{i_0} \ldots \hat f_{i_j} \ldots f_{i_p})^{-e}
\end{align*}
Les deux formules coïncident, ce qui achève la démonstration.
\end{proof}

\noindent
Soient un complexe borné $I^\bullet$ de $A$-modules et un
complexe de $A$-modules $M^\bullet$. On dispose du
complexe Hom correspondant $\Hom^\bullet(I^\bullet, M^\bullet)$. Son terme
en degré $n$ est
$$
\bigoplus\nolimits_{p + q = n} \Hom_A(I^{-q}, M^p)
$$
et sa différentielle est décrite dans Compléments sur l'algèbre, section
\ref{more-algebra-section-hom-complexes}.
Puisque le complexe $I^\bullet$ n'a qu'un nombre fini de termes non nuls, la
somme directe ci-dessus est finie.
Les conventions pour le complexe total associé au
complexe de {\v C}ech d'un complexe sont celles de
Cohomologie, section \ref{cohomology-section-cech-cohomology-of-complexes}.

\begin{lemma}
\label{lemma-alternating-cech-complex-complex}
Dans la situation \ref{situation-complex}, soit $M^\bullet$ un
complexe de $A$-modules, et
notons $\mathcal{F}^\bullet$ le complexe associé de
$\mathcal{O}_X$-modules. Alors
il existe un isomorphisme canonique de complexes
$$
\colim_e \Hom^\bullet(I^\bullet(f_1^e, \ldots, f_r^e), M^\bullet)
\longrightarrow
\text{Tot}(\check{\mathcal{C}}_{alt}^\bullet(\mathcal{U}, \mathcal{F}^\bullet))
$$
fonctoriel en $M^\bullet$.
\end{lemma}

\begin{proof}
Considérons le complexe double $F^{\bullet, \bullet}$ de termes
$F^{p, q} = \mathcal{C}_{alt}^p(\mathcal{U}, \mathcal{F}^q)$
étudié dans
Cohomologie, section \ref{cohomology-section-cech-cohomology-of-complexes}.
Considérons le complexe double $G^{\bullet, \bullet}$ de termes
$G^{p, q} = \colim_e \Hom_A(I^{-p}(f_1^e, \ldots, f_r^e), M^q)$ et
dont les différentielles sont données par fonctorialité (sans intervention
de signes). Les morphismes $\psi^{p, q} : G^{p, q} \to F^{p, q}$
construits dans la démonstration du lemme \ref{lemma-alternating-cech-complex}
sont des isomorphismes compatibles aux différentielles $d_1$ (d'après le lemme)
et $d_2$ (ce qui est clair). Cependant, les différentielles $d$ des
complexes figurant à gauche et à droite de la flèche de l'énoncé
ont des signes différents. Pour $g \in G^{p, q}$, on a en effet
$$
d(g) =  d_2(g) - (-1)^{p + q} d_1(g) 
$$
(voir Compléments sur l'algèbre, section \ref{more-algebra-section-hom-complexes}),
tandis que la différentielle, pour $f \in F^{p, q}$, est donnée par
$$
d(f) = d_1(f) + (-1)^p d_2(f)
$$
On peut donc corriger les signes en multipliant $\psi^{p, q}$ par
$(-1)^{pq + p(p - 1)/2}$.
\end{proof}

\begin{lemma}
\label{lemma-alternating-cech-complex-complex-computes-cohomology}
Dans la situation \ref{situation-complex}, soit $\mathcal{F}^\bullet$
un complexe de $\mathcal{O}_X$-modules quasi-cohérents. Alors
il existe un isomorphisme canonique
$$
\text{Tot}(\check{\mathcal{C}}_{alt}^\bullet(\mathcal{U}, \mathcal{F}^\bullet))
\longrightarrow
R\Gamma(U, \mathcal{F}^\bullet)
$$
dans $D(A)$, fonctoriel en $\mathcal{F}^\bullet$.
\end{lemma}

\begin{proof}
Soit $\mathcal{B}$ l'ensemble des ouverts affines de $U$. Puisque les groupes de
cohomologie supérieure d'un module quasi-cohérent sur un schéma affine sont nuls
(Cohomologie des schémas, lemme
\ref{coherent-lemma-quasi-coherent-affine-cohomology-zero}),
il s'agit d'un cas particulier de
Cohomologie, lemme \ref{cohomology-lemma-alternating-cech-complex-complex-ss}.
\end{proof}

\noindent
Dans la situation \ref{situation-complex}, notons $I_e$ l'objet de
$D(\mathcal{O}_X)$ correspondant au complexe de $A$-modules
$I^\bullet(f_1^e, \ldots, f_r^e)$ par l'équivalence du
lemme \ref{lemma-affine-compare-bounded}. Les morphismes
(\ref{equation-system}) donnent un système
$$
I_1 \leftarrow
I_2 \leftarrow
I_3 \leftarrow \ldots
$$
De plus, on dispose d'un système compatible de morphismes $I_e \to \mathcal{O}_X$
qui deviennent des isomorphismes après restriction à $U$. Ainsi, pour
tout objet $E$ de $D(\mathcal{O}_X)$, il existe un morphisme canonique
\begin{equation}
\label{equation-comparison}
\colim_e \Hom_{D(\mathcal{O}_X)}(I_e, E) \longrightarrow H^0(U, E)
\end{equation}
obtenu en envoyant un morphisme $I_e \to E$ sur sa restriction à $U$
et en utilisant l'égalité
$\Hom_{D(\mathcal{O}_U)}(\mathcal{O}_U, E|_U) = H^0(U, E)$.

\begin{proposition}
\label{proposition-represent-cohomology-class-on-open}
Dans la situation \ref{situation-complex}, pour tout objet $E$
de $D_\QCoh(\mathcal{O}_X)$, le morphisme
(\ref{equation-comparison}) est un isomorphisme.
\end{proposition}

\begin{proof}
D'après le lemme \ref{lemma-affine-compare-bounded}, on peut supposer que $E$
est donné par un complexe de faisceaux quasi-cohérents $\mathcal{F}^\bullet$.
Soit $M^\bullet = \Gamma(X, \mathcal{F}^\bullet)$ le complexe correspondant
de $A$-modules. D'après
les lemmes \ref{lemma-alternating-cech-complex-complex} et
\ref{lemma-alternating-cech-complex-complex-computes-cohomology},
on a des quasi-isomorphismes
$$
\colim_e \Hom^\bullet(I^\bullet(f_1^e, \ldots, f_r^e), M^\bullet)
\longrightarrow
\text{Tot}(\check{\mathcal{C}}_{alt}^\bullet(\mathcal{U}, \mathcal{F}^\bullet))
\longrightarrow
R\Gamma(U, \mathcal{F}^\bullet)
$$
D'après Compléments sur l'algèbre, lemme \ref{more-algebra-lemma-RHom-out-of-projective} et
équation (\ref{more-algebra-equation-h0-RHom}),
le $H^0$ du complexe
$\Hom^\bullet(I^\bullet(f_1^e, \ldots, f_r^e), M^\bullet)$
calcule le $\Hom$ dans $D(A)$. En prenant $H^0$ des deux côtés, on obtient donc
$$
\colim_e \Hom_{D(A)}(I^\bullet(f_1^e, \ldots, f_r^e), M^\bullet)
=
H^0(U, E)
$$
Puisque $\Hom_{D(A)}(I^\bullet(f_1^e, \ldots, f_r^e), M^\bullet) =
\Hom_{D(\mathcal{O}_X)}(I_e, E)$ d'après
le lemme \ref{lemma-affine-compare-bounded}, le résultat en découle.
\end{proof}

\noindent
Dans la situation \ref{situation-complex}, notons $K_e$ l'objet de
$D(\mathcal{O}_X)$ correspondant au complexe de $A$-modules
$K^\bullet(f_1^e, \ldots, f_r^e)$ par l'équivalence du
lemme \ref{lemma-affine-compare-bounded}. On a ainsi des triangles
distingués
$$
I_e \to \mathcal{O}_X \to K_e \to I_e[1]
$$
et un système
$$
K_1 \leftarrow
K_2 \leftarrow
K_3 \leftarrow \ldots
$$
compatible au système $(I_e)$.
De plus, on dispose d'un système compatible de morphismes
$$
K_e \to H^0(K_e) = \mathcal{O}_X/(f_1^e, \ldots, f_r^e)
$$

\begin{lemma}
\label{lemma-represent-cohomology-class-on-closed}
Dans la situation \ref{situation-complex}, soit $E$ un objet de
$D_\QCoh(\mathcal{O}_X)$.
Supposons que $H^i(E)|_U = 0$ pour $i = - r + 1, \ldots, 0$.
Alors, étant donné $s \in H^0(X, E)$, il existe un $e \geq 0$ et
un morphisme $K_e \to E$ tels que $s$ appartienne à l'image de
$H^0(X, K_e) \to H^0(X, E)$.
\end{lemma}

\begin{proof}
Puisque $U$ est recouvert par $r$ ouverts affines, on a $H^j(U, \mathcal{F}) = 0$
pour $j \geq r$ et tout module quasi-cohérent
(Cohomologie des schémas, lemme \ref{coherent-lemma-vanishing-nr-affines}).
D'après le lemme \ref{lemma-application-nice-K-injective}, $H^0(U, E)$
est égal à $H^0(U, \tau_{\geq -r + 1}E)$. On dispose d'une
suite spectrale
$$
H^j(U, H^i(\tau_{\geq -r + 1}E)) \Rightarrow H^{i + j}(U, \tau_{\geq -N}E)
$$
voir Catégories dérivées, lemme \ref{derived-lemma-two-ss-complex-functor}.
Ainsi $H^0(U, E) = 0$ par l'hypothèse d'annulation des faisceaux de cohomologie de $E$.
On en déduit que $s|_U = 0$.
Considérons $s$ comme un morphisme $\mathcal{O}_X \to E$ dans $D(\mathcal{O}_X)$.
D'après la proposition \ref{proposition-represent-cohomology-class-on-open},
le composé $I_e \to \mathcal{O}_X \to E$ est nul pour un certain $e$.
Le triangle distingué $I_e \to \mathcal{O}_X \to K_e \to I_e[1]$
fournit alors un morphisme $K_e \to E$ tel que $s$ soit le composé
$\mathcal{O}_X \to K_e \to E$.
\end{proof}

\section{Complexes pseudo-cohérents et parfaits}
\label{section-spell-out}

\noindent
Dans cette section, nous établissons le lien entre les notions
générales définies dans
Cohomologie, sections \ref{cohomology-section-strictly-perfect},
\ref{cohomology-section-pseudo-coherent},
\ref{cohomology-section-tor} et
\ref{cohomology-section-perfect},
et les notions correspondantes pour les complexes de modules, dans
Compléments sur l'algèbre, sections
\ref{more-algebra-section-pseudo-coherent},
\ref{more-algebra-section-tor} et
\ref{more-algebra-section-perfect}.

\begin{lemma}
\label{lemma-pseudo-coherent}
Soit $X$ un schéma. Si $E$ est un objet $m$-pseudo-cohérent
de $D(\mathcal{O}_X)$, alors $H^i(E)$ est un
$\mathcal{O}_X$-module quasi-cohérent pour $i > m$, et $H^m(E)$ est un quotient
d'un $\mathcal{O}_X$-module quasi-cohérent.
Si $E$ est pseudo-cohérent, alors $E$ est un objet de
$D_\QCoh(\mathcal{O}_X)$.
\end{lemma}

\begin{proof}
Localement sur $X$, il existe un complexe strictement parfait $\mathcal{E}^\bullet$
tel que $H^i(E)$ soit isomorphe à $H^i(\mathcal{E}^\bullet)$ pour $i > m$
et que $H^m(E)$ soit un quotient de $H^m(\mathcal{E}^\bullet)$. Les faisceaux
$\mathcal{E}^i$ sont des facteurs directs de modules libres de rang fini,
donc sont quasi-cohérents. Le lemme en résulte.
\end{proof}

\begin{lemma}
\label{lemma-pseudo-coherent-affine}
Soit $X = \Spec(A)$ un schéma affine. Soit $M^\bullet$ un
complexe de $A$-modules et soit $E$ l'objet correspondant
de $D(\mathcal{O}_X)$. Alors $E$ est $m$-pseudo-cohérent
(resp.\ pseudo-cohérent) en tant qu'objet de $D(\mathcal{O}_X)$
si et seulement si $M^\bullet$ est $m$-pseudo-cohérent (resp.\ pseudo-cohérent)
en tant que complexe de $A$-modules.
\end{lemma}

\begin{proof}
Il résulte immédiatement des définitions que, si $M^\bullet$ est
$m$-pseudo-cohérent, il en est de même de $E$. Pour établir la réciproque, supposons
$E$ $m$-pseudo-cohérent. Puisque $X = \Spec(A)$ est quasi-compact et que
les ouverts principaux forment une base de sa topologie, on peut trouver un
recouvrement par des ouverts principaux $X = D(f_1) \cup \ldots \cup D(f_n)$, des complexes strictement
parfaits $\mathcal{E}_i^\bullet$ sur $D(f_i)$ et des
morphismes $\alpha_i : \mathcal{E}_i^\bullet \to E|_{U_i}$ induisant
des isomorphismes sur $H^j$ pour $j > m$ et des surjections sur $H^m$.
D'après Cohomologie, lemme \ref{cohomology-lemma-local-actual},
quitte à raffiner le recouvrement ouvert,
on peut supposer que $\alpha_i$ est donné par un morphisme de complexes
$\mathcal{E}_i^\bullet \to \widetilde{M^\bullet}|_{U_i}$
pour tout $i$. D'après Modules, lemme
\ref{modules-lemma-direct-summand-of-locally-free-is-locally-free},
les termes $\mathcal{E}_i^n$ sont des modules localement libres de rang fini.
Ainsi, quitte à raffiner le recouvrement ouvert, on peut supposer chaque
$\mathcal{E}_i^n$ libre de rang fini sur $\mathcal{O}_{U_i}$.
Il résulte de la définition que $M^\bullet_{f_i}$ est
un complexe $m$-pseudo-cohérent de $A_{f_i}$-modules.
On conclut en appliquant
Compléments sur l'algèbre, lemme \ref{more-algebra-lemma-glue-pseudo-coherent}.

\medskip\noindent
Le cas « pseudo-cohérent » résulte du fait que $E$ est
pseudo-cohérent si et seulement si $E$ est $m$-pseudo-cohérent pour
tout $m$ (par définition), et qu'il en est de même de $M^\bullet$
d'après Compléments sur l'algèbre, lemme \ref{more-algebra-lemma-pseudo-coherent}.
\end{proof}

\begin{lemma}
\label{lemma-identify-pseudo-coherent-noetherian}
Soit $X$ un schéma noethérien. Soit $E$ un objet de
$D_\QCoh(\mathcal{O}_X)$. Pour $m \in \mathbf{Z}$, les
conditions suivantes sont équivalentes :
\begin{enumerate}
\item $H^i(E)$ est cohérent pour $i \geq m$ et nul pour $i \gg 0$,
\item $E$ est $m$-pseudo-cohérent.
\end{enumerate}
En particulier, $E$ est pseudo-cohérent si et seulement si $E$ est un objet
de $D^-_{\textit{Coh}}(\mathcal{O}_X)$.
\end{lemma}

\begin{proof}
Puisque $X$ est quasi-compact, dans chacun des cas (1) et (2), l'objet $E$
est borné supérieurement. La question est donc locale sur $X$, et l'on peut supposer
$X$ affine. Écrivons $X = \Spec(A)$ pour un anneau noethérien $A$.
Dans ce cas, $E$ correspond à un complexe de $A$-modules $M^\bullet$
d'après le lemme \ref{lemma-affine-compare-bounded}. D'après
le lemme \ref{lemma-pseudo-coherent-affine},
$E$ est $m$-pseudo-cohérent si et seulement si $M^\bullet$
est $m$-pseudo-cohérent. D'autre part, $H^i(E)$ est cohérent
si et seulement si $H^i(M^\bullet)$ est un $A$-module de type fini
(Propriétés, lemme \ref{properties-lemma-finite-type-module}).
Le résultat découle donc de Compléments sur l'algèbre, lemme
\ref{more-algebra-lemma-Noetherian-pseudo-coherent}.
\end{proof}

\begin{lemma}
\label{lemma-tor-dimension-affine}
Soit $X = \Spec(A)$ un schéma affine. Soit $M^\bullet$ un
complexe de $A$-modules et soit $E$ l'objet correspondant
de $D(\mathcal{O}_X)$. Alors
\begin{enumerate}
\item $E$ a une amplitude de Tor contenue dans $[a, b]$ si et seulement si $M^\bullet$
a une amplitude de Tor contenue dans $[a, b]$.
\item $E$ est de dimension de Tor finie si et seulement si $M^\bullet$
est de dimension de Tor finie.
\end{enumerate}
\end{lemma}

\begin{proof}
L'assertion (2) résulte immédiatement de (1). Pour démontrer (1), nous
utiliserons l'équivalence $D(A) = D_\QCoh(X)$ du
lemme \ref{lemma-affine-compare-bounded}
sans autre mention.
Supposons que $M^\bullet$ ait une amplitude de Tor contenue dans $[a, b]$. Alors $K^\bullet$
est isomorphe dans $D(A)$ à un complexe $K^\bullet$ de $A$-modules plats
tel que $K^i = 0$ pour $i \not \in [a, b]$ ; voir
Compléments sur l'algèbre, lemme \ref{more-algebra-lemma-tor-amplitude}.
Alors $E$ est isomorphe à $\widetilde{K^\bullet}$. Puisque chaque
$\widetilde{K^i}$ est un $\mathcal{O}_X$-module plat,
$E$ a une amplitude de Tor contenue dans $[a, b]$ d'après
Cohomologie, lemme \ref{cohomology-lemma-tor-amplitude}.

\medskip\noindent
Supposons que $E$ ait une amplitude de Tor contenue dans $[a, b]$. Alors $E$ est borné,
donc $M^\bullet$ appartient à $K^-(A)$. On peut ainsi remplacer $M^\bullet$
par un complexe borné supérieurement de $A$-modules. On peut même choisir
une résolution projective et supposer que $M^\bullet$ est un complexe borné supérieurement
de $A$-modules libres. Pour tout $A$-module $N$, on a alors
$$
E \otimes_{\mathcal{O}_X}^\mathbf{L} \widetilde{N}
\cong
\widetilde{M^\bullet} \otimes_{\mathcal{O}_X}^\mathbf{L} \widetilde{N}
\cong
\widetilde{M^\bullet \otimes_A N}
$$
dans $D(\mathcal{O}_X)$. L'annulation des faisceaux de cohomologie du
membre de gauche entraîne donc que $M^\bullet$ a une amplitude de Tor contenue dans $[a, b]$.
\end{proof}

\begin{lemma}
\label{lemma-tor-dimension-rel-affine}
Soit $f : X \to S$ un morphisme de schémas affines correspondant
au morphisme d'anneaux $R \to A$. Soit $M^\bullet$ un
complexe de $A$-modules et soit $E$ l'objet correspondant
de $D(\mathcal{O}_X)$. Alors
\begin{enumerate}
\item $E$, en tant qu'objet de $D(f^{-1}\mathcal{O}_S)$, a une amplitude de Tor contenue dans
$[a, b]$ si et seulement si $M^\bullet$ a une amplitude de Tor contenue dans $[a, b]$
en tant qu'objet de $D(R)$.
\item $E$ est localement de dimension de Tor finie en tant qu'objet de
$D(f^{-1}\mathcal{O}_S)$ si et seulement si $M^\bullet$
est de dimension de Tor finie en tant qu'objet de $D(R)$.
\end{enumerate}
\end{lemma}

\begin{proof}
Considérons un idéal premier $\mathfrak q \subset A$ au-dessus de $\mathfrak p \subset R$.
Soient $x \in X$ et $s = f(x) \in S$ les points correspondants.
Alors $(f^{-1}\mathcal{O}_S)_x = \mathcal{O}_{S, s} = R_\mathfrak p$
et $E_x = M^\bullet_\mathfrak q$. Compte tenu de ces identifications, on établit
l'équivalence comme suit.

\medskip\noindent
Si $M^\bullet$ a une amplitude de Tor contenue dans $[a, b]$ en tant que complexe de $R$-modules,
il en est de même du localisé de $M^\bullet$ en tout idéal premier de $A$.
On déduit alors de
Cohomologie, lemme \ref{cohomology-lemma-tor-amplitude-stalk},
que $E$ a une amplitude de Tor contenue dans $[a, b]$ en tant que complexe de faisceaux
de $f^{-1}\mathcal{O}_S$-modules.
Réciproquement, supposons que $E$ ait une amplitude de Tor contenue dans $[a, b]$
en tant qu'objet de $D(f^{-1}\mathcal{O}_S)$.
On en déduit, à l'aide du lemme qui vient d'être cité, que
$M^\bullet_\mathfrak q$ a une amplitude de Tor contenue dans $[a, b]$
en tant que complexe de $R_\mathfrak p$-modules pour tout
idéal premier $\mathfrak q \subset A$ au-dessus de $\mathfrak p \subset R$.
D'après Compléments sur l'algèbre, lemme \ref{more-algebra-lemma-tor-amplitude-localization},
il en résulte que $M^\bullet$ a une amplitude de Tor contenue dans $[a, b]$
en tant que complexe de $R$-modules.
Cela achève la démonstration de (1).

\medskip\noindent
Puisque $X$ est quasi-compact, si $E$ est localement de dimension de Tor finie
en tant que complexe de $f^{-1}\mathcal{O}_S$-modules, alors $E$
a en fait une amplitude de Tor contenue dans $[a, b]$ pour certains $a, b$, en tant que complexe de
$f^{-1}\mathcal{O}_S$-modules. Ainsi (2) résulte de (1).
\end{proof}

\begin{lemma}
\label{lemma-tor-qc-qs}
Soit $X$ un schéma quasi-séparé. Soit $E$ un objet
de $D_\QCoh(\mathcal{O}_X)$. Soient $a \leq b$. Les
conditions suivantes sont équivalentes :
\begin{enumerate}
\item $E$ a une amplitude de Tor contenue dans $[a, b]$,
\item pour tout $\mathcal{F}$ dans $\QCoh(\mathcal{O}_X)$,
on a $H^i(E \otimes_{\mathcal{O}_X}^\mathbf{L} \mathcal{F}) = 0$
pour $i \not \in [a, b]$.
\end{enumerate}
\end{lemma}

\begin{proof}
Il est clair que (1) entraîne (2). Supposons (2). Soit $U \subset X$
un ouvert affine. Puisque $X$ est quasi-séparé, le morphisme $j : U \to X$
est quasi-compact et séparé ; ainsi $j_*$ transforme les modules quasi-cohérents
en modules quasi-cohérents
(Schémas, lemme \ref{schemes-lemma-push-forward-quasi-coherent}).
Le foncteur
$\QCoh(\mathcal{O}_X) \to \QCoh(\mathcal{O}_U)$
est donc essentiellement surjectif. Il en résulte que la condition (2)
entraîne l'annulation de
$H^i(E|_U \otimes_{\mathcal{O}_U}^\mathbf{L} \mathcal{G})$
pour $i \not \in [a, b]$, pour tout $\mathcal{O}_U$-module quasi-cohérent
$\mathcal{G}$. Écrivons $U = \Spec(A)$, et soit $M^\bullet$ le
complexe de $A$-modules correspondant à $E|_U$ d'après
le lemme \ref{lemma-affine-compare-bounded}.
Nous venons de montrer que $M^\bullet \otimes_A^\mathbf{L} N$
a des groupes de cohomologie nuls en dehors de l'intervalle $[a, b]$ ;
autrement dit, $M^\bullet$ a une amplitude de Tor contenue dans $[a, b]$.
D'après le lemme \ref{lemma-tor-dimension-affine},
on en déduit que $E|_U$ a une amplitude de Tor contenue dans $[a, b]$.
Cela démontre le lemme.
\end{proof}

\begin{lemma}
\label{lemma-perfect-affine}
Soit $X = \Spec(A)$ un schéma affine. Soit $M^\bullet$ un
complexe de $A$-modules et soit $E$ l'objet correspondant
de $D(\mathcal{O}_X)$. Alors $E$ est un objet parfait de $D(\mathcal{O}_X)$
si et seulement si $M^\bullet$ est parfait en tant qu'objet de $D(A)$.
\end{lemma}

\begin{proof}
C'est une conséquence formelle des
lemmes \ref{lemma-pseudo-coherent-affine} et
\ref{lemma-tor-dimension-affine},
de Cohomologie, lemme \ref{cohomology-lemma-perfect}, et de
Compléments sur l'algèbre, lemme \ref{more-algebra-lemma-perfect}.
\end{proof}

\noindent
Notre description des complexes pseudo-cohérents
sur les schémas permet de montrer que certains Hom internes
sont quasi-cohérents.

\begin{lemma}
\label{lemma-quasi-coherence-internal-hom}
Soit $X$ un schéma.
\begin{enumerate}
\item Si $L$ appartient à $D^+_\QCoh(\mathcal{O}_X)$ et si
$K$ dans $D(\mathcal{O}_X)$ est pseudo-cohérent, alors
$R\SheafHom(K, L)$ appartient à $D_\QCoh(\mathcal{O}_X)$
et est localement borné inférieurement.
\item Si $L$ appartient à $D_\QCoh(\mathcal{O}_X)$ et si
$K$ dans $D(\mathcal{O}_X)$ est parfait, alors
$R\SheafHom(K, L)$ appartient à $D_\QCoh(\mathcal{O}_X)$.
\item Si $X = \Spec(A)$ est affine et si $K, L \in D(A)$, alors
$$
R\SheafHom(\widetilde{K}, \widetilde{L}) = \widetilde{R\Hom_A(K, L)}
$$
dans les deux cas suivants :
\begin{enumerate}
\item $K$ est pseudo-cohérent et $L$ est borné inférieurement,
\item $K$ est parfait et $L$ est quelconque.
\end{enumerate}
\item Si $X = \Spec(A)$ et si $K, L$ appartiennent à $D(A)$, le $n$-ième
faisceau de cohomologie de $R\SheafHom(\widetilde{K}, \widetilde{L})$
est le faisceau associé au préfaisceau
$$
X \supset D(f) \longmapsto \Ext^n_{A_f}(K \otimes_A A_f, L \otimes_A A_f)
$$
pour $f \in A$.
\end{enumerate}
\end{lemma}

\begin{proof}
La construction du Hom interne dans la catégorie dérivée de
$\mathcal{O}_X$ commute à la localisation (voir
Cohomologie, section \ref{cohomology-section-internal-hom}).
Pour démontrer (1) et (2), on peut donc remplacer $X$ par un ouvert affine.
D'après les lemmes \ref{lemma-affine-compare-bounded},
\ref{lemma-pseudo-coherent-affine} et
\ref{lemma-perfect-affine},
il suffit alors de démontrer (3).

\medskip\noindent
L'assertion (3) résulte du calcul du
Hom interne de Cohomologie, lemme
\ref{cohomology-lemma-Rhom-complex-of-direct-summands-finite-free},
en représentant $K$ par un complexe borné supérieurement (resp.\ borné) de
$A$-modules projectifs de type fini et $L$ par un complexe borné inférieurement
(resp.\ quelconque) de $A$-modules.

\medskip\noindent
Pour démontrer (4), rappelons que, sur tout espace annelé, le $n$-ième faisceau de cohomologie de
$R\SheafHom(A, B)$ est le faisceau associé au préfaisceau
$$
U \mapsto \Hom_{D(U)}(A|_U, B|_U[n]) =
\Ext^n_{D(\mathcal{O}_U)}(A|_U, B|_U)
$$
Voir Cohomologie, section \ref{cohomology-section-internal-hom}.
D'autre part, la restriction de $\widetilde{K}$ à un ouvert principal
$D(f)$ est l'image de $K \otimes_A A_f$, et de même pour $L$.
L'assertion (4) résulte donc de l'équivalence de catégories du
lemme \ref{lemma-affine-compare-bounded}.
\end{proof}

\begin{lemma}
\label{lemma-internal-hom-evaluate-tensor-isomorphism}
Soit $X$ un schéma. Soient $K, L, M$ des objets de $D_\QCoh(\mathcal{O}_X)$.
Le morphisme
$$
K \otimes_{\mathcal{O}_X}^\mathbf{L} R\SheafHom(M, L)
\longrightarrow
R\SheafHom(M, K \otimes_{\mathcal{O}_X}^\mathbf{L} L)
$$
de Cohomologie, lemme \ref{cohomology-lemma-internal-hom-diagonal-better},
est un isomorphisme dans les cas suivants :
\begin{enumerate}
\item $M$ est parfait, ou
\item $K$ est parfait, ou
\item $M$ est pseudo-cohérent, $L \in D^+(\mathcal{O}_X)$, et $K$ est de dimension
de Tor finie.
\end{enumerate}
\end{lemma}

\begin{proof}
Le lemme \ref{lemma-quasi-coherence-internal-hom}
ramène les cas (1) et (3) au cas affine, traité dans
Compléments sur l'algèbre, lemme
\ref{more-algebra-lemma-internal-hom-evaluate-tensor-isomorphism}.
(Il faut aussi utiliser les lemmes \ref{lemma-pseudo-coherent-affine},
\ref{lemma-perfect-affine} et \ref{lemma-tor-dimension-affine}
pour passer à l'énoncé algébrique.)
Si $K$ est parfait, sans autre hypothèse, on
ne sait pas si les deux termes de la flèche appartiennent à $D_\QCoh(\mathcal{O}_X)$,
mais le résultat reste vrai, car on peut travailler localement et se ramener
au cas où $K$ est un complexe borné de modules libres de rang fini,
auquel cas il est clair.
\end{proof}





\section{Catégorie dérivée des modules cohérents}
\label{section-derived-coherent}

\noindent
Soit $X$ un schéma localement noethérien. Dans ce cas, la catégorie
$\textit{Coh}(\mathcal{O}_X) \subset \textit{Mod}(\mathcal{O}_X)$
des $\mathcal{O}_X$-modules cohérents est une sous-catégorie de Serre faible ; voir
Homologie, section \ref{homology-section-serre-subcategories}
et
Cohomologie des schémas, lemme \ref{coherent-lemma-coherent-abelian-Noetherian}.
Notons
$$
D_{\textit{Coh}}(\mathcal{O}_X) \subset D(\mathcal{O}_X)
$$
la sous-catégorie des complexes dont les faisceaux de cohomologie sont cohérents ; voir
Catégories dérivées, section \ref{derived-section-triangulated-sub}.
On obtient ainsi un foncteur canonique
\begin{equation}
\label{equation-compare-coherent}
D(\textit{Coh}(\mathcal{O}_X))
\longrightarrow
D_{\textit{Coh}}(\mathcal{O}_X)
\end{equation}
voir Catégories dérivées, équation (\ref{derived-equation-compare}).

\begin{lemma}
\label{lemma-coh-to-qcoh}
Soit $X$ un schéma noethérien. Alors le foncteur
$$
D^-(\textit{Coh}(\mathcal{O}_X))
\longrightarrow
D^-_{\textit{Coh}(\mathcal{O}_X)}(\QCoh(\mathcal{O}_X))
$$
est une équivalence.
\end{lemma}

\begin{proof}
Observons que $\textit{Coh}(\mathcal{O}_X) \subset \QCoh(\mathcal{O}_X)$
est une sous-catégorie de Serre ; voir
Homologie, définition \ref{homology-definition-serre-subcategory} et
lemme \ref{homology-lemma-characterize-serre-subcategory}, ainsi que
Cohomologie des schémas, lemmes
\ref{coherent-lemma-coherent-abelian-Noetherian} et
\ref{coherent-lemma-coherent-Noetherian-quasi-coherent-sub-quotient}.
D'autre part, si $\mathcal{G} \to \mathcal{F}$ est une surjection
d'un $\mathcal{O}_X$-module quasi-cohérent sur un
$\mathcal{O}_X$-module cohérent, il existe un sous-module cohérent
$\mathcal{G}' \subset \mathcal{G}$ qui s'envoie surjectivement sur $\mathcal{F}$.
En effet, on peut écrire $\mathcal{G}$ comme réunion filtrante de ses sous-modules
cohérents d'après
Propriétés, lemme \ref{properties-lemma-quasi-coherent-colimit-finite-type},
et l'un d'entre eux convient alors.
Le lemme résulte donc de
Catégories dérivées, lemme \ref{derived-lemma-fully-faithful-embedding}.
\end{proof}

\begin{proposition}
\label{proposition-DCoh}
Soit $X$ un schéma noethérien. Alors les foncteurs
$$
D^-(\textit{Coh}(\mathcal{O}_X))
\longrightarrow
D^-_{\textit{Coh}}(\mathcal{O}_X)
\quad\text{et}\quad
D^b(\textit{Coh}(\mathcal{O}_X))
\longrightarrow
D^b_{\textit{Coh}}(\mathcal{O}_X)
$$
sont des équivalences.
\end{proposition}

\begin{proof}
Considérons le diagramme commutatif
$$
\xymatrix{
D^-(\textit{Coh}(\mathcal{O}_X)) \ar[r] \ar[d] &
D^-_{\textit{Coh}}(\mathcal{O}_X) \ar[d] \\
D^-(\QCoh(\mathcal{O}_X)) \ar[r] &
D^-_\QCoh(\mathcal{O}_X)
}
$$
D'après le lemme \ref{lemma-coh-to-qcoh}, la flèche verticale de gauche est pleinement fidèle.
D'après la proposition \ref{proposition-Noetherian}, la flèche inférieure est une équivalence.
Par construction, la flèche verticale de droite est pleinement fidèle.
On en déduit que la flèche horizontale supérieure est pleinement fidèle.
Si $K$ est un objet de $D^-_{\textit{Coh}}(\mathcal{O}_X)$,
l'objet $K'$ de $D^-(\QCoh(\mathcal{O}_X))$ qui lui correspond
par la proposition \ref{proposition-Noetherian} a des
faisceaux de cohomologie cohérents. Ainsi $K'$ appartient à l'image essentielle
de la flèche verticale de gauche, d'après le lemme \ref{lemma-coh-to-qcoh},
et la flèche horizontale supérieure est donc essentiellement surjective.
Cela achève la démonstration dans le cas borné supérieurement. Le cas
borné en découle immédiatement.
\end{proof}

\begin{lemma}
\label{lemma-direct-image-coherent}
Soit $S$ un schéma noethérien. Soit $f : X \to S$ un morphisme de schémas
localement de type fini. Soit $E$ un objet de
$D^b_{\textit{Coh}}(\mathcal{O}_X)$ tel que le support de $H^i(E)$
soit propre sur $S$ pour tout $i$.
Alors $Rf_*E$ est un objet de $D^b_{\textit{Coh}}(\mathcal{O}_S)$.
\end{lemma}

\begin{proof}
Considérons la suite spectrale
$$
R^pf_*H^q(E) \Rightarrow R^{p + q}f_*E
$$
voir Catégories dérivées, lemme \ref{derived-lemma-two-ss-complex-functor}.
Par hypothèse et d'après
Cohomologie des schémas, lemme
\ref{coherent-lemma-support-proper-over-base-pushforward},
les faisceaux $R^pf_*H^q(E)$ sont cohérents. Ainsi
$R^{p + q}f_*E$ est cohérent, c'est-à-dire que $Rf_*E \in D_{\textit{Coh}}(\mathcal{O}_S)$.
La bornitude inférieure est immédiate. La bornitude supérieure
résulte de
Cohomologie des schémas, lemme
\ref{coherent-lemma-quasi-coherence-higher-direct-images},
ou du
lemme \ref{lemma-quasi-coherence-direct-image}.
\end{proof}

\begin{lemma}
\label{lemma-direct-image-coherent-bdd-below}
Soit $S$ un schéma noethérien. Soit $f : X \to S$ un morphisme de schémas
localement de type fini. Soit $E$ un objet de
$D^+_{\textit{Coh}}(\mathcal{O}_X)$ tel que le support de $H^i(E)$
soit propre sur $S$ pour tout $i$.
Alors $Rf_*E$ est un objet de $D^+_{\textit{Coh}}(\mathcal{O}_S)$.
\end{lemma}

\begin{proof}
La démonstration est la même que celle du
lemme \ref{lemma-direct-image-coherent}.
On peut aussi le déduire du
lemme \ref{lemma-direct-image-coherent}
en appliquant le foncteur exact $Rf_*$ aux
triangles distingués
$\tau_{\leq a}E \to E \to \tau_{\geq a + 1}E \to \tau_{\leq a}E[1]$.
\end{proof}

\begin{lemma}
\label{lemma-coherent-internal-hom}
Soit $X$ un schéma localement noethérien. Si $L$ appartient à
$D^+_{\textit{Coh}}(\mathcal{O}_X)$ et si $K$ appartient à
$D^-_{\textit{Coh}}(\mathcal{O}_X)$, alors
$R\SheafHom(K, L)$ appartient à $D^+_{\textit{Coh}}(\mathcal{O}_X)$.
\end{lemma}

\begin{proof}
Il suffit de le démontrer lorsque $X$ est le spectre
d'un anneau noethérien $A$.
D'après le lemme \ref{lemma-identify-pseudo-coherent-noetherian},
$K$ est pseudo-cohérent.
On peut alors utiliser le lemme \ref{lemma-quasi-coherence-internal-hom}
pour ramener la question à l'énoncé algébrique suivant :
pour $L \in D^+_{\textit{Coh}}(A)$ et $K$ dans $D^-_{\textit{Coh}}(A)$,
$R\Hom_A(K, L)$ appartient à $D^+_{\textit{Coh}}(A)$.
Puisque $L$ est borné inférieurement et $K$ borné supérieurement, on dispose d'une
suite spectrale convergente
$$
\Ext^p_A(K, H^q(L)) \Rightarrow \text{Ext}^{p + q}_A(K, L)
$$
ainsi que de suites spectrales convergentes
$$
\Ext^i_A(H^{-j}(K), H^q(L)) \Rightarrow \text{Ext}^{i + j}_A(K, H^q(L))
$$
Voir Injectifs, remarques \ref{injectives-remark-spectral-sequences-ext}
et \ref{injectives-remark-spectral-sequences-ext-variant}.
Cela achève la démonstration, car les modules $\Ext^p_A(M, N)$
sont de type fini pour des $A$-modules de type fini $M$, $N$, d'après
Algèbre, lemme \ref{algebra-lemma-ext-noetherian}.
\end{proof}

\begin{lemma}
\label{lemma-perfect-on-noetherian}
Soit $X$ un schéma noethérien. Soit $E$ dans $D(\mathcal{O}_X)$ un objet parfait.
Alors
\begin{enumerate}
\item $E$ appartient à $D^b_{\textit{Coh}}(\mathcal{O}_X)$,
\item si $L$ appartient à $D_{\textit{Coh}}(\mathcal{O}_X)$, alors
$E \otimes_{\mathcal{O}_X}^\mathbf{L} L$ et
$R\SheafHom_{\mathcal{O}_X}(E, L)$ appartiennent à
$D_{\textit{Coh}}(\mathcal{O}_X)$,
\item si $L$ appartient à $D^b_{\textit{Coh}}(\mathcal{O}_X)$, alors
$E \otimes_{\mathcal{O}_X}^\mathbf{L} L$ et
$R\SheafHom_{\mathcal{O}_X}(E, L)$ appartiennent à
$D^b_{\textit{Coh}}(\mathcal{O}_X)$,
\item si $L$ appartient à $D^+_{\textit{Coh}}(\mathcal{O}_X)$, alors
$E \otimes_{\mathcal{O}_X}^\mathbf{L} L$ et
$R\SheafHom_{\mathcal{O}_X}(E, L)$ appartiennent à
$D^+_{\textit{Coh}}(\mathcal{O}_X)$,
\item si $L$ appartient à $D^-_{\textit{Coh}}(\mathcal{O}_X)$, alors
$E \otimes_{\mathcal{O}_X}^\mathbf{L} L$ et
$R\SheafHom_{\mathcal{O}_X}(E, L)$ appartiennent à
$D^-_{\textit{Coh}}(\mathcal{O}_X)$.
\end{enumerate}
\end{lemma}

\begin{proof}
Puisque $X$ est quasi-compact, chacune de ces assertions se vérifie
sur les membres d'un recouvrement ouvert quelconque de $X$.
On peut donc supposer $E$ représenté par
un complexe borné $\mathcal{E}^\bullet$ de modules libres de rang fini ; voir
Cohomologie, lemme \ref{cohomology-lemma-perfect-on-locally-ringed}.
Dans ce cas, chacune des assertions est claire, car
$R\SheafHom_{\mathcal{O}_X}(E, L)$
et $E \otimes_{\mathcal{O}_X}^\mathbf{L} L$ se calculent au
niveau des complexes à l'aide de $\mathcal{E}^\bullet$ ; voir
Cohomologie, lemmes \ref{cohomology-lemma-Rhom-strictly-perfect} et
\ref{cohomology-lemma-bounded-flat-K-flat}. Nous omettons quelques détails.
\end{proof}

\begin{lemma}
\label{lemma-ext-finite}
Soit $A$ un anneau noethérien. Soit $X$ un schéma propre sur $A$.
Pour $L$ dans
$D^+_{\textit{Coh}}(\mathcal{O}_X)$ et $K$ dans
$D^-_{\textit{Coh}}(\mathcal{O}_X)$, les $A$-modules
$\Ext_{\mathcal{O}_X}^n(K, L)$ sont de type fini.
\end{lemma}

\begin{proof}
Rappelons que
$$
\Ext_{\mathcal{O}_X}^n(K, L) =
H^n(X, R\SheafHom_{\mathcal{O}_X}(K, L)) =
H^n(\Spec(A), Rf_*R\SheafHom_{\mathcal{O}_X}(K, L))
$$
voir Cohomologie, lemme \ref{cohomology-lemma-section-RHom-over-U}
et Cohomologie, section \ref{cohomology-section-Leray}.
Le résultat découle donc des
lemmes \ref{lemma-coherent-internal-hom} et
\ref{lemma-direct-image-coherent-bdd-below}.
\end{proof}

\begin{lemma}
\label{lemma-perfect-on-regular}
Soit $X$ un schéma régulier. Alors
tout objet de $D^b_{\textit{Coh}}(\mathcal{O}_X)$ est parfait.
Si $X$ est quasi-compact, c'est-à-dire noethérien et régulier,
alors, réciproquement, tout objet parfait de
$D(\mathcal{O}_X)$ appartient à $D^b_{\textit{Coh}}(\mathcal{O}_X)$.
\end{lemma}

\begin{proof}
Soit $K$ un objet de $D^b_{\textit{Coh}}(\mathcal{O}_X)$.
Pour vérifier que $K$ est parfait, on peut travailler sur les ouverts affines de $X$
(voir Cohomologie, section \ref{cohomology-section-perfect}).
Alors $K$ est parfait d'après le lemme \ref{lemma-perfect-affine} et
Compléments sur l'algèbre, lemme \ref{more-algebra-lemma-regular-perfect}.
La réciproque est le lemme \ref{lemma-perfect-on-noetherian}.
\end{proof}



\section{Descente des propriétés de finitude des complexes}
\label{section-descent-finiteness}

\noindent
Cette section est l'analogue de
Descente, section \ref{descent-section-descent-finiteness},
pour les objets de la catégorie dérivée d'un schéma.
Le résultat le plus simple de ce type est sans doute le suivant.

\begin{lemma}
\label{lemma-tor-amplitude-descends}
Soit $f : X \to Y$ un morphisme surjectif et plat de schémas
(ou, plus généralement, d'espaces localement annelés).
Soit $E \in D(\mathcal{O}_Y)$. Soient $a, b \in \mathbf{Z}$.
Alors $E$ a une amplitude de Tor contenue dans $[a, b]$ si et seulement si
$Lf^*E$ a une amplitude de Tor contenue dans $[a, b]$.
\end{lemma}

\begin{proof}
L'image réciproque préserve toujours l'amplitude de Tor ; voir
Cohomologie, lemme \ref{cohomology-lemma-tor-amplitude-pullback}.
On peut vérifier sur les fibres que l'amplitude de Tor est contenue dans $[a, b]$ ; voir
Cohomologie, lemme \ref{cohomology-lemma-tor-amplitude-stalk}.
Un homomorphisme local et plat d'anneaux est fidèlement plat d'après
Algèbre, lemme \ref{algebra-lemma-local-flat-ff}.
Le résultat découle donc de
Compléments sur l'algèbre, lemme
\ref{more-algebra-lemma-flat-descent-tor-amplitude}.
\end{proof}

\begin{lemma}
\label{lemma-pseudo-coherent-descends-fpqc}
Soit $\{f_i : X_i \to X\}$ un recouvrement fpqc de schémas. Soit
$E \in D_\QCoh(\mathcal{O}_X)$. Soit $m \in \mathbf{Z}$.
Alors $E$ est $m$-pseudo-cohérent si et seulement si chaque
$Lf_i^*E$ est $m$-pseudo-cohérent.
\end{lemma}

\begin{proof}
L'image réciproque préserve toujours la $m$-pseudo-cohérence ; voir
Cohomologie, lemme \ref{cohomology-lemma-pseudo-coherent-pullback}.
Réciproquement, supposons que $Lf_i^*E$ soit $m$-pseudo-cohérent pour tout $i$.
Soit $U \subset X$ un ouvert affine. Il suffit de démontrer que
$E|_U$ est $m$-pseudo-cohérent. Puisque $\{f_i : X_i \to X\}$ est un
recouvrement fpqc, on peut trouver un nombre fini d'ouverts affines $V_j \subset X_{a(j)}$
tels que $f_{a(j)}(V_j) \subset U$ et $U = \bigcup f_{a(j)}(V_j)$.
Posons $V = \coprod V_i$.
On peut donc remplacer $X$ par $U$ et $\{f_i : X_i \to X\}$ par
$\{V \to U\}$, et supposer que $X$ est affine et que le recouvrement
est donné par un unique morphisme surjectif et plat $\{f : Y \to X\}$
de schémas affines. Dans ce cas, le résultat découle de
Compléments sur l'algèbre, lemme \ref{more-algebra-lemma-flat-descent-pseudo-coherent},
à l'aide des lemmes \ref{lemma-affine-compare-bounded} et
\ref{lemma-pseudo-coherent-affine}.
\end{proof}

\begin{lemma}
\label{lemma-pseudo-coherent-descends-fppf}
Soit $\{f_i : X_i \to X\}$ un recouvrement fppf de schémas. Soit
$E \in D(\mathcal{O}_X)$. Soit $m \in \mathbf{Z}$.
Alors $E$ est $m$-pseudo-cohérent si et seulement si chaque
$Lf_i^*E$ est $m$-pseudo-cohérent.
\end{lemma}

\begin{proof}
L'image réciproque préserve toujours la $m$-pseudo-cohérence ; voir
Cohomologie, lemme \ref{cohomology-lemma-pseudo-coherent-pullback}.
Réciproquement, supposons que $Lf_i^*E$ soit $m$-pseudo-cohérent pour tout $i$.
Soit $U \subset X$ un ouvert affine. Il suffit de démontrer que
$E|_U$ est $m$-pseudo-cohérent. Puisque $\{f_i : X_i \to X\}$ est un
recouvrement fppf, on peut trouver un nombre fini d'ouverts affines $V_j \subset X_{a(j)}$
tels que $f_{a(j)}(V_j) \subset U$ et $U = \bigcup f_{a(j)}(V_j)$.
Posons $V = \coprod V_i$.
On peut donc remplacer $X$ par $U$ et $\{f_i : X_i \to X\}$ par
$\{V \to U\}$, et supposer que $X$ est affine et que le recouvrement
est donné par un unique morphisme surjectif et plat $\{f : Y \to X\}$
de présentation finie.

\medskip\noindent
Puisque $f$ est plat, le foncteur dérivé $Lf^*$ est donné par $f^*$, et $f^*$
est exact. Ainsi $H^i(Lf^*E) = f^*H^i(E)$. Puisque $Lf^*E$ est $m$-pseudo-cohérent,
on a $Lf^*E \in D^-(\mathcal{O}_Y)$. Puisque $f$ est surjectif et plat,
on a $E \in D^-(\mathcal{O}_X)$. Soit $i \in \mathbf{Z}$ le plus grand
entier tel que $H^i(E)$ soit non nul. Si $i < m$, la conclusion est acquise. Sinon,
$f^*H^i(E)$ est un $\mathcal{O}_Y$-module de type fini d'après
Cohomologie, lemme \ref{cohomology-lemma-finite-cohomology}.
D'après Descente, lemme \ref{descent-lemma-finite-type-descends-fppf},
le $\mathcal{O}_X$-module $H^i(E)$ est alors de type fini.
Ainsi, quitte à remplacer $X$ par les membres d'un recouvrement ouvert affine fini,
on peut supposer qu'il existe un morphisme
$$
\alpha : \mathcal{O}_X^{\oplus n}[-i] \longrightarrow E
$$
tel que $H^i(\alpha)$ soit surjectif. Soit $C$ le cône de $\alpha$
dans $D(\mathcal{O}_X)$. En prenant l'image réciproque sur $Y$ et en utilisant
Cohomologie, lemme \ref{cohomology-lemma-cone-pseudo-coherent},
on voit que $Lf^*C$ est $m$-pseudo-cohérent. De plus, $H^j(C) = 0$
pour $j \geq i$. Par récurrence sur $i$, on en déduit que $C$ est
$m$-pseudo-cohérent. En appliquant de nouveau
Cohomologie, lemme \ref{cohomology-lemma-cone-pseudo-coherent},
on conclut.
\end{proof}

\begin{lemma}
\label{lemma-perfect-descends-fpqc}
Soit $\{f_i : X_i \to X\}$ un recouvrement fpqc de schémas. Soit
$E \in D(\mathcal{O}_X)$. Alors $E$ est parfait
si et seulement si chaque $Lf_i^*E$ est parfait.
\end{lemma}

\begin{proof}
L'image réciproque préserve toujours les complexes parfaits ; voir
Cohomologie, lemme \ref{cohomology-lemma-perfect-pullback}.
Réciproquement, supposons que $Lf_i^*E$ soit parfait pour tout $i$.
Les faisceaux de cohomologie de chaque $Lf_i^*E$ sont alors quasi-cohérents ; voir
le lemme \ref{lemma-pseudo-coherent}
et
Cohomologie, lemme \ref{cohomology-lemma-perfect}.
Puisque les morphismes $f_i$ sont plats, on a $H^p(Lf_i^*E) = f_i^*H^p(E)$.
Les faisceaux de cohomologie de $E$ sont donc quasi-cohérents d'après
Descente, proposition \ref{descent-proposition-fpqc-descent-quasi-coherent}.
Cela étant, le lemme résulte formellement de
Cohomologie, lemme \ref{cohomology-lemma-perfect},
et des
lemmes \ref{lemma-tor-amplitude-descends} et
\ref{lemma-pseudo-coherent-descends-fpqc}.
\end{proof}

\begin{lemma}
\label{lemma-closed-push-pseudo-coherent}
Soit $i : Z \to X$ un morphisme d'espaces annelés tel que
$i$ soit une immersion fermée des espaces topologiques sous-jacents et que
$i_*\mathcal{O}_Z$ soit pseudo-cohérent en tant que $\mathcal{O}_X$-module.
Soit $E \in D(\mathcal{O}_Z)$. Alors $E$ est $m$-pseudo-cohérent
si et seulement si $Ri_*E$ est $m$-pseudo-cohérent.
\end{lemma}

\begin{proof}
Tout au long de cette démonstration, on utilisera l'exactitude du foncteur $i_*$ et,
par suite, l'égalité $Ri_* = i_*$ ; voir Modules, lemme \ref{modules-lemma-i-star-exact}.

\medskip\noindent
Supposons $E$ $m$-pseudo-cohérent. Soit $x \in X$. Nous allons trouver un voisinage
de $x$ sur lequel $i_*E$ soit $m$-pseudo-cohérent. Si $x \not \in Z$,
c'est clair. On peut donc supposer $x \in Z$. Nous utiliserons le fait
que les $U \cap Z$, pour $x \in U \subset X$ ouvert, forment un système fondamental de
voisinages de $x$ dans $Z$. Quitte à restreindre $X$, on peut supposer $E$
borné supérieurement. Nous raisonnons par récurrence sur
le plus grand entier $p$ tel que $H^p(E)$ soit non nul. Si $p < m$,
il n'y a rien à démontrer. Si $p \geq m$, alors $H^p(E)$ est un
$\mathcal{O}_Z$-module de type fini ; voir
Cohomologie, lemme \ref{cohomology-lemma-finite-cohomology}.
On peut donc choisir, quitte à restreindre $X$, un morphisme
$\mathcal{O}_Z^{\oplus n}[-p] \to E$ qui induit une surjection
$\mathcal{O}_Z^{\oplus n} \to H^p(E)$. Choisissons un triangle distingué
$$
\mathcal{O}_Z^{\oplus n}[-p] \to E \to C \to \mathcal{O}_Z^{\oplus n}[-p + 1]
$$
On voit que $H^j(C) = 0$ pour $j \geq p$ et que $C$ est $m$-pseudo-cohérent
d'après Cohomologie, lemme \ref{cohomology-lemma-cone-pseudo-coherent}.
Par récurrence, $i_*C$ est $m$-pseudo-cohérent sur $X$.
Puisque $i_*\mathcal{O}_Z$ est aussi $m$-pseudo-cohérent sur $X$, on déduit
du triangle distingué
$$
i_*\mathcal{O}_Z^{\oplus n}[-p] \to i_*E \to i_*C \to
i_*\mathcal{O}_Z^{\oplus n}[-p + 1]
$$
et de
Cohomologie, lemme \ref{cohomology-lemma-cone-pseudo-coherent},
que $i_*E$ est $m$-pseudo-cohérent.

\medskip\noindent
Supposons que $i_*E$ soit $m$-pseudo-cohérent. Soit $z \in Z$.
Nous allons trouver un voisinage de $z$ sur lequel $E$
soit $m$-pseudo-cohérent. Nous utiliserons le fait
que les $U \cap Z$, pour $z \in U \subset X$ ouvert, forment un système fondamental de
voisinages de $z$ dans $Z$. Quitte à restreindre $X$, on peut supposer $i_*E$
et donc $E$ bornés supérieurement. Nous raisonnons par récurrence sur
le plus grand entier $p$ tel que $H^p(E)$ soit non nul. Si $p < m$,
il n'y a rien à démontrer. Si $p \geq m$, alors $H^p(i_*E) = i_*H^p(E)$
est un $\mathcal{O}_X$-module de type fini ; voir
Cohomologie, lemme \ref{cohomology-lemma-finite-cohomology}.
Choisissons un complexe $\mathcal{E}^\bullet$ de $\mathcal{O}_Z$-modules
représentant $E$. On peut choisir, quitte à restreindre $X$,
un morphisme $\alpha : \mathcal{O}_X^{\oplus n}[-p] \to i_*\mathcal{E}^\bullet$
qui induit une surjection
$\mathcal{O}_X^{\oplus n} \to i_*H^p(\mathcal{E}^\bullet)$.
Par adjonction, on obtient un morphisme
$\alpha : \mathcal{O}_Z^{\oplus n}[-p] \to \mathcal{E}^\bullet$
qui induit une surjection
$\mathcal{O}_Z^{\oplus n} \to H^p(\mathcal{E}^\bullet)$.
Choisissons un triangle distingué
$$
\mathcal{O}_Z^{\oplus n}[-p] \to E \to C \to \mathcal{O}_Z^{\oplus n}[-p + 1]
$$
On voit que $H^j(C) = 0$ pour $j \geq p$. Du triangle distingué
$$
i_*\mathcal{O}_Z^{\oplus n}[-p] \to i_*E \to i_*C \to
i_*\mathcal{O}_Z^{\oplus n}[-p + 1]
$$
du fait que $i_*\mathcal{O}_Z$ est pseudo-cohérent
et de
Cohomologie, lemme \ref{cohomology-lemma-cone-pseudo-coherent},
on déduit que $i_*C$ est $m$-pseudo-cohérent.
Par récurrence, $C$ est $m$-pseudo-cohérent.
En appliquant de nouveau Cohomologie, lemme \ref{cohomology-lemma-cone-pseudo-coherent},
on en déduit que $E$ est $m$-pseudo-cohérent.
\end{proof}

\begin{lemma}
\label{lemma-finite-push-pseudo-coherent}
Soit $f : X \to Y$ un morphisme fini de schémas tel que
$f_*\mathcal{O}_X$ soit pseudo-cohérent en tant que
$\mathcal{O}_Y$-module\footnote{Cela signifie que $f$ est pseudo-cohérent ; voir
Compléments sur les morphismes, lemme
\ref{more-morphisms-lemma-finite-pseudo-coherent}.}.
Soit $E \in D_\QCoh(\mathcal{O}_X)$. Alors $E$ est $m$-pseudo-cohérent
si et seulement si $Rf_*E$ est $m$-pseudo-cohérent.
\end{lemma}

\begin{proof}
Il s'agit de la traduction de
Compléments sur l'algèbre, lemme \ref{more-algebra-lemma-finite-push-pseudo-coherent},
dans le langage des schémas. Pour effectuer ce passage, utiliser les
lemmes \ref{lemma-affine-compare-bounded} et
\ref{lemma-pseudo-coherent-affine}.
\end{proof}


\section{Relèvement des complexes}
\label{section-lift}

\noindent
Soit $U \subset X$ un sous-espace ouvert d'un espace annelé,
et notons $j : U \to X$ le morphisme d'inclusion. Le foncteur
$D(\mathcal{O}_X) \to D(\mathcal{O}_U)$ est essentiellement surjectif car
$Rj_*$ est un inverse à droite de la restriction.
Dans cette section, nous étendons ce résultat aux complexes dont les faisceaux
de cohomologie sont quasi-cohérents, etc.

\begin{lemma}
\label{lemma-lift-quasi-coherent}
Soient $X$ un schéma et $j : U \to X$ une immersion ouverte
quasi-compacte. Les foncteurs
$$
D_\QCoh(\mathcal{O}_X) \to D_\QCoh(\mathcal{O}_U)
\quad\text{et}\quad
D^+_\QCoh(\mathcal{O}_X) \to D^+_\QCoh(\mathcal{O}_U)
$$
sont essentiellement surjectifs. Si $X$ est quasi-compact, les foncteurs
$$
D^-_\QCoh(\mathcal{O}_X) \to D^-_\QCoh(\mathcal{O}_U)
\quad\text{et}\quad
D^b_\QCoh(\mathcal{O}_X) \to D^b_\QCoh(\mathcal{O}_U)
$$
sont essentiellement surjectifs.
\end{lemma}

\begin{proof}
L'argument précédant le lemme s'applique au premier cas, car $Rj_*$
envoie $D_\QCoh(\mathcal{O}_U)$ dans $D_\QCoh(\mathcal{O}_X)$
d'après le lemme \ref{lemma-quasi-coherence-direct-image}.
Il est clair que $Rj_*$ envoie
$D^+_\QCoh(\mathcal{O}_U)$ dans
$D^+_\QCoh(\mathcal{O}_X)$,
ce qui donne l'assertion relative aux complexes bornés inférieurement.
Enfin, le lemme \ref{lemma-quasi-coherence-direct-image}
assure que $Rj_*$ envoie
$D^-_\QCoh(\mathcal{O}_U)$ dans
$D^-_\QCoh(\mathcal{O}_X)$
si $X$ est quasi-compact. La dernière assertion résulte de ces deux faits.
\end{proof}

\begin{lemma}
\label{lemma-lift-coherent}
Soient $X$ un schéma noethérien et $j : U \to X$ une immersion ouverte.
Le foncteur
$D^b_{\textit{Coh}}(\mathcal{O}_X) \to D^b_{\textit{Coh}}(\mathcal{O}_U)$
est essentiellement surjectif.
\end{lemma}

\begin{proof}
Soit $K$ un objet de $D^b_{\textit{Coh}}(\mathcal{O}_U)$.
D'après la proposition \ref{proposition-DCoh}, on peut représenter $K$ par un complexe
borné $\mathcal{F}^\bullet$ de $\mathcal{O}_U$-modules cohérents.
Écrivons $\mathcal{F}^i = 0$ pour $i \not \in [a, b]$, avec $a \leq b$.
Puisque $j$ est quasi-compact et séparé, les termes du complexe borné
$j_*\mathcal{F}^\bullet$ sont des modules quasi-cohérents sur $X$, voir
Schémas, lemme \ref{schemes-lemma-push-forward-quasi-coherent}.
Choisissons par récurrence un sous-module cohérent
$\mathcal{G}^i \subset j_*\mathcal{F}^i$ de la façon suivante.
Pour $i = a$, choisissons un sous-module cohérent quelconque
$\mathcal{G}^a \subset j_*\mathcal{F}^a$ dont la restriction
à $U$ soit $\mathcal{F}^a$. C'est possible d'après
Propriétés, lemme \ref{properties-lemma-extend}.
Pour $i > a$, choisissons d'abord un sous-module cohérent quelconque
$\mathcal{H}^i \subset j_*\mathcal{F}^i$
dont la restriction à $U$ soit $\mathcal{F}^i$,
puis posons
$\mathcal{G}^i = \Im(\mathcal{H}^i \oplus \mathcal{G}^{i - 1}
\to j_*\mathcal{F}^i)$. Il est clair que
$\mathcal{G}^\bullet \subset j_*\mathcal{F}^\bullet$
est un complexe borné de $\mathcal{O}_X$-modules cohérents
dont la restriction à $U$ est $\mathcal{F}^\bullet$, comme voulu.
\end{proof}

\begin{lemma}
\label{lemma-lift-pseudo-coherent}
Soient $X$ un schéma affine et $U \subset X$ un sous-schéma ouvert
quasi-compact. Pour tout objet pseudo-cohérent $E$ de $D(\mathcal{O}_U)$,
il existe un complexe borné supérieurement de $\mathcal{O}_X$-modules libres de rang fini
dont la restriction à $U$ est isomorphe à $E$.
\end{lemma}

\begin{proof}
D'après le lemme \ref{lemma-pseudo-coherent}, $E$ est un objet de
$D_\QCoh(\mathcal{O}_U)$. D'après le
lemme \ref{lemma-lift-quasi-coherent},
on peut supposer que $E = E'|U$ pour un objet $E'$ de
$D_\QCoh(\mathcal{O}_X)$.
Écrivons $X = \Spec(A)$. D'après le lemme \ref{lemma-affine-compare-bounded},
on peut trouver un complexe $M^\bullet$ de $A$-modules dont le
complexe de $\mathcal{O}_X$-modules associé représente $E'$.

\medskip\noindent
Choisissons $f_1, \ldots, f_r \in A$ tels que $U = D(f_1) \cup \ldots \cup D(f_r)$.
D'après le lemme \ref{lemma-pseudo-coherent-affine}, les complexes
$M^\bullet_{f_j}$ sont des complexes pseudo-cohérents de $A_{f_j}$-modules.
Soit $n$ un entier. Supposons donné un morphisme de complexes
$\alpha : F^\bullet \to M^\bullet$, où $F^\bullet$ est
borné supérieurement, $F^i = 0$ pour $i < n$, et chaque $F^i$ est un
$R$-module libre de rang fini, tel que
$$
H^i(\alpha_{f_j}) : H^i(F^\bullet_{f_j}) \to H^i(M^\bullet_{f_j})
$$
soit un isomorphisme pour $i > n$ et soit surjectif pour $i = n$. On a le diagramme
$$
\xymatrix{
& F^n \ar[r] \ar[d]^\alpha & F^{n + 1} \ar[d]^\alpha \ar[r] & \ldots \\
M^{n-1} \ar[r] & M^n \ar[r] & M^{n + 1} \ar[r] & \ldots
}
$$
Puisque chaque $M^\bullet_{f_j}$ a sa cohomologie nulle
en degrés assez grands, on peut trouver un tel morphisme pour $n \gg 0$.
Par récurrence sur $n$, nous allons le prolonger en un morphisme
de complexes $F^\bullet \to M^\bullet$
tel que $H^i(\alpha_{f_j})$ soit un isomorphisme
pour tout $i$. Le lemme en résultera en prenant $\widetilde{F^\bullet}$.

\medskip\noindent
L'étape de récurrence consiste à prolonger le diagramme
ci-dessus en ajoutant $F^{n - 1}$. Soit $C^\bullet$ le cône de $\alpha$
(Catégories dérivées, définition \ref{derived-definition-cone}).
La suite exacte longue de cohomologie montre que
$H^i(C^\bullet_{f_j}) = 0$ pour $i \geq n$. D'après
Compléments sur l'algèbre, lemme \ref{more-algebra-lemma-cone-pseudo-coherent},
$C^\bullet_{f_j}$ est $(n - 1)$-pseudo-cohérent. D'après
Compléments sur l'algèbre, lemme \ref{more-algebra-lemma-finite-cohomology},
$H^{n - 1}(C^\bullet_{f_j})$ est un $A_{f_j}$-module de type fini.
Choisissons un $A$-module libre de rang fini $F^{n - 1}$ et un morphisme de $A$-modules
$\beta : F^{n - 1} \to C^{n - 1}$ tels que le composé
$F^{n - 1} \to C^{n - 1} \to C^n$ soit nul et que
$F^{n - 1}_{f_j}$ se surjecte sur $H^{n - 1}(C^\bullet_{f_j})$.
(Certains détails sont omis ; indication : chasser les dénominateurs.)
Puisque $C^{n - 1} = M^{n - 1} \oplus F^n$,
on peut écrire $\beta = (\alpha^{n - 1}, -d^{n - 1})$. La nullité du
composé $F^{n - 1} \to C^{n - 1} \to C^n$ implique
que ces applications s'insèrent dans un morphisme de complexes
$$
\xymatrix{
& F^{n - 1} \ar[d]^{\alpha^{n - 1}} \ar[r]_{d^{n - 1}} &
F^n \ar[r] \ar[d]^\alpha &
F^{n + 1} \ar[d]^\alpha \ar[r] & \ldots \\
\ldots \ar[r] &
M^{n - 1} \ar[r] & M^n \ar[r] & M^{n + 1} \ar[r] & \ldots
}
$$
De plus, ces applications définissent un morphisme de triangles distingués
$$
\xymatrix{
(F^n \to \ldots) \ar[r] \ar[d] &
(F^{n-1} \to \ldots) \ar[r] \ar[d] &
F^{n-1} \ar[r] \ar[d]_\beta &
(F^n \to \ldots)[1] \ar[d] \\
(F^n \to \ldots) \ar[r] &
M^\bullet \ar[r] &
C^\bullet \ar[r] &
(F^n \to \ldots)[1]
}
$$
Le choix de $\beta$ implique donc que le morphisme de complexes
$(F^{-1} \to \ldots) \to M^\bullet$ induit un isomorphisme sur la
cohomologie localisée en $f_j$ en degrés $\geq n$ et une surjection
en degré $n - 1$. Ceci achève la démonstration du lemme.
\end{proof}

\noindent
Les deux lemmes suivants devraient probablement être placés ailleurs.

\begin{lemma}
\label{lemma-vanishing-ext}
Soit $X$ un schéma quasi-compact et quasi-séparé.
Soit $E \in D^b_\QCoh(\mathcal{O}_X)$.
Il existe un entier $n_0 > 0$ tel que
$\Ext^n_{D(\mathcal{O}_X)}(\mathcal{E}, E) = 0$
pour tout
$\mathcal{O}_X$-module localement libre de rang fini $\mathcal{E}$ et tout $n \geq n_0$.
\end{lemma}

\begin{proof}
Rappelons que $\Ext^n_{D(\mathcal{O}_X)}(\mathcal{E}, E) =
\Hom_{D(\mathcal{O}_X)}(\mathcal{E}, E[n])$. Nous disposons de
Mayer-Vietoris pour les morphismes dans la catégorie dérivée, voir
Cohomologie, lemme \ref{cohomology-lemma-mayer-vietoris-hom}.
Ainsi, si $X = U \cup V$ et si la conclusion du lemme vaut
pour $E|_U$, $E|_V$ et $E|_{U \cap V}$ avec une borne $n_0$,
elle vaut pour $E$ avec la borne $n_0 + 1$.
Il suffit donc de démontrer le lemme lorsque $X$ est affine, voir
Cohomologie des schémas, lemme \ref{coherent-lemma-induction-principle}.

\medskip\noindent
Supposons que $X = \Spec(A)$ soit affine. Choisissons un complexe de $A$-modules
$M^\bullet$ dont le complexe de modules quasi-cohérents associé
représente $E$, voir le lemme \ref{lemma-affine-compare-bounded}.
Écrivons $\mathcal{E} = \widetilde{P}$ pour un $A$-module $P$.
Puisque $\mathcal{E}$ est localement libre de rang fini, $P$
est un $A$-module projectif de type fini. On a
\begin{align*}
\Hom_{D(\mathcal{O}_X)}(\mathcal{E}, E[n])
& = 
\Hom_{D(A)}(P, M^\bullet[n]) \\
& =
\Hom_{K(A)}(P, M^\bullet[n]) \\
& =
\Hom_A(P, H^n(M^\bullet))
\end{align*}
La première égalité résulte du lemme \ref{lemma-affine-compare-bounded},
la deuxième de
Catégories dérivées, lemme
\ref{derived-lemma-morphisms-from-projective-complex}, et
la dernière de l'exactitude du foncteur $\Hom_A(P, -)$.
Comme $E$, et donc $M^\bullet$, est borné,
on obtient zéro pour tout $n$ assez grand.
\end{proof}

\noindent
On peut renforcer le lemme suivant (l'annulation est uniforme
pour tous les $L$ dont les faisceaux de cohomologie non nuls
sont concentrés dans un intervalle fixé).

\begin{lemma}
\label{lemma-ext-from-perfect-into-bounded-QCoh}
Soit $X$ un schéma quasi-compact et quasi-séparé.
Soit $K$ un objet parfait de $D(\mathcal{O}_X)$. Alors
\begin{enumerate}
\item il existe des entiers $a \leq b$ tels que, pour tout
$L \in D_\QCoh(\mathcal{O}_X)$ avec $H^i(L) = 0$ pour $i \in [a, b]$,
on ait $\Hom_{D(\mathcal{O}_X)}(K, L) = 0$, et
\item si $L$ est borné, alors $\Ext^n_{D(\mathcal{O}_X)}(K, L)$
est nul sauf pour un nombre fini de valeurs de $n$.
\end{enumerate}
\end{lemma}

\begin{proof}
L'assertion (2) résulte de (1), car $\Ext^n_{D(\mathcal{O}_X)}(K, L) =
\Hom_{D(\mathcal{O}_X)}(K, L[n])$. Démontrons (1).
Puisque $K$ est parfait, on a
$$
\Hom_{D(\mathcal{O}_X)}(K, L) =
H^0(X, K^\vee \otimes_{\mathcal{O}_X}^\mathbf{L} L)
$$
où $K^\vee$ est le complexe parfait « dual » de $K$, voir
Cohomologie, lemme \ref{cohomology-lemma-dual-perfect-complex}.
Notons que $K^\vee \otimes_{\mathcal{O}_X}^\mathbf{L} L$
appartient à $D_\QCoh(X)$ d'après les
lemmes \ref{lemma-quasi-coherence-tensor-product} et
\ref{lemma-pseudo-coherent} (ce dernier montre qu'un complexe parfait
a des faisceaux de cohomologie quasi-cohérents). Supposons que $K^\vee$ ait
une amplitude de Tor contenue dans $[a, b]$. Alors la suite spectrale
$$
E_1^{p, q} = H^p(K^\vee \otimes_{\mathcal{O}_X}^\mathbf{L} H^q(L))
\Rightarrow
H^{p + q}(K^\vee \otimes_{\mathcal{O}_X}^\mathbf{L} L)
$$
montre que $H^j(K^\vee \otimes_{\mathcal{O}_X}^\mathbf{L} L)$
est nul si $H^q(L) = 0$ pour $q \in [j - b, j - a]$.
Soit $N$ l'entier $d$ de Cohomologie des schémas,
lemme \ref{coherent-lemma-vanishing-nr-affines-quasi-separated}.
Alors $H^0(X, K^\vee \otimes_{\mathcal{O}_X}^\mathbf{L} L)$
est nul si les faisceaux de cohomologie
$$
H^{-N}(K^\vee \otimes_{\mathcal{O}_X}^\mathbf{L} L),
\ H^{-N + 1}(K^\vee \otimes_{\mathcal{O}_X}^\mathbf{L} L),
\ \ldots,
\ H^0(K^\vee \otimes_{\mathcal{O}_X}^\mathbf{L} L)
$$
sont nuls. En effet, d'après le lemme cité et le
lemme \ref{lemma-application-nice-K-injective}, on a
$$
H^0(X, K^\vee \otimes_{\mathcal{O}_X}^\mathbf{L} L) =
H^0(X, \tau_{\geq -N}(K^\vee \otimes_{\mathcal{O}_X}^\mathbf{L} L))
$$
et, par annulation des faisceaux de cohomologie, ceci est égal à
$H^0(X, \tau_{\geq 1}(K^\vee \otimes_{\mathcal{O}_X}^\mathbf{L} L))$,
qui est nul d'après Catégories dérivées, lemme
\ref{derived-lemma-negative-vanishing}.
Il en résulte que $\Hom_{D(\mathcal{O}_X)}(K, L)$ est nul si
$H^i(L) = 0$ pour $i \in [-b - N, -a]$.
\end{proof}

\begin{lemma}
\label{lemma-lift-perfect-complex-plus-locally-free}
Soit $X$ un schéma affine. Soit $U \subset X$ un ouvert quasi-compact.
Pour tout objet parfait $E$ de $D(\mathcal{O}_U)$, il existe un entier
$r$ et un faisceau localement libre de rang fini $\mathcal{F}$ sur $U$ tels que
$\mathcal{F}[-r] \oplus E$ soit la restriction d'un objet parfait de
$D(\mathcal{O}_X)$.
\end{lemma}

\begin{proof}
Écrivons $X = \Spec(A)$. Rappelons qu'un complexe parfait est
pseudo-cohérent, voir
Cohomologie, lemme \ref{cohomology-lemma-perfect}.
D'après le lemme \ref{lemma-lift-pseudo-coherent}, on peut trouver un complexe borné supérieurement
$\mathcal{F}^\bullet$ de $A$-modules libres de rang fini tel que $E$ soit
isomorphe à $\mathcal{F}^\bullet|_U$ dans $D(\mathcal{O}_U)$.
D'après Cohomologie, lemme \ref{cohomology-lemma-perfect}, et puisque
$U$ est quasi-compact, $E$ est de dimension de Tor finie ; supposons que
$E$ ait une amplitude de Tor contenue dans $[a, b]$. Choisissons $r < a$ et posons
$$
\mathcal{K} = \Ker(\mathcal{F}^{r} \to \mathcal{F}^{r + 1})
= \Im(\mathcal{F}^{r - 1} \to \mathcal{F}^r).
$$
Puisque $E$ a une amplitude de Tor contenue dans $[a, b]$,
$\mathcal{F} = \mathcal{K}|_U$ est
plat (Cohomologie, lemme \ref{cohomology-lemma-last-one-flat}).
Ainsi $\mathcal{F}$ est plat et de présentation finie, donc localement
libre de rang fini (Propriétés, lemme \ref{properties-lemma-finite-locally-free}).
Il en résulte que
$$
\mathcal{F} \to \mathcal{F}^r|_U \to \mathcal{F}^{r + 1}|_U \to \ldots
$$
est un complexe strictement parfait sur $U$ représentant $E$. D'autre part,
le complexe $P = (\mathcal{F}^r \to \mathcal{F}^{r + 1} \to \ldots )$
est un complexe parfait sur $X$. Les
troncatures bêtes donnent un triangle distingué
$$
P|_U \to E \to \mathcal{F}[-r - 1] \to (P|_U)[1]
$$
Si le morphisme $E \to \mathcal{F}[-r - 1]$ est nul dans $D(\mathcal{O}_U)$,
alors $P|_U = \mathcal{F}[-r - 2] \oplus E$, voir
Catégories dérivées, lemme \ref{derived-lemma-split}.
C'est le cas pour $r \ll 0$, par exemple d'après le
lemme \ref{lemma-ext-from-perfect-into-bounded-QCoh}.
\end{proof}

\begin{lemma}
\label{lemma-lift-map}
Soit $X$ un schéma affine. Soit $U \subset X$ un ouvert quasi-compact.
Soient $E, E'$ des objets de $D_\QCoh(\mathcal{O}_X)$, avec $E$ parfait.
Pour tout morphisme $\alpha : E|_U \to E'|_U$, il existe des morphismes
$$
E \xleftarrow{\beta} E_1 \xrightarrow{\gamma} E'
$$
de complexes sur $X$, avec $E_1$ parfait, tels que $\beta : E_1 \to E$
se restreigne en un isomorphisme sur $U$ et que
$\alpha = \gamma|_U \circ \beta|_U^{-1}$.
De plus, on peut supposer que $E_1 = E \otimes_{\mathcal{O}_X}^\mathbf{L} I$
pour un complexe parfait $I$ sur $X$.
\end{lemma}

\begin{proof}
Écrivons $X = \Spec(A)$. Écrivons $U = D(f_1) \cup \ldots \cup D(f_r)$. Choisissons
un complexe borné de $A$-modules projectifs de type fini $M^\bullet$ représentant
$E$ (lemme \ref{lemma-perfect-affine}). Choisissons un complexe de $A$-modules
$(M')^\bullet$ représentant $E'$ (lemme \ref{lemma-affine-compare-bounded}).
Alors le complexe $H^\bullet = \Hom_A(M^\bullet, (M')^\bullet)$
est un complexe de $A$-modules dont le complexe de
$\mathcal{O}_X$-modules quasi-cohérents associé représente $R\SheafHom(E, E')$, voir
Cohomologie, lemme \ref{cohomology-lemma-Rhom-strictly-perfect}.
Le morphisme $\alpha$ détermine un élément $s$ de $H^0(U, R\SheafHom(E, E'))$, voir
Cohomologie, lemme \ref{cohomology-lemma-section-RHom-over-U}.
Il existe un $e$ et un morphisme
$$
\xi : I^\bullet(f_1^e, \ldots, f_r^e) \to \Hom_A(M^\bullet, (M')^\bullet)
$$
correspondant à $s$, voir la
proposition \ref{proposition-represent-cohomology-class-on-open}.
En prenant pour $E_1$ l'objet correspondant au
complexe de $\mathcal{O}_X$-modules quasi-cohérents
associé à
$$
\text{Tot}(I^\bullet(f_1^e, \ldots, f_r^e) \otimes_A M^\bullet)
$$
on obtient $E_1 \to E$ à l'aide du morphisme canonique
$I^\bullet(f_1^e, \ldots, f_r^e) \to A$, et $E_1 \to E'$
à l'aide de $\xi$ et de
Cohomologie, lemme \ref{cohomology-lemma-section-RHom-over-U}.
\end{proof}

\begin{lemma}
\label{lemma-lift-perfect-complex-plus-shift}
Soit $X$ un schéma affine. Soit $U \subset X$ un ouvert quasi-compact.
Pour tout objet parfait $F$ de $D(\mathcal{O}_U)$,
l'objet $F \oplus F[1]$ est la restriction
d'un objet parfait de $D(\mathcal{O}_X)$.
\end{lemma}

\begin{proof}
D'après le lemme \ref{lemma-lift-perfect-complex-plus-locally-free},
on peut trouver un objet parfait $E$ de $D(\mathcal{O}_X)$
tel que $E|_U = \mathcal{F}[r] \oplus F$ pour un
$\mathcal{O}_U$-module localement libre de rang fini $\mathcal{F}$.
D'après le lemme \ref{lemma-lift-map}, on peut trouver un morphisme de
complexes parfaits $\alpha : E_1 \to E$ tel que $(E_1)|_U \cong E|_U$
et que $\alpha|_U$ soit le morphisme
$$
\left(
\begin{matrix}
\text{id}_{\mathcal{F}[r]} & 0 \\
0 & 0
\end{matrix}
\right)
:
\mathcal{F}[r] \oplus F \to \mathcal{F}[r] \oplus F
$$
Le cône de $\alpha$ convient alors.
\end{proof}

\begin{lemma}
\label{lemma-perfect-into-support-on-T}
Soit $X$ un schéma quasi-compact et quasi-séparé.
Soit $f \in \Gamma(X, \mathcal{O}_X)$.
Pour tout morphisme $\alpha : E \to E'$ dans
$D_\QCoh(\mathcal{O}_X)$ tel que
\begin{enumerate}
\item $E$ soit parfait, et
\item $E'$ soit à support dans $T = V(f)$,
\end{enumerate}
il existe un $n \geq 0$ tel que $f^n \alpha  = 0$.
\end{lemma}

\begin{proof}
Nous disposons de Mayer-Vietoris pour les morphismes dans la catégorie dérivée, voir
Cohomologie, lemme \ref{cohomology-lemma-mayer-vietoris-hom}.
Ainsi, si $X = U \cup V$ et si la conclusion du lemme vaut
pour $f|_U$, $f|_V$ et $f|_{U \cap V}$, elle vaut pour $f$.
Il suffit donc de démontrer le lemme lorsque $X$ est affine, voir
Cohomologie des schémas, lemme \ref{coherent-lemma-induction-principle}.

\medskip\noindent
Soit $X = \Spec(A)$. Alors $f \in A$. Nous utiliserons
l'équivalence $D(A) = D_\QCoh(X)$ du
lemme \ref{lemma-affine-compare-bounded}
sans la mentionner davantage.
Représentons $E$ par un complexe borné de $A$-modules projectifs de type fini
$P^\bullet$. C'est possible d'après le lemme \ref{lemma-perfect-affine}.
Soit $t$ le plus grand entier tel que $P^t$ soit non nul.
Le triangle distingué
$$
P^t[-t] \to P^\bullet \to \sigma_{\leq t - 1}P^\bullet \to P^t[-t + 1]
$$
montre que, par récurrence sur la longueur du complexe $P^\bullet$,
on peut se ramener au cas où $P^\bullet$ n'a qu'un seul terme non nul.
Ceci, joint au foncteur de décalage, nous ramène au cas où $P^\bullet$
est constitué d'un seul $A$-module projectif de type fini $P$ en degré $0$.
Représentons $E'$ par un complexe $M^\bullet$ de $A$-modules.
Alors $\alpha$ correspond à un morphisme $P \to H^0(M^\bullet)$.
Puisque le module $H^0(M^\bullet)$ est à support dans $V(f)$ par l'hypothèse (2),
tout élément de $H^0(M^\bullet)$ est annulé par une puissance
de $f$. Comme $P$ est un $A$-module de type fini, le morphisme
$f^n\alpha : P \to H^0(M^\bullet)$ est nul pour un certain $n$, comme voulu.
\end{proof}

\begin{lemma}
\label{lemma-lift-perfect-complex-plus-shift-support}
Soit $X$ un schéma affine. Soit $T \subset X$ un sous-ensemble fermé
tel que $X \setminus T$ soit quasi-compact. Soit $U \subset X$ un
ouvert quasi-compact. Pour tout objet parfait $F$ de $D(\mathcal{O}_U)$
à support dans $T \cap U$, l'objet $F \oplus F[1]$ est la restriction
d'un objet parfait $E$ de $D(\mathcal{O}_X)$ à support dans $T$.
\end{lemma}

\begin{proof}
Écrivons $T = V(g_1, \ldots, g_s)$. En remplaçant $g_j$ par une puissance,
on peut supposer que la multiplication par $g_j$ est nulle sur $F$, voir le
lemme \ref{lemma-perfect-into-support-on-T}. Choisissons $E$ comme dans le
lemme \ref{lemma-lift-perfect-complex-plus-shift}.
Notons que $g_j : E \to E$ se restreint à zéro sur $U$.
Choisissons un triangle distingué
$$
E \xrightarrow{g_1} E \to C_1 \to E[1]
$$
D'après Catégories dérivées, lemme \ref{derived-lemma-split},
l'objet $C_1$ se restreint à
$F \oplus F[1] \oplus F[1] \oplus F[2]$ sur $U$.
De plus, $g_1 : C_1 \to C_1$ est de carré nul d'après
Catégories dérivées, lemme \ref{derived-lemma-third-map-square-zero}.
En effet, le diagramme
$$
\xymatrix{
E \ar[r] \ar[d]_0 & C_1 \ar[d]_{g_1} \ar[r] & E[1] \ar[d]_0 \\
E \ar[r] & C_1 \ar[r] & E[1]
}
$$
est commutatif, car les composés $E \xrightarrow{g_1} E \to C_1$ et
$C_1 \to E[1] \xrightarrow{g_1} E[1]$ sont nuls. En poursuivant, et en prenant
pour $C_{i + 1}$ le cône du morphisme $g_i : C_i \to C_i$, on obtient
un complexe parfait $C_s$ sur $X$ à support dans $T$
dont la restriction à $U$ est
$$
F \oplus F[1]^{\oplus s} \oplus F[2]^{\oplus {s \choose 2}}
\oplus \ldots \oplus F[s]
$$
Choisissons des morphismes de complexes parfaits $\beta : C' \to C_s$
et $\gamma : C' \to C_s$ comme dans le lemme \ref{lemma-lift-map},
tels que $\beta|_U$ soit un isomorphisme et que
$\gamma|_U \circ \beta|_U^{-1}$ soit le morphisme
$$
F \oplus F[1]^{\oplus s} \oplus F[2]^{\oplus {s \choose 2}}
\oplus \ldots \oplus F[s]
\to
F \oplus F[1]^{\oplus s} \oplus F[2]^{\oplus {s \choose 2}}
\oplus \ldots \oplus F[s]
$$
qui est l'identité sur tous les facteurs directs sauf sur $F$, où il est nul.
D'après le lemme \ref{lemma-lift-map}, on a aussi
$C' = C_s \otimes^\mathbf{L} I$ pour un complexe parfait
$I$ sur $X$. La nullité de $g_j^2\text{id}_{C_s}$ entraîne donc la
même propriété pour $C'$. Ainsi $C'$ est également à support dans $T$.
Alors $\text{Cone}(\gamma)$ convient.
\end{proof}

\noindent
On trouvera un cas particulier du lemme suivant dans
\cite{Neeman-Grothendieck}.

\begin{lemma}
\label{lemma-lift-map-from-perfect-complex-with-support}
Soit $X$ un schéma quasi-compact et quasi-séparé.
Soit $U \subset X$ un ouvert quasi-compact. Soit $T \subset X$
un sous-ensemble fermé tel que $X \setminus T$ soit rétrocompact dans $X$.
Soit $E$ un objet de $D_\QCoh(\mathcal{O}_X)$.
Soit $\alpha : P \to E|_U$ un morphisme, où $P$ est un objet parfait de
$D(\mathcal{O}_U)$ à support dans $T \cap U$. Il existe alors un morphisme
$\beta : R \to E$, où $R$ est un objet parfait de $D(\mathcal{O}_X)$
à support dans $T$, tel que $P$ soit facteur direct de $R|_U$ dans
$D(\mathcal{O}_U)$, de façon compatible avec $\alpha$ et $\beta|_U$.
\end{lemma}

\begin{proof}
Puisque $X$ est quasi-compact, il existe un entier $m$ tel que
$X = U \cup V_1 \cup \ldots \cup V_m$ pour des ouverts affines $V_j$ de $X$.
Par récurrence sur $m$, on voit que l'on peut supposer $m = 1$. Autrement
dit, on peut supposer que $X = U \cup V$, avec $V$ affine. D'après le
lemme \ref{lemma-lift-perfect-complex-plus-shift-support},
on peut choisir un objet parfait $Q$ de $D(\mathcal{O}_V)$
à support dans $T \cap V$ et un isomorphisme
$Q|_{U \cap V} \to (P \oplus P[1])|_{U \cap V}$.
D'après le lemme \ref{lemma-lift-map}, on peut remplacer $Q$ par
$Q \otimes^\mathbf{L} I$ (toujours à support dans $T \cap V$)
et supposer que le morphisme
$$
Q|_{U \cap V} \to (P \oplus P[1])|_{U \cap V}
\longrightarrow P|_{U \cap V}
\longrightarrow
E|_{U \cap V}
$$
se relève en $Q \to E|_V$. D'après
Cohomologie, lemme \ref{cohomology-lemma-glue},
on trouve un morphisme $a : R \to E$ de $D(\mathcal{O}_X)$
tel que $a|_U$ soit isomorphe à $P \oplus P[1] \to E|_U$
et $a|_V$ soit isomorphe à $Q \to E|_V$.
Ainsi $R$ est parfait et à support dans $T$, comme voulu.
\end{proof}

\begin{remark}
\label{remark-addendum}
La démonstration du lemme \ref{lemma-lift-map-from-perfect-complex-with-support}
montre que
$$
R|_U = P \oplus P^{\oplus n_1}[1] \oplus \ldots \oplus P^{\oplus n_m}[m]
$$
pour des entiers $m \geq 0$ et $n_j \geq 0$. Ainsi le faisceau de cohomologie
de plus haut degré de $R|_U$ est celui de $P$. En répétant la construction pour le morphisme
$P^{\oplus n_1}[1] \oplus \ldots \oplus P^{\oplus n_m}[m] \to R|_U$, en prenant
des cônes et en procédant par récurrence, on peut obtenir l'égalité des faisceaux de cohomologie
de $R|_U$ et de $P$ au-dessus de tout degré donné.
\end{remark}



\section{Approximation par des complexes parfaits}
\label{section-approximation}

\noindent
Dans cette section, nous étudions l'observation, due à Neeman et Lipman,
selon laquelle un complexe pseudo-cohérent peut être « approché » par des complexes parfaits.

\begin{definition}
\label{definition-approximation-holds}
Soit $X$ un schéma. Considérons les triplets $(T, E, m)$ où
\begin{enumerate}
\item $T \subset X$ est un sous-ensemble fermé,
\item $E$ est un objet de $D_\QCoh(\mathcal{O}_X)$, et
\item $m \in \mathbf{Z}$.
\end{enumerate}
On dit que {\it l'approximation est possible pour le triplet} $(T, E, m)$ s'il
existe un objet parfait $P$ de $D(\mathcal{O}_X)$ à support dans $T$
et un morphisme $\alpha : P \to E$ qui induit des isomorphismes $H^i(P) \to H^i(E)$
pour $i > m$ et une surjection $H^m(P) \to H^m(E)$.
\end{definition}

\noindent
L'approximation ne peut pas être possible pour tout triplet. En effet, il est clair que,
si l'approximation est possible pour le triplet $(T, E, m)$, alors
\begin{enumerate}
\item $E$ est $m$-pseudo-cohérent, voir
Cohomologie, définition \ref{cohomology-definition-pseudo-coherent}, et
\item les faisceaux de cohomologie $H^i(E)$ sont à support dans $T$ pour $i \geq m$.
\end{enumerate}
De plus, le « support » d'un complexe parfait est un sous-schéma fermé
dont le complémentaire est rétrocompact dans $X$ (détails omis). On ne peut donc
espérer que l'approximation soit possible sans cette hypothèse sur $T$.
Ceci explique en partie les conditions de la définition suivante.

\begin{definition}
\label{definition-approximation}
Soit $X$ un schéma. On dit que {\it l'approximation par des complexes parfaits est possible}
sur $X$ si, pour tout sous-ensemble fermé $T \subset X$ tel que $X \setminus T$
soit rétrocompact dans $X$, il existe un entier $r$ tel que,
pour tout triplet $(T, E, m)$ comme dans la
définition \ref{definition-approximation-holds} vérifiant
\begin{enumerate}
\item $E$ est $(m - r)$-pseudo-cohérent, et
\item $H^i(E)$ est à support dans $T$ pour $i \geq m - r$,
\end{enumerate}
l'approximation soit possible.
\end{definition}

\noindent
Nous démontrerons que l'approximation par des complexes parfaits est possible pour les
schémas quasi-compacts et quasi-séparés. Il semble que la deuxième
condition soit nécessaire à notre méthode de démonstration. Il est possible que la
première condition puisse être affaiblie en « $E$ est $m$-pseudo-cohérent »
par une analyse soigneuse des arguments ci-dessous.

\begin{lemma}
\label{lemma-open}
Soit $X$ un schéma. Soit $U \subset X$ un sous-schéma ouvert.
Soit $(T, E, m)$ un triplet comme dans la
définition \ref{definition-approximation-holds}.
Si
\begin{enumerate}
\item $T \subset U$,
\item l'approximation est possible pour $(T, E|_U, m)$, et
\item les faisceaux $H^i(E)$ pour $i \geq m$ sont à support dans $T$,
\end{enumerate}
alors l'approximation est possible pour $(T, E, m)$.
\end{lemma}

\begin{proof}
Soit $j : U \to X$ le morphisme d'inclusion.
Si $P \to E|_U$ est une approximation du triplet $(T, E|_U, m)$
sur $U$, alors $j_!P = Rj_*P \to j_!(E|_U) \to E$ est une approximation
de $(T, E, m)$ sur $X$.
Voir Cohomologie, lemmes \ref{cohomology-lemma-pushforward-restriction} et
\ref{cohomology-lemma-pushforward-perfect}.
\end{proof}

\begin{lemma}
\label{lemma-approximation-affine}
Soit $X$ un schéma affine. Alors l'approximation est possible pour tout
triplet $(T, E, m)$ comme dans la définition \ref{definition-approximation-holds}
tel qu'il existe un entier $r \geq 0$ vérifiant
\begin{enumerate}
\item $E$ est $m$-pseudo-cohérent,
\item $H^i(E)$ est à support dans $T$ pour $i \geq m - r + 1$,
\item $X \setminus T$ est la réunion de $r$ ouverts affines.
\end{enumerate}
En particulier, l'approximation par des complexes parfaits est possible pour les schémas affines.
\end{lemma}

\begin{proof}
Écrivons $X = \Spec(A)$. Écrivons $T = V(f_1, \ldots, f_r)$.
(Le cas $r = 0$, c'est-à-dire $T = X$, résulte immédiatement du
lemme \ref{lemma-pseudo-coherent-affine} et des définitions.)
Soit $(T, E, m)$ un triplet comme dans le lemme.
Soit $t$ le plus grand entier tel que $H^t(E)$ soit non nul.
Nous procéderons par récurrence sur $t$. Le cas initial est $t < m$ ;
le résultat est alors trivial. Supposons maintenant que $t \geq m$. D'après
Cohomologie, lemme \ref{cohomology-lemma-finite-cohomology},
le faisceau $H^t(E)$ est de type fini. Comme il est quasi-cohérent,
il est engendré par un nombre fini de sections
(Propriétés, lemme \ref{properties-lemma-finite-type-module}).
Pour tout $s \in \Gamma(X, H^t(E)) = H^t(X, E)$
(voir la démonstration du lemme \ref{lemma-affine-compare-bounded}),
on peut trouver un $e > 0$ et un morphisme $K_e[-t] \to E$
tels que $s$ appartienne à l'image de
$H^0(K_e) = H^t(K_e[-t]) \to H^t(E)$, voir le
lemme \ref{lemma-represent-cohomology-class-on-closed}.
Une somme directe finie de ces morphismes donne un morphisme
$P \to E$, où $P$ est un complexe parfait à support dans $T$,
où $H^i(P) = 0$ pour $i > t$, et où $H^t(P) \to E$ est
surjectif. Choisissons un triangle distingué
$$
P \to E \to E' \to P[1]
$$
Alors $E'$ est $m$-pseudo-cohérent
(Cohomologie, lemme \ref{cohomology-lemma-cone-pseudo-coherent}),
$H^i(E') = 0$ pour $i \geq t$, et
$H^i(E')$ est à support dans $T$ pour $i \geq m - r + 1$.
Par récurrence, on trouve une approximation $P' \to E'$
de $(T, E', m)$. Insérons le composé $P' \to E' \to P[1]$
dans un triangle distingué $P \to P'' \to P' \to P[1]$
et prolongeons les morphismes $P' \to E'$ et $P[1] \to P[1]$ en
un morphisme de triangles distingués
$$
\xymatrix{
P \ar[r] \ar[d] & P'' \ar[d] \ar[r] & P' \ar[d] \ar[r] & P[1] \ar[d] \\
P \ar[r] &  E \ar[r] & E' \ar[r] & P[1]
}
$$
à l'aide de TR3. Alors $P''$ est un complexe parfait
(Cohomologie, lemme \ref{cohomology-lemma-two-out-of-three-perfect})
à support dans $T$.
Une chasse au diagramme immédiate montre que $P'' \to E$ est l'approximation
cherchée.
\end{proof}

\begin{lemma}
\label{lemma-induction-step}
Soit $X$ un schéma. Soit $X = U \cup V$ un recouvrement ouvert,
avec $U$ quasi-compact, $V$ affine et $U \cap V$ quasi-compact.
Si l'approximation par des complexes parfaits est possible sur $U$,
elle est possible sur $X$.
\end{lemma}

\begin{proof}
Soit $T \subset X$ un sous-ensemble fermé tel que $X \setminus T$ soit rétrocompact
dans $X$. Soit $r_U$ l'entier de la définition \ref{definition-approximation}
adapté au couple $(U, T \cap U)$.
Posons $T' = T \setminus U$. Notons que
$T' \subset V$ et que $V \setminus T' = (X \setminus T) \cap U \cap V$
est quasi-compact par notre hypothèse sur $T$.
Soit $r'$ le nombre d'ouverts affines nécessaires pour recouvrir $V \setminus T'$.
Nous affirmons que $r = \max(r_U, r')$ convient pour le couple $(X, T)$.

\medskip\noindent
Pour le voir, choisissons un triplet $(T, E, m)$ tel que $E$ soit
$(m - r)$-pseudo-cohérent et que $H^i(E)$ soit à support dans $T$ pour
$i \geq m - r$. Soit $t$ le plus grand entier tel que
$H^t(E)|_U$ soit non nul. (Un tel entier existe, car $U$ est quasi-compact
et $E|_U$ est $(m - r)$-pseudo-cohérent.)
Nous démontrerons que $E$ peut être approché en procédant par récurrence sur $t$.

\medskip\noindent
Cas initial : $t \leq m - r'$. Cela signifie que $H^i(E)$ est à support
dans $T'$ pour $i \geq m - r'$. Le
lemme \ref{lemma-approximation-affine}
assure donc l'existence d'une approximation
$P \to E|_V$ de $(T', E|_V, m)$ sur $V$.
En appliquant le lemme \ref{lemma-open}, on voit que
$(T', E, m)$ peut être approché. Une telle approximation
est aussi une approximation de $(T, E, m)$.

\medskip\noindent
Étape de récurrence. Choisissons une approximation $P \to E|_U$
de $(T \cap U, E|_U, m)$. Elle donne en particulier une surjection
$H^t(P) \to H^t(E|_U)$. D'après le
lemme \ref{lemma-lift-perfect-complex-plus-shift-support},
on peut choisir un objet parfait $Q$ de $D(\mathcal{O}_V)$
à support dans $T \cap V$ et un isomorphisme
$Q|_{U \cap V} \to (P \oplus P[1])|_{U \cap V}$.
D'après le lemme \ref{lemma-lift-map}, on peut remplacer $Q$ par
$Q \otimes^\mathbf{L} I$
et supposer que le morphisme
$$
Q|_{U \cap V} \to (P \oplus P[1])|_{U \cap V}
\longrightarrow P|_{U \cap V}
\longrightarrow
E|_{U \cap V}
$$
se relève en $Q \to E|_V$. D'après
Cohomologie, lemme \ref{cohomology-lemma-glue},
on trouve un morphisme $a : R \to E$ de $D(\mathcal{O}_X)$
tel que $a|_U$ soit isomorphe à $P \oplus P[1] \to E|_U$
et $a|_V$ soit isomorphe à $Q \to E|_V$.
Ainsi $R$ est parfait et à support dans $T$,
et le morphisme $H^t(R) \to H^t(E)$ est surjectif après restriction à $U$.
Choisissons un triangle distingué
$$
R \to E \to E' \to R[1]
$$
Alors $E'$ est $(m - r)$-pseudo-cohérent
(Cohomologie, lemme \ref{cohomology-lemma-cone-pseudo-coherent}),
$H^i(E')|_U = 0$ pour $i \geq t$, et
$H^i(E')$ est à support dans $T$ pour $i \geq m - r$.
Par récurrence, on trouve une approximation $R' \to E'$
de $(T, E', m)$. Insérons le composé $R' \to E' \to R[1]$
dans un triangle distingué $R \to R'' \to R' \to R[1]$
et prolongeons les morphismes $R' \to E'$ et $R[1] \to R[1]$ en
un morphisme de triangles distingués
$$
\xymatrix{
R \ar[r] \ar[d] & R'' \ar[d] \ar[r] & R' \ar[d] \ar[r] & R[1] \ar[d] \\
R \ar[r] &  E \ar[r] & E' \ar[r] & R[1]
}
$$
à l'aide de TR3. Alors $R''$ est un complexe parfait
(Cohomologie, lemme \ref{cohomology-lemma-two-out-of-three-perfect})
à support dans $T$.
Une chasse au diagramme immédiate montre que $R'' \to E$ est l'approximation
cherchée.
\end{proof}

\begin{theorem}
\label{theorem-approximation}
Soit $X$ un schéma quasi-compact et quasi-séparé.
Alors l'approximation par des complexes parfaits est possible sur $X$.
\end{theorem}

\begin{proof}
Ceci résulte du principe de récurrence de
Cohomologie des schémas, lemme \ref{coherent-lemma-induction-principle},
et des lemmes \ref{lemma-induction-step} et \ref{lemma-approximation-affine}.
\end{proof}





\section{Générateurs des catégories dérivées}
\label{section-generating}

\noindent
Dans cette section, nous démontrons que la catégorie dérivée
$D_\QCoh(\mathcal{O}_X)$ d'un schéma quasi-compact
et quasi-séparé peut être engendrée par un seul objet parfait.
Nous invitons vivement le lecteur à lire la démonstration de ce résultat dans le remarquable article de
Bondal et van den Bergh, voir \cite{BvdB}.

\begin{lemma}
\label{lemma-direct-summand-of-a-restriction}
Soit $X$ un schéma quasi-compact et quasi-séparé.
Soit $U$ un sous-schéma ouvert quasi-compact.
Soit $P$ un objet parfait de $D(\mathcal{O}_U)$.
Alors $P$ est facteur direct de la restriction d'un objet parfait
de $D(\mathcal{O}_X)$.
\end{lemma}

\begin{proof}
C'est un cas particulier du lemme \ref{lemma-lift-map-from-perfect-complex-with-support}.
\end{proof}

\begin{lemma}
\label{lemma-orthogonal-koszul-complex}
\begin{reference}
\cite[proposition 6.1]{Bokstedt-Neeman}
\end{reference}
Dans la situation \ref{situation-complex}, notons $j : U \to X$ l'immersion
ouverte et soit $K$ l'objet parfait de $D(\mathcal{O}_X)$
correspondant au complexe de Koszul de $f_1, \ldots, f_r$ sur $A$.
Pour $E \in D_\QCoh(\mathcal{O}_X)$, les conditions suivantes sont équivalentes :
\begin{enumerate}
\item $E = Rj_*(E|_U)$, et
\item $\Hom_{D(\mathcal{O}_X)}(K[n], E) = 0$ pour tout $n \in \mathbf{Z}$.
\end{enumerate}
\end{lemma}

\begin{proof}
Choisissons un triangle distingué $E \to Rj_*(E|_U) \to N \to E[1]$.
Observons que
$$
\Hom_{D(\mathcal{O}_X)}(K[n], Rj_*(E|_U)) =
\Hom_{D(\mathcal{O}_U)}(K|_U[n], E) = 0
$$
pour tout $n$, car $K|_U = 0$. Il suffit donc de démontrer le résultat pour
$N$. Autrement dit, on peut supposer que $E$ se restreint à zéro sur $U$.
Observons qu'il existe des triangles distingués
$$
\begin{aligned}
K^\bullet(f_1^{e_1}, \ldots, f_i^{e'_i}, \ldots, f_r^{e_r}) &\to
K^\bullet(f_1^{e_1}, \ldots, f_i^{e'_i + e''_i}, \ldots, f_r^{e_r}) \\
{}&\to
K^\bullet(f_1^{e_1}, \ldots, f_i^{e''_i}, \ldots, f_r^{e_r}) \to \ldots
\end{aligned}
$$
de complexes de Koszul, voir
Compléments sur l'algèbre, lemme \ref{more-algebra-lemma-koszul-mult}.
Ainsi, si $\Hom_{D(\mathcal{O}_X)}(K[n], E) = 0$ pour tout $n \in \mathbf{Z}$,
la même propriété vaut en remplaçant $K$ par
$K_e$ comme dans le lemme \ref{lemma-represent-cohomology-class-on-closed}.
Notre lemme résulte donc immédiatement de ce dernier et du fait que $E$
est déterminé par le complexe de $A$-modules $R\Gamma(X, E)$, voir le
lemme \ref{lemma-affine-compare-bounded}.
\end{proof}

\begin{theorem}
\label{theorem-bondal-van-den-Bergh}
Soit $X$ un schéma quasi-compact et quasi-séparé. La catégorie
$D_\QCoh(\mathcal{O}_X)$ peut être engendrée par un seul
objet parfait. Plus précisément, il existe un objet parfait
$P$ de $D(\mathcal{O}_X)$ tel que, pour
$E \in D_\QCoh(\mathcal{O}_X)$, les conditions suivantes soient équivalentes :
\begin{enumerate}
\item $E = 0$, et
\item $\Hom_{D(\mathcal{O}_X)}(P[n], E) = 0$ pour tout $n \in \mathbf{Z}$.
\end{enumerate}
\end{theorem}

\begin{proof}
Nous utiliserons le principe de récurrence de
Cohomologie des schémas, lemme \ref{coherent-lemma-induction-principle}.

\medskip\noindent
Si $X$ est affine, alors $\mathcal{O}_X$ est un générateur parfait.
Ceci résulte du lemme \ref{lemma-affine-compare-bounded}.

\medskip\noindent
Supposons que $X = U \cup V$ soit un recouvrement ouvert, avec $U$ quasi-compact,
tel que le théorème soit vrai pour $U$, et que $V$ soit un ouvert affine.
Soit $P$ un objet parfait de $D(\mathcal{O}_U)$ qui engendre
$D_\QCoh(\mathcal{O}_U)$. À l'aide du
lemme \ref{lemma-direct-summand-of-a-restriction}, on peut
choisir un objet parfait
$Q$ de $D(\mathcal{O}_X)$ dont la restriction à $U$ est une somme directe dont
l'un des facteurs est $P$. Écrivons $V = \Spec(A)$. Soit $Z = X \setminus U$.
C'est un sous-ensemble fermé de $V$ tel que $V \setminus Z$ soit quasi-compact.
Choisissons $f_1, \ldots, f_r \in A$ tels que
$Z = V(f_1, \ldots, f_r)$. Soit $K \in D(\mathcal{O}_V)$ l'objet parfait
correspondant au complexe de Koszul de $f_1, \ldots, f_r$ sur $A$.
Comme $K$ est à support dans le fermé $Z \subset V$, l'image directe
$K' = R(V \to X)_*K$ est un objet parfait de $D(\mathcal{O}_X)$ dont la
restriction à $V$ est $K$ (voir
Cohomologie, lemme \ref{cohomology-lemma-pushforward-perfect}).
Nous affirmons que $Q \oplus K'$ engendre
$D_\QCoh(\mathcal{O}_X)$.

\medskip\noindent
Soit $E$ un objet de $D_\QCoh(\mathcal{O}_X)$ tel qu'il
n'existe aucun morphisme non nul d'un décalé quelconque de $Q \oplus K'$ vers $E$.
D'après Cohomologie, lemme \ref{cohomology-lemma-pushforward-restriction},
on a $K' =  R(V \to X)_! K$, et donc
$$
\Hom_{D(\mathcal{O}_X)}(K'[n], E) = \Hom_{D(\mathcal{O}_V)}(K[n], E|_V)
$$
D'après le lemme \ref{lemma-orthogonal-koszul-complex}, l'annulation de
ces groupes implique que $E|_V$ est isomorphe à
$R(U \cap V \to V)_*E|_{U \cap V}$. Il en résulte que $E = R(U \to X)_*E|_U$
(un petit détail est omis). Dans ce cas,
$$
\Hom_{D(\mathcal{O}_X)}(Q[n], E) = \Hom_{D(\mathcal{O}_U)}(Q|_U[n], E|_U)
$$
contient $\Hom_{D(\mathcal{O}_U)}(P[n], E|_U)$ comme facteur direct.
Par le choix de $P$, l'annulation de ces groupes implique donc que $E|_U$
est nul. Par suite, $E$ est nul.
\end{proof}

\noindent
Le résultat suivant renforce le
théorème \ref{theorem-bondal-van-den-Bergh}
et se démontre par exactement les mêmes méthodes.
Rappelons que, pour un sous-ensemble fermé $T$ d'un schéma $X$, on note
$D_T(\mathcal{O}_X)$ la sous-catégorie strictement pleine, saturée
et triangulée de $D(\mathcal{O}_X)$ formée des objets
à support dans $T$ (définition \ref{definition-supported-on}).
De même, on note $D_{\QCoh, T}(\mathcal{O}_X)$ la sous-catégorie strictement pleine, saturée
et triangulée de $D(\mathcal{O}_X)$ formée des
complexes dont les faisceaux de cohomologie sont quasi-cohérents et à support
dans $T$.

\begin{lemma}
\label{lemma-generator-with-support}
\begin{reference}
\cite[théorème 6.8]{Rouquier-dimensions}
\end{reference}
Soit $X$ un schéma quasi-compact et quasi-séparé. Soit $T \subset X$ un
sous-ensemble fermé tel que $X \setminus T$ soit quasi-compact. Avec les notations
ci-dessus, la catégorie $D_{\QCoh, T}(\mathcal{O}_X)$ est engendrée par un
seul objet parfait.
\end{lemma}

\begin{proof}
Nous utiliserons le principe de récurrence de
Cohomologie des schémas, lemme \ref{coherent-lemma-induction-principle}.

\medskip\noindent
Supposons que $X = \Spec(A)$ soit affine. Il existe alors
$f_1, \ldots, f_r \in A$ tels que $T = V(f_1, \ldots, f_r)$.
Soit $K$ le complexe de Koszul de $f_1, \ldots, f_r$ comme dans le
lemme \ref{lemma-orthogonal-koszul-complex}.
Alors $K$ est un objet parfait dont la cohomologie est à support dans
$T$, donc un objet parfait de $D_{\QCoh, T}(\mathcal{O}_X)$.
D'autre part, si $E \in D_{\QCoh, T}(\mathcal{O}_X)$ et si
$\Hom(K, E[n]) = 0$ pour tout $n$, le
lemme \ref{lemma-orthogonal-koszul-complex}
donne $E = Rj_*(E|_{X \setminus T}) = 0$.
Ainsi $K$ engendre $D_{\QCoh, T}(\mathcal{O}_X)$
(selon notre définition des générateurs de catégories triangulées dans
Catégories dérivées, définition \ref{derived-definition-generators}).

\medskip\noindent
Supposons que $X = U \cup V$ soit un recouvrement ouvert, avec $V$ affine et
$U$ quasi-compact, tel que le lemme soit vrai pour $U$.
Soit $P$ un objet parfait de $D(\mathcal{O}_U)$ à support dans $T \cap U$
qui engendre $D_{\QCoh, T \cap U}(\mathcal{O}_U)$. À l'aide du
lemme \ref{lemma-lift-map-from-perfect-complex-with-support},
on peut choisir un objet parfait $Q$ de $D(\mathcal{O}_X)$ à support dans $T$
dont la restriction à $U$ est une somme directe dont l'un des facteurs est $P$.
Écrivons $V = \Spec(B)$. Soit $Z = X \setminus U$. Alors $Z$ est un sous-ensemble fermé
de $V$ tel que $V \setminus Z$ soit quasi-compact. Comme $X$ est quasi-séparé,
il en résulte que $Z \cap T$ est un sous-ensemble fermé de $V$ tel que
$W = V \setminus (Z \cap T)$ soit quasi-compact. On peut donc choisir
$g_1, \ldots, g_s \in B$ tels que $Z \cap T = V(g_1, \ldots, g_r)$.
Soit $K \in D(\mathcal{O}_V)$ l'objet parfait correspondant au
complexe de Koszul de $g_1, \ldots, g_s$ sur $B$. Comme $K$ est
à support dans le fermé $(Z \cap T) \subset V$, l'image directe
$K' = R(V \to X)_*K$ est un objet parfait de $D(\mathcal{O}_X)$ dont la
restriction à $V$ est $K$ (voir
Cohomologie, lemme \ref{cohomology-lemma-pushforward-perfect}).
Nous affirmons que $Q \oplus K'$ engendre
$D_{\QCoh, T}(\mathcal{O}_X)$.

\medskip\noindent
Soit $E$ un objet de $D_{\QCoh, T}(\mathcal{O}_X)$ tel qu'il
n'existe aucun morphisme non nul d'un décalé quelconque de $Q \oplus K'$ vers $E$.
D'après Cohomologie, lemme \ref{cohomology-lemma-pushforward-restriction},
on a $K' =  R(V \to X)_! K$, et donc
$$
\Hom_{D(\mathcal{O}_X)}(K'[n], E) = \Hom_{D(\mathcal{O}_V)}(K[n], E|_V)
$$
D'après le lemme \ref{lemma-orthogonal-koszul-complex}, on a donc
$E|_V = Rj_*E|_W$, où $j : W \to V$ est l'inclusion. On a le diagramme
$$
\xymatrix{
W \ar[r]_j & V & Z \cap T \ar[l] \ar[d] \\
U \cap V \ar[u]^{j'} \ar[ru]_{j''} & & Z \ar[lu]
}
$$
Puisque $E$ est à support dans $T$, $E|_W$ est à support dans
$T \cap W = T \cap U \cap V$, qui est fermé dans $W$.
On en déduit que
$$
E|_V = Rj_*(E|_W) = Rj_*(Rj'_*(E|_{U \cap V})) = Rj''_*(E|_{U \cap V})
$$
où la deuxième égalité résulte de l'assertion (1) de
Cohomologie, lemme \ref{cohomology-lemma-pushforward-restriction}.
Il en résulte que $E = R(U \to X)_*E|_U$ (un petit détail est omis).
Dans ce cas,
$$
\Hom_{D(\mathcal{O}_X)}(Q[n], E) = \Hom_{D(\mathcal{O}_U)}(Q|_U[n], E|_U)
$$
contient $\Hom_{D(\mathcal{O}_U)}(P[n], E|_U)$ comme facteur direct.
Par le choix de $P$, l'annulation de ces groupes implique donc que $E|_U$
est nul. Par suite, $E$ est nul.
\end{proof}

\section{Un exemple de générateur}
\label{section-example-generator}

\noindent
Dans cette section, nous démontrons que la catégorie dérivée de l'espace projectif
sur un anneau est engendrée par un fibré vectoriel, et même par une somme
directe de tordus du faisceau structural.

\medskip\noindent
Le lemme suivant affirme que $\bigoplus_{n \geq 0} \mathcal{L}^{\otimes -n}$
est un générateur si $\mathcal{L}$ est ample.

\begin{lemma}
\label{lemma-nonzero-some-cohomology}
Soient $X$ un schéma et $\mathcal{L}$ un
$\mathcal{O}_X$-module inversible ample. Si $K$ est un objet non nul de
$D_\QCoh(\mathcal{O}_X)$, alors, pour un certain $n \geq 0$ et un certain $p \in \mathbf{Z}$,
le groupe de cohomologie
$H^p(X, K \otimes_{\mathcal{O}_X}^\mathbf{L} \mathcal{L}^{\otimes n})$
est non nul.
\end{lemma}

\begin{proof}
Rappelons que, puisque $X$ possède un faisceau inversible ample, il est quasi-compact
et séparé (Propriétés, définition \ref{properties-definition-ample} et
lemme \ref{properties-lemma-affine-s-opens-cover-quasi-separated}).
On peut donc appliquer la
proposition \ref{proposition-quasi-compact-affine-diagonal}
et représenter $K$ par un complexe $\mathcal{F}^\bullet$ de
modules quasi-cohérents. Choisissons un $p$ tel que
$\mathcal{H}^p = \Ker(\mathcal{F}^p \to \mathcal{F}^{p + 1})/
\Im(\mathcal{F}^{p - 1} \to \mathcal{F}^p)$ soit non nul.
Choisissons un point $x \in X$ tel que la fibre $\mathcal{H}^p_x$ soit
non nulle. Choisissons un $n \geq 0$ et un $s \in \Gamma(X, \mathcal{L}^{\otimes n})$
tels que $X_s$ soit un voisinage ouvert affine de $x$.
Choisissons $\tau \in \mathcal{H}^p(X_s)$ dont l'image soit un élément non nul
de la fibre $\mathcal{H}^p_x$ ; c'est possible,
car $\mathcal{H}^p$ est quasi-cohérent et $X_s$ est affine.
Puisque le foncteur des sections sur $X_s$ est exact sur les
modules quasi-cohérents, on peut trouver une section $\tau' \in \mathcal{F}^p(X_s)$
d'image nulle dans $\mathcal{F}^{p + 1}(X_s)$ et d'image
$\tau$ dans $\mathcal{H}^p(X_s)$. D'après
Propriétés, lemme \ref{properties-lemma-invert-s-sections},
il existe un $m$ tel que $\tau' \otimes s^{\otimes m}$
soit l'image d'une section
$\tau'' \in \Gamma(X, \mathcal{F}^p \otimes \mathcal{L}^{\otimes mn})$.
En appliquant encore une fois le même lemme, on trouve $l \geq 0$ tel que
$\tau'' \otimes s^{\otimes l}$ ait une image nulle dans
$\mathcal{F}^{p + 1} \otimes \mathcal{L}^{\otimes (m + l)n}$.
Alors $\tau''$ donne une classe non nulle dans
$H^p(X, K \otimes^\mathbf{L}_{\mathcal{O}_X} \mathcal{L}^{(m + l)n})$,
comme voulu.
\end{proof}

\begin{lemma}
\label{lemma-construct-the-next-one}
Soit $A$ un anneau. Soit $X = \mathbf{P}^n_A$. Pour tout $a \in \mathbf{Z}$,
il existe un complexe exact
$$
0 \to \mathcal{O}_X(a) \to \ldots
\to \mathcal{O}_X(a + i)^{\oplus {n + 1 \choose i}} \to
\ldots \to \mathcal{O}_X(a + n + 1) \to 0
$$
de fibrés vectoriels sur $X$.
\end{lemma}

\begin{proof}
Rappelons que $\mathbf{P}^n_A$ est $\text{Proj}(A[X_0, \ldots, X_n])$, voir
Constructions, définition \ref{constructions-definition-projective-space}.
Considérons le complexe de Koszul
$$
K_\bullet = K_\bullet(A[X_0, \ldots, X_n], X_0, \ldots, X_n)
$$
sur $S = A[X_0, \ldots, X_n]$ associé à $X_0, \ldots, X_n$.
Puisque $X_0, \ldots, X_n$ est manifestement une suite régulière dans
l'anneau de polynômes $S$, on voit
(Compléments sur l'algèbre, lemme \ref{more-algebra-lemma-regular-koszul-regular})
que le complexe de Koszul $K_\bullet$ est exact, sauf en degré $0$,
où la cohomologie est $S/(X_0, \ldots, X_n)$.
Notons que $K_\bullet$ devient un complexe de modules gradués si l'on
place les générateurs de $K_i$ en degré $+i$. On obtient donc un
complexe exact
$$
0 \to S(-n - 1) \to \ldots \to S(-n - 1 + i)^{\oplus {n \choose i}} \to \ldots
\to S \to S/(X_0, \ldots, X_n) \to 0
$$
En appliquant le foncteur exact $\tilde{\ }$ de Constructions,
lemme \ref{constructions-lemma-proj-sheaves}, et en utilisant le fait que
le dernier terme appartient au noyau de ce foncteur,
on obtient le complexe exact
$$
0 \to \mathcal{O}_X(-n - 1) \to \ldots
\to \mathcal{O}_X(-n - 1 + i)^{\oplus {n + 1 \choose i}} \to
\ldots \to \mathcal{O}_X \to 0
$$
En tensorisant par les faisceaux inversibles $\mathcal{O}_X(n + a + 1)$,
on obtient les complexes exacts du lemme.
\end{proof}

\begin{lemma}
\label{lemma-generator-P1}
Soit $A$ un anneau. Soit $X = \mathbf{P}^n_A$. Alors
$$
E =
\mathcal{O}_X \oplus \mathcal{O}_X(-1) \oplus \ldots \oplus \mathcal{O}_X(-n)
$$
est un générateur
(Catégories dérivées, définition \ref{derived-definition-generators})
de $D_\QCoh(X)$.
\end{lemma}

\begin{proof}
Soit $K \in D_\QCoh(\mathcal{O}_X)$. Supposons que
$\Hom(E, K[p]) = 0$ pour tout $p \in \mathbf{Z}$.
Nous devons montrer que $K = 0$.
D'après Catégories dérivées, lemme
\ref{derived-lemma-right-orthogonal},
$\Hom(E', K[p])$ est nul pour tout $E' \in \langle E \rangle$
et tout $p \in \mathbf{Z}$.
D'après le lemme \ref{lemma-construct-the-next-one}
appliqué avec $a = -n - 1$,
on a $\mathcal{O}_X(-n - 1) \in \langle E \rangle$,
car ce faisceau est quasi-isomorphe à un complexe borné
dont les termes sont des sommes directes finies de facteurs directs de $E$.
En répétant l'argument avec $a = -n - 2$, on obtient
$\mathcal{O}_X(-n - 2) \in \langle E \rangle$.
Par récurrence, on trouve $\mathcal{O}_X(-m) \in \langle E \rangle$
pour tout $m \geq 0$.
Puisque
$$
\Hom(\mathcal{O}_X(-m), K[p]) =
H^p(X, K \otimes_{\mathcal{O}_X}^\mathbf{L} \mathcal{O}_X(m)) =
H^p(X, K \otimes_{\mathcal{O}_X}^\mathbf{L} \mathcal{O}_X(1)^{\otimes m})
$$
on en déduit que $K = 0$ d'après le lemme \ref{lemma-nonzero-some-cohomology}.
(On utilise aussi le fait que $\mathcal{O}_X(1)$ est un
faisceau inversible ample sur $X$, ce qui résulte de
Propriétés, lemme \ref{properties-lemma-open-in-proj-ample}.)
\end{proof}

\begin{remark}
\label{remark-pullback-generator}
Soit $f : X \to Y$ un morphisme de schémas quasi-compacts et quasi-séparés.
Soit $E \in D_\QCoh(\mathcal{O}_Y)$ un générateur
(voir le théorème \ref{theorem-bondal-van-den-Bergh}).
Alors les conditions suivantes sont équivalentes :
\begin{enumerate}
\item pour $K \in D_\QCoh(\mathcal{O}_X)$, on a
$Rf_*K = 0$ si et seulement si $K = 0$,
\item $Rf_* : D_\QCoh(\mathcal{O}_X) \to D_\QCoh(\mathcal{O}_Y)$
est conservatif, et
\item $Lf^*E$ engendre $D_\QCoh(\mathcal{O}_X)$.
\end{enumerate}
L'équivalence de (1) et (2) est une conséquence formelle du fait que
$Rf_* : D_\QCoh(\mathcal{O}_X) \to D_\QCoh(\mathcal{O}_Y)$ est un
foncteur exact de catégories triangulées. De même, l'équivalence
de (1) et (3) résulte formellement du fait que $Lf^*$
est adjoint à gauche de $Rf_*$.
Ces conditions sont vérifiées si $f$ est affine (lemme \ref{lemma-affine-morphism}),
ou si $f$ est une immersion ouverte, ou si $f$ est un composé de tels morphismes.
On en déduit que
\begin{enumerate}
\item si $X$ est un schéma quasi-affine, alors $\mathcal{O}_X$ engendre
$D_\QCoh(\mathcal{O}_X)$,
\item si $X \subset \mathbf{P}^n_A$ est un
sous-schéma localement fermé quasi-compact, alors
$\mathcal{O}_X \oplus \mathcal{O}_X(-1) \oplus \ldots \oplus \mathcal{O}_X(-n)$
engendre $D_\QCoh(\mathcal{O}_X)$ d'après le
lemme \ref{lemma-generator-P1}.
\end{enumerate}
\end{remark}




\section{Objets compacts et parfaits}
\label{section-compact}

\noindent
Soit $X$ un schéma noethérien de dimension finie. D'après
Cohomologie, proposition \ref{cohomology-proposition-vanishing-Noetherian},
et
Cohomologie sur les sites, lemme \ref{sites-cohomology-lemma-when-jshriek-compact},
les faisceaux de modules $j_!\mathcal{O}_U$ sont des objets compacts
de $D(\mathcal{O}_X)$ pour tout ouvert $U \subset X$.
Ces faisceaux ne sont en général pas quasi-cohérents et ne donnent donc pas
des objets parfaits de la catégorie dérivée $D(\mathcal{O}_X)$.
Cependant, si l'on se restreint aux complexes dont les faisceaux de cohomologie
sont quasi-cohérents, cette situation ne se présente pas.
Voici l'énoncé précis.

\begin{proposition}
\label{proposition-compact-is-perfect}
Soit $X$ un schéma quasi-compact et quasi-séparé.
Un objet de $D_\QCoh(\mathcal{O}_X)$ est compact
si et seulement s'il est parfait.
\end{proposition}

\begin{proof}
Si $K$ est un objet parfait de $D(\mathcal{O}_X)$ de dual
$K^\vee$ (Cohomologie, lemme \ref{cohomology-lemma-dual-perfect-complex}),
on a
$$
\Hom_{D(\mathcal{O}_X)}(K, M) =
H^0(X, K^\vee \otimes_{\mathcal{O}_X}^\mathbf{L} M)
$$
fonctoriellement en $M$. Puisque $K^\vee \otimes_{\mathcal{O}_X}^\mathbf{L} -$
commute aux sommes directes et que $H^0(X, -)$ commute aux sommes
directes sur $D_\QCoh(\mathcal{O}_X)$ d'après le
lemme \ref{lemma-quasi-coherence-pushforward-direct-sums},
on en déduit que $K$ est compact dans $D_\QCoh(\mathcal{O}_X)$.

\medskip\noindent
Réciproquement, soit $K$ un objet compact de $D_\QCoh(\mathcal{O}_X)$.
Pour montrer que $K$ est parfait, il suffit de montrer que
$K|_U$ est parfait pour tout ouvert affine $U \subset X$, voir
Cohomologie, lemme \ref{cohomology-lemma-perfect-independent-representative}.
Observons que $j : U \to X$ est un morphisme quasi-compact et séparé.
Par conséquent,
$Rj_* : D_\QCoh(\mathcal{O}_U) \to D_\QCoh(\mathcal{O}_X)$
commute aux sommes directes, voir le
lemme \ref{lemma-quasi-coherence-pushforward-direct-sums}.
L'adjonction entre la restriction à $U$ et $Rj_*$ implique donc que
$K|_U$ est un objet compact de $D_\QCoh(\mathcal{O}_U)$.
On se ramène ainsi au cas où $X$ est affine.

\medskip\noindent
Supposons que $X = \Spec(A)$ soit affine. D'après le lemme \ref{lemma-affine-compare-bounded},
le problème se traduit en la même question pour $D(A)$.
Pour $D(A)$, le résultat est
Compléments sur l'algèbre, proposition \ref{more-algebra-proposition-perfect-is-compact}.
\end{proof}

\begin{remark}
\label{remark-classical-generator}
Soit $X$ un schéma quasi-compact et quasi-séparé. Soit $G$ un
objet parfait de $D(\mathcal{O}_X)$ qui engendre
$D_\QCoh(\mathcal{O}_X)$. D'après le théorème \ref{theorem-bondal-van-den-Bergh},
il en existe au moins un. En combinant le
lemme \ref{lemma-quasi-coherence-direct-sums} avec la
proposition \ref{proposition-compact-is-perfect} et avec
Catégories dérivées, proposition
\ref{derived-proposition-generator-versus-classical-generator},
on voit que $G$ est un générateur classique de $D_{perf}(\mathcal{O}_X)$.
\end{remark}

\noindent
Le résultat suivant renforce la
proposition \ref{proposition-compact-is-perfect}.
Soit $T \subset X$ un sous-ensemble fermé d'un schéma $X$. Comme précédemment,
$D_T(\mathcal{O}_X)$ désigne la sous-catégorie strictement pleine, saturée
et triangulée de $D(\mathcal{O}_X)$ formée des objets
à support dans $T$ (définition \ref{definition-supported-on}).
Puisque les sommes directes commutent aux faisceaux de cohomologie, il en résulte
que $D_T(\mathcal{O}_X)$ possède des sommes directes, qui coïncident
avec celles de $D(\mathcal{O}_X)$.

\begin{lemma}
\label{lemma-compact-is-perfect-with-support}
Soit $X$ un schéma quasi-compact et quasi-séparé.
Soit $T \subset X$ un sous-ensemble fermé tel que $X \setminus T$
soit quasi-compact. Un objet de $D_{\QCoh, T}(\mathcal{O}_X)$ est compact
si et seulement s'il est parfait en tant qu'objet de $D(\mathcal{O}_X)$.
\end{lemma}

\begin{proof}
Observons que $D_{\QCoh, T}(\mathcal{O}_X)$ est une catégorie triangulée
possédant des sommes directes, d'après la remarque précédant le lemme.
D'après la proposition \ref{proposition-compact-is-perfect},
les objets parfaits définissent des objets compacts de $D(\mathcal{O}_X)$,
donc a fortiori de toute sous-catégorie stable par sommes
directes. Pour la réciproque, nous utiliserons l'existence d'un générateur
$E \in D_{\QCoh, T}(\mathcal{O}_X)$ qui est un complexe parfait
de $\mathcal{O}_X$-modules, voir le
lemme \ref{lemma-generator-with-support}.
D'après ce qui précède, $E$ est donc compact. Il résulte alors de
Catégories dérivées, proposition
\ref{derived-proposition-generator-versus-classical-generator},
que $E$ est un générateur classique de la sous-catégorie pleine
des objets compacts de $D_{\QCoh, T}(\mathcal{O}_X)$.
Ainsi, tout objet compact se construit à partir de $E$ par
une suite finie d'opérations consistant à prendre
(a) des décalés, (b) des sommes directes finies, (c) des cônes, et
(d) des facteurs directs. Chacune de ces opérations conserve
la propriété d'être parfait, d'où le résultat.
\end{proof}

\begin{remark}
\label{remark-classical-generator-with-support}
Soit $X$ un schéma quasi-compact et quasi-séparé.
Soit $T \subset X$ un sous-ensemble fermé tel que $X \setminus T$
soit quasi-compact. Soit $G$ un
objet parfait de $D_{\QCoh, T}(\mathcal{O}_X)$ qui engendre
$D_{\QCoh, T}(\mathcal{O}_X)$. D'après le lemme \ref{lemma-generator-with-support},
il en existe au moins un. En combinant l'existence de sommes directes dans
$D_{\QCoh, T}(\mathcal{O}_X)$ avec le
lemme \ref{lemma-compact-is-perfect-with-support} et avec
Catégories dérivées, proposition
\ref{derived-proposition-generator-versus-classical-generator},
on voit que $G$ est un générateur classique de $D_{perf, T}(\mathcal{O}_X)$.
\end{remark}

\noindent
Le lemme suivant est une application des idées intervenant dans
la démonstration du lemme précédent.

\begin{lemma}
\label{lemma-map-from-pseudo-coherent-to-complex-with-support}
Soit $X$ un schéma quasi-compact et quasi-séparé. Soit $T \subset X$
un sous-ensemble fermé tel que $U = X \setminus T$ soit quasi-compact.
Soit $\alpha : P \to E$ un morphisme de $D_\QCoh(\mathcal{O}_X)$ vérifiant
l'une des conditions suivantes :
\begin{enumerate}
\item $P$ est parfait et $E$ est à support dans $T$, ou
\item $P$ est pseudo-cohérent, $E$ est à support dans $T$, et $E$ est borné inférieurement.
\end{enumerate}
Alors il existe un complexe parfait de $\mathcal{O}_X$-modules $I$
et un morphisme $I \to \mathcal{O}_X[0]$ tels que
$I \otimes^\mathbf{L} P \to E$ soit nul et que
$I|_U \to \mathcal{O}_U[0]$ soit un
isomorphisme.
\end{lemma}

\begin{proof}
Posons $\mathcal{D} = D_{\QCoh, T}(\mathcal{O}_X)$. Dans les deux cas, le complexe
$K = R\SheafHom(P, E)$ est un objet de $\mathcal{D}$. Voir le
lemme \ref{lemma-quasi-coherence-internal-hom} pour la quasi-cohérence.
Il est clair que $K$ est à support dans $T$, car la formation de $R\SheafHom$
commute à la restriction aux ouverts.
Le morphisme $\alpha$ définit un élément de
$H^0(K) = \Hom_{D(\mathcal{O}_X)}(\mathcal{O}_X[0], K)$.
Il suffit alors de démontrer le résultat pour le morphisme
$\alpha : \mathcal{O}_X[0] \to K$.

\medskip\noindent
Soit $E \in \mathcal{D}$ un générateur parfait, voir le
lemme \ref{lemma-generator-with-support}. Écrivons
$$
K = \text{hocolim} K_n
$$
comme dans Catégories dérivées, lemme \ref{derived-lemma-write-as-colimit},
en utilisant le générateur $E$. Puisque le foncteur $\mathcal{D} \to D(\mathcal{O}_X)$
commute aux sommes directes, on voit que $K = \text{hocolim} K_n$
vaut dans $D(\mathcal{O}_X)$. Puisque $\mathcal{O}_X$ est un objet compact
de $D(\mathcal{O}_X)$, on trouve un $n$ et un morphisme
$\alpha_n : \mathcal{O}_X \to K_n$ qui induit $\alpha$, voir
Catégories dérivées, lemme \ref{derived-lemma-commutes-with-countable-sums}.
D'après Catégories dérivées, lemme \ref{derived-lemma-factor-through},
appliqué au morphisme $\mathcal{O}_X[0] \to K_n$ dans la catégorie
ambiante $D(\mathcal{O}_X)$, on voit que $\alpha_n$ se factorise en
$\mathcal{O}_X[0] \to Q \to K_n$, où $Q$ est un objet
de $\langle E \rangle$. On en déduit que $Q$ est un complexe parfait
à support dans $T$.

\medskip\noindent
Choisissons un triangle distingué
$$
I \to \mathcal{O}_X[0] \to Q \to I[1]
$$
Par construction, $I$ est parfait, le morphisme $I \to \mathcal{O}_X[0]$
se restreint en un isomorphisme sur $U$, et le composé
$I \to K$ est nul, car $\alpha$ se factorise par $Q$.
Ceci démontre le lemme.
\end{proof}






\section{Catégories dérivées comme catégories de modules}
\label{section-derived-is-dga}

\noindent
Dans cette section, nous tirons quelques conséquences de ce qui précède.
Nous avons d'abord besoin de quelques lemmes supplémentaires.

\begin{lemma}
\label{lemma-tensor-with-QCoh-complex}
Soit $X$ un schéma. Soit $K^\bullet$ un complexe de $\mathcal{O}_X$-modules
dont les faisceaux de cohomologie sont quasi-cohérents.\newline Soit
$(E, d) = \Hom_{\text{Comp}^{dg}(\mathcal{O}_X)}(K^\bullet, K^\bullet)$
l'algèbre différentielle graduée des endomorphismes. Alors le foncteur
$$
- \otimes_E^\mathbf{L} K^\bullet :
D(E, \text{d}) \longrightarrow D(\mathcal{O}_X)
$$
de
Algèbre différentielle graduée, lemme
\ref{dga-lemma-tensor-with-complex-derived},
a son image contenue dans $D_\QCoh(\mathcal{O}_X)$.
\end{lemma}

\begin{proof}
Soit $P$ un $E$-module différentiel gradué ayant la propriété (P),
et soit $F_\bullet$ une filtration de $P$ comme dans
Algèbre différentielle graduée, section \ref{dga-section-P-resolutions}.
On a alors
$$
P \otimes_E K^\bullet = \text{hocolim}\ F_iP \otimes_E K^\bullet
$$
Chacun des $F_iP$ possède une filtration finie dont les quotients successifs
sont des sommes directes de $E[k]$. Le résultat s'ensuit aisément.
\end{proof}

\noindent
Le résultat suivant est tiré de \cite{BvdB}.

\begin{theorem}
\label{theorem-DQCoh-is-Ddga}
Soit $X$ un schéma quasi-compact et quasi-séparé.
Il existe alors une algèbre différentielle graduée $(E, \text{d})$
dont seul un nombre fini de groupes de cohomologie $H^i(E)$ sont non nuls,
telle que $D_\QCoh(\mathcal{O}_X)$ soit équivalente
à $D(E, \text{d})$.
\end{theorem}

\begin{proof}
Soit $K^\bullet$ un complexe K-injectif de $\mathcal{O}$-modules qui
soit parfait et engendre $D_\QCoh(\mathcal{O}_X)$. Un tel complexe
existe d'après le théorème \ref{theorem-bondal-van-den-Bergh}
et l'existence des résolutions K-injectives. Nous démontrerons que le
théorème est vérifié avec
$$
(E, \text{d}) = \Hom_{\text{Comp}^{dg}(\mathcal{O}_X)}(K^\bullet, K^\bullet)
$$
où $\text{Comp}^{dg}(\mathcal{O}_X)$ est la catégorie différentielle graduée
des complexes de $\mathcal{O}$-modules. Voir
Algèbre différentielle graduée, section \ref{dga-section-variant-base-change}.
Puisque $K^\bullet$ est K-injectif, on
a
\begin{equation}
\label{equation-E-is-OK}
H^n(E) = \Ext^n_{D(\mathcal{O}_X)}(K^\bullet, K^\bullet)
\end{equation}
pour tout $n \in \mathbf{Z}$. Seul un nombre fini de ces groupes Ext
sont non nuls, d'après le lemme \ref{lemma-ext-from-perfect-into-bounded-QCoh}.
Considérons le foncteur
$$
- \otimes_E^\mathbf{L} K^\bullet :
D(E, \text{d}) \longrightarrow D(\mathcal{O}_X)
$$
de
Algèbre différentielle graduée, lemme
\ref{dga-lemma-tensor-with-complex-derived}.
Puisque $K^\bullet$ est parfait, il définit un objet compact de
$D(\mathcal{O}_X)$, voir la proposition \ref{proposition-compact-is-perfect}.
En combinant ceci avec (\ref{equation-E-is-OK}), on voit que le foncteur ci-dessus est pleinement
fidèle, d'après
Algèbre différentielle graduée, lemmes
\ref{dga-lemma-fully-faithful-in-compact-case}. Il possède un adjoint à droite
$$
R\Hom(K^\bullet, - ) : D(\mathcal{O}_X) \longrightarrow D(E, \text{d})
$$
d'après Algèbre différentielle graduée, lemmes
\ref{dga-lemma-tensor-with-complex-hom-adjoint},
qui est un foncteur quasi-inverse à gauche par les propriétés générales des foncteurs
adjoints. D'autre part, il résulte du
lemme \ref{lemma-tensor-with-QCoh-complex} que l'on obtient
$$
- \otimes_E^\mathbf{L} K^\bullet :
D(E, \text{d}) \longrightarrow D_\QCoh(\mathcal{O}_X)
$$
et, par le choix de $K^\bullet$ comme générateur de
$D_\QCoh(\mathcal{O}_X)$, le noyau de l'adjoint
restreint à $D_\QCoh(\mathcal{O}_X)$ est nul.
Un argument formel montre que l'on obtient l'équivalence voulue, voir
Catégories dérivées, lemme
\ref{derived-lemma-fully-faithful-adjoint-kernel-zero}.
\end{proof}

\begin{remark}[Variante avec support]
\label{remark-DQCoh-is-Ddga-with-support}
Soit $X$ un schéma quasi-compact et quasi-séparé. Soit $T \subset X$
un sous-ensemble fermé tel que $X \setminus T$ soit quasi-compact.
L'analogue du théorème \ref{theorem-DQCoh-is-Ddga} est vrai
pour $D_{\QCoh, T}(\mathcal{O}_X)$.
Ceci résulte exactement du même argument que dans la démonstration
du théorème, en utilisant les
lemmes \ref{lemma-generator-with-support} et
\ref{lemma-compact-is-perfect-with-support}
et une variante avec supports du lemme \ref{lemma-tensor-with-QCoh-complex}.

Si nous en avons besoin, nous énoncerons ici précisément le
résultat et en donnerons une démonstration détaillée.
\end{remark}

\begin{remark}[Unicité de l'algèbre différentielle graduée]
\label{remark-independence-choice}
Soit $X$ un schéma quasi-compact et quasi-séparé sur un anneau $R$.
Par la construction de la démonstration du
théorème \ref{theorem-DQCoh-is-Ddga},
il existe une algèbre différentielle graduée $(A, \text{d})$ sur $R$
telle que $D_\QCoh(X)$ soit $R$-linéairement équivalente à
$D(A, \text{d})$ comme catégorie triangulée.
On peut se demander : dans quelle mesure $(A, \text{d})$ est-elle unique ?
La réponse n'est que légèrement plus forte que l'affirmation selon laquelle
$(A, \text{d})$ est bien définie à équivalence dérivée près.
Supposons en effet que $(B, \text{d})$ soit un second tel couple.
On a alors
$$
(A, \text{d}) = \Hom_{\text{Comp}^{dg}(\mathcal{O}_X)}(K^\bullet, K^\bullet)
$$
et
$$
(B, \text{d}) = \Hom_{\text{Comp}^{dg}(\mathcal{O}_X)}(L^\bullet, L^\bullet)
$$
pour des complexes K-injectifs $K^\bullet$ et $L^\bullet$
de $\mathcal{O}_X$-modules correspondant à des générateurs parfaits
de $D_\QCoh(\mathcal{O}_X)$. Posons
$$
\Omega = \Hom_{\text{Comp}^{dg}(\mathcal{O}_X)}(K^\bullet, L^\bullet)
\quad
\Omega' = \Hom_{\text{Comp}^{dg}(\mathcal{O}_X)}(L^\bullet, K^\bullet)
$$
Alors $\Omega$ est un $B^{opp} \otimes_R A$-module différentiel gradué,
et $\Omega'$ est un $A^{opp} \otimes_R B$-module différentiel gradué.
De plus, l'équivalence
$$
D(A, \text{d}) \to D_\QCoh(\mathcal{O}_X) \to
D(B, \text{d})
$$
est donnée par le foncteur $- \otimes_A^\mathbf{L} \Omega'$, et il en va
de même pour le quasi-inverse. Nous sommes donc dans la situation
de Algèbre différentielle graduée, remarque \ref{dga-remark-hochschild-cohomology}.
Si nous avons besoin de cette remarque, nous donnerons ici un énoncé précis
et une démonstration détaillée.
\end{remark}




\section{Caractérisation des complexes pseudo-cohérents, I}
\label{section-pseudo-coherent-hocolim}

\noindent
Les méthodes précédentes permettent de caractériser les objets pseudo-cohérents
comme limites inductives homotopiques d'approximations par des objets parfaits.

\begin{lemma}
\label{lemma-pseudo-coherent-hocolim}
Soit $X$ un schéma quasi-compact et quasi-séparé.
Soit $K \in D(\mathcal{O}_X)$. Les conditions suivantes sont équivalentes :
\begin{enumerate}
\item $K$ est pseudo-cohérent, et
\item $K = \text{hocolim} K_n$, où
$K_n$ est parfait et $\tau_{\geq -n}K_n \to \tau_{\geq -n}K$
est un isomorphisme pour tout $n$.
\end{enumerate}
\end{lemma}

\begin{proof}
L'implication (2) $\Rightarrow$ (1) est vraie sur tout espace annelé.
Supposons en effet (2) vérifiée. Rappelons qu'un objet parfait de la catégorie
dérivée est pseudo-cohérent, voir
Cohomologie, lemme \ref{cohomology-lemma-perfect}.
Il résulte alors des définitions que
$\tau_{\geq -n}K_n$ est $(-n + 1)$-pseudo-cohérent,
donc que $\tau_{\geq -n}K$ est $(-n + 1)$-pseudo-cohérent,
et par suite que $K$ est $(-n + 1)$-pseudo-cohérent. Ceci est vrai pour
tout $n$, donc $K$ est pseudo-cohérent, voir
Cohomologie, définition \ref{cohomology-definition-pseudo-coherent}.

\medskip\noindent
Supposons (1) vérifiée. Commençons par choisir une approximation
$K_1 \to K$ de $(X, K, -2)$ par un complexe parfait $K_1$, voir les
définitions \ref{definition-approximation-holds} et
\ref{definition-approximation} et le
théorème \ref{theorem-approximation}.
Supposons avoir construit par récurrence
$$
K_1 \to K_2 \to \ldots \to K_n \to K
$$
avec $K_i$ parfait et tel que
$\tau_{\geq -i}K_i \to \tau_{\geq -i}K$ soit un isomorphisme
pour tout $1 \leq i \leq n$. Choisissons alors $a \leq b$ comme dans le
lemme \ref{lemma-ext-from-perfect-into-bounded-QCoh}
pour l'objet parfait $K_n$. Choisissons une approximation
$K_{n + 1} \to K$ de $(X, K, \min(a - 1, -n - 1))$.
Choisissons un triangle distingué
$$
K_{n + 1} \to K \to C \to K_{n + 1}[1]
$$
On voit alors que $C \in D_\QCoh(\mathcal{O}_X)$ vérifie
$H^i(C) = 0$ pour $i \geq a$. Par le choix de $a, b$,
on a donc $\Hom_{D(\mathcal{O}_X)}(K_n, C) = 0$.
Le composé $K_n \to K \to C$ est donc nul. D'après
Catégories dérivées, lemme \ref{derived-lemma-representable-homological},
on peut donc factoriser $K_n \to K$ par $K_{n + 1}$,
ce qui démontre l'étape de récurrence.

\medskip\noindent
Il reste à démontrer que $K = \text{hocolim} K_n$.
Ceci résulte de l'application de
Catégories dérivées, lemme \ref{derived-lemma-cohomology-of-hocolim},
aux foncteurs
$H^i( - ) : D(\mathcal{O}_X) \to \textit{Mod}(\mathcal{O}_X)$
et du choix des $K_n$.
\end{proof}

\begin{lemma}
\label{lemma-pseudo-coherent-hocolim-with-support}
Soit $X$ un schéma quasi-compact et quasi-séparé.
Soit $T \subset X$ un sous-ensemble fermé tel que $X \setminus T$
soit quasi-compact. Soit $K \in D(\mathcal{O}_X)$ à support dans $T$.
Les conditions suivantes sont équivalentes :
\begin{enumerate}
\item $K$ est pseudo-cohérent, et
\item $K = \text{hocolim} K_n$, où
$K_n$ est parfait, à support dans $T$, et
$\tau_{\geq -n}K_n \to \tau_{\geq -n}K$ est un isomorphisme pour tout $n$.
\end{enumerate}
\end{lemma}

\begin{proof}
La démonstration de ce lemme est exactement celle du
lemme \ref{lemma-pseudo-coherent-hocolim},
à ceci près que, pour choisir les approximations, on utilise
les triplets $(T, K, m)$.
\end{proof}









\section{Un exemple d'équivalence}
\label{section-example-equivalence}

\noindent
Dans la section \ref{section-example-generator},
nous avons démontré que la catégorie dérivée de l'espace projectif $\mathbf{P}^n_A$
sur un anneau $A$ est engendrée par un fibré vectoriel, et même par une somme
directe de tordus du faisceau structural. Dans cette section, nous démontrons que ceci
détermine une équivalence entre $D_\QCoh(\mathcal{O}_{\mathbf{P}^n_A})$
et la catégorie dérivée d'une $A$-algèbre.

\medskip\noindent
Avant d'énoncer le résultat, nous avons besoin de quelques notations.
Soit $A$ un anneau. Soit $X = \mathbf{P}^n_A = \text{Proj}(S)$,
où $S = A[X_0, \ldots, X_n]$. D'après le lemme \ref{lemma-generator-P1},
nous savons que
\begin{equation}
\label{equation-generator-Pn}
P =
\mathcal{O}_X \oplus \mathcal{O}_X(-1) \oplus \ldots \oplus \mathcal{O}_X(-n)
\end{equation}
est un générateur parfait de $D_\QCoh(\mathcal{O}_X)$.
Considérons la $A$-algèbre (non commutative)
\begin{equation}
\label{equation-algebra-for-Pn}
R = \Hom_{\mathcal{O}_X}(P, P) =
\left(
\begin{matrix}
S_0 & S_1 & S_2 & \ldots & \ldots \\
0 & S_0 & S_1 & \ldots & \ldots\\
0 & 0 & S_0 & \ldots & \ldots \\
\ldots & \ldots & \ldots & \ldots & \ldots \\
0 & \ldots & \ldots & \ldots & S_0
\end{matrix}
\right)
\end{equation}
munie de la multiplication et de l'addition évidentes. Si l'on considère $P$ comme un complexe
de $\mathcal{O}_X$-modules de la manière habituelle (c'est-à-dire avec $P$ en degré $0$
et zéro dans tous les autres degrés), on a
$$
R = \Hom_{\text{Comp}^{dg}(\mathcal{O}_X)}(P, P)
$$
où, dans le membre de droite, on considère $R$ comme une algèbre différentielle graduée
sur $A$
à différentielle nulle (c'est-à-dire avec $R$ en degré $0$ et zéro dans tous
les autres degrés). D'après la discussion de
Algèbre différentielle graduée, section \ref{dga-section-variant-base-change},
on obtient un foncteur dérivé
$$
- \otimes_R^\mathbf{L} P : D(R) \longrightarrow D(\mathcal{O}_X),
$$
voir en particulier Algèbre différentielle graduée, lemme
\ref{dga-lemma-tensor-with-complex-derived}.
D'après le lemme \ref{lemma-tensor-with-QCoh-complex},
l'image essentielle de ce foncteur
est contenue dans $D_\QCoh(\mathcal{O}_X)$.

\begin{lemma}
\label{lemma-Pn-module-category}
\begin{reference}
\cite{Beilinson}
\end{reference}
Soit $A$ un anneau. Soit $X = \mathbf{P}^n_A = \text{Proj}(S)$,
où $S = A[X_0, \ldots, X_n]$. Avec
$P$ comme dans (\ref{equation-generator-Pn}) et
$R$ comme dans (\ref{equation-algebra-for-Pn}),
le foncteur
$$
- \otimes_R^\mathbf{L} P : D(R) \longrightarrow D_\QCoh(\mathcal{O}_X)
$$
est une équivalence $A$-linéaire de catégories triangulées envoyant $R$
sur $P$.
\end{lemma}

\noindent
Autrement dit, la catégorie dérivée des modules quasi-cohérents sur
l'espace projectif est équivalente à la catégorie dérivée des modules
sur une algèbre (non commutative).
Cette propriété de l'espace projectif semble assez exceptionnelle parmi
tous les schémas projectifs sur $A$.

\begin{proof}
Pour démontrer que notre foncteur est pleinement fidèle, il suffit de démontrer
que $\Ext^i_X(P, P)$ est nul pour $i \not = 0$ et égal
à $R$ pour $i = 0$, voir
Algèbre différentielle graduée, lemme
\ref{dga-lemma-fully-faithful-in-compact-case}.
Comme dans la démonstration du
lemme \ref{lemma-ext-from-perfect-into-bounded-QCoh},
on voit que
$$
\Ext^i_X(P, P) = H^i(X, P^\wedge \otimes P) =
\bigoplus\nolimits_{0 \leq a, b \leq n} H^i(X, \mathcal{O}_X(a - b))
$$
Le calcul de la cohomologie de l'espace projectif
(Cohomologie des schémas, lemme
\ref{coherent-lemma-cohomology-projective-space-over-ring})
montre que
ces groupes $\Ext$ sont nuls sauf si $i = 0$.
Pour $i = 0$, on retrouve $R$, puisque c'est ainsi que nous avons défini $R$
dans (\ref{equation-algebra-for-Pn}).
D'après Algèbre différentielle graduée, lemme
\ref{dga-lemma-tensor-with-complex-hom-adjoint},
notre foncteur possède un adjoint à droite, à savoir
$R\Hom(P, -) : D_\QCoh(\mathcal{O}_X) \to D(R)$.
Puisque $P$ engendre $D_\QCoh(\mathcal{O}_X)$ d'après le
lemme \ref{lemma-generator-P1},
on voit que le noyau de $R\Hom(P, -)$ est nul.
Notre foncteur est donc une équivalence de catégories
triangulées, d'après Catégories dérivées, lemme
\ref{derived-lemma-fully-faithful-adjoint-kernel-zero}.
\end{proof}






\section{Nouvelle étude du cohérateur}
\label{section-better-coherator}

\noindent
Dans la section \ref{section-coherator}, nous avons construit et étudié
l'adjoint à droite $RQ_X$ du foncteur canonique
$D(\QCoh(\mathcal{O}_X)) \to D(\mathcal{O}_X)$.
Il a été construit comme le prolongement dérivé à droite du cohérateur
$Q_X : \textit{Mod}(\mathcal{O}_X) \to \QCoh(\mathcal{O}_X)$.
Dans cette section, nous étudions quand le foncteur d'inclusion
$$
D_\QCoh(\mathcal{O}_X) \longrightarrow D(\mathcal{O}_X)
$$
possède un adjoint à droite. Si cet adjoint à droite existe, nous le
noterons\footnote{Cette notation est probablement inhabituelle. Cependant, nous avons déjà
employé $Q_X$ pour le cohérateur et $RQ_X$ pour son prolongement dérivé.}
$$
DQ_X :
D(\mathcal{O}_X) \longrightarrow D_\QCoh(\mathcal{O}_X)
$$
Il se trouve que les schémas quasi-compacts et quasi-séparés
possèdent un tel adjoint à droite.

\begin{lemma}
\label{lemma-better-coherator}
Soit $X$ un schéma quasi-compact et quasi-séparé.
Le foncteur d'inclusion $D_\QCoh(\mathcal{O}_X) \to D(\mathcal{O}_X)$
possède un adjoint à droite $DQ_X$.
\end{lemma}

\begin{proof}[Première démonstration]
Nous utiliserons le principe de récurrence de
Cohomologie des schémas, lemme \ref{coherent-lemma-induction-principle},
pour démontrer ceci. Si $D(\QCoh(\mathcal{O}_X)) \to D_\QCoh(\mathcal{O}_X)$
est une équivalence, le lemme est vrai, car le foncteur
$RQ_X$ de la section \ref{section-coherator} est adjoint à droite du foncteur
$D(\QCoh(\mathcal{O}_X)) \to D(\mathcal{O}_X)$.
En particulier, notre lemme est vrai pour les schémas affines, voir le
lemme \ref{lemma-affine-coherator}.
Il suffit donc de montrer que, si $X = U \cup V$
est une réunion de deux ouverts quasi-compacts et si le
lemme est vrai pour $U$, $V$ et $U \cap V$, alors il est vrai pour $X$.

\medskip\noindent
L'adjoint existe si et seulement si, pour tout objet $K$ de
$D(\mathcal{O}_X)$, on peut trouver un triangle distingué
$$
E' \to E \to K \to E'[1]
$$
dans $D(\mathcal{O}_X)$
tel que $E'$ appartienne à $D_\QCoh(\mathcal{O}_X)$ et que
$\Hom(M, K) = 0$ pour tout $M$ de $D_\QCoh(\mathcal{O}_X)$. Voir
Catégories dérivées, lemme \ref{derived-lemma-right-adjoint}.
Considérons le triangle distingué
$$
E \to Rj_{U, *}E|_U \oplus Rj_{V, *}E|_V \to
Rj_{U \cap V, *}E|_{U \cap V} \to E[1]
$$
dans $D(\mathcal{O}_X)$ de
Cohomologie, lemme \ref{cohomology-lemma-exact-sequence-j-star}.
D'après Catégories dérivées, lemme \ref{derived-lemma-prepare-adjoint},
il suffit de construire les triangles distingués voulus
pour $Rj_{U, *}E|_U$, $Rj_{V, *}E|_V$ et
$Rj_{U \cap V, *}E|_{U \cap V}$. Ceci nous ramène à l'assertion
étudiée dans le paragraphe suivant.

\medskip\noindent
Soit $j : U \to X$ l'immersion ouverte correspondant à un ouvert quasi-compact $U$
pour lequel le lemme est vrai. Soit $L$ un objet de $D(\mathcal{O}_U)$.
Il existe alors un triangle distingué
$$
E' \to Rj_*L \to K \to E'[1]
$$
dans $D(\mathcal{O}_X)$
tel que $E'$ appartienne à $D_\QCoh(\mathcal{O}_X)$ et que
$\Hom(M, K) = 0$ pour tout $M$ de $D_\QCoh(\mathcal{O}_X)$.
Pour le voir, choisissons un triangle distingué
$$
L' \to L \to Q \to L'[1]
$$
dans $D(\mathcal{O}_U)$ tel que $L'$ appartienne à $D_\QCoh(\mathcal{O}_U)$
et que $\Hom(N, Q) = 0$ pour tout $N$ de $D_\QCoh(\mathcal{O}_U)$.
C'est possible, car l'énoncé de
Catégories dérivées, lemme \ref{derived-lemma-right-adjoint},
est une équivalence.
On obtient un triangle distingué
$$
Rj_*L' \to Rj_*L \to Rj_*Q \to Rj_*L'[1]
$$
dans $D(\mathcal{O}_X)$. Observons que $Rj_*L'$ appartient à $D_\QCoh(\mathcal{O}_X)$
d'après le lemme \ref{lemma-quasi-coherence-direct-image}.
D'autre part, si $M$ appartient à $D_\QCoh(\mathcal{O}_X)$, alors
$$
\Hom(M, Rj_*Q) = \Hom(Lj^*M, Q) = 0
$$
car $Lj^*M$ appartient à $D_\QCoh(\mathcal{O}_U)$ d'après le
lemme \ref{lemma-quasi-coherence-pullback}.
Ceci achève la démonstration.
\end{proof}

\begin{proof}[Deuxième démonstration]
L'adjoint existe d'après Catégories dérivées, proposition
\ref{derived-proposition-brown}. Les hypothèses sont vérifiées :
tout d'abord, $D_\QCoh(\mathcal{O}_X)$ possède des sommes directes,
et celles-ci commutent au foncteur d'inclusion
(lemme \ref{lemma-quasi-coherence-direct-sums}).
D'autre part, $D_\QCoh(\mathcal{O}_X)$
est engendrée par des objets compacts, car elle possède un générateur
parfait d'après le théorème \ref{theorem-bondal-van-den-Bergh},
et les objets parfaits sont compacts d'après la
proposition \ref{proposition-compact-is-perfect}.
\end{proof}

\begin{lemma}
\label{lemma-pushforward-better-coherator}
Soit $f : X \to Y$ un morphisme quasi-compact et quasi-séparé
de schémas. Si les adjoints à droite $DQ_X$ et $DQ_Y$
des foncteurs d'inclusion $D_\QCoh \to D$ existent pour $X$ et $Y$, alors
$$
Rf_* \circ DQ_X = DQ_Y \circ Rf_*
$$
\end{lemma}

\begin{proof}
L'énoncé a un sens, car $Rf_*$ envoie
$D_\QCoh(\mathcal{O}_X)$ dans $D_\QCoh(\mathcal{O}_Y)$ d'après le
lemme \ref{lemma-quasi-coherence-direct-image}.
Il est vrai, car $Lf^*$ envoie de même
$D_\QCoh(\mathcal{O}_Y)$ dans $D_\QCoh(\mathcal{O}_X)$
(lemme \ref{lemma-quasi-coherence-pullback}),
de sorte que $Rf_* \circ DQ_X$ et $DQ_Y \circ Rf_*$
sont tous deux adjoints à droite de $Lf^* : D_\QCoh(\mathcal{O}_Y) \to D(\mathcal{O}_X)$.
\end{proof}

\begin{remark}
\label{remark-explain-consequence}
Soit $X$ un schéma quasi-compact et quasi-séparé.
Soit $X = U \cup V$, avec $U$ et $V$ ouverts quasi-compacts.
D'après le lemme \ref{lemma-better-coherator}, les foncteurs
$DQ_X$, $DQ_U$, $DQ_V$, $DQ_{U \cap V}$ existent. De plus, il existe un
triangle distingué canonique
$$
DQ_X(K) \to Rj_{U, *}DQ_U(K|_U) \oplus Rj_{V, *}DQ_V(K|_V)
\to Rj_{U \cap V, *}DQ_{U \cap V}(K|_{U \cap V}) \to
$$
pour tout $K \in D(\mathcal{O}_X)$. Ceci résulte de l'application du
foncteur exact $DQ_X$ au triangle distingué de
Cohomologie, lemme \ref{cohomology-lemma-exact-sequence-j-star},
et de trois applications du lemme \ref{lemma-pushforward-better-coherator}.
\end{remark}

\begin{lemma}
\label{lemma-boundedness-better-coherator}
Soit $X$ un schéma quasi-compact et quasi-séparé. Le foncteur
$DQ_X$ du lemme \ref{lemma-better-coherator}
possède la propriété de bornitude suivante :
il existe un entier $N = N(X)$ tel que, si
$K$ appartient à $D(\mathcal{O}_X)$ et vérifie
$H^i(U, K) = 0$ pour $U$ ouvert affine de $X$ et $i \not \in [a, b]$, alors
les faisceaux de cohomologie $H^i(DQ_X(K))$ sont nuls pour
$i \not \in [a, b + N]$.
\end{lemma}

\begin{proof}
Nous utiliserons le principe de récurrence de
Cohomologie des schémas, lemme \ref{coherent-lemma-induction-principle}.

\medskip\noindent
Si $X$ est affine, le lemme est vrai avec $N = 0$, car alors
$RQ_X = DQ_X$ s'obtient en prenant le complexe de faisceaux
quasi-cohérents associé à $R\Gamma(X, K)$.
Voir les lemmes \ref{lemma-affine-compare-bounded} et \ref{lemma-affine-coherator}.

\medskip\noindent
Supposons que $U, V$ soient des ouverts quasi-compacts de $X$
et que le lemme soit vrai pour $U$, $V$ et $U \cap V$,
avec des entiers $N(U)$, $N(V)$ et $N(U \cap V)$.
Supposons maintenant que $K$ appartienne à $D(\mathcal{O}_X)$ et vérifie
$H^i(W, K) = 0$ pour tout ouvert affine $W \subset X$ et tout $i \not \in [a, b]$.
Alors $K|_U$, $K|_V$, $K|_{U \cap V}$ ont la même propriété.
On voit donc que $RQ_U(K|_U)$, $RQ_V(K|_V)$ et
$RQ_{U \cap V}(K|_{U \cap V})$ ont leurs faisceaux de cohomologie
nuls hors de l'intervalle $[a, b + \max(N(U), N(V), N(U \cap V))$.
Puisque les foncteurs $Rj_{U, *}$, $Rj_{V, *}$, $Rj_{U \cap V, *}$
sont de dimension cohomologique finie sur $D_\QCoh$ d'après le
lemme \ref{lemma-quasi-coherence-direct-image},
il existe un $N$ tel que
$Rj_{U, *}DQ_U(K|_U)$, $Rj_{V, *}DQ_V(K|_V)$ et
$Rj_{U \cap V, *}DQ_{U \cap V}(K|_{U \cap V})$ aient leurs faisceaux
de cohomologie nuls hors de l'intervalle $[a, b + N]$.
On conclut alors à l'aide du triangle distingué
de la remarque \ref{remark-explain-consequence}.
\end{proof}

\begin{example}
\label{example-inverse-limit-quasi-coherent}
Soit $X$ un schéma quasi-compact et quasi-séparé. Soit $(\mathcal{F}_n)$
un système projectif de faisceaux quasi-cohérents. Puisque $DQ_X$ est un adjoint à
droite, il commute aux produits, donc aux limites dérivées.
On a donc
$$
DQ_X(R\lim \mathcal{F}_n) =
(R\lim\text{ dans }D_\QCoh(\mathcal{O}_X))(\mathcal{F}_n)
$$
où la première limite $R\lim$ est prise dans $D(\mathcal{O}_X)$.
Posons $K = R\lim \mathcal{F}_n$ pour celle-ci.
Pour tout ouvert affine $U \subset X$, on a
$$
H^i(U, K) =
H^i(R\Gamma(U, R\lim \mathcal{F}_n)) =
H^i(R\lim R\Gamma(U, \mathcal{F}_n)) =
H^i(R\lim \Gamma(U, \mathcal{F}_n))
$$
car la cohomologie commute aux limites dérivées et les
faisceaux quasi-cohérents
$\mathcal{F}_n$ n'ont pas de cohomologie supérieure sur les affines.
Le calcul de $R\lim$ dans la catégorie des groupes
abéliens montre que $H^i(U, K) = 0$
sauf si $i \in [0, 1]$. On en déduit finalement que
la limite $R\lim$ dans $D_\QCoh(\mathcal{O}_X)$, qui est
$DQ_X(K)$ d'après ce qui précède, appartient à $D^b_\QCoh(\mathcal{O}_X)$
d'après le lemme \ref{lemma-boundedness-better-coherator}.
\end{example}













\section{Cohomologie et changement de base, IV}
\label{section-cohomology-and-base-change-perfect}

\noindent
Cette section poursuit la discussion de
Cohomologie des schémas, section
\ref{coherent-section-cohomology-and-base-change-perfect}.
Nous disposons d'abord d'une version très générale de la formule
de projection pour les morphismes quasi-compacts et quasi-séparés de schémas
et les complexes dont les faisceaux de cohomologie sont quasi-cohérents.

\begin{lemma}
\label{lemma-cohomology-base-change}
Soit $f : X \to Y$ un morphisme quasi-compact et quasi-séparé
de schémas. Pour $E$ dans $D_\QCoh(\mathcal{O}_X)$ et
$K$ dans $D_\QCoh(\mathcal{O}_Y)$, le morphisme
$$
Rf_*(E) \otimes_{\mathcal{O}_Y}^\mathbf{L} K
\longrightarrow
Rf_*(E \otimes_{\mathcal{O}_X}^\mathbf{L} Lf^*K)
$$
défini dans
Cohomologie, équation (\ref{cohomology-equation-projection-formula-map}),
est un isomorphisme.
\end{lemma}

\begin{proof}
Pour vérifier que ce morphisme est un isomorphisme, on peut travailler localement sur $Y$.
On se ramène donc au cas où $Y$ est affine.

\medskip\noindent
Supposons que $K = \bigoplus K_i$ soit une somme
directe de complexes $K_i \in D_\QCoh(\mathcal{O}_Y)$.
Si l'énoncé est vrai pour chaque $K_i$, il est vrai pour $K$.
En effet, les foncteurs $Lf^*$ et $\otimes^\mathbf{L}$ conservent
les sommes directes par construction, et $Rf_*$ commute aux sommes directes
(pour les complexes dont les faisceaux de cohomologie sont quasi-cohérents) d'après le
lemme \ref{lemma-quasi-coherence-pushforward-direct-sums}.
De plus, supposons que $K \to L \to M \to K[1]$ soit un triangle
distingué dans $D_\QCoh(Y)$. Si l'énoncé du
lemme est vrai pour deux des objets $K, L, M$, il est vrai pour le troisième
(car les foncteurs considérés sont des foncteurs exacts de catégories triangulées).

\medskip\noindent
Supposons $Y$ affine, et écrivons $Y = \Spec(A)$. Le foncteur
$\widetilde{\ } : D(A) \to D_\QCoh(\mathcal{O}_Y)$ est une équivalence
(lemme \ref{lemma-affine-compare-bounded}).
Soit $T$ la propriété, pour $K \in D(A)$, selon laquelle
l'énoncé du lemme est vrai pour $\widetilde{K}$.
La discussion ci-dessus et
Compléments sur l'algèbre, remarque \ref{more-algebra-remark-P-resolution},
montrent qu'il suffit de démontrer que $T$ est vraie pour $A[k]$.
Ceci achève la démonstration, car l'énoncé du lemme
est clair pour les décalés du faisceau structural.
\end{proof}

\begin{definition}
\label{definition-tor-independent}
Soit $S$ un schéma. Soient $X$, $Y$ des schémas sur $S$. On dit que $X$ et
$Y$ sont {\it Tor-indépendants sur $S$} si, pour tous $x \in X$ et
$y \in Y$ d'image commune $s \in S$, les anneaux
$\mathcal{O}_{X, x}$ et $\mathcal{O}_{Y, y}$ sont Tor-indépendants
sur $\mathcal{O}_{S, s}$ (voir
Compléments sur l'algèbre, définition \ref{more-algebra-definition-tor-independent}).
\end{definition}

\begin{lemma}
\label{lemma-tor-independent}
Soient $f : X \to S$ et $g : Y \to S$ des morphismes de schémas.
Les conditions suivantes sont équivalentes :
\begin{enumerate}
\item $X$ et $Y$ sont Tor-indépendants sur $S$, et
\item pour tous ouverts affines $U \subset X$, $V \subset Y$, $W \subset S$
avec $f(U) \subset W$ et $g(V) \subset W$, les anneaux
$\mathcal{O}_X(U)$ et $\mathcal{O}_Y(V)$ sont Tor-indépendants sur
$\mathcal{O}_S(W)$ ;
\item il existe un recouvrement ouvert affine $S = \bigcup W_i$ et,
pour chaque $i$, des recouvrements ouverts affines $f^{-1}(W_i) = \bigcup U_{ij}$
et $g^{-1}(W_i) = \bigcup V_{ik}$ tels que les anneaux
$\mathcal{O}_X(U_{ij})$ et $\mathcal{O}_Y(V_{ik})$ soient Tor-indépendants sur
$\mathcal{O}_S(W_i)$ pour tous $i, j, k$.
\end{enumerate}
\end{lemma}

\begin{proof}
Omise. Indication : utiliser Compléments sur l'algèbre, lemme
\ref{more-algebra-lemma-tor-independent}.
\end{proof}

\begin{lemma}
\label{lemma-flat-base-change-tor-independent}
Soient $X \to S$ et $Y \to S$ des morphismes de schémas. Soit $S' \to S$ un
morphisme de schémas et notons $X' = X \times_S S'$
et $Y' = Y \times_S S'$.
Si $X$ et $Y$ sont Tor-indépendants sur $S$ et si $S' \to S$ est plat,
alors $X'$ et $Y'$ sont Tor-indépendants sur $S'$.
\end{lemma}

\begin{proof}
Omise. Indication : utiliser le lemme \ref{lemma-tor-independent} et,
sur les ouverts affines, Compléments sur l'algèbre, lemme
\ref{more-algebra-lemma-flat-base-change-tor-independent}.
\end{proof}

\begin{lemma}
\label{lemma-compare-base-change}
Soit $g : S' \to S$ un morphisme de schémas.
Soit $f : X \to S$ quasi-compact et quasi-séparé.
Considérons le diagramme de changement de base
$$
\xymatrix{
X' \ar[r]_{g'} \ar[d]_{f'} &
X \ar[d]^f \\
S' \ar[r]^g &
S
}
$$
Si $X$ et $S'$ sont Tor-indépendants sur $S$, alors, pour tout
$E \in D_\QCoh(\mathcal{O}_X)$, la flèche canonique
$Lg^*Rf_*E \to Rf'_*L(g')^*E$ est un isomorphisme.
\end{lemma}

\begin{proof}
Pour tout objet $E$ de $D(\mathcal{O}_X)$, on peut utiliser
Cohomologie, remarque \ref{cohomology-remark-base-change}, pour obtenir un morphisme
canonique de changement de base $Lg^*Rf_*E \to Rf'_*L(g')^*E$. Pour vérifier qu'il
est un isomorphisme, on peut travailler localement sur $S'$. On peut donc supposer que
$g : S' \to S$ est un morphisme de schémas affines. En particulier, $g$
est affine, et il suffit de montrer que
$$
Rg_*Lg^*Rf_*E \to Rg_*Rf'_*L(g')^*E = Rf_*(Rg'_* L(g')^* E)
$$
est un isomorphisme, voir le lemme \ref{lemma-affine-morphism}
(et utiliser les lemmes \ref{lemma-quasi-coherence-pullback},
\ref{lemma-quasi-coherence-tensor-product} et
\ref{lemma-quasi-coherence-direct-image}
pour voir que les objets $Rf'_*L(g')^*E$ et $Lg^*Rf_*E$
ont des faisceaux de cohomologie quasi-cohérents). Notons que $g'$ est
également affine (Morphismes, lemme \ref{morphisms-lemma-base-change-affine}).
D'après le lemme \ref{lemma-affine-morphism-pull-push}, le morphisme devient
$$
Rf_*E \otimes_{\mathcal{O}_S}^\mathbf{L} g_*\mathcal{O}_{S'}
\longrightarrow
Rf_*(E \otimes_{\mathcal{O}_X}^\mathbf{L} g'_*\mathcal{O}_{X'})
$$
Observons que $g'_*\mathcal{O}_{X'} = f^*g_*\mathcal{O}_{S'}$ (par changement
de base affine, voir Cohomologie des schémas, lemme
\ref{coherent-lemma-affine-base-change}).
Ainsi, d'après le
lemme \ref{lemma-cohomology-base-change}, il suffit de démontrer que
$Lf^*g_*\mathcal{O}_{S'} = f^*g_*\mathcal{O}_{S'}$. Ceci résulte de notre
hypothèse que $X$ et $S'$ sont Tor-indépendants sur $S$. En effet, pour
le vérifier, on peut travailler localement sur $X$, donc aussi supposer $X$ affine.
Écrivons $X = \Spec(A)$, $S = \Spec(R)$ et $S' = \Spec(R')$. Notre hypothèse
implique que $A$ et $R'$ sont Tor-indépendants sur $R$
(Compléments sur l'algèbre, lemme \ref{more-algebra-lemma-tor-independent}), c'est-à-dire que
$\text{Tor}_i^R(A, R') = 0$ pour $i > 0$. Autrement dit,
$A \otimes_R^\mathbf{L} R' = A \otimes_R R'$, ce qui signifie exactement
que $Lf^*g_*\mathcal{O}_{S'} = f^*g_*\mathcal{O}_{S'}$
(utiliser le lemme \ref{lemma-quasi-coherence-pullback}).
\end{proof}

\noindent
Le lemme suivant sera utilisé dans le chapitre sur les complexes dualisants.

\begin{lemma}
\label{lemma-affine-morphism-and-hom-out-of-perfect}
Considérons un carré cartésien
$$
\xymatrix{
X' \ar[r]_{g'} \ar[d]_{f'} & X \ar[d]^f \\
S' \ar[r]^g & S
}
$$
de schémas quasi-compacts et quasi-séparés. Supposons que $g$ et $f$ soient
Tor-indépendants et que $S = \Spec(R)$, $S' = \Spec(R')$ soient affines. Pour
$M, K \in D(\mathcal{O}_X)$, le morphisme canonique
$$
R\Hom_X(M, K) \otimes^\mathbf{L}_R R'
\longrightarrow
R\Hom_{X'}(L(g')^*M, L(g')^*K)
$$
dans $D(R')$ est un isomorphisme dans les deux cas suivants :
\begin{enumerate}
\item $M \in D(\mathcal{O}_X)$ est parfait et $K \in D_\QCoh(X)$, ou
\item $M \in D(\mathcal{O}_X)$ est pseudo-cohérent,
$K \in D_\QCoh^+(X)$, et $R'$ est de dimension de Tor finie sur $R$.
\end{enumerate}
\end{lemma}

\begin{proof}
Il existe un morphisme canonique
$R\Hom_X(M, K) \to R\Hom_{X'}(L(g')^*M, L(g')^*K)$
dans $D(\Gamma(X, \mathcal{O}_X))$ entre complexes Hom globaux, voir
Cohomologie, section \ref{cohomology-section-global-RHom}.
Par restriction des scalaires, on peut le considérer comme un morphisme dans $D(R)$.
On utilise alors l'adjonction entre la restriction et
$- \otimes_R^\mathbf{L} R'$ pour obtenir le morphisme affiché dans le lemme.
Ce morphisme étant défini, il suffit de démontrer que c'est un isomorphisme
dans la catégorie dérivée des groupes abéliens.

\medskip\noindent
Le membre de droite est égal à
$$
R\Hom_X(M, R(g')_*L(g')^*K) =
R\Hom_X(M, K \otimes_{\mathcal{O}_X}^\mathbf{L} g'_*\mathcal{O}_{X'})
$$
d'après le lemme \ref{lemma-affine-morphism-pull-push}. Dans les deux cas, le complexe
$R\SheafHom(M, K)$ est un objet de $D_\QCoh(\mathcal{O}_X)$ d'après le
lemme \ref{lemma-quasi-coherence-internal-hom}. Il existe un morphisme naturel
$$
R\SheafHom(M, K) \otimes_{\mathcal{O}_X}^\mathbf{L} g'_*\mathcal{O}_{X'}
\longrightarrow
R\SheafHom(M, K \otimes_{\mathcal{O}_X}^\mathbf{L} g'_*\mathcal{O}_{X'})
$$
qui est un isomorphisme dans les deux cas d'après le
lemme \ref{lemma-internal-hom-evaluate-tensor-isomorphism}.
Pour voir que ce lemme s'applique dans le cas (2), notons que
$g'_*\mathcal{O}_{X'} = Rg'_*\mathcal{O}_{X'} =
Lf^*g_*\mathcal{O}_{S'}$, la deuxième égalité résultant du
lemme \ref{lemma-compare-base-change}.
En utilisant le lemme \ref{lemma-tor-dimension-affine} et
Cohomologie, lemme \ref{cohomology-lemma-tor-amplitude-pullback},
on en déduit que $g'_*\mathcal{O}_{X'}$ est de dimension de Tor finie.
Ainsi, dans les deux cas, en remplaçant $K$ par $R\SheafHom(M, K)$, on se ramène
à démontrer que
$$
R\Gamma(X, K) \otimes^\mathbf{L}_A A' \longrightarrow
R\Gamma(X, K \otimes^\mathbf{L}_{\mathcal{O}_X} g'_*\mathcal{O}_{X'})
$$
est un isomorphisme.
Notons que le membre de gauche est égal à $R\Gamma(X', L(g')^*K)$
d'après le lemme \ref{lemma-affine-morphism-pull-push}.
Le résultat découle donc du
lemme \ref{lemma-compare-base-change}.
\end{proof}

\begin{remark}
\label{remark-multiplication-map}
Reprenons les notations du lemme \ref{lemma-affine-morphism-and-hom-out-of-perfect}.
Le diagramme
$$
\xymatrix{
R\Hom_X(M, Rg'_*L) \otimes_R^\mathbf{L} R' \ar[r] \ar[d]_\mu &
R\Hom_{X'}(L(g')^*M, L(g')^*Rg'_*L) \ar[d]^a \\
R\Hom_X(M, R(g')_*L) \ar@{=}[r] &
R\Hom_{X'}(L(g')^*M, L)
}
$$
est commutatif ; la flèche horizontale supérieure est le morphisme du lemme,
$\mu$ est le morphisme de multiplication, et $a$ provient du morphisme d'adjonction
$L(g')^*Rg'_*L \to L$. Le morphisme de multiplication est le morphisme d'adjonction
$K' \otimes_R^\mathbf{L} R' \to K'$ pour tout $K' \in D(R')$.
\end{remark}

\begin{lemma}
\label{lemma-tor-independence-and-tor-amplitude}
Considérons un carré cartésien de schémas
$$
\xymatrix{
X' \ar[r]_{g'} \ar[d]_{f'} & X \ar[d]^f \\
S' \ar[r]^g & S
}
$$
Supposons que $g$ et $f$ soient Tor-indépendants.
\begin{enumerate}
\item Si $E \in D(\mathcal{O}_X)$ a une amplitude de Tor
contenue dans $[a, b]$ en tant que complexe de $f^{-1}\mathcal{O}_S$-modules,
alors $L(g')^*E$ a une amplitude de Tor
contenue dans $[a, b]$ en tant que complexe de $f^{-1}\mathcal{O}_{S'}$-modules.
\item Si $\mathcal{G}$ est un $\mathcal{O}_X$-module plat
sur $S$, alors $L(g')^*\mathcal{G} = (g')^*\mathcal{G}$.
\end{enumerate}
\end{lemma}

\begin{proof}
On peut calculer la dimension de Tor sur les fibres, voir
Cohomologie, lemme \ref{cohomology-lemma-tor-amplitude-stalk}.
Si $x' \in X'$ a pour image $x \in X$, alors
$$
(L(g')^*E)_{x'} =
E_x \otimes_{\mathcal{O}_{X, x}}^\mathbf{L} \mathcal{O}_{X', x'}
$$
Soient $s' \in S'$ et $s \in S$ les images de $x'$ et $x$.
Puisque $X$ et $S'$ sont Tor-indépendants sur $S$, on peut appliquer
Compléments sur l'algèbre, lemme \ref{more-algebra-lemma-base-change-comparison},
pour voir que le membre de droite de la formule affichée est égal à
$E_x \otimes_{\mathcal{O}_{S, s}}^\mathbf{L} \mathcal{O}_{S', s'}$
dans $D(\mathcal{O}_{S', s'})$.
L'assertion (1) résulte donc de
Compléments sur l'algèbre, lemme \ref{more-algebra-lemma-pull-tor-amplitude}.
Pour voir (2), observons que la platitude de $\mathcal{G}$ équivaut à
la condition que $\mathcal{G}[0]$ ait une amplitude de Tor contenue dans $[0, 0]$.
On conclut en appliquant (1).
\end{proof}

\begin{lemma}
\label{lemma-compare-base-change-closed-immersion}
Considérons un diagramme cartésien de schémas
$$
\xymatrix{
Z' \ar[r]_{i'} \ar[d]_g & X' \ar[d]^f \\
Z \ar[r]^i & X
}
$$
où $i$ est une immersion fermée. Si $Z$ et $X'$ sont
Tor-indépendants sur $X$, alors $Ri'_* \circ Lg^* = Lf^* \circ Ri_*$
comme foncteurs $D(\mathcal{O}_Z) \to D(\mathcal{O}_{X'})$.
\end{lemma}

\begin{proof}
Notons que le lemme doit être vrai pour tout $K \in D(\mathcal{O}_Z)$.
Observons que $i_*$ et $i'_*$ sont des foncteurs exacts ; par conséquent,
$Ri_*$ et $Ri'_*$ se calculent en appliquant $i_*$ et $i'_*$
à des représentants quelconques. Ainsi le morphisme de changement de base
$$
Lf^*(Ri_*(K)) \longrightarrow Ri'_*(Lg^*(K))
$$
sur les fibres en un point $z' \in Z'$ d'image $z \in Z$ est donné par
$$
K_z \otimes_{\mathcal{O}_{X, z}}^\mathbf{L} \mathcal{O}_{X', z'}
\longrightarrow
K_z \otimes_{\mathcal{O}_{Z, z}}^\mathbf{L} \mathcal{O}_{Z', z'}
$$
Ce morphisme est un isomorphisme d'après
Compléments sur l'algèbre, lemme \ref{more-algebra-lemma-base-change-comparison},
et la Tor-indépendance supposée.
\end{proof}








\section{Formule de K\"unneth, II}
\label{section-kunneth}

\noindent
Pour le cas où la base est un corps, voir
Variétés, section \ref{varieties-section-kunneth}.
Considérons un diagramme cartésien de schémas
$$
\xymatrix{
& X \times_S Y \ar[ld]^p \ar[rd]_q \ar[dd]^f \\
X \ar[rd]_a & & Y \ar[ld]^b \\
& S
}
$$
Soient $K \in D(\mathcal{O}_X)$ et $M \in D(\mathcal{O}_Y)$.
Il existe un morphisme canonique
\begin{equation}
\label{equation-kunneth}
Ra_*K \otimes_{\mathcal{O}_S}^\mathbf{L} Rb_*M
\longrightarrow
Rf_*(Lp^*K \otimes_{\mathcal{O}_{X \times_S Y}}^\mathbf{L} Lq^*M)
\end{equation}
En effet, on peut utiliser les morphismes
$Ra_*K \to Ra_*Rp_* Lp^*K = Rf_*Lp^*K$ et
$Rb_*M \to Rb_*Rq_* Lq^*M = Rf_*Lq^*M$, puis le produit cup
relatif (Cohomologie, remarque \ref{cohomology-remark-cup-product}).

\medskip\noindent
Posons $A = \Gamma(S, \mathcal{O}_S)$. Il existe un morphisme de K\"unneth global
\begin{equation}
\label{equation-kunneth-global}
R\Gamma(X, K) \otimes_A^\mathbf{L} R\Gamma(Y, M)
\longrightarrow
R\Gamma(X \times_S Y,
Lp^*K \otimes_{\mathcal{O}_{X \times_S Y}}^\mathbf{L} Lq^*M)
\end{equation}
dans $D(A)$. Ce morphisme se construit à l'aide des morphismes de rappel
$R\Gamma(X, K) \to R\Gamma(X \times_S Y, Lp^*K)$ et
$R\Gamma(Y, M) \to R\Gamma(X \times_S Y, Lq^*M)$ et
du produit cup construit dans
Cohomologie, section \ref{cohomology-section-cup-product}.

\begin{lemma}
\label{lemma-kunneth}
Dans la situation ci-dessus, si $a$ et $b$ sont quasi-compacts et quasi-séparés
et si $X$ et $Y$ sont Tor-indépendants sur $S$, alors (\ref{equation-kunneth})
est un isomorphisme pour $K \in D_\QCoh(\mathcal{O}_X)$ et
$M \in D_\QCoh(\mathcal{O}_Y)$. Si de plus $S = \Spec(A)$ est affine,
alors le morphisme (\ref{equation-kunneth-global}) est un isomorphisme.
\end{lemma}

\begin{proof}[Première démonstration]
Cela résulte de la suite d'isomorphismes suivante
\begin{align*}
Rf_*(Lp^*K \otimes_{\mathcal{O}_{X \times_S Y}}^\mathbf{L} Lq^*M)
& =
Ra_*Rp_*(Lp^*K \otimes_{\mathcal{O}_{X \times_S Y}}^\mathbf{L} Lq^*M) \\
& =
Ra_*(K \otimes_{\mathcal{O}_X}^\mathbf{L} Rp_*Lq^*M) \\
& =
Ra_*(K \otimes_{\mathcal{O}_X}^\mathbf{L} La^*Rb_*M) \\
& =
Ra_*K \otimes_{\mathcal{O}_S}^\mathbf{L} Rb_*M
\end{align*}
La première égalité vaut parce que $f = a \circ p$. La deuxième égalité
résulte du lemme \ref{lemma-cohomology-base-change}. La troisième, du
lemme \ref{lemma-compare-base-change}. La quatrième, du
lemme \ref{lemma-cohomology-base-change}.
Nous omettons la vérification que la composée de ces isomorphismes
est le morphisme (\ref{equation-kunneth}).
Si $S$ est affine, alors la source et le but de la flèche
(\ref{equation-kunneth-global}) s'obtiennent en appliquant
$R\Gamma(S, -)$ à la source et au but de (\ref{equation-kunneth}),
ce qui donne la dernière assertion ; nous omettons les détails.
\end{proof}

\begin{proof}[Deuxième démonstration]
La construction de la flèche (\ref{equation-kunneth}) est compatible
avec la restriction aux sous-schémas ouverts de $S$, comme il ressort
immédiatement de la construction du produit cup relatif. Il suffit donc de montrer
que (\ref{equation-kunneth}) est un isomorphisme lorsque $S$ est affine.

\medskip\noindent
Supposons $S = \Spec(A)$ affine. Par Leray, on a
$R\Gamma(S, Rf_*K) = R\Gamma(X, K)$, et de même dans les autres cas.
D'après Cohomologie, lemme \ref{cohomology-lemma-compose-cup-product},
le morphisme (\ref{equation-kunneth}) induit le morphisme
(\ref{equation-kunneth-global}) après application de $R\Gamma(S, -)$.
D'autre part, d'après les lemmes \ref{lemma-quasi-coherence-direct-image} et
\ref{lemma-quasi-coherence-tensor-product},
la source et le but du morphisme
(\ref{equation-kunneth}) appartiennent à $D_\QCoh(\mathcal{O}_S)$.
Ainsi, d'après le lemme \ref{lemma-affine-compare-bounded}, il suffit de montrer que
(\ref{equation-kunneth-global}) est un isomorphisme.

\medskip\noindent
Supposons $S = \Spec(A)$, $X = \Spec(B)$ et $Y = \Spec(C)$ tous affines.
Nous utiliserons le lemme \ref{lemma-affine-compare-bounded} sans autre mention.
Dans ce cas, on peut choisir un complexe K-plat $K^\bullet$ de $B$-modules
dont les termes sont plats et tel que $K$ soit représenté par $\widetilde{K}^\bullet$.
De même, on peut choisir un complexe K-plat $M^\bullet$ de $C$-modules
dont les termes sont plats et tel que $M$ soit représenté par $\widetilde{M}^\bullet$.
Voir Compléments sur l'algèbre, lemme \ref{more-algebra-lemma-K-flat-resolution}.
Alors $\widetilde{K}^\bullet$ est un complexe K-plat de $\mathcal{O}_X$-modules,
et de même pour $\widetilde{M}^\bullet$ ; voir
le lemme \ref{lemma-affine-K-flat}. Ainsi,
$La^*K$ est représenté par
$$
a^*\widetilde{K}^\bullet = \widetilde{K^\bullet \otimes_A C}
$$
et de même pour $Lb^*M$. À son tour, c'est un complexe K-plat
de $\mathcal{O}_{X \times_S Y}$-modules d'après le lemme cité ci-dessus
et Compléments sur l'algèbre, lemme \ref{more-algebra-lemma-base-change-K-flat}.
On voit donc finalement que le complexe de
$\mathcal{O}_{X \times_S Y}$-modules associé à
$$
\text{Tot}((K^\bullet \otimes_A C) \otimes_{B \otimes_A C}
(B \otimes_A M^\bullet)) =
\text{Tot}(K^\bullet \otimes_A M^\bullet)
$$
représente $La^*K \otimes_{\mathcal{O}_{X \times_S Y}}^\mathbf{L} Lb^*M$
dans la catégorie dérivée de $X \times_S Y$. En prenant les sections globales,
on obtient $\text{Tot}(K^\bullet \otimes_A M^\bullet)$, qui est bien entendu
également le complexe représentant
$R\Gamma(X, K) \otimes_A^\mathbf{L} R\Gamma(Y, M)$.
Le fait que l'isomorphisme soit donné par le produit cup résulte de la
relation entre le véritable produit cup et le produit cup naïf
dans Cohomologie, section \ref{cohomology-section-cup-product}
(et en particulier de
Cohomologie, lemme \ref{cohomology-lemma-cup-compatible-with-naive},
ainsi que de la discussion qui le suit).

\medskip\noindent
Supposons $S = \Spec(A)$ et $Y$ affines. Nous allons utiliser le principe
d'induction de
Cohomologie des schémas, lemme \ref{coherent-lemma-induction-principle},
pour démontrer l'assertion. Pour cela, il suffit de montrer ceci :
si $X = U \cup V$ est un recouvrement ouvert avec $U$ et $V$ quasi-compacts
et si le morphisme (\ref{equation-kunneth-global})
$$
R\Gamma(U, K) \otimes_A^\mathbf{L} R\Gamma(Y, M)
\longrightarrow
R\Gamma(U \times_S Y,
Lp^*K \otimes_{\mathcal{O}_{X \times_S Y}}^\mathbf{L} Lq^*M)
$$
pour $U$ et $Y$ sur $S$, le morphisme (\ref{equation-kunneth-global})
$$
R\Gamma(V, K) \otimes_A^\mathbf{L} R\Gamma(Y, M)
\longrightarrow
R\Gamma(V \times_S Y,
Lp^*K \otimes_{\mathcal{O}_{X \times_S Y}}^\mathbf{L} Lq^*M)
$$
pour $V$ et $Y$ sur $S$, et le morphisme (\ref{equation-kunneth-global})
$$
R\Gamma(U \cap V, K) \otimes_A^\mathbf{L} R\Gamma(Y, M)
\longrightarrow
R\Gamma((U \cap V) \times_S Y,
Lp^*K \otimes_{\mathcal{O}_{X \times_S Y}}^\mathbf{L} Lq^*M)
$$
pour $U \cap V$ et $Y$ sur $S$ sont des isomorphismes, alors le morphisme
(\ref{equation-kunneth-global}) pour $X$ et $Y$ sur $S$ en est aussi un.
Cependant, d'après Cohomologie, lemme \ref{cohomology-lemma-mayer-vietoris-cup},
ces morphismes s'insèrent dans un morphisme de triangles distingués dont
(\ref{equation-kunneth-global}) est la dernière branche ; on conclut donc
par Catégories dérivées, lemme \ref{derived-lemma-third-isomorphism-triangle}.

\medskip\noindent
Supposons $S = \Spec(A)$ affine. Pour achever la démonstration, on peut appliquer
le principe d'induction de
Cohomologie des schémas, lemme \ref{coherent-lemma-induction-principle},
à $Y$. En effet, ce qui précède montre déjà que notre morphisme est
un isomorphisme lorsque $Y$ est affine. La suite de l'argument est
exactement la même que dans le paragraphe précédent, en échangeant les rôles de
$X$ et $Y$.
\end{proof}

\begin{lemma}
\label{lemma-cohomology-de-rham-base-change}
Soit $a : X \to S$ un morphisme quasi-compact et quasi-séparé
de schémas. Soit $\mathcal{F}^\bullet$ un complexe localement borné
de $a^{-1}\mathcal{O}_S$-modules. Supposons que, pour tout $n \in \mathbf{Z}$,
le faisceau $\mathcal{F}^n$ soit un $a^{-1}\mathcal{O}_S$-module plat et que
$\mathcal{F}^n$ possède une structure de $\mathcal{O}_X$-module quasi-cohérent
compatible avec la structure de $a^{-1}\mathcal{O}_S$-module donnée (mais les
différentielles du complexe $\mathcal{F}^\bullet$ ne sont pas nécessairement
$\mathcal{O}_X$-linéaires). Alors les assertions suivantes sont vraies :
\begin{enumerate}
\item $Ra_*\mathcal{F}^\bullet$ est localement borné,
\item $Ra_*\mathcal{F}^\bullet$ appartient à $D_\QCoh(\mathcal{O}_S)$,
\item $Ra_*\mathcal{F}^\bullet$ est localement de dimension de Tor finie,
\item $\mathcal{G} \otimes_{\mathcal{O}_S}^\mathbf{L} Ra_*\mathcal{F}^\bullet =
Ra_*(a^{-1}\mathcal{G} \otimes_{a^{-1}\mathcal{O}_S} \mathcal{F}^\bullet)$
pour $\mathcal{G} \in \QCoh(\mathcal{O}_S)$, et
\item $K \otimes_{\mathcal{O}_S}^\mathbf{L} Ra_*\mathcal{F}^\bullet =
Ra_*(a^{-1}K \otimes_{a^{-1}\mathcal{O}_S}^\mathbf{L} \mathcal{F}^\bullet)$
pour $K \in D_\QCoh(\mathcal{O}_S)$.
\end{enumerate}
\end{lemma}

\begin{proof}
Les parties (1), (2) et (3) sont locales sur $S$ ; on peut donc supposer, et l'on suppose,
que $S$ est affine. Puisque $a$ est quasi-compact, on en déduit que $X$ est quasi-compact.
Comme $\mathcal{F}^\bullet$ est localement borné, on en déduit que
$\mathcal{F}^\bullet$ est borné.

\medskip\noindent
Pour (1) et (2), on peut utiliser la première suite spectrale
$R^pa_*\mathcal{F}^q \Rightarrow R^{p + q}a_*\mathcal{F}^\bullet$ de
Catégories dérivées, lemme \ref{derived-lemma-two-ss-complex-functor}.
En combinant Cohomologie des schémas, lemme
\ref{coherent-lemma-quasi-coherence-higher-direct-images},
et Homologie, lemme \ref{homology-lemma-biregular-ss-converges},
on conclut.

\medskip\noindent
Démontrons (3) au moyen du principe d'induction de
Cohomologie des schémas, lemme \ref{coherent-lemma-induction-principle}.
Plus précisément, pour un ouvert quasi-compact $U$ de $X$, considérons la
condition que $R(a|_U)_*(\mathcal{F}^\bullet|_U)$ soit de
dimension de Tor finie. Si $U, V$ sont des ouverts quasi-compacts de
$X$, alors on dispose d'un triangle distingué de Mayer--Vietoris relatif
$$
R(a|_{U \cup V})_*\mathcal{F}^\bullet|_{U \cup V} \to
R(a|_U)_*\mathcal{F}^\bullet|_U \oplus
R(a|_V)_*\mathcal{F}^\bullet|_V \to
R(a|_{U \cap V})_*\mathcal{F}^\bullet|_{U \cap V} \to
$$
d'après Cohomologie, lemme \ref{cohomology-lemma-unbounded-relative-mayer-vietoris}.
Le comportement de l'amplitude de Tor dans les triangles distingués
(voir Cohomologie, lemme \ref{cohomology-lemma-cone-tor-amplitude})
montre que, si le résultat est connu pour $U$, $V$ et $U \cap V$, alors
il vaut pour $U \cup V$. Cela nous ramène au cas où
$X$ est affine. Dans ce cas, on a
$$
Ra_*\mathcal{F}^\bullet = a_*\mathcal{F}^\bullet
$$
d'après le lemme d'acyclicité de Leray
(Catégories dérivées, lemme \ref{derived-lemma-leray-acyclicity})
et l'annulation des images directes supérieures des modules quasi-cohérents
par un morphisme affine
(Cohomologie des schémas, lemme \ref{coherent-lemma-relative-affine-vanishing}).
Puisque $\mathcal{F}^n$ est plat sur $S$ par hypothèse et que $X$ est affine, les modules
$a_*\mathcal{F}^n$ sont plats pour tout $n$. Par conséquent, $a_*\mathcal{F}^\bullet$
est un complexe borné de $\mathcal{O}_S$-modules plats, donc il est de
dimension de Tor finie.

\medskip\noindent
Démonstration de la partie (5). Notons
$a' : (X, a^{-1}\mathcal{O}_S) \to (S, \mathcal{O}_S)$
le morphisme plat évident d'espaces annelés. La partie (5) affirme que
$$
K \otimes_{\mathcal{O}_S}^\mathbf{L} Ra'_*\mathcal{F}^\bullet =
Ra'_*(L(a')^*K \otimes_{a^{-1}\mathcal{O}_S}^\mathbf{L}
\mathcal{F}^\bullet)
$$
Ainsi,
Cohomologie, équation (\ref{cohomology-equation-projection-formula-map}),
fournit un morphisme fonctoriel du membre de gauche vers celui de droite, et nous voulons
montrer que ce morphisme est un isomorphisme.
Cette question est locale sur $S$ ; on peut donc supposer, et l'on suppose, que $S$
est affine. Le reste de la démonstration est {\it exactement} identique à la
démonstration du lemme \ref{lemma-cohomology-base-change}, sauf qu'il faut
montrer que le foncteur
$K \mapsto Ra'_*(L(a')^*K \otimes_{a^{-1}\mathcal{O}_S}^\mathbf{L}
\mathcal{F}^\bullet)$ commute aux sommes directes.
C'est ici que nous utilisons le fait que $\mathcal{F}^n$ possède une structure
de $\mathcal{O}_X$-module quasi-cohérent. En effet, observons que
$K \mapsto L(a')^*K
\otimes_{a^{-1}\mathcal{O}_S}^\mathbf{L} \mathcal{F}^\bullet$
commute aux sommes directes arbitraires. Ensuite, si
$\mathcal{F}^\bullet$ consiste en un seul $\mathcal{O}_X$-module
quasi-cohérent $\mathcal{F}^\bullet = \mathcal{F}^n[-n]$,
alors $L(a')^*G
\otimes_{a^{-1}\mathcal{O}_S}^\mathbf{L} \mathcal{F}^\bullet =
La^*K \otimes_{\mathcal{O}_X}^\mathbf{L} \mathcal{F}^n[-n]$ ; voir
Cohomologie, lemme \ref{cohomology-lemma-variant-derived-pullback}.
Dans ce cas, la commutation aux sommes directes résulte donc du
lemme \ref{lemma-quasi-coherence-pushforward-direct-sums}.
Dans le cas général, puisque $S$ est affine (donc $X$ quasi-compact)
et que $\mathcal{F}^\bullet$ est localement borné, on voit que
$$
\mathcal{F}^\bullet = (\mathcal{F}^a \to \ldots \to \mathcal{F}^b)
$$
est borné. En raisonnant par récurrence sur $b - a$ et en considérant le
triangle distingué
$$
\mathcal{F}^b[-b] \to (\mathcal{F}^a \to \ldots \to \mathcal{F}^b)
\to (\mathcal{F}^a \to \ldots \to \mathcal{F}^{b - 1}) \to
\mathcal{F}^b[-b + 1]
$$
on achève la démonstration de cette partie. Certains détails sont omis.

\medskip\noindent
Démonstration de la partie (4). Soit $a' : (X, a^{-1}\mathcal{O}_S) \to (S, \mathcal{O}_S)$
comme ci-dessus. Puisque $\mathcal{F}^\bullet$ est un complexe localement borné
de $a^{-1}\mathcal{O}_S$-modules plats, on voit que le complexe
$a^{-1}\mathcal{G} \otimes_{a^{-1}\mathcal{O}_S} \mathcal{F}^\bullet$
représente $L(a')^*\mathcal{G}
\otimes_{a^{-1}\mathcal{O}_S}^\mathbf{L}
\mathcal{F}^\bullet$ dans $D(a^{-1}\mathcal{O}_S)$. L'assertion (4)
résulte donc de (5).
\end{proof}

\begin{lemma}
\label{lemma-K-flat}
Soit $f : X \to Y$ un morphisme de schémas, avec $Y = \Spec(A)$ affine.
Soit $\mathcal{U} : X = \bigcup_{i \in I} U_i$ un recouvrement ouvert affine fini
tel que toutes les intersections finies
$U_{i_0 \ldots i_p} = U_{i_0} \cap \ldots \cap U_{i_p}$
soient affines. Soit $\mathcal{F}^\bullet$ un complexe borné de
$f^{-1}\mathcal{O}_Y$-modules. Supposons que, pour tout $n \in \mathbf{Z}$,
le faisceau $\mathcal{F}^n$ soit un $f^{-1}\mathcal{O}_Y$-module plat et que
$\mathcal{F}^n$ possède une structure de $\mathcal{O}_X$-module quasi-cohérent
compatible avec la structure de $p^{-1}\mathcal{O}_Y$-module donnée (mais les
différentielles du complexe $\mathcal{F}^\bullet$ ne sont pas nécessairement
$\mathcal{O}_X$-linéaires). Alors le complexe
$\text{Tot}(\check{\mathcal{C}}^\bullet(\mathcal{U}, \mathcal{F}^\bullet))$
est K-plat comme complexe de $A$-modules.
\end{lemma}

\begin{proof}
On peut écrire
$$
\mathcal{F}^\bullet = (\mathcal{F}^a \to \ldots \to \mathcal{F}^b)
$$
En raisonnant par récurrence sur $b - a$ et en considérant le triangle distingué
$$
\mathcal{F}^b[-b] \to (\mathcal{F}^a \to \ldots \to \mathcal{F}^b)
\to (\mathcal{F}^a \to \ldots \to \mathcal{F}^{b - 1}) \to
\mathcal{F}^b[-b + 1]
$$
et en utilisant
Compléments sur l'algèbre, lemme \ref{more-algebra-lemma-K-flat-two-out-of-three},
on se ramène au cas où $\mathcal{F}^\bullet$ consiste en un seul
$\mathcal{O}_X$-module quasi-cohérent $\mathcal{F}$
placé en degré $0$. Dans ce cas, le complexe de {\v C}ech
de $\mathcal{F}$ relativement à $\mathcal{U}$ est homotopiquement équivalent au
complexe de {\v C}ech alterné ; voir
Cohomologie, lemme \ref{cohomology-lemma-alternating-usual}.
Comme $U_{i_0 \ldots i_p}$ est toujours affine, on voit que
$\mathcal{F}(U_{i_0 \ldots i_p})$ est plat sur $A$.
Par conséquent,
$\check{\mathcal{C}}_{alt}^\bullet(\mathcal{U}, \mathcal{F})$
est un complexe borné de $A$-modules plats, donc K-plat
d'après Compléments sur l'algèbre, lemme
\ref{more-algebra-lemma-derived-tor-quasi-isomorphism}.
\end{proof}

\noindent
Soient $X, Y, S, a, b, p, q, f$ comme dans l'introduction de cette section.
Soit $\mathcal{F}$ un $\mathcal{O}_X$-module.
Soit $\mathcal{G}$ un $\mathcal{O}_Y$-module.
Posons $A = \Gamma(S, \mathcal{O}_S)$.
Considérons le morphisme
\begin{equation}
\label{equation-kunneth-single-sheaves}
R\Gamma(X, \mathcal{F})
\otimes_A^\mathbf{L}
R\Gamma(Y, \mathcal{G})
\longrightarrow
R\Gamma(X \times_S Y,
p^*\mathcal{F}
\otimes_{\mathcal{O}_{X \times_S Y}}
q^*\mathcal{G})
\end{equation}
dans $D(A)$. Ce morphisme se construit à l'aide des morphismes de rappel
$R\Gamma(X, \mathcal{F}) \to R\Gamma(X \times_S Y, p^*\mathcal{F})$ et
$R\Gamma(Y, \mathcal{G}) \to R\Gamma(X \times_S Y, q^*\mathcal{G})$,
du produit cup construit dans
Cohomologie, section \ref{cohomology-section-cup-product}, et
du morphisme canonique
$p^*\mathcal{F}
\otimes_{\mathcal{O}_{X \times_S Y}}^\mathbf{L}
q^*\mathcal{G}
\to
p^*\mathcal{F}
\otimes_{\mathcal{O}_{X \times_S Y}}
q^*\mathcal{G}$.

\begin{lemma}
\label{lemma-kunneth-single-sheaf}
Dans la situation ci-dessus, le morphisme (\ref{equation-kunneth-single-sheaves}) est un
isomorphisme si $S$ est affine, si $\mathcal{F}$ et $\mathcal{G}$ sont plats sur $S$ et
quasi-cohérents, et si $X$ et $Y$ sont quasi-compacts à diagonale affine.
\end{lemma}

\begin{proof}
Nous recommandons vivement au lecteur de lire d'abord la démonstration de
Variétés, lemme \ref{varieties-lemma-kunneth}.
Choisissons des recouvrements ouverts affines finis
$\mathcal{U} : X = \bigcup_{i \in I} U_i$ et
$\mathcal{V} : Y = \bigcup_{j \in J} V_j$.
Ils déterminent un recouvrement ouvert affine
$\mathcal{W} : X \times_S Y = \bigcup_{(i, j) \in I \times J} U_i \times_S V_j$.
Notons que $\mathcal{W}$ raffine
$\text{pr}_1^{-1}\mathcal{U}$ et $\text{pr}_2^{-1}\mathcal{V}$.
Ainsi, d'après la discussion de Cohomologie, section
\ref{cohomology-section-cech-cohomology-of-complexes},
on obtient des morphismes
$$
\check{\mathcal{C}}^\bullet(\mathcal{U}, \mathcal{F})
\to
\check{\mathcal{C}}^\bullet(\mathcal{W}, p^*\mathcal{F})
\quad\text{et}\quad
\check{\mathcal{C}}^\bullet(\mathcal{V}, \mathcal{G})
\to
\check{\mathcal{C}}^\bullet(\mathcal{W}, q^*\mathcal{G})
$$
bien définis à homotopie près et compatibles avec les morphismes de rappel en cohomologie.
Dans Cohomologie, équation (\ref{cohomology-equation-needs-signs}),
nous avons construit un morphisme de complexes
$$
\text{Tot}(
\check{\mathcal{C}}^\bullet(\mathcal{W}, p^*\mathcal{F})
\otimes_A
\check{\mathcal{C}}^\bullet(\mathcal{W}, q^*\mathcal{G}))
\longrightarrow
\check{\mathcal{C}}^\bullet(\mathcal{W},
p^*\mathcal{F} \otimes_{\mathcal{O}_{X \times_S Y}}
q^*\mathcal{G})
$$
qui est compatible avec le produit cup en cohomologie d'après
Cohomologie, lemme \ref{cohomology-lemma-diagrams-commute}.
En combinant ce qui précède, on obtient un morphisme de complexes
\begin{equation}
\label{equation-kunneth-on-cech}
\text{Tot}(
\check{\mathcal{C}}^\bullet(\mathcal{U}, \mathcal{F})
\otimes_A
\check{\mathcal{C}}^\bullet(\mathcal{V}, \mathcal{G}))
\to
\check{\mathcal{C}}^\bullet(\mathcal{W},
p^*\mathcal{F}
\otimes_{\mathcal{O}_{X \times_S Y}}
q^*\mathcal{G})
\end{equation}
Nous affirmons qu'il s'agit du morphisme de l'énoncé du lemme, c'est-à-dire que
la source et le but de cette flèche s'identifient à la source
et au but de (\ref{equation-kunneth-single-sheaves}). En effet, d'après
Cohomologie des schémas, lemme
\ref{coherent-lemma-quasi-coherent-affine-cohomology-zero},
et
Cohomologie, lemme \ref{cohomology-lemma-cech-complex-complex-computes},
les morphismes canoniques
$$
\check{\mathcal{C}}^\bullet(\mathcal{U}, \mathcal{F})
\to
R\Gamma(X, \mathcal{F}),
\quad
\check{\mathcal{C}}^\bullet(\mathcal{V}, \mathcal{G})
\to
R\Gamma(Y, \mathcal{G})
$$
et
$$
\check{\mathcal{C}}^\bullet(\mathcal{W},
p^*\mathcal{F} \otimes_{\mathcal{O}_{X \times_S Y}} q^*\mathcal{G})
\to
R\Gamma(X \times_S Y,
p^*\mathcal{F} \otimes_{\mathcal{O}_{X \times_S Y}}
q^*\mathcal{G})
$$
sont des isomorphismes. D'autre part, le complexe
$\check{\mathcal{C}}^\bullet(\mathcal{U}, \mathcal{F})$
est K-plat d'après le lemme \ref{lemma-K-flat}, et on en déduit que
$\text{Tot}(
\check{\mathcal{C}}^\bullet(\mathcal{U}, \mathcal{F})
\otimes_A
\check{\mathcal{C}}^\bullet(\mathcal{V}, \mathcal{G}))$
représente le produit tensoriel dérivé
$R\Gamma(X, \mathcal{F}) \otimes_A^\mathbf{L} R\Gamma(Y, \mathcal{G})$,
comme annoncé.

\medskip\noindent
Il reste à montrer que (\ref{equation-kunneth-on-cech})
est un quasi-isomorphisme. Nous allons le faire par décalage de dimension.
Posons $d(\mathcal{F}) = \max \{d \mid H^d(X, \mathcal{F}) \not = 0\}$.
Supposons $d(\mathcal{F}) > 0$. Posons $U = \coprod\nolimits_{i \in I} U_i$.
C'est un schéma affine puisque $I$ est fini. Notons
$j : U \to X$ le morphisme qui est l'inclusion $U_i \to X$
sur chaque $U_i$. Comme la diagonale de $X$ est affine, le morphisme
$j$ est affine ; voir
Morphismes, lemme \ref{morphisms-lemma-affine-permanence}.
Il s'ensuit que $\mathcal{F}' = j_*j^*\mathcal{F}$ est plat sur $S$ ; voir
Morphismes, lemme \ref{morphisms-lemma-pushforward-flat-affine}.
Il s'ensuit aussi que $d(\mathcal{F}') = 0$ en combinant
Cohomologie des schémas, lemmes
\ref{coherent-lemma-relative-affine-cohomology} et
\ref{coherent-lemma-quasi-coherent-affine-cohomology-zero}.
Pour tout $x \in X$, le morphisme $\mathcal{F}_x  \to \mathcal{F}'_x$
est l'inclusion d'un facteur direct : si $x \in U_i$,
alors $\mathcal{F}' \to (U_i \to X)_*\mathcal{F}|_{U_i}$
fournit une rétraction. On en déduit que
$\mathcal{F} \to \mathcal{F}'$ est injectif et que
$\mathcal{F}'' = \mathcal{F}'/\mathcal{F}$
est lui aussi plat sur $S$. La suite exacte courte
$0 \to \mathcal{F} \to \mathcal{F}' \to \mathcal{F}'' \to 0$
de $\mathcal{O}_X$-modules quasi-cohérents plats
donne une suite exacte courte de complexes
$$
\begin{aligned}
0 &\to
\text{Tot}(
\check{\mathcal{C}}^\bullet(\mathcal{U}, \mathcal{F})
\otimes_A
\check{\mathcal{C}}^\bullet(\mathcal{V}, \mathcal{G})) \\
&\to
\text{Tot}(
\check{\mathcal{C}}^\bullet(\mathcal{U}, \mathcal{F}')
\otimes_A
\check{\mathcal{C}}^\bullet(\mathcal{V}, \mathcal{G})) \\
&\to
\text{Tot}(
\check{\mathcal{C}}^\bullet(\mathcal{U}, \mathcal{F}'')
\otimes_A
\check{\mathcal{C}}^\bullet(\mathcal{V}, \mathcal{G})) \to 0
\end{aligned}
$$
et une suite exacte courte de complexes
$$
\begin{aligned}
0 &\to
\check{\mathcal{C}}^\bullet(\mathcal{W},
p^*\mathcal{F}
\otimes_{\mathcal{O}_{X \times_S Y}}
q^*\mathcal{G}) \\
&\to
\check{\mathcal{C}}^\bullet(\mathcal{W},
p^*\mathcal{F}'
\otimes_{\mathcal{O}_{X \times_S Y}}
q^*\mathcal{G}) \\
&\to
\check{\mathcal{C}}^\bullet(\mathcal{W},
p^*\mathcal{F}''
\otimes_{\mathcal{O}_{X \times_S Y}}
q^*\mathcal{G}) \to 0
\end{aligned}
$$
De plus, les morphismes (\ref{equation-kunneth-on-cech}) entre celles-ci
sont compatibles avec ces suites exactes courtes. Il suffit donc de montrer que
(\ref{equation-kunneth-on-cech})
est un isomorphisme pour $\mathcal{F}'$ et $\mathcal{F}''$. Enfin,
on a $d(\mathcal{F}'') < d(\mathcal{F})$.
On se ramène ainsi au cas $d(\mathcal{F}) = 0$.

\medskip\noindent
En raisonnant de la même manière pour $\mathcal{G}$, on voit que l'on
peut supposer que $\mathcal{F}$ et $\mathcal{G}$
n'ont de cohomologie non nulle qu'en degré $0$.
Observons que cela signifie que $\Gamma(X, \mathcal{F})$
est quasi-isomorphe au complexe $K$-plat
$\check{\mathcal{C}}^\bullet(\mathcal{U}, \mathcal{F})$
de $A$-modules situé en degrés $\geq 0$.
Il s'ensuit que $\Gamma(X, \mathcal{F})$ est un $A$-module plat
(car on peut calculer les Tor supérieurs contre ce module
en tensorisant par le complexe de Cech).
Soit $V \subset Y$ un ouvert affine. Considérons le recouvrement ouvert affine
$\mathcal{U}_V : X \times_S V = \bigcup_{i \in I} U_i \times_S V$.
Il est immédiat que
$$
\check{\mathcal{C}}^\bullet(\mathcal{U}, \mathcal{F})
\otimes_A \mathcal{G}(V) =
\check{\mathcal{C}}^\bullet(\mathcal{U}_V,
p^*\mathcal{F} \otimes_{\mathcal{O}_{X \times Y}}
q^*\mathcal{G})
$$
(égalité de complexes). Comme $\mathcal{G}(V)$ est plat
sur $A$, on voit que
$\Gamma(X, \mathcal{F}) \otimes_A \mathcal{G}(V) \to
\check{\mathcal{C}}^\bullet(\mathcal{U}, \mathcal{F})
\otimes_A \mathcal{G}(V)$ est un quasi-isomorphisme.
Puisque le faisceau associé à
$V \mapsto \check{\mathcal{C}}^\bullet(\mathcal{U}_V,
p^*\mathcal{F} \otimes_{\mathcal{O}_{X \times Y}}
q^*\mathcal{G})$ représente
$Rq_*(p^*\mathcal{F} \otimes_{\mathcal{O}_{X \times Y}} q^*\mathcal{G})$
d'après Cohomologie des schémas, lemme
\ref{coherent-lemma-separated-case-relative-cech},
on en déduit que
$$
Rq_*(p^*\mathcal{F} \otimes_{\mathcal{O}_{X \times Y}} q^*\mathcal{G})
\cong
\Gamma(X, \mathcal{F}) \otimes_A \mathcal{G}
$$
sur $Y$, où la notation du membre de droite désigne le module
$$
b^*\widetilde{\Gamma(X, \mathcal{F})} \otimes_{\mathcal{O}_Y} \mathcal{G}
$$
En utilisant la suite spectrale de Leray pour $q$, on obtient
$$
H^n(X \times_S Y, p^*\mathcal{F} \otimes_{\mathcal{O}_{X \times Y}}
q^*\mathcal{G}) =
H^n(Y,
b^*\widetilde{\Gamma(X, \mathcal{F})} \otimes_{\mathcal{O}_Y} \mathcal{G})
$$
En appliquant le lemme \ref{lemma-cohomology-base-change} au morphisme
$b : Y \to S = \Spec(A)$ et en utilisant le fait que $\Gamma(X, \mathcal{F})$
est plat sur $A$, on conclut que
$H^n(X \times_S Y, p^*\mathcal{F} \otimes_{\mathcal{O}_{X \times Y}}
q^*\mathcal{G})$ est nul pour $n > 0$ et isomorphe à
$H^0(X, \mathcal{F}) \otimes_A H^0(Y, \mathcal{G})$ pour $n = 0$.
Bien entendu, nous utilisons ici aussi le fait que $\mathcal{G}$ n'a de
cohomologie qu'en degré $0$.
Cela achève la démonstration (à ceci près qu'il faudrait vérifier que
l'isomorphisme est bien donné par le produit cup en degré $0$ ; nous omettons
cette vérification).
\end{proof}

\begin{remark}
\label{remark-annoying-compatibility}
Soit $S = \Spec(A)$ un schéma affine. Soient $a : X \to S$
et $b : Y \to S$ des morphismes de schémas. Soient $\mathcal{F}$ et $\mathcal{G}$
des $\mathcal{O}_X$-modules quasi-cohérents, et soit $\mathcal{E}$
un $\mathcal{O}_Y$-module quasi-cohérent. Soit
$\xi \in H^i(X, \mathcal{G})$, dont le rappel est
$p^*\xi \in H^i(X \times_S Y, p^*\mathcal{G})$.
Alors le diagramme suivant est commutatif
$$
\xymatrix{
R\Gamma(X, \mathcal{F})[-i] \otimes_A^\mathbf{L} R\Gamma(Y, \mathcal{E})
\ar[d] \ar[rr]_-{\xi \otimes \text{id}} & &
R\Gamma(X, \mathcal{G} \otimes_{\mathcal{O}_X} \mathcal{F})
\otimes_A^\mathbf{L} R\Gamma(Y, \mathcal{E}) \ar[d] \\
R\Gamma(X \times_S Y, p^*\mathcal{F} \otimes q^*\mathcal{E})[-i]
\ar[rr]^-{p^*\xi} & &
R\Gamma(X \times_S Y,
p^*(\mathcal{G} \otimes_{\mathcal{O}_X} \mathcal{F}) \otimes q^*\mathcal{E})
}
$$
où les produits tensoriels non décorés sont pris sur $\mathcal{O}_{X \times_S Y}$.
Les flèches horizontales proviennent de Cohomologie, remarque
\ref{cohomology-remark-cup-with-element-map-total-cohomology},
et les flèches verticales sont (\ref{equation-kunneth-global}) ; elles sont donc
données par le rappel suivi du produit cup sur $X \times_S Y$.
Le diagramme commute parce que le produit cup global (sur $X \times_S Y$
avec les faisceaux $p^*\mathcal{G}$, $p^*\mathcal{F}$ et $q^*\mathcal{E}$)
est associatif ; voir
Cohomologie, lemme \ref{cohomology-lemma-cup-product-associative}.
\end{remark}


\section{Formule de K\"unneth, III}
\label{section-kunneth-complexes}

\noindent
Soient $X, Y, S, a, b, p, q, f$ comme dans l'introduction de la
section \ref{section-kunneth}. Dans cette section, étant donnés un
$\mathcal{O}_X$-module $\mathcal{F}$ et un $\mathcal{O}_Y$-module
$\mathcal{G}$, posons
$$
\mathcal{F} \boxtimes \mathcal{G} =
p^*\mathcal{F} \otimes_{\mathcal{O}_{X \times_S Y}} q^*\mathcal{G}
$$
Notons que, contrairement à ce qui se produira dans une section ultérieure, nous prenons
ici le produit tensoriel non dérivé.

\medskip\noindent
Sur $X$, soit $\mathcal{F}^\bullet$ un complexe de faisceaux de groupes abéliens
dont les termes sont des $\mathcal{O}_X$-modules quasi-cohérents
et tel que les différentielles
$d^i_\mathcal{F} : \mathcal{F}^i \to \mathcal{F}^{i + 1}$
soient des opérateurs différentiels sur $X/S$ d'ordre fini ; voir
Morphismes, section \ref{morphisms-section-differential-operators}.
De même, sur $Y$, soit $\mathcal{G}^\bullet$ un complexe de faisceaux
de groupes abéliens dont les termes sont des $\mathcal{O}_Y$-modules quasi-cohérents
et tel que les différentielles
$d^j_\mathcal{G} : \mathcal{G}^j \to \mathcal{G}^{j + 1}$
soient des opérateurs différentiels sur $Y/S$ d'ordre fini.
En appliquant la construction de
Morphismes, lemme \ref{morphisms-lemma-base-change-differential-operators},
on obtient un complexe double
$$
\xymatrix{
\ldots &
\ldots &
\ldots &
\ldots \\
\ldots \ar[r] &
\mathcal{F}^i \boxtimes \mathcal{G}^{j + 1}
\ar[r]^{d_1^{i, j + 1}} \ar[u] &
\mathcal{F}^{i + 1} \boxtimes \mathcal{G}^{j + 1} \ar[r] \ar[u] &
\ldots \\
\ldots \ar[r] &
\mathcal{F}^i \boxtimes \mathcal{G}^j
\ar[r]^{d_1^{i, j}} \ar[u]^{d_2^{i, j}} &
\mathcal{F}^{i + 1} \boxtimes \mathcal{G}^j \ar[r] \ar[u]_{d_2^{i + 1, j}} &
\ldots \\
\ldots &
\ldots \ar[u] &
\ldots \ar[u] &
\ldots
}
$$
de modules quasi-cohérents dont les morphismes sont des opérateurs différentiels
d'ordre fini sur $X \times_S Y / S$. Voir la discussion de
Morphismes, remarque \ref{morphisms-remark-base-change-differential-operators},
et
Homologie, exemple \ref{homology-example-double-complex-as-tensor-product-of}.
Explicitement, on pose
$$
d_1^{i, j} = d^i_\mathcal{F} \boxtimes 1
\quad\text{et}\quad
d_2^{i, j} = 1 \boxtimes d^j_\mathcal{G}
$$
Dans la discussion ci-dessous, la notation
$$
\text{Tot}(\mathcal{F}^\bullet \boxtimes \mathcal{G}^\bullet)
$$
désigne le complexe total associé à ce complexe double.
Les termes de ce complexe sont des
$\mathcal{O}_{X \times_S Y}$-modules quasi-cohérents, et ses différentielles
sont des opérateurs différentiels d'ordre fini sur $X \times_S Y / S$.

\medskip\noindent
Dans la situation ci-dessus, il existe un morphisme de « produit cup relatif »
\begin{equation}
\label{equation-relative-de-rham-kunneth}
Ra_*(\mathcal{F}^\bullet)
\otimes_{\mathcal{O}_S}^\mathbf{L}
Rb_*(\mathcal{G}^\bullet)
\longrightarrow
Rf_*\left(\text{Tot}(\mathcal{F}^\bullet \boxtimes \mathcal{G}^\bullet)\right)
\end{equation}
On peut en effet construire ce morphisme en combinant
\begin{enumerate}
\item $Ra_*(\mathcal{F}^\bullet) \to Rf_*(p^{-1}\mathcal{F}^\bullet)$,
\item $Rb_*(\mathcal{G}^\bullet) \to Rf_*(q^{-1}\mathcal{G}^\bullet)$,
\item
$Rf_*(p^{-1}\mathcal{F}^\bullet)
\otimes_{\mathcal{O}_S}^\mathbf{L}
Rf_*(q^{-1}\mathcal{G}^\bullet)
\to
Rf_*(p^{-1}\mathcal{F}^\bullet \otimes_{f^{-1}\mathcal{O}_S}^\mathbf{L}
q^{-1}\mathcal{G}^\bullet)$,
\item
$p^{-1}\mathcal{F}^\bullet \otimes_{f^{-1}\mathcal{O}_S}^\mathbf{L}
q^{-1}\mathcal{G}^\bullet \to
\text{Tot}(p^{-1}\mathcal{F}^\bullet \otimes_{f^{-1}\mathcal{O}_S}
q^{-1}\mathcal{G}^\bullet)$
\item
$\text{Tot}(p^{-1}\mathcal{F}^\bullet \otimes_{f^{-1}\mathcal{O}_S}
q^{-1}\mathcal{G}^\bullet) \to \text{Tot}(\mathcal{F}^\bullet \boxtimes
\mathcal{G}^\bullet)$.
\end{enumerate}
Les morphismes (1) et (2) sont des morphismes de rappel, le morphisme (3) est
le produit cup relatif ; voir
Cohomologie, remarque \ref{cohomology-remark-cup-product} ;
le morphisme (4) compare les produits tensoriels dérivé et non dérivé, et
le morphisme (5) est donné par les morphismes évidents
$p^{-1}\mathcal{F}^i \otimes_{f^{-1}\mathcal{O}_S} q^{-1}\mathcal{G}^j
\to \mathcal{F}^i \boxtimes \mathcal{G}^j$ sur les complexes doubles sous-jacents.

\medskip\noindent
Posons $A = \Gamma(S, \mathcal{O}_S)$. Il existe un morphisme de « produit cup global »
\begin{equation}
\label{equation-de-rham-kunneth}
R\Gamma(X, \mathcal{F}^\bullet)
\otimes_A^\mathbf{L}
R\Gamma(Y, \mathcal{G}^\bullet)
\longrightarrow
R\Gamma(X \times_S Y,
\text{Tot}(\mathcal{F}^\bullet \boxtimes \mathcal{G}^\bullet))
\end{equation}
dans $D(A)$. Ce morphisme se construit comme le produit cup relatif ci-dessus,
en utilisant
\begin{enumerate}
\item $R\Gamma(X, \mathcal{F}^\bullet) \to
R\Gamma(X \times_S Y, p^{-1}\mathcal{F}^\bullet)$
\item $R\Gamma(Y, \mathcal{G}^\bullet) \to
R\Gamma(X \times_S Y, q^{-1}\mathcal{G}^\bullet)$,
\item $R\Gamma(X \times_S Y, p^{-1}\mathcal{F}^\bullet)
\otimes_A^\mathbf{L} R\Gamma(X \times_S Y, q^{-1}\mathcal{G}^\bullet) \to
R\Gamma(X \times_S Y,
p^{-1}\mathcal{F}^\bullet \otimes_{f^{-1}\mathcal{O}_S}^\mathbf{L}
q^{-1}\mathcal{G}^\bullet)$,
\item
$p^{-1}\mathcal{F}^\bullet \otimes_{f^{-1}\mathcal{O}_S}^\mathbf{L}
q^{-1}\mathcal{G}^\bullet \to
\text{Tot}(p^{-1}\mathcal{F}^\bullet \otimes_{f^{-1}\mathcal{O}_S}
q^{-1}\mathcal{G}^\bullet)$
\item
$\text{Tot}(p^{-1}\mathcal{F}^\bullet \otimes_{f^{-1}\mathcal{O}_S}
q^{-1}\mathcal{G}^\bullet) \to \text{Tot}(\mathcal{F}^\bullet \boxtimes
\mathcal{G}^\bullet)$.
\end{enumerate}
Ici, les morphismes (1) et (2) sont les morphismes de rappel, et le morphisme (3)
est le produit cup construit dans Cohomologie, section \ref{cohomology-section-cup-product}.
Les morphismes (4) et (5) sont ceux indiqués au paragraphe précédent.

\begin{lemma}
\label{lemma-kunneth-special}
Dans la situation ci-dessus, le produit cup (\ref{equation-de-rham-kunneth})
est un isomorphisme dans $D(A)$ si les hypothèses suivantes sont satisfaites :
\begin{enumerate}
\item $S = \Spec(A)$ est affine,
\item $X$ et $Y$ sont quasi-compacts à diagonale affine,
\item $\mathcal{F}^\bullet$ est borné,
\item $\mathcal{G}^\bullet$ est borné inférieurement,
\item $\mathcal{F}^n$ est plat sur $S$, et
\item $\mathcal{G}^m$ est plat sur $S$.
\end{enumerate}
\end{lemma}

\begin{proof}
Nous utiliserons les notations $\mathcal{A}_{X/S}$ et $\mathcal{A}_{Y/S}$
introduites dans Morphismes, remarque
\ref{morphisms-remark-base-change-differential-operators}.
Supposons que l'on ait des morphismes de complexes
$$
\mathcal{F}_1^\bullet \to
\mathcal{F}_2^\bullet \to
\mathcal{F}_3^\bullet \to
\mathcal{F}_1^\bullet[1]
$$
dans la catégorie $\mathcal{A}_{X/S}$.
La fonctorialité du produit cup fournit alors un diagramme commutatif
$$
\xymatrix{
R\Gamma(X, \mathcal{F}_1^\bullet)
\otimes_A^\mathbf{L}
R\Gamma(Y, \mathcal{G}^\bullet)
\ar[r] \ar[d] &
R\Gamma(X \times_S Y,
\text{Tot}(\mathcal{F}_1^\bullet \boxtimes \mathcal{G}^\bullet)) \ar[d] \\
R\Gamma(X, \mathcal{F}_2^\bullet)
\otimes_A^\mathbf{L}
R\Gamma(Y, \mathcal{G}^\bullet)
\ar[r] \ar[d] &
R\Gamma(X \times_S Y,
\text{Tot}(\mathcal{F}_2^\bullet \boxtimes \mathcal{G}^\bullet)) \ar[d] \\
R\Gamma(X, \mathcal{F}_3^\bullet)
\otimes_A^\mathbf{L}
R\Gamma(Y, \mathcal{G}^\bullet)
\ar[r] \ar[d] &
R\Gamma(X \times_S Y,
\text{Tot}(\mathcal{F}_3^\bullet \boxtimes \mathcal{G}^\bullet)) \ar[d] \\
R\Gamma(X, \mathcal{F}_1^\bullet[1])
\otimes_A^\mathbf{L}
R\Gamma(Y, \mathcal{G}^\bullet)
\ar[r] &
R\Gamma(X \times_S Y,
\text{Tot}(\mathcal{F}_1^\bullet[1] \boxtimes \mathcal{G}^\bullet))
}
$$
Si les morphismes initiaux forment un triangle distingué dans la catégorie homotopique
de $\mathcal{A}_{X/S}$, alors
les colonnes de ce diagramme forment des triangles distingués dans $D(A)$.

\medskip\noindent
Dans la situation du lemme,
supposons $\mathcal{F}^n = 0$ pour $n < i$. On peut alors considérer la
suite exacte courte de complexes, scindée terme à terme,
$$
0 \to \sigma_{\geq i + 1}\mathcal{F}^\bullet \to
\mathcal{F}^\bullet \to \mathcal{F}^i[-i] \to 0
$$
où la troncature est celle de
Homologie, section \ref{homology-section-truncations}.
Elle donne le triangle distingué
$$
\sigma_{\geq i + 1}\mathcal{F}^\bullet \to
\mathcal{F}^\bullet \to
\mathcal{F}^i[-i] \to
(\sigma_{\geq i + 1}\mathcal{F}^\bullet)[1]
$$
dans la catégorie homotopique de $\mathcal{A}_{X/S}$,
où la dernière flèche est donnée par le morphisme de bord
$\mathcal{F}^i \to \mathcal{F}^{i + 1}$.
La discussion qui précède montre qu'il suffit de démontrer le lemme pour
$\mathcal{F}^i[-i]$ et $\sigma_{\geq i + 1}\mathcal{F}^\bullet$.
Puisque $\sigma_{\geq i + 1}\mathcal{F}^\bullet$ a moins de termes non nuls,
une récurrence montre que, si l'on peut démontrer le lemme lorsque $\mathcal{F}^\bullet$
n'est non nul qu'en un seul degré, alors le lemme en résulte.
On peut donc supposer que $\mathcal{F}^\bullet$ n'est non nul qu'en un seul degré.

\medskip\noindent
Supposons que $\mathcal{F}^\bullet$ soit le complexe ayant pour terme, plat sur $S$,
le $\mathcal{O}_X$-module quasi-cohérent $\mathcal{F}$ placé en degré $0$,
et nul dans les autres degrés. Observons que $R\Gamma(X, \mathcal{F})$ est de
dimension de Tor finie, par exemple d'après le lemme \ref{lemma-cohomology-de-rham-base-change}.
Disons que son amplitude de Tor est contenue dans $[i, j]$.
Choisissons $N \gg 0$ et considérons le triangle distingué
$$
\sigma_{\geq N + 1}\mathcal{G}^\bullet \to
\mathcal{G}^\bullet \to
\sigma_{\leq N}\mathcal{G}^\bullet \to
(\sigma_{\geq N + 1}\mathcal{G}^\bullet)[1]
$$
dans la catégorie homotopique de $\mathcal{A}_{Y/S}$. Observons maintenant que les deux complexes
$$
R\Gamma(X, \mathcal{F})
\otimes_A^\mathbf{L}
R\Gamma(Y, \sigma_{\geq N + 1}\mathcal{G}^\bullet)
\quad\text{et}\quad
R\Gamma(X \times_S Y,
\text{Tot}(\mathcal{F} \boxtimes \sigma_{\geq N + 1}\mathcal{G}^\bullet))
$$
ont une cohomologie nulle en degrés $\leq N + i$. Ainsi, en utilisant les arguments
donnés ci-dessus, si l'on veut démontrer notre assertion en un degré donné, on peut
supposer que $\mathcal{G}^\bullet$ est borné.
En répétant encore une fois les arguments précédents, on peut aussi supposer
que $\mathcal{G}^\bullet$ n'est non nul qu'en un seul degré.
Ce cas est traité par le lemme \ref{lemma-kunneth-single-sheaf}.
\end{proof}






\section{Formule de K\"unneth pour Ext}
\label{section-kunneth-Ext}

\noindent
Considérons un diagramme cartésien de schémas
$$
\xymatrix{
& X \times_S Y \ar[ld]^p \ar[rd]_q \ar[dd]^f \\
X \ar[rd]_a & & Y \ar[ld]^b \\
& S
}
$$
Pour $K \in D(\mathcal{O}_X)$ et $M \in D(\mathcal{O}_Y)$,
posons dans cette section
$$
K \boxtimes M =
Lp^*K \otimes_{\mathcal{O}_{X \times_S Y}}^\mathbf{L} Lq^*M
$$
Nous affirmons qu'il existe un morphisme canonique
\begin{equation}
\label{equation-kunneth-ext}
Ra_*R\SheafHom(K, K')
\otimes_{\mathcal{O}_S}^\mathbf{L}
Rb_*R\SheafHom(M, M')
\longrightarrow
Rf_*(R\SheafHom(K \boxtimes M, K' \boxtimes M'))
\end{equation}
pour $K, K' \in D(\mathcal{O}_X)$ et $M, M' \in D(\mathcal{O}_Y)$.
On peut en effet prendre le morphisme adjoint au morphisme
$$
\begin{matrix}
Lf^*\left(Ra_*R\SheafHom(K, K')
\otimes_{\mathcal{O}_S}^\mathbf{L}
Rb_* R\SheafHom(M, M')\right) = \\
Lf^* Ra_* R\SheafHom(K, K')
\otimes_{\mathcal{O}_{X \times_S Y}}^\mathbf{L}
Lf^* Rb_* R\SheafHom(M, M') = \\
Lp^* La^* Ra_* R\SheafHom(K, K')
\otimes_{\mathcal{O}_{X \times_S Y}}^\mathbf{L}
Lq^* Lb^* Rb_* R\SheafHom(M, M') \to \\
Lp^* R\SheafHom(K, K')
\otimes_{\mathcal{O}_{X \times_S Y}}^\mathbf{L}
Lq^* R\SheafHom(M, M') \to \\
R\SheafHom(Lp^*K, Lp^*K')
\otimes_{\mathcal{O}_{X \times_S Y}}^\mathbf{L}
R\SheafHom(Lq^*M, Lq^*M') \to \\
R\SheafHom(K \boxtimes M, K' \boxtimes M')
\end{matrix}
$$
Ici, la première égalité exprime la compatibilité des rappels avec les produits tensoriels,
Cohomologie, lemme \ref{cohomology-lemma-pullback-tensor-product}.
La deuxième égalité utilise $f = a \circ p = b \circ q$ et
la composition des rappels,
Cohomologie, lemme \ref{cohomology-lemma-derived-pullback-composition}.
La première flèche est donnée par les morphismes d'adjonction
$La^* Ra_* \to \text{id}$ et
$Lb^* Rb_* \to \text{id}$, car l'image directe et le rappel sont adjoints,
Cohomologie, lemme \ref{cohomology-lemma-adjoint}.
La deuxième flèche est donnée par
Cohomologie, remarque \ref{cohomology-remark-prepare-fancy-base-change}.
La troisième et dernière flèche est celle de
Cohomologie, remarque \ref{cohomology-remark-tensor-internal-hom}.
Le résultat suivant en est un cas particulier simple.

\begin{lemma}
\label{lemma-kunneth-Ext}
Dans la situation ci-dessus, supposons que $a$ et $b$ soient quasi-compacts et
quasi-séparés et que $X$ et $Y$ soient Tor-indépendants sur $S$.
Si $K$ est parfait, $K' \in D_\QCoh(\mathcal{O}_X)$, $M$ est parfait, et si
$M' \in D_\QCoh(\mathcal{O}_Y)$, alors (\ref{equation-kunneth-ext})
est un isomorphisme.
\end{lemma}

\begin{proof}
Dans ce cas, on a $R\SheafHom(K, K') = K' \otimes^\mathbf{L} K^\vee$,
$R\SheafHom(M, M') = M' \otimes^\mathbf{L} M^\vee$, et
$$
R\SheafHom(K \boxtimes M, K' \boxtimes M') =
(K' \otimes^\mathbf{L} K^\vee) \boxtimes
(M' \otimes^\mathbf{L} M^\vee)
$$
Voir Cohomologie, lemme \ref{cohomology-lemma-dual-perfect-complex} ;
on utilise également le fait que la perfection est préservée par rappel
et par produit tensoriel.
Ce cas résulte donc du lemme \ref{lemma-kunneth}.
(Nous omettons de vérifier que ces identifications
donnent le même morphisme.)
\end{proof}







\section{Cohomologie et changement de base, V}
\label{section-cohomology-base-change}

\noindent
Dans la section \ref{section-cohomology-and-base-change-perfect},
nous avons vu qu'un théorème de changement de base vaut lorsque les morphismes sont Tor-indépendants.
Même dans le cas affine, il ne peut exister de théorème de changement de base sans une telle
condition ; voir
Compléments sur l'algèbre, section \ref{more-algebra-section-tor-independence}.
Dans cette section, nous analysons quand on peut obtenir un résultat de changement de base
« un complexe à la fois ».

\medskip\noindent
Pour que cela fonctionne, supposons donné un diagramme commutatif
$$
\xymatrix{
X' \ar[r]_{g'} \ar[d]_{f'} &
X \ar[d]^f \\
S' \ar[r]^g &
S
}
$$
de schémas (en général, nous supposerons qu'il est cartésien).
Soit $K \in D_\QCoh(\mathcal{O}_X)$,
et soit $L(g')^*K \to K'$ un morphisme dans $D_\QCoh(\mathcal{O}_{X'})$.
Pour un point $x' \in X'$, posons $x = g'(x') \in X$,
$s' = f'(x') \in S'$ et $s = f(x) = g(s')$.
On peut alors considérer les morphismes
$$
K_x \otimes_{\mathcal{O}_{S, s}}^\mathbf{L} \mathcal{O}_{S', s'} \to
K_x \otimes_{\mathcal{O}_{X, x}}^\mathbf{L} \mathcal{O}_{X', x'} \to
K'_{x'}
$$
où la première flèche est celle de Compléments sur l'algèbre,
équation (\ref{more-algebra-equation-comparison-map}),
et la seconde provient de
$(L(g')^*K)_{x'} =
K_x \otimes_{\mathcal{O}_{X, x}}^\mathbf{L} \mathcal{O}_{X', x'}$
et du morphisme donné $L(g')^*K \to K'$.
Pour chaque $i \in \mathbf{Z}$, on obtient une structure de
$\mathcal{O}_{X, x} \otimes_{\mathcal{O}_{S, s}} \mathcal{O}_{S', s'}$-module
sur
$H^i(K_x \otimes_{\mathcal{O}_{S, s}}^\mathbf{L} \mathcal{O}_{S', s'})$.
En rassemblant le tout, on obtient des morphismes canoniques
\begin{equation}
\label{equation-bc}
H^i(K_x \otimes_{\mathcal{O}_{S, s}}^\mathbf{L} \mathcal{O}_{S', s'})
\otimes_{(\mathcal{O}_{X, x} \otimes_{\mathcal{O}_{S, s}} \mathcal{O}_{S', s'})}
\mathcal{O}_{X', x'}
\longrightarrow
H^i(K'_{x'})
\end{equation}
de $\mathcal{O}_{X', x'}$-modules.

\begin{lemma}
\label{lemma-single-complex-base-change-condition}
Soit
$$
\xymatrix{
X' \ar[r]_{g'} \ar[d]_{f'} &
X \ar[d]^f \\
S' \ar[r]^g &
S
}
$$
un diagramme cartésien de schémas. Soit $K \in D_\QCoh(\mathcal{O}_X)$,
et soit $L(g')^*K \to K'$ un morphisme dans $D_\QCoh(\mathcal{O}_{X'})$.
Les assertions suivantes sont équivalentes :
\begin{enumerate}
\item pour tous $x' \in X'$ et $i \in \mathbf{Z}$, le morphisme (\ref{equation-bc})
est un isomorphisme,
\item si $U \subset X$ et $V' \subset S'$ sont des ouverts affines dont les images sont
contenues dans l'ouvert affine $V \subset S$, et si $U' = V' \times_V U$,
la composée
$$
R\Gamma(U, K) \otimes_{\mathcal{O}_S(U)}^\mathbf{L} \mathcal{O}_{S'}(V')
\to
R\Gamma(U, K) \otimes_{\mathcal{O}_X(U)}^\mathbf{L} \mathcal{O}_{X'}(U')
\to
R\Gamma(U', K')
$$
est un isomorphisme dans $D(\mathcal{O}_{S'}(V'))$, et
\item il existe un ensemble $I$ de quadruplets $U_i, V_i', V_i, U_i'$, $i \in I$,
comme en (2), tels que $X' = \bigcup U'_i$.
\end{enumerate}
\end{lemma}

\begin{proof}
La deuxième flèche de (2) provient de l'égalité
$$
R\Gamma(U, K) \otimes_{\mathcal{O}_X(U)}^\mathbf{L} \mathcal{O}_{X'}(U') =
R\Gamma(U', L(g')^*K)
$$
du lemme \ref{lemma-quasi-coherence-pullback} et du morphisme donné
$L(g')^*K \to K'$. La première flèche de (2) est celle de
Compléments sur l'algèbre, équation (\ref{more-algebra-equation-comparison-map}).
Il est clair que (2) implique (3). Observons que (1) est local sur $X'$.
Il suffit donc de montrer que, si $X$, $S$, $S'$ et $X'$ sont affines, alors
(1) équivaut à la condition que
$$
R\Gamma(X, K) \otimes_{\mathcal{O}_S(S)}^\mathbf{L} \mathcal{O}_{S'}(S')
\to
R\Gamma(X, K) \otimes_{\mathcal{O}_X(X)}^\mathbf{L} \mathcal{O}_{X'}(X')
\to
R\Gamma(X', K')
$$
soit un isomorphisme dans $D(\mathcal{O}_{S'}(S'))$. Écrivons
$S = \Spec(R)$, $X = \Spec(A)$, $S' = \Spec(R')$, $X' = \Spec(A')$,
et supposons que $K$ corresponde au complexe $M^\bullet$ de $A$-modules, et que
$K'$ corresponde au complexe $N^\bullet$ de $A'$-modules.
Notons que $A' = A \otimes_R R'$. La condition ci-dessus signifie que la composée
$$
M^\bullet \otimes_R^\mathbf{L} R' \to
M^\bullet \otimes_A^\mathbf{L} A' \to
N^\bullet
$$
est un isomorphisme dans $D(R')$. De manière équivalente, pour tout
$i \in \mathbf{Z}$, le morphisme
$$
H^i(M^\bullet \otimes_R^\mathbf{L} R') \to
H^i(M^\bullet \otimes_A^\mathbf{L} A') \to
H^i(N^\bullet)
$$
est un isomorphisme. Observons qu'il s'agit d'un morphisme de $A \otimes_R R'$-modules,
c'est-à-dire de $A'$-modules. D'autre part, (1) exige
que, pour des idéaux premiers compatibles
$\mathfrak q' \subset A'$, $\mathfrak q \subset A$,
$\mathfrak p' \subset R'$, $\mathfrak p \subset R$,
la composée
$$
H^i(M^\bullet_\mathfrak q \otimes_{R_\mathfrak p}^\mathbf{L} R'_{\mathfrak p'})
\otimes_{(A_\mathfrak q \otimes_{R_\mathfrak p} R'_{\mathfrak p'})}
A'_{\mathfrak q'} \to
H^i(M^\bullet_{\mathfrak q}
\otimes_{A_\mathfrak q}^\mathbf{L} A'_{\mathfrak q'})
\to H^i(N^\bullet_{\mathfrak q'})
$$
soit un isomorphisme. Comme
$$
H^i(M^\bullet_\mathfrak q \otimes_{R_\mathfrak p}^\mathbf{L} R'_{\mathfrak p'})
\otimes_{(A_\mathfrak q \otimes_{R_\mathfrak p} R'_{\mathfrak p'})}
A'_{\mathfrak q'} =
H^i(M^\bullet \otimes_R^\mathbf{L} R') \otimes_{A'} A'_{\mathfrak q'}
$$
est la localisation en $\mathfrak q'$,
on voit que ces deux conditions sont équivalentes d'après
Algèbre, lemme \ref{algebra-lemma-characterize-zero-local}.
\end{proof}

\begin{lemma}
\label{lemma-single-complex-base-change}
Soit
$$
\xymatrix{
X' \ar[r]_{g'} \ar[d]_{f'} &
X \ar[d]^f \\
S' \ar[r]^g &
S
}
$$
un diagramme cartésien de schémas. Soit $K \in D_\QCoh(\mathcal{O}_X)$,
et soit $L(g')^*K \to K'$ un morphisme dans $D_\QCoh(\mathcal{O}_{X'})$.
Si
\begin{enumerate}
\item les conditions équivalentes du
lemme \ref{lemma-single-complex-base-change-condition} sont satisfaites, et
\item $f$ est quasi-compact et quasi-séparé,
\end{enumerate}
alors la composée $Lg^*Rf_*K \to Rf'_*L(g')^*K \to Rf'_*K'$
est un isomorphisme.
\end{lemma}

\begin{proof}
On pourrait le démontrer par la même méthode que dans la démonstration du
lemme \ref{lemma-compare-base-change}, mais nous allons plutôt utiliser
le principe d'induction et Mayer--Vietoris relatif.

\medskip\noindent
Pour vérifier que le morphisme est un isomorphisme, on peut travailler localement sur $S'$.
On peut donc supposer que $g : S' \to S$ est un morphisme de schémas affines.
En particulier, $X$ est un schéma quasi-compact et quasi-séparé.
Nous utiliserons le principe d'induction de
Cohomologie des schémas, lemme \ref{coherent-lemma-induction-principle},
pour démontrer que, pour tout ouvert quasi-compact $U \subset X$, le morphisme construit
de la même manière $Lg^*R(U \to S)_*K|_U \to R(U' \to S')_*K'|_{U'}$
est un isomorphisme. Ici $U' = (g')^{-1}(U)$.

\medskip\noindent
Si $U \subset X$ est un ouvert affine, alors le résultat est
vrai par hypothèse ; voir
le lemme \ref{lemma-single-complex-base-change-condition}, partie (2),
et la traduction algébrique fournie par les
lemmes \ref{lemma-affine-compare-bounded} et
\ref{lemma-quasi-coherence-pullback}.

\medskip\noindent
L'étape d'induction. Supposons que $X = U \cup V$ soit un recouvrement ouvert
avec $U$, $V$ et $U \cap V$
quasi-compacts, tel que le résultat soit vrai pour $U$, $V$ et $U \cap V$.
Notons $a = f|_U$, $b = f|_V$ et $c = f|_{U \cap V}$.
Soient $a' : U' \to S'$, $b' : V' \to S'$ et $c' : U' \cap V' \to S'$
les changements de base de $a$, $b$ et $c$.
En utilisant les triangles distingués de Mayer--Vietoris relatif
(Cohomologie, lemme \ref{cohomology-lemma-unbounded-relative-mayer-vietoris}),
on obtient un diagramme commutatif
$$
\xymatrix{
Lg^*Rf_*K \ar[r] \ar[d] &
Rf'_* K' \ar[d] \\
Lg^*Ra_* K|_U \oplus
Lg^*Rb_* K|_V \ar[r] \ar[d] &
Ra'_* K'|_{U'} \oplus
Rb'_* K'|_{V'} \ar[d] \\
Lg^*Rc_* K|_{U \cap V} \ar[r] \ar[d] &
Rc'_* K'|_{U' \cap V'} \ar[d] \\
Lg^*Rf_* K[1] \ar[r] &
Rf'_* K'[1]
}
$$
Puisque les deuxième et troisième flèches horizontales sont des isomorphismes, la première
l'est aussi (Catégories dérivées, lemme \ref{derived-lemma-third-isomorphism-triangle}),
ce qui achève la démonstration du lemme.
\end{proof}

\begin{lemma}
\label{lemma-single-complex-base-change-condition-inherited}
Soit
$$
\xymatrix{
X' \ar[r]_{g'} \ar[d]_{f'} &
X \ar[d]^f \\
S' \ar[r]^g &
S
}
$$
un diagramme cartésien de schémas. Soit $K \in D_\QCoh(\mathcal{O}_X)$,
et soit $L(g')^*K \to K'$ un morphisme dans $D_\QCoh(\mathcal{O}_{X'})$.
Si les conditions équivalentes du
lemme \ref{lemma-single-complex-base-change-condition} sont satisfaites, alors
\begin{enumerate}
\item pour $E \in D_\QCoh(\mathcal{O}_X)$, les conditions
équivalentes du lemme \ref{lemma-single-complex-base-change-condition} sont satisfaites
pour $L(g')^*(E \otimes^\mathbf{L} K) \to L(g')^*E \otimes^\mathbf{L} K'$,
\item si $E$ dans $D(\mathcal{O}_X)$ est parfait, les conditions équivalentes du
lemme \ref{lemma-single-complex-base-change-condition} sont satisfaites pour
$L(g')^*R\SheafHom(E, K) \to R\SheafHom(L(g')^*E, K')$, et
\item si $K$ est borné inférieurement et si $E$ dans $D(\mathcal{O}_X)$ est
pseudo-cohérent, les conditions équivalentes du
lemme \ref{lemma-single-complex-base-change-condition} sont satisfaites pour
$L(g')^*R\SheafHom(E, K) \to R\SheafHom(L(g')^*E, K')$.
\end{enumerate}
\end{lemma}

\begin{proof}
L'énoncé a un sens, car les complexes en jeu ont des faisceaux de cohomologie
quasi-cohérents d'après les lemmes
\ref{lemma-quasi-coherence-pullback},
\ref{lemma-quasi-coherence-tensor-product} et
\ref{lemma-quasi-coherence-internal-hom}, ainsi que les
lemmes de Cohomologie \ref{cohomology-lemma-pseudo-coherent-pullback} et
\ref{cohomology-lemma-perfect-pullback}.
Ceci étant dit, on peut vérifier que les morphismes (\ref{equation-bc})
sont des isomorphismes dans le cas (1) en calculant la source et le but
de (\ref{equation-bc}) au moyen de la transitivité du produit tensoriel ; voir
Compléments sur l'algèbre, lemme \ref{more-algebra-lemma-triple-tensor-product}.
Le morphisme dans (2) et (3) est la composée
$$
L(g')^*R\SheafHom(E, K) \to R\SheafHom(L(g')^*E, L(g')^*K)
\to R\SheafHom(L(g')^*E, K')
$$
où la première flèche est celle de la
remarque de Cohomologie \ref{cohomology-remark-prepare-fancy-base-change},
et la seconde provient du morphisme donné $L(g')^*K \to K'$.
Pour démontrer que les morphismes (\ref{equation-bc}) sont des isomorphismes, on représente
$E_x$ par un complexe borné de $\mathcal{O}_{X. x}$-modules projectifs de type fini
dans le cas (2), ou par un complexe borné supérieurement de modules libres de type fini dans le cas (3),
et l'on calcule la source et le but de la flèche.
Certains détails sont omis.
\end{proof}

\begin{lemma}
\label{lemma-base-change-tensor}
Soit $f : X \to S$ un morphisme quasi-compact et quasi-séparé de
schémas. Soit $E \in D_\QCoh(\mathcal{O}_X)$. Soit $\mathcal{G}^\bullet$
un complexe borné supérieurement de $\mathcal{O}_X$-modules
quasi-cohérents et plats sur $S$. Alors la formation de
$$
Rf_*(E \otimes^\mathbf{L}_{\mathcal{O}_X} \mathcal{G}^\bullet)
$$
commute aux changements de base arbitraires (voir la démonstration pour un énoncé précis).
\end{lemma}

\begin{proof}
L'énoncé signifie ce qui suit. Soit $g : S' \to S$ un morphisme de
schémas et considérons le diagramme de changement de base
$$
\xymatrix{
X' \ar[r]_{g'} \ar[d]_{f'} &
X \ar[d]^f \\
S' \ar[r]^g &
S
}
$$
autrement dit $X' = S' \times_S X$. Le lemme affirme que
$$
Lg^*Rf_*(E \otimes^\mathbf{L}_{\mathcal{O}_X} \mathcal{G}^\bullet)
\longrightarrow
Rf'_*\left(
L(g')^*E \otimes^\mathbf{L}_{\mathcal{O}_{X'}} (g')^*\mathcal{G}^\bullet
\right)
$$
est un isomorphisme. Observons que, dans le membre de droite, nous n'utilisons
{\bf pas} l'image inverse dérivée pour $\mathcal{G}^\bullet$.
Pour le démontrer, appliquons les lemmes \ref{lemma-single-complex-base-change} et
\ref{lemma-single-complex-base-change-condition-inherited} : il suffit alors de démontrer que le morphisme canonique
$$
L(g')^*\mathcal{G}^\bullet \to (g')^*\mathcal{G}^\bullet
$$
satisfait les conditions équivalentes du
lemme \ref{lemma-single-complex-base-change-condition}.
Cela résulte de la vérification de la condition sur les fibres, où l'assertion
découle immédiatement du fait que
$\mathcal{G}^\bullet_x \otimes_{\mathcal{O}_{S, s}} \mathcal{O}_{S', s'}$
calcule le produit tensoriel dérivé, par les hypothèses faites sur le complexe
$\mathcal{G}^\bullet$.
\end{proof}

\begin{lemma}
\label{lemma-base-change-RHom}
Soit $f : X \to S$ un morphisme quasi-compact et quasi-séparé de schémas.
Soit $E$ un objet de $D(\mathcal{O}_X)$.
Soit $\mathcal{G}^\bullet$ un complexe de
$\mathcal{O}_X$-modules quasi-cohérents. Si
\begin{enumerate}
\item $E$ est parfait, $\mathcal{G}^\bullet$ est borné supérieurement,
et $\mathcal{G}^n$ est plat sur $S$, ou bien
\item $E$ est pseudo-cohérent, $\mathcal{G}^\bullet$ est borné,
et $\mathcal{G}^n$ est plat sur $S$,
\end{enumerate}
alors la formation de
$$
Rf_*R\SheafHom(E, \mathcal{G}^\bullet)
$$
commute aux changements de base arbitraires (voir la démonstration pour un énoncé précis).
\end{lemma}

\begin{proof}
L'énoncé signifie ce qui suit. Soit $g : S' \to S$ un morphisme de
schémas et considérons le diagramme de changement de base
$$
\xymatrix{
X' \ar[r]_{g'} \ar[d]_{f'} &
X \ar[d]^f \\
S' \ar[r]^g &
S
}
$$
autrement dit $X' = S' \times_S X$. Le lemme affirme que
$$
Lg^*Rf_*R\SheafHom(E, \mathcal{G}^\bullet)
\longrightarrow
R(f')_*R\SheafHom(L(g')^*E, (g')^*\mathcal{G}^\bullet)
$$
est un isomorphisme. Observons que, dans le membre de droite, nous n'utilisons
{\bf pas} l'image inverse dérivée pour $\mathcal{G}^\bullet$. Pour le démontrer, appliquons
les lemmes \ref{lemma-single-complex-base-change} et
\ref{lemma-single-complex-base-change-condition-inherited} : il suffit alors de démontrer que le morphisme canonique
$$
L(g')^*\mathcal{G}^\bullet \to (g')^*\mathcal{G}^\bullet
$$
satisfait les conditions équivalentes du
lemme \ref{lemma-single-complex-base-change-condition}.
Cela a été démontré dans la démonstration du lemme \ref{lemma-base-change-tensor}.
\end{proof}




\section{Production de complexes parfaits}
\label{section-producing-perfect}

\noindent
Le lemme suivant est notre principal outil technique pour produire des
complexes parfaits. Des versions ultérieures de ce résultat s'y ramèneront
par approximation noethérienne ; voir la
section \ref{section-cohomology-and-base-change-final}.

\begin{lemma}
\label{lemma-perfect-direct-image}
Soit $S$ un schéma noethérien. Soit $f : X \to S$ un morphisme de schémas
localement de type fini. Soit $E \in D(\mathcal{O}_X)$ tel que
\begin{enumerate}
\item $E \in D^b_{\textit{Coh}}(\mathcal{O}_X)$,
\item le support de $H^i(E)$ est propre sur $S$ pour tout $i$, et
\item $E$ est de dimension Tor finie comme objet de $D(f^{-1}\mathcal{O}_S)$.
\end{enumerate}
Alors $Rf_*E$ est un objet parfait de $D(\mathcal{O}_S)$.
\end{lemma}

\begin{proof}
Le lemme \ref{lemma-direct-image-coherent} montre que $Rf_*E$ est un objet de
$D^b_{\textit{Coh}}(\mathcal{O}_S)$. Ainsi $Rf_*E$ est pseudo-cohérent
(lemme \ref{lemma-identify-pseudo-coherent-noetherian}).
Il suffit donc de montrer que $Rf_*E$ est de dimension Tor finie ; voir
Cohomologie, lemme \ref{cohomology-lemma-perfect}.
D'après le lemme \ref{lemma-tor-qc-qs}, il suffit de vérifier que
$Rf_*(E) \otimes_{\mathcal{O}_S}^\mathbf{L} \mathcal{F}$
est à cohomologie uniformément bornée pour tout faisceau quasi-cohérent
$\mathcal{F}$ sur $S$. Le fait d'être borné supérieurement est clair puisque $Rf_*(E)$
est borné supérieurement. Soit $T \subset X$ la réunion des supports
des $H^i(E)$ pour tous les $i$. Alors $T$ est propre sur $S$ en vertu des hypothèses
(1) et (2) ; voir Cohomologie des schémas, lemme
\ref{coherent-lemma-union-closed-proper-over-base}.
En particulier, il existe un ouvert quasi-compact
$X' \subset X$ contenant $T$. En posant $f' = f|_{X'}$, on a
$Rf_*(E) = Rf'_*(E|_{X'})$, car la restriction de $E$ à $X \setminus T$ est nulle.
On peut donc remplacer $X$ par $X'$ et supposer $f$ quasi-compact.
De plus, $f$ est quasi-séparé d'après Morphismes, lemme
\ref{morphisms-lemma-finite-type-Noetherian-quasi-separated}. Or
$$
Rf_*(E) \otimes_{\mathcal{O}_S}^\mathbf{L} \mathcal{F} =
Rf_*\left(E \otimes_{\mathcal{O}_X}^\mathbf{L} Lf^*\mathcal{F}\right) =
Rf_*\left(E \otimes_{f^{-1}\mathcal{O}_S}^\mathbf{L} f^{-1}\mathcal{F}\right)
$$
d'après
le lemme \ref{lemma-cohomology-base-change}
et le
lemme de Cohomologie \ref{cohomology-lemma-variant-derived-pullback}.
Par l'hypothèse (3), le complexe
$E \otimes_{f^{-1}\mathcal{O}_S}^\mathbf{L} f^{-1}\mathcal{F}$
a ses faisceaux de cohomologie dans un
intervalle fini fixé, disons $[a, b]$. Son image par $Rf_*$
a donc sa cohomologie dans l'intervalle $[a, \infty)$, ce qui conclut.
\end{proof}

\begin{lemma}
\label{lemma-tensor-perfect}
Soit $S$ un schéma noethérien. Soit $f : X \to S$ un morphisme de schémas
localement de type fini. Soit $E \in D(\mathcal{O}_X)$ parfait.
Soit $\mathcal{G}^\bullet$ un complexe borné de
$\mathcal{O}_X$-modules cohérents, plats sur $S$ et à support propre sur $S$.
Alors $K = Rf_*(E \otimes_{\mathcal{O}_X}^\mathbf{L} \mathcal{G}^\bullet)$
est un objet parfait de $D(\mathcal{O}_S)$.
\end{lemma}

\begin{proof}
L'objet $K$ est parfait d'après le lemme \ref{lemma-perfect-direct-image}.
Vérifions que ce lemme s'applique. Localement, $E$ est isomorphe à un complexe fini
de $\mathcal{O}_X$-modules libres de type fini. Ainsi, localement,
$E \otimes^\mathbf{L}_{\mathcal{O}_X} \mathcal{G}^\bullet$ est isomorphe
à un complexe fini dont les termes sont de la forme
$$
\bigoplus\nolimits_{i = a, \ldots, b} (\mathcal{G}^i)^{\oplus r_i}
$$
pour certains entiers $a, b, r_a, \ldots, r_b$. Il en résulte immédiatement que les
faisceaux de cohomologie $H^i(E \otimes^\mathbf{L}_{\mathcal{O}_X} \mathcal{G})$
sont cohérents. L'hypothèse sur la dimension Tor est également satisfaite puisque
$\mathcal{G}^i$ est plat sur $f^{-1}\mathcal{O}_S$.
\end{proof}

\begin{lemma}
\label{lemma-ext-perfect}
Soit $S$ un schéma noethérien. Soit $f : X \to S$ un morphisme de schémas
localement de type fini. Soit $E \in D(\mathcal{O}_X)$ parfait.
Soit $\mathcal{G}^\bullet$ un complexe borné de
$\mathcal{O}_X$-modules cohérents, plats sur $S$ et à support propre sur $S$.
Alors $K = Rf_*R\SheafHom(E, \mathcal{G}^\bullet)$ est un objet parfait de
$D(\mathcal{O}_S)$.
\end{lemma}

\begin{proof}
Puisque $E$ est un complexe parfait, il existe un complexe parfait dual
$E^\vee$ ; voir Cohomologie, lemme \ref{cohomology-lemma-dual-perfect-complex}.
Observons que $R\SheafHom(E, \mathcal{G}^\bullet) =
E^\vee \otimes^\mathbf{L}_{\mathcal{O}_X} \mathcal{G}^\bullet$.
Le caractère parfait de $K$ résulte donc du lemme \ref{lemma-tensor-perfect}.
\end{proof}

\noindent
Nous généraliserons le lemme suivant aux morphismes plats et propres
sur des bases quelconques dans le
lemme \ref{lemma-flat-proper-perfect-direct-image-general},
et aux morphismes parfaits propres dans
Compléments sur les morphismes, lemme
\ref{more-morphisms-lemma-perfect-proper-perfect-direct-image}.

\begin{lemma}
\label{lemma-flat-proper-perfect-direct-image}
Soit $S$ un schéma noethérien. Soit $f : X \to S$ un morphisme plat et propre
de schémas. Soit $E \in D(\mathcal{O}_X)$ parfait. Alors
$Rf_*E$ est un objet parfait de $D(\mathcal{O}_S)$.
\end{lemma}

\begin{proof}
Affirmons que le lemme \ref{lemma-perfect-direct-image} s'applique.
Les conditions (1) et (2) sont immédiates. La condition (3) est locale
sur $X$. On peut donc supposer $X$ et $S$ affines et $E$
représenté par un complexe strictement parfait de $\mathcal{O}_X$-modules.
Puisque $\mathcal{O}_X$ est plat comme faisceau de $f^{-1}\mathcal{O}_S$-modules,
la condition (3) est satisfaite.
\end{proof}






\section{Une formule de projection pour Ext}
\label{section-ext}

\noindent
Le lemme \ref{lemma-compute-ext} (ou des résultats analogues dans la littérature)
est parfois utilisé pour vérifier l'un des critères d'Artin pour les
foncteurs de Quot, les schémas de Hilbert et d'autres problèmes de modules.
Supposons que $f : X \to S$ soit un morphisme propre, plat et de présentation finie
de schémas et que $E \in D(\mathcal{O}_X)$ soit parfait.
Le lemme affirme alors que
$$
\Ext^i_X(E, f^*\mathcal{F}) =
\Ext^i_S((Rf_*E^\vee)^\vee, \mathcal{F})
$$
pour $\mathcal{F}$ quasi-cohérent sur $S$.
Sous cette forme, l'énoncé ressemble à une formule de projection
pour Ext ; de fait, le résultat découle assez
facilement du lemme \ref{lemma-cohomology-base-change}.

\begin{lemma}
\label{lemma-compute-tensor-perfect}
Sous les hypothèses et avec les notations du lemme \ref{lemma-tensor-perfect},
il existe des isomorphismes fonctoriels
$$
H^i(S, K \otimes^\mathbf{L}_{\mathcal{O}_S} \mathcal{F})
\longrightarrow
H^i(X, E \otimes_{\mathcal{O}_X}^\mathbf{L}
(\mathcal{G}^\bullet \otimes_{\mathcal{O}_X} f^*\mathcal{F}))
$$
pour $\mathcal{F}$ quasi-cohérent sur $S$,
compatibles aux morphismes de connexion (voir la démonstration).
\end{lemma}

\begin{proof}
On a
$$
\mathcal{G}^\bullet \otimes_{\mathcal{O}_X}^\mathbf{L} Lf^*\mathcal{F} =
\mathcal{G}^\bullet \otimes_{f^{-1}\mathcal{O}_S}^\mathbf{L} f^{-1}\mathcal{F} =
\mathcal{G}^\bullet \otimes_{f^{-1}\mathcal{O}_S} f^{-1}\mathcal{F} =
\mathcal{G}^\bullet \otimes_{\mathcal{O}_X} f^*\mathcal{F}
$$
où la première égalité résulte du
lemme de Cohomologie \ref{cohomology-lemma-variant-derived-pullback},
la deuxième de ce que $\mathcal{G}^n$ est un $f^{-1}\mathcal{O}_S$-module plat, et
la troisième de la définition des images inverses. On obtient donc
\begin{align*}
H^i(X, E \otimes^\mathbf{L}_{\mathcal{O}_X}
(\mathcal{G}^\bullet \otimes_{\mathcal{O}_X} f^*\mathcal{F}))
& =
H^i(X, E \otimes^\mathbf{L}_{\mathcal{O}_X} \mathcal{G}^\bullet
\otimes_{\mathcal{O}_X}^\mathbf{L} Lf^*\mathcal{F}) \\
& =
H^i(S,
Rf_*(E \otimes^\mathbf{L}_{\mathcal{O}_X} \mathcal{G}^\bullet
\otimes^\mathbf{L}_{\mathcal{O}_X} Lf^*\mathcal{F})) \\
& =
H^i(S, Rf_*(E \otimes^\mathbf{L}_{\mathcal{O}_X} \mathcal{G}^\bullet)
\otimes^\mathbf{L}_{\mathcal{O}_S} \mathcal{F}) \\
& =
H^i(S, K \otimes^\mathbf{L}_{\mathcal{O}_S} \mathcal{F}) 
\end{align*}
La première égalité résulte de ce qui précède, la deuxième de Leray
(Cohomologie, lemme \ref{cohomology-lemma-before-Leray}), et
la troisième du lemme \ref{lemma-cohomology-base-change}.
L'assertion relative aux morphismes de connexion signifie ce qui suit. Étant donnée une suite
exacte courte $0 \to \mathcal{F}_1 \to \mathcal{F}_2 \to \mathcal{F}_3 \to 0$
de $\mathcal{O}_S$-modules quasi-cohérents, les isomorphismes s'insèrent dans des
diagrammes commutatifs
$$
\xymatrix{
H^i(S, K \otimes^\mathbf{L}_{\mathcal{O}_S} \mathcal{F}_3)
\ar[r] \ar[d]_\delta &
H^i(X, E \otimes^\mathbf{L}_{\mathcal{O}_X}
(\mathcal{G}^\bullet \otimes_{\mathcal{O}_X} f^*\mathcal{F}_3))
\ar[d]^\delta \\
H^{i + 1}(S, K \otimes^\mathbf{L}_{\mathcal{O}_S} \mathcal{F}_1)
\ar[r] &
H^{i + 1}(X, E \otimes^\mathbf{L}_{\mathcal{O}_X}
(\mathcal{G}^\bullet \otimes_{\mathcal{O}_X} f^*\mathcal{F}_1))
}
$$
où les morphismes de connexion proviennent du triangle distingué
$$
K \otimes^\mathbf{L}_{\mathcal{O}_S} \mathcal{F}_1 \to
K \otimes^\mathbf{L}_{\mathcal{O}_S} \mathcal{F}_2 \to
K \otimes^\mathbf{L}_{\mathcal{O}_S} \mathcal{F}_3 \to
K \otimes^\mathbf{L}_{\mathcal{O}_S} \mathcal{F}_1[1]
$$
et du triangle distingué de $D(\mathcal{O}_X)$ associé à
la suite exacte courte
$$
0 \to
\mathcal{G}^\bullet \otimes_{\mathcal{O}_X} f^*\mathcal{F}_1 \to
\mathcal{G}^\bullet \otimes_{\mathcal{O}_X} f^*\mathcal{F}_2 \to
\mathcal{G}^\bullet \otimes_{\mathcal{O}_X} f^*\mathcal{F}_3 \to 0
$$
de complexes de $\mathcal{O}_X$-modules.
Cette suite est exacte parce que $\mathcal{G}^n$ est plat sur $S$.
Nous omettons la vérification de la commutativité du diagramme affiché.
\end{proof}

\begin{lemma}
\label{lemma-compute-ext-perfect}
Sous les hypothèses et avec les notations du lemme \ref{lemma-ext-perfect},
il existe des isomorphismes fonctoriels
$$
H^i(S, K \otimes^\mathbf{L}_{\mathcal{O}_S} \mathcal{F})
\longrightarrow
\Ext^i_{\mathcal{O}_X}(E,
\mathcal{G}^\bullet \otimes_{\mathcal{O}_X} f^*\mathcal{F})
$$
pour $\mathcal{F}$ quasi-cohérent sur $S$,
compatibles aux morphismes de connexion (voir la démonstration).
\end{lemma}

\begin{proof}
Comme dans la démonstration du lemme \ref{lemma-ext-perfect}, soit $E^\vee$ le
complexe parfait dual et rappelons que
$K = Rf_*(E^\vee \otimes_{\mathcal{O}_X}^\mathbf{L} \mathcal{G}^\bullet)$.
Puisque l'on a aussi
$$
\Ext^i_{\mathcal{O}_X}(E,
\mathcal{G}^\bullet \otimes_{\mathcal{O}_X} f^*\mathcal{F})
=
H^i(X, E^\vee \otimes^\mathbf{L}_{\mathcal{O}_X}
(\mathcal{G}^\bullet \otimes_{\mathcal{O}_X} f^*\mathcal{F}))
$$
par construction de $E^\vee$, l'existence des isomorphismes résulte
du lemme \ref{lemma-compute-tensor-perfect} appliqué à $E^\vee$
et $\mathcal{G}^\bullet$.
L'assertion relative aux morphismes de connexion signifie ce qui suit. Étant donnée une suite
exacte courte $0 \to \mathcal{F}_1 \to \mathcal{F}_2 \to \mathcal{F}_3 \to 0$,
les isomorphismes s'insèrent dans des diagrammes commutatifs
$$
\xymatrix{
H^i(S, K \otimes^\mathbf{L}_{\mathcal{O}_S} \mathcal{F}_3)
\ar[r] \ar[d]_\delta &
\Ext^i_{\mathcal{O}_X}(E,
\mathcal{G}^\bullet \otimes_{\mathcal{O}_X} f^*\mathcal{F}_3) \ar[d]^\delta \\
H^{i + 1}(S, K \otimes^\mathbf{L}_{\mathcal{O}_S} \mathcal{F}_1)
\ar[r] &
\Ext^{i + 1}_{\mathcal{O}_X}(E,
\mathcal{G}^\bullet \otimes_{\mathcal{O}_X} f^*\mathcal{F}_1)
}
$$
où les morphismes de connexion proviennent du triangle distingué
$$
K \otimes^\mathbf{L}_{\mathcal{O}_S} \mathcal{F}_1 \to
K \otimes^\mathbf{L}_{\mathcal{O}_S} \mathcal{F}_2 \to
K \otimes^\mathbf{L}_{\mathcal{O}_S} \mathcal{F}_3 \to
K \otimes^\mathbf{L}_{\mathcal{O}_S} \mathcal{F}_1[1]
$$
et du triangle distingué de $D(\mathcal{O}_X)$ associé à
la suite exacte courte
$$
0 \to
\mathcal{G}^\bullet \otimes_{\mathcal{O}_X} f^*\mathcal{F}_1 \to
\mathcal{G}^\bullet \otimes_{\mathcal{O}_X} f^*\mathcal{F}_2 \to
\mathcal{G}^\bullet \otimes_{\mathcal{O}_X} f^*\mathcal{F}_3 \to 0
$$
de complexes.
Cette suite est exacte parce que $\mathcal{G}$ est plat sur $S$.
Nous omettons la vérification de la commutativité du diagramme affiché.
\end{proof}

\begin{lemma}
\label{lemma-compute-ext}
Soient $f : X \to S$ un morphisme de schémas, $E \in D(\mathcal{O}_X)$
et $\mathcal{G}^\bullet$ un complexe de $\mathcal{O}_X$-modules.
Supposons que
\begin{enumerate}
\item $S$ soit noethérien,
\item $f$ soit localement de type fini,
\item $E \in D^-_{\textit{Coh}}(\mathcal{O}_X)$,
\item $\mathcal{G}^\bullet$ soit un complexe borné de
$\mathcal{O}_X$-modules cohérents, plats sur $S$ et à support propre sur $S$.
\end{enumerate}
Alors les deux assertions suivantes sont vraies :
\begin{enumerate}
\item[(A)] pour tout $m \in \mathbf{Z}$, il existe un objet parfait $K$
de $D(\mathcal{O}_S)$ et des morphismes fonctoriels
$$
\alpha^i_\mathcal{F} :
\Ext^i_{\mathcal{O}_X}(E,
\mathcal{G}^\bullet \otimes_{\mathcal{O}_X} f^*\mathcal{F})
\longrightarrow
H^i(S, K \otimes^\mathbf{L}_{\mathcal{O}_S} \mathcal{F})
$$
pour $\mathcal{F}$ quasi-cohérent sur $S$, compatibles aux morphismes de connexion
(voir la démonstration), tels que $\alpha^i_\mathcal{F}$ soit un isomorphisme pour $i \leq m$ ;
\item[(B)] il existe un $L \in D(\mathcal{O}_S)$ pseudo-cohérent
et des isomorphismes fonctoriels
$$
\Ext^i_{\mathcal{O}_S}(L, \mathcal{F}) \longrightarrow
\Ext^i_{\mathcal{O}_X}(E,
\mathcal{G}^\bullet \otimes_{\mathcal{O}_X} f^*\mathcal{F})
$$
pour $\mathcal{F}$ quasi-cohérent sur $S$, compatibles aux morphismes de connexion.
\end{enumerate}
\end{lemma}

\begin{proof}
Démonstration de (A).
Supposons $\mathcal{G}^i$ non nul seulement pour $i \in [a, b]$.
On peut remplacer $X$ par un voisinage ouvert quasi-compact de
la réunion des supports des $\mathcal{G}^i$.
On peut donc supposer $X$ noethérien.
Dans ce cas, $X$ et $f$ sont quasi-compacts et quasi-séparés.
Choisissons une approximation $P \to E$ par un complexe parfait $P$ de
$(X, E, -m - 1 + a)$
(ce qui est possible d'après le théorème \ref{theorem-approximation}).
Alors le morphisme induit
$$
\Ext^i_{\mathcal{O}_X}(E,
\mathcal{G}^\bullet \otimes_{\mathcal{O}_X} f^*\mathcal{F})
\longrightarrow
\Ext^i_{\mathcal{O}_X}(P,
\mathcal{G}^\bullet \otimes_{\mathcal{O}_X} f^*\mathcal{F})
$$
est un isomorphisme pour $i \leq m$. En effet, le noyau, resp.\ le conoyau, de ce
morphisme est un quotient, resp.\ un sous-module, de
$$
\Ext^i_{\mathcal{O}_X}(C,
\mathcal{G}^\bullet \otimes_{\mathcal{O}_X} f^*\mathcal{F})
\quad\text{resp.}\quad
\Ext^{i + 1}_{\mathcal{O}_X}(C,
\mathcal{G}^\bullet \otimes_{\mathcal{O}_X} f^*\mathcal{F})
$$
où $C$ est le cône de $P \to E$. Puisque les faisceaux de cohomologie de $C$
sont nuls en degrés $\geq -m - 1 + a$, ces groupes $\Ext$ sont nuls
pour $i \leq m + 1$ d'après
Catégories dérivées, lemme \ref{derived-lemma-negative-exts}.
On est ainsi ramené au cas où
$E$ est un complexe parfait, qui est le lemme \ref{lemma-compute-ext-perfect}.
L'assertion relative aux morphismes de connexion est expliquée dans la démonstration du
lemme \ref{lemma-compute-ext-perfect}.

\medskip\noindent
Démonstration de (B). Comme dans la démonstration de (A), on peut supposer $X$ noethérien.
Observons que $E$ est pseudo-cohérent d'après le
lemme \ref{lemma-identify-pseudo-coherent-noetherian}.
D'après le lemme \ref{lemma-pseudo-coherent-hocolim}, on peut écrire
$E = \text{hocolim} E_n$, où $E_n$ est parfait et où $E_n \to E$ induit
un isomorphisme sur les troncatures $\tau_{\geq -n}$. Soit $E_n^\vee$
le complexe parfait dual
(Cohomologie, lemme \ref{cohomology-lemma-dual-perfect-complex}).
On obtient un système projectif $\ldots \to E_3^\vee \to E_2^\vee \to E_1^\vee$
d'objets parfaits. Celui-ci donne à son tour un système projectif
$$
\ldots \to K_3 \to K_2 \to K_1\quad\text{avec}\quad
K_n = Rf_*(E_n^\vee \otimes_{\mathcal{O}_X}^\mathbf{L} \mathcal{G}^\bullet)
$$
d'objets parfaits sur $S$ ; voir le lemme \ref{lemma-tensor-perfect}.
D'après le lemme \ref{lemma-compute-ext-perfect} et sa démonstration, ainsi que
les arguments du paragraphe précédent (avec $P = E_n$),
pour tout $\mathcal{F}$ quasi-cohérent sur $S$, on dispose de
morphismes canoniques fonctoriels
$$
\xymatrix{
& \Ext^i_{\mathcal{O}_X}(E,
\mathcal{G}^\bullet \otimes_{\mathcal{O}_X} f^*\mathcal{F})
\ar[ld] \ar[rd] \\
H^i(S, K_{n + 1} \otimes_{\mathcal{O}_S}^\mathbf{L} \mathcal{F})
\ar[rr] & &
H^i(S, K_n \otimes_{\mathcal{O}_S}^\mathbf{L} \mathcal{F})
}
$$
qui sont des isomorphismes pour $i \leq n + a$.
Soit $L_n = K_n^\vee$ le complexe parfait dual.
On voit alors que $L_1 \to L_2 \to L_3 \to \ldots$
est un système d'objets parfaits de $D(\mathcal{O}_S)$
tel que, pour tout $\mathcal{F}$ quasi-cohérent sur $S$,
les morphismes
$$
\Ext^i_{\mathcal{O}_S}(L_{n + 1}, \mathcal{F})
\longrightarrow
\Ext^i_{\mathcal{O}_S}(L_n, \mathcal{F})
$$
sont des isomorphismes pour $i \leq n + a - 1$. Il en résulte que
$L_n \to L_{n + 1}$ induit un isomorphisme sur les troncatures
$\tau_{\geq -n - a + 2}$ (indication : prendre le cône de $L_n \to L_{n + 1}$
et considérer son dernier faisceau de cohomologie non nul).
Ainsi $L = \text{hocolim} L_n$ est pseudo-cohérent ; voir le
lemme \ref{lemma-pseudo-coherent-hocolim}. La propriété universelle
des colimites homotopiques donne
$\Ext^i_{\mathcal{O}_S}(L, \mathcal{F}) =
\Ext^i_{\mathcal{O}_S}(L_n, \mathcal{F})$
pour $i \leq n + a - 3$, ce qui achève la démonstration.
\end{proof}

\begin{remark}
\label{remark-base-change-of-L}
Le complexe pseudo-cohérent $L$ de la partie (B) du lemme \ref{lemma-compute-ext}
est canoniquement associé à la situation. Par exemple,
la formation de $L$ comme en (B) est compatible au changement de base.
Autrement dit, étant donné un diagramme cartésien
$$
\xymatrix{
X' \ar[r]_{g'} \ar[d]_{f'} &
X \ar[d]^f \\
S' \ar[r]^g &
S
}
$$
de schémas, on dispose d'isomorphismes canoniques fonctoriels
$$
\Ext^i_{\mathcal{O}_{S'}}(Lg^*L, \mathcal{F}') \longrightarrow
\Ext^i_{\mathcal{O}_X}(L(g')^*E,
(g')^*\mathcal{G}^\bullet \otimes_{\mathcal{O}_{X'}} (f')^*\mathcal{F}')
$$
pour $\mathcal{F}'$ quasi-cohérent sur $S'$. Observons que nous n'utilisons {\bf pas}
l'image inverse dérivée de $\mathcal{G}^\bullet$ dans le membre de droite.
Si nous avons un jour besoin de ce fait, nous
formulerons ici un résultat précis et en donnerons une démonstration détaillée.
\end{remark}





\section{Limites et catégories dérivées}
\label{section-limits}

\noindent
Dans cette section, nous rassemblons quelques résultats sur la catégorie dérivée
d'un schéma qui est la limite d'un système projectif de schémas.
Plus précisément, nous travaillerons dans la situation suivante.

\begin{situation}
\label{situation-descent}
Soit $S = \lim_{i \in I} S_i$ la limite d'un système filtrant de schémas
dont les morphismes de transition $f_{i'i} : S_{i'} \to S_i$ sont affines.
Nous supposons que $S_i$ est quasi-compact et quasi-séparé pour tout $i \in I$.
Nous notons $f_i : S \to S_i$ la projection. Fixons également un élément $0 \in I$.
\end{situation}

\begin{lemma}
\label{lemma-descend-homomorphisms}
Dans la situation \ref{situation-descent},
soient $E_0$ et $K_0$ des objets de
$D(\mathcal{O}_{S_0})$.
Posons $E_i = Lf_{i0}^*E_0$ et $K_i = Lf_{i0}^*K_0$ pour $i \geq 0$,
et posons $E = Lf_0^*E_0$ et $K = Lf_0^*K_0$. Alors le morphisme
$$
\colim_{i \geq 0} \Hom_{D(\mathcal{O}_{S_i})}(E_i, K_i)
\longrightarrow
\Hom_{D(\mathcal{O}_S)}(E, K)
$$
est un isomorphisme dans chacun des deux cas suivants :
\begin{enumerate}
\item $E_0$ est parfait et $K_0 \in D_\QCoh(\mathcal{O}_{S_0})$ ;
\item $E_0$ est pseudo-cohérent et
$K_0 \in D_\QCoh(\mathcal{O}_{S_0})$ est de dimension Tor finie.
\end{enumerate}
\end{lemma}

\begin{proof}
Pour tout ouvert $U_0 \subset S_0$, considérons la condition $P$ selon laquelle le morphisme canonique
$$
\colim_{i \geq 0} \Hom_{D(\mathcal{O}_{U_i})}(E_i|_{U_i}, K_i|_{U_i})
\longrightarrow
\Hom_{D(\mathcal{O}_U)}(E|_U, K|_U)
$$
est un isomorphisme, où $U = f_0^{-1}(U_0)$ et $U_i = f_{i0}^{-1}(U_0)$.
Nous allons démontrer que $P$ est satisfaite pour tout ouvert quasi-compact $U_0$
par le principe d'induction de
Cohomologie des schémas, lemme \ref{coherent-lemma-induction-principle}.
La condition (2) de ce lemme résulte immédiatement de Mayer--Vietoris
pour les Hom de la catégorie dérivée ; voir
Cohomologie, lemme \ref{cohomology-lemma-mayer-vietoris-hom}.
Il suffit donc de démontrer le lemme lorsque $S_0$ est affine.

\medskip\noindent
Supposons $S_0$ affine. Écrivons $S_0 = \Spec(A_0)$, $S_i = \Spec(A_i)$ et
$S = \Spec(A)$. Nous utiliserons le lemme \ref{lemma-affine-compare-bounded}
sans le mentionner davantage.

\medskip\noindent
Dans le cas (1), l'objet $E_0^\bullet$ correspond à un complexe fini
de $A_0$-modules projectifs de type fini ; voir le lemme \ref{lemma-perfect-affine}.
On peut représenter l'objet $K_0$ par un complexe K-plat $K_0^\bullet$
de $A_0$-modules. Dans cette situation, il s'agit de démontrer que
$$
\colim_{i \geq 0} \Hom_{D(A_i)}(E_0^\bullet \otimes_{A_0} A_i,
K_0^\bullet \otimes_{A_0} A_i)
\longrightarrow
\Hom_{D(A)}(E_0^\bullet \otimes_{A_0} A, K_0^\bullet \otimes_{A_0} A)
$$
est un isomorphisme. Comme $E_0^\bullet$ est un complexe borné supérieurement de modules projectifs,
on peut réécrire ce morphisme sous la forme
$$
\colim_{i \geq 0} \Hom_{K(A_0)}(E_0^\bullet,
K_0^\bullet \otimes_{A_0} A_i)
\longrightarrow
\Hom_{K(A_0)}(E_0^\bullet, K_0^\bullet \otimes_{A_0} A)
$$
Puisqu'il n'y a qu'un nombre fini de modules non nuls
$E_0^n$ et que ceux-ci sont tous de présentation finie, ce
morphisme est un isomorphisme.

\medskip\noindent
Dans le cas (2), l'objet $E_0$ correspond à un
complexe borné supérieurement $E_0^\bullet$ de $A_0$-modules libres de type fini ;
voir le lemme \ref{lemma-pseudo-coherent-affine}.
On peut représenter $K_0$ par un complexe fini $K_0^\bullet$
de $A_0$-modules plats ; voir le lemme \ref{lemma-tor-dimension-affine}
et
Compléments sur l'algèbre, lemme \ref{more-algebra-lemma-tor-amplitude}.
En particulier, $K_0^\bullet$ est K-plat, et l'on peut raisonner comme précédemment
pour obtenir le morphisme
$$
\colim_{i \geq 0} \Hom_{K(A_0)}(E_0^\bullet,
K_0^\bullet \otimes_{A_0} A_i)
\longrightarrow
\Hom_{K(A_0)}(E_0^\bullet, K_0^\bullet \otimes_{A_0} A)
$$
Il est clair que ce morphisme est un isomorphisme (seul un nombre fini de
termes intervient puisque $K_0^\bullet$ est borné).
\end{proof}

\begin{lemma}
\label{lemma-descend-perfect}
Dans la situation \ref{situation-descent}, la catégorie des objets parfaits
de $D(\mathcal{O}_S)$ est la colimite des catégories
d'objets parfaits de $D(\mathcal{O}_{S_i})$.
\end{lemma}

\begin{proof}
Pour tout ouvert $U_0 \subset S_0$, considérons la condition $P$ selon laquelle
le foncteur
$$
\colim_{i \geq 0} D_{perf}(\mathcal{O}_{U_i})
\longrightarrow
D_{perf}(\mathcal{O}_U)
$$
est une équivalence, où ${}_{perf}$ désigne la sous-catégorie pleine des
objets parfaits et où $U = f_0^{-1}(U_0)$ et $U_i = f_{i0}^{-1}(U_0)$.
Nous allons démontrer que $P$ est satisfaite pour tout ouvert quasi-compact $U_0$
par le principe d'induction de
Cohomologie des schémas, lemme \ref{coherent-lemma-induction-principle}.
D'abord, observons que le foncteur est déjà pleinement fidèle
d'après le lemme \ref{lemma-descend-homomorphisms}. Il suffit donc de démontrer
qu'il est essentiellement surjectif.

\medskip\noindent
Vérifions d'abord la condition (2) du principe d'induction. Supposons donc
que $S_0 = U_0 \cup V_0$ et que $P$ soit satisfaite pour
$U_0$, $V_0$ et $U_0 \cap V_0$. Soit $E$ un objet parfait
de $D(\mathcal{O}_S)$. On peut trouver $i \geq 0$, un objet $E_{U, i}$ parfait sur $U_i$
et un objet $E_{V, i}$ parfait sur $V_i$, dont les images inverses sur $U$ et $V$ sont isomorphes
à $E|_U$ et $E|_V$. Notons
$$
a : E_{U, i} \to (Rf_{i, *}E)|_{U_i}
\quad\text{et}\quad
b : E_{V, i} \to (Rf_{i, *}E)|_{V_i}
$$
les morphismes adjoints aux isomorphismes $Lf_i^*E_{U, i} \to E|_U$
et $Lf_i^*E_{V, i} \to E|_V$.
Par pleine fidélité, quitte à augmenter $i$,
on peut trouver un isomorphisme
$c : E_{U, i}|_{U_i \cap V_i} \to E_{V, i}|_{U_i \cap V_i}$
qui, par image inverse, donne les identifications
$$
Lf_i^*E_{U, i}|_{U \cap V} \to E|_{U \cap V} \to Lf_i^*E_{V, i}|_{U \cap V}.
$$
Appliquons Cohomologie, lemme \ref{cohomology-lemma-glue},
pour obtenir un objet $E_i$ sur $S_i$ et un morphisme $d : E_i \to Rf_{i, *}E$
dont les restrictions sont les morphismes $a$ et $b$ sur $U_i$ et $V_i$.
Il est alors clair que $E_i$ est parfait et que
$d$ est adjoint à un isomorphisme $Lf_i^*E_i \to E$.

\medskip\noindent
Enfin, vérifions la condition (1) du principe d'induction, autrement
dit vérifions le lemme lorsque $S_0$ est affine.
Écrivons $S_0 = \Spec(A_0)$, $S_i = \Spec(A_i)$ et
$S = \Spec(A)$. Les lemmes \ref{lemma-affine-compare-bounded}
et \ref{lemma-perfect-affine} montrent qu'il faut établir que
$$
D_{perf}(A) = \colim D_{perf}(A_i)
$$
Cela résulte du fait que les complexes parfaits sur des anneaux sont
donnés par des complexes finis de modules projectifs de type fini (donc de présentation finie).
Voir Compléments sur l'algèbre, lemme
\ref{more-algebra-lemma-colimit-perfect-complexes}, pour les détails.
\end{proof}





\section{Cohomologie et changement de base, VI}
\label{section-cohomology-and-base-change-final}

\noindent
Une dernière section sur la cohomologie et le changement de base, qui poursuit
la discussion des sections
\ref{section-cohomology-and-base-change-perfect},
\ref{section-cohomology-base-change} et
\ref{section-producing-perfect}.
Un cas particulier facile à saisir est donné dans la
remarque \ref{remark-explain-perfect-direct-image}.

\begin{lemma}
\label{lemma-base-change-tensor-perfect}
Soit $f : X \to S$ un morphisme de présentation finie.
Soit $E \in D(\mathcal{O}_X)$ un objet parfait. Soit $\mathcal{G}^\bullet$
un complexe borné de $\mathcal{O}_X$-modules de présentation finie,
plats sur $S$ et à support propre sur $S$. Alors
$$
K = Rf_*(E \otimes_{\mathcal{O}_X}^\mathbf{L} \mathcal{G}^\bullet)
$$
est un objet parfait de $D(\mathcal{O}_S)$ et sa formation
commute à tout changement de base.
\end{lemma}

\begin{proof}
L'assertion relative au changement de base est le lemme \ref{lemma-base-change-tensor}.
Il suffit donc de montrer que $K$ est un objet parfait. Si $S$ est
noethérien, cela résulte du
lemme \ref{lemma-tensor-perfect}.
Nous allons nous ramener à ce cas par approximation noethérienne.
Nous conseillons au lecteur de sauter le reste de cette démonstration.

\medskip\noindent
La question est locale sur $S$ ; on peut donc supposer $S$ affine.
Écrivons $S = \Spec(R)$. Écrivons $R = \colim R_i$ comme colimite filtrante
d'anneaux noethériens $R_i$. D'après Limites, lemme
\ref{limits-lemma-descend-finite-presentation},
il existe un $i$ et un schéma $X_i$ de présentation finie sur $R_i$
dont le changement de base à $R$ est $X$. D'après
Limites, lemme \ref{limits-lemma-descend-modules-finite-presentation},
quitte à augmenter $i$, on peut supposer qu'il existe un complexe borné
de $\mathcal{O}_{X_i}$-modules de présentation finie
$\mathcal{G}_i^\bullet$ dont l'image inverse sur $X$ est $\mathcal{G}^\bullet$.
Quitte à augmenter $i$, on peut supposer $\mathcal{G}_i^n$ plat sur $R_i$ ; voir
Limites, lemme \ref{limits-lemma-descend-module-flat-finite-presentation}.
Quitte à augmenter $i$, on peut supposer que le support de $\mathcal{G}_i^n$
est propre sur $R_i$ ; voir
Limites, lemme \ref{limits-lemma-eventually-proper-support},
et Cohomologie des schémas, lemme
\ref{coherent-lemma-module-support-proper-over-base}.
Enfin, d'après le lemme \ref{lemma-descend-perfect},
quitte à augmenter $i$, on peut supposer qu'il existe un objet parfait
$E_i$ de $D(\mathcal{O}_{X_i})$ dont l'image inverse sur
$X$ est $E$. En appliquant le lemme \ref{lemma-tensor-perfect}
à $X_i \to \Spec(R_i)$, $E_i$, $\mathcal{G}_i^\bullet$, et en utilisant la
propriété de changement de base déjà démontrée, on obtient le résultat.
\end{proof}

\begin{remark}
\label{remark-explain-perfect-direct-image}
Soit $R$ un anneau. Soit $X$ un schéma de présentation finie sur
$R$. Soit $\mathcal{G}$ un $\mathcal{O}_X$-module de présentation finie,
plat sur $R$ et à support propre sur $R$. D'après le
lemme \ref{lemma-base-change-tensor-perfect},
il existe un complexe fini de $R$-modules projectifs de type fini
$M^\bullet$ tel que l'on ait
$$
R\Gamma(X_{R'}, \mathcal{G}_{R'}) = M^\bullet \otimes_R R'
$$
fonctoriellement en la $R$-algèbre $R'$.
\end{remark}

\begin{lemma}
\label{lemma-base-change-tensor-pseudo-coherent}
Soit $f : X \to S$ un morphisme de présentation finie.
Soit $E \in D(\mathcal{O}_X)$ un objet pseudo-cohérent.
Soit $\mathcal{G}^\bullet$ un complexe borné supérieurement de
$\mathcal{O}_X$-modules de présentation finie, plats sur $S$
et à support propre sur $S$. Alors
$$
K = Rf_*(E \otimes_{\mathcal{O}_X}^\mathbf{L} \mathcal{G}^\bullet)
$$
est un objet pseudo-cohérent de $D(\mathcal{O}_S)$ et sa formation
commute à tout changement de base.
\end{lemma}

\begin{proof}
L'assertion relative au changement de base est le lemme \ref{lemma-base-change-tensor}.
Il suffit donc de montrer que $K$ est un objet pseudo-cohérent.
Cela résultera du lemme \ref{lemma-base-change-tensor-perfect}
par approximation par des complexes parfaits. Nous conseillons au lecteur de
sauter le reste de la démonstration.

\medskip\noindent
La question est locale sur $S$ ; on peut donc supposer $S$ affine.
Alors $X$ est quasi-compact et quasi-séparé. De plus, il
existe un entier $N$ tel que l'image directe totale
$Rf_* : D_\QCoh(\mathcal{O}_X) \to D_\QCoh(\mathcal{O}_S)$
soit de dimension cohomologique $N$, comme expliqué dans le
lemme \ref{lemma-quasi-coherence-direct-image}.
Choisissons un entier $b$ tel que $\mathcal{G}^i = 0$ pour $i > b$.
Il suffit de montrer que $K$ est $m$-pseudo-cohérent pour
tout $m$. Choisissons une approximation $P \to E$ par un complexe parfait $P$
de $(X, E, m - N - 1 - b)$. Cela est possible d'après le
théorème \ref{theorem-approximation}.
Choisissons un triangle distingué
$$
P \to E \to C \to P[1]
$$
dans $D_\QCoh(\mathcal{O}_X)$. Les faisceaux de cohomologie de $C$ sont nuls
en degrés $\geq m - N - 1 - b$. Les faisceaux de cohomologie de
$C \otimes^\mathbf{L} \mathcal{G}^\bullet$ sont donc nuls en degrés
$\geq m - N - 1$. Ainsi, les faisceaux de cohomologie de
$Rf_*(C \otimes^\mathbf{L} \mathcal{G}^\bullet)$
sont nuls en degrés $\geq m - 1$.
Par conséquent,
$$
Rf_*(P \otimes^\mathbf{L} \mathcal{G}^\bullet) \to
Rf_*(E \otimes^\mathbf{L} \mathcal{G}^\bullet)
$$
est un isomorphisme sur les faisceaux de cohomologie en degrés $\geq m$.
Supposons ensuite que $H^i(P) = 0$ pour $i > a$. Alors
$
P \otimes^\mathbf{L} \sigma_{\geq m - N - 1 - a}\mathcal{G}^\bullet
\longrightarrow
P \otimes^\mathbf{L} \mathcal{G}^\bullet
$
est un isomorphisme sur les faisceaux de cohomologie en degrés $\geq m - N - 1$.
On trouve donc, de nouveau, que
$$
Rf_*(P \otimes^\mathbf{L} \sigma_{\geq m - N - 1 - a}\mathcal{G}^\bullet) \to
Rf_*(P \otimes^\mathbf{L} \mathcal{G}^\bullet)
$$
est un isomorphisme sur les faisceaux de cohomologie en degrés $\geq m$.
D'après le lemme \ref{lemma-base-change-tensor-perfect}, la source
est un complexe parfait.
On conclut que $K$ est $m$-pseudo-cohérent, comme voulu.
\end{proof}

\begin{lemma}
\label{lemma-flat-proper-perfect-direct-image-general}
Soit $S$ un schéma. Soit $f : X \to S$ un morphisme propre
de présentation finie.
\begin{enumerate}
\item Soit $E \in D(\mathcal{O}_X)$ parfait et supposons $f$ plat. Alors
$Rf_*E$ est un objet parfait de $D(\mathcal{O}_S)$ et sa formation
commute à tout changement de base.
\item Soit $\mathcal{G}$ un $\mathcal{O}_X$-module de présentation finie,
plat sur $S$. Alors $Rf_*\mathcal{G}$ est un objet parfait de
$D(\mathcal{O}_S)$ et sa formation commute à tout changement de base.
\end{enumerate}
\end{lemma}

\begin{proof}
Ce sont des cas particuliers du
lemme \ref{lemma-base-change-tensor-perfect}, appliqué avec
(1) $\mathcal{G}^\bullet$ égal à $\mathcal{O}_X$ placé en degré $0$
et (2) $E = \mathcal{O}_X$ et $\mathcal{G}^\bullet$ constitué de
$\mathcal{G}$ placé en degré $0$.
\end{proof}

\begin{lemma}
\label{lemma-flat-proper-pseudo-coherent-direct-image-general}
Soit $S$ un schéma. Soit $f : X \to S$ un morphisme plat et propre
de présentation finie. Soit $E \in D(\mathcal{O}_X)$
pseudo-cohérent. Alors $Rf_*E$ est un objet pseudo-cohérent de
$D(\mathcal{O}_S)$ et sa formation commute à tout changement de base.
\end{lemma}

\noindent
Plus généralement, si $f : X \to S$ est propre et si $E$ sur $X$ est pseudo-cohérent
relativement à $S$ (Compléments sur les morphismes, définition
\ref{more-morphisms-definition-relative-pseudo-coherence}),
alors $Rf_*E$ est pseudo-cohérent
(mais sa formation ne commute pas au changement de base dans cette généralité).
Voir \cite{Kiehl}.

\begin{proof}
C'est un cas particulier du
lemme \ref{lemma-base-change-tensor-pseudo-coherent}, appliqué avec
$\mathcal{G}^\bullet$ égal à $\mathcal{O}_X$ placé en degré $0$.
\end{proof}

\begin{lemma}
\label{lemma-pullback-and-limits}
Soit $R$ un anneau. Soit $X$ un schéma, et soit
$f : X \to \Spec(R)$ un morphisme propre, plat et
de présentation finie. Soit $(M_n)$ un système projectif
de $R$-modules dont les morphismes de transition sont surjectifs.
Alors le morphisme canonique
$$
\mathcal{O}_X \otimes_R (\lim M_n)
\longrightarrow
\lim \mathcal{O}_X \otimes_R M_n
$$
induit un isomorphisme de la source vers l'image du but par $DQ_X$.
\end{lemma}

\begin{proof}
L'énoncé signifie que, pour tout objet $E$ de
$D_\QCoh(\mathcal{O}_X)$, le morphisme induit
$$
\Hom(E, \mathcal{O}_X \otimes_R (\lim M_n))
\longrightarrow
\Hom(E, \lim \mathcal{O}_X \otimes_R M_n)
$$
est un isomorphisme. Puisque $D_\QCoh(\mathcal{O}_X)$ possède
un générateur parfait (théorème \ref{theorem-bondal-van-den-Bergh}),
il suffit de le vérifier pour $E$ parfait.
D'après le lemme \ref{lemma-Rlim-quasi-coherent}, on a
$\lim \mathcal{O}_X \otimes_R M_n = R\lim \mathcal{O}_X \otimes_R M_n$.
Le foncteur exact
$R\Hom_X(E, -) : D_\QCoh(\mathcal{O}_X) \to D(R)$
de Cohomologie, section \ref{cohomology-section-global-RHom},
commute aux produits et donc aux limites dérivées, d'où
$$
R\Hom_X(E, \lim \mathcal{O}_X \otimes_R M_n) =
R\lim R\Hom_X(E, \mathcal{O}_X \otimes_R M_n)
$$
Soit $E^\vee$ le complexe parfait dual ; voir
Cohomologie, lemme \ref{cohomology-lemma-dual-perfect-complex}.
On a
$$
R\Hom_X(E, \mathcal{O}_X \otimes_R M_n) =
R\Gamma(X, E^\vee \otimes_{\mathcal{O}_X}^\mathbf{L} Lf^*M_n) =
R\Gamma(X, E^\vee) \otimes_R^\mathbf{L} M_n
$$
d'après le lemme \ref{lemma-cohomology-base-change}.
Le lemme \ref{lemma-flat-proper-perfect-direct-image-general}
montre que $R\Gamma(X, E^\vee)$ est un complexe parfait de $R$-modules.
En particulier, c'est un complexe pseudo-cohérent, et le
lemme de Compléments sur l'algèbre \ref{more-algebra-lemma-pseudo-coherent-tensor-limit}
donne
$$
R\lim R\Gamma(X, E^\vee) \otimes_R^\mathbf{L} M_n =
R\Gamma(X, E^\vee) \otimes_R^\mathbf{L} \lim M_n
$$
comme voulu.
\end{proof}

\begin{lemma}
\label{lemma-base-change-RHom-perfect}
Soit $f : X \to S$ un morphisme de présentation finie.
Soit $E \in D(\mathcal{O}_X)$ un objet parfait. Soit $\mathcal{G}^\bullet$
un complexe borné de $\mathcal{O}_X$-modules de présentation finie,
plats sur $S$ et à support propre sur $S$. Alors
$$
K = Rf_*R\SheafHom(E, \mathcal{G}^\bullet)
$$
est un objet parfait de $D(\mathcal{O}_S)$ et sa formation
commute à tout changement de base.
\end{lemma}

\begin{proof}
L'assertion relative au changement de base est le lemme \ref{lemma-base-change-RHom}.
Il suffit donc de montrer que $K$ est un objet parfait. Si $S$ est
noethérien, cela résulte du
lemme \ref{lemma-ext-perfect}.
Nous allons nous ramener à ce cas par approximation noethérienne.
Nous conseillons au lecteur de sauter le reste de cette démonstration.

\medskip\noindent
La question est locale sur $S$ ; on peut donc supposer $S$ affine.
Écrivons $S = \Spec(R)$. Écrivons $R = \colim R_i$ comme colimite filtrante
d'anneaux noethériens $R_i$. D'après Limites, lemme
\ref{limits-lemma-descend-finite-presentation},
il existe un $i$ et un schéma $X_i$ de présentation finie sur $R_i$
dont le changement de base à $R$ est $X$. D'après
Limites, lemme \ref{limits-lemma-descend-modules-finite-presentation},
quitte à augmenter $i$, on peut supposer qu'il existe un complexe borné
de $\mathcal{O}_{X_i}$-modules de présentation finie $\mathcal{G}_i^\bullet$
dont l'image inverse sur $X$ est $\mathcal{G}^\bullet$. Quitte à augmenter $i$,
on peut supposer $\mathcal{G}_i^n$ plat sur $R_i$ ; voir
Limites, lemme \ref{limits-lemma-descend-module-flat-finite-presentation}.
Quitte à augmenter $i$, on peut supposer que le support de $\mathcal{G}_i^n$
est propre sur $R_i$ ; voir
Limites, lemme \ref{limits-lemma-eventually-proper-support}, et
Cohomologie des schémas, lemme
\ref{coherent-lemma-module-support-proper-over-base}.
Enfin, d'après le lemme \ref{lemma-descend-perfect},
quitte à augmenter $i$, on peut supposer qu'il existe un objet parfait
$E_i$ de $D(\mathcal{O}_{X_i})$ dont l'image inverse sur
$X$ est $E$. En appliquant le lemme \ref{lemma-ext-perfect}
à $X_i \to \Spec(R_i)$, $E_i$, $\mathcal{G}_i^\bullet$, et en utilisant la
propriété de changement de base déjà démontrée, on obtient le résultat.
\end{proof}









\section{Complexes parfaits}
\label{section-perfect-complexes}

\noindent
Nous parlons d'abord des lieux de saut des nombres de Betti des complexes parfaits.
Étant donnés un complexe $E$ sur un schéma $X$ et un point $x$ de $X$, nous écrivons souvent
$E \otimes_{\mathcal{O}_X}^\mathbf{L} \kappa(x)$ au lieu de l'expression plus correcte
$Li_x^*E$, où $i_x : x \to X$ est le morphisme canonique.

\begin{lemma}
\label{lemma-jump-loci}
Soit $X$ un schéma. Soit $E \in D(\mathcal{O}_X)$ pseudo-cohérent
(par exemple parfait). Pour tout $i \in \mathbf{Z}$, considérons la fonction
$$
\beta_i : X \longrightarrow \{0, 1, 2, \ldots\},\quad
x \longmapsto
\dim_{\kappa(x)}
H^i(E \otimes_{\mathcal{O}_X}^\mathbf{L} \kappa(x))
$$
Alors :
\begin{enumerate}
\item la formation de $\beta_i$ commute à tout changement de base ;
\item les fonctions $\beta_i$ sont semi-continues supérieurement ;
\item les ensembles de niveau de $\beta_i$ sont localement constructibles dans $X$.
\end{enumerate}
\end{lemma}

\begin{proof}
Considérons un morphisme de schémas $f : Y \to X$ et un point $y \in Y$.
Soit $x$ l'image de $y$, et considérons le diagramme commutatif
$$
\xymatrix{
y \ar[r]_j \ar[d]_g & Y \ar[d]^f \\
x \ar[r]^i & X
}
$$
On voit que $Lg^* \circ Li^* = Lj^* \circ Lf^*$. Il en résulte que
la fonction $\beta'_i$ associée au complexe pseudo-cohérent $Lf^*E$
est l'image inverse de la fonction $\beta_i$, c'est-à-dire
$\beta'_i = \beta_i \circ f$. Tel est le sens de (1).

\medskip\noindent
Fixons $i$ et soit $x \in X$. Il suffit de démontrer (2) et (3)
dans un voisinage ouvert de $x$ ; on peut donc supposer $X$ affine.
On peut alors représenter $E$ par un complexe borné supérieurement $\mathcal{F}^\bullet$
de modules libres de type fini (lemme \ref{lemma-lift-pseudo-coherent}).
Alors $P = \sigma_{\geq i - 1}\mathcal{F}^\bullet$ est un objet parfait,
et $P \to E$ induit un isomorphisme
$$
H^i(P \otimes_{\mathcal{O}_X}^\mathbf{L} \kappa(x')) \to
H^i(E \otimes_{\mathcal{O}_X}^\mathbf{L} \kappa(x'))
$$
pour tout $x' \in X$. On peut donc supposer $E$ parfait. Dans ce cas,
d'après Compléments sur l'algèbre, lemme
\ref{more-algebra-lemma-lift-perfect-from-residue-field},
il existe un voisinage ouvert affine $U$ de $x$ et
$a \leq b$ tels que $E|_U$ soit représenté par un complexe
$$
\ldots \to 0 \to \mathcal{O}_U^{\oplus \beta_a(x)}
\to \mathcal{O}_U^{\oplus \beta_{a + 1}(x)} \to
\ldots \to
\mathcal{O}_U^{\oplus \beta_{b - 1}(x)} \to
\mathcal{O}_U^{\oplus \beta_b(x)} \to 0
\to \ldots
$$
(On utilise également ici des résultats antérieurs pour ramener le problème à l'algèbre, par exemple les
lemmes \ref{lemma-affine-compare-bounded} et
\ref{lemma-perfect-affine}.)
Il en résulte immédiatement que $\beta_i(x') \leq \beta_i(x)$
pour tout $x' \in U$. Cela démontre que $\beta_i$ est semi-continue
supérieurement.

\medskip\noindent
Pour démontrer (3), on peut supposer que $X$ est affine et que
$E$ est donné par un complexe de
$\mathcal{O}_X$-modules libres de type fini (par exemple en raisonnant comme au paragraphe
précédent, ou en utilisant Cohomologie, lemme
\ref{cohomology-lemma-perfect-on-locally-ringed}).
Il faut donc montrer qu'étant donné un complexe
$$
\mathcal{O}_X^{\oplus a} \to
\mathcal{O}_X^{\oplus b} \to
\mathcal{O}_X^{\oplus c}
$$
la fonction qui associe à un point $x \in X$ la dimension de la cohomologie
au milieu de $\kappa_x^{\oplus a} \to \kappa_x^{\oplus b} \to \kappa_x^{\oplus c}$
a des ensembles de niveau constructibles. Soit
$A \in \text{Mat}(a \times b, \Gamma(X, \mathcal{O}_X))$ la matrice
de la première flèche. Le rang de l'image de $A$ dans
$\text{Mat}(a \times b, \kappa(x))$ est égal à $r$ si tous les
mineurs $(r + 1) \times (r + 1)$ de $A$ s'annulent en $x$ et s'il existe un
mineur $r \times r$ de $A$ qui ne s'annule pas en $x$. L'ensemble
des points où le rang vaut $r$ est donc un ensemble localement fermé constructible.
En raisonnant de même pour la seconde flèche et en réunissant le tout,
on obtient le résultat voulu.
\end{proof}

\begin{lemma}
\label{lemma-chi-locally-constant}
Soit $X$ un schéma. Soit $E \in D(\mathcal{O}_X)$ parfait.
La fonction
$$
\chi_E : X \longrightarrow \mathbf{Z},\quad
x \longmapsto \sum (-1)^i
\dim_{\kappa(x)} H^i(E \otimes_{\mathcal{O}_X}^\mathbf{L} \kappa(x))
$$
est localement constante sur $X$.
\end{lemma}

\begin{proof}
D'après Cohomologie, lemme
\ref{cohomology-lemma-perfect-on-locally-ringed},
on peut, localement sur $X$, représenter $E$ par un complexe fini
$\mathcal{E}^\bullet$ de $\mathcal{O}_X$-modules libres de type fini.
Sur un tel ouvert, la fonction $\chi_E$ est constante de valeur
$\sum (-1)^i \text{rank}(\mathcal{E}^i)$.
\end{proof}

\begin{lemma}
\label{lemma-open-where-cohomology-in-degree-i-rank-r}
Soit $X$ un schéma. Soit $E \in D(\mathcal{O}_X)$ parfait.
Étant donnés $i, r \in \mathbf{Z}$, il existe un
sous-schéma ouvert $U \subset X$ caractérisé par les propriétés suivantes :
\begin{enumerate}
\item $E|_U \cong H^i(E|_U)[-i]$ et $H^i(E|_U)$ est un
$\mathcal{O}_U$-module localement libre de rang $r$ ;
\item un morphisme $f : Y \to X$ se factorise par $U$ si et seulement si
$Lf^*E$ est isomorphe à un module localement libre de rang $r$
placé en degré $i$.
\end{enumerate}
\end{lemma}

\begin{proof}
Soient $\beta_j : X \to \{0, 1, 2, \ldots\}$, pour $j \in \mathbf{Z}$,
les fonctions du lemme \ref{lemma-jump-loci}. Alors l'ensemble
$$
W = \{x \in X \mid \beta_j(x) \leq 0\text{ pour tout }j \not = i\}
$$
est ouvert dans $X$, et sa formation commute à l'image inverse sur tout
$Y$ au-dessus de $X$. Cela résulte du lemme et du fait que,
a priori, dans un voisinage de tout point, seul un nombre fini
des $\beta_j$ est non nul. On peut donc remplacer $X$ par $W$
et supposer que $\beta_j(x) = 0$ pour tout $x \in X$ et tout $j \not = i$.
Dans ce cas, $H^i(E)$ est un module localement libre de type fini et
$E \cong H^i(E)[-i]$ ; voir par exemple
Compléments sur l'algèbre, lemme
\ref{more-algebra-lemma-lift-perfect-from-residue-field}.
Ainsi, $X$ est la réunion disjointe des sous-schémas ouverts sur lesquels le
rang de $H^i(E)$ est fixé, ce qui conclut.
\end{proof}

\begin{lemma}
\label{lemma-locally-closed-where-H0-locally-free}
Soit $X$ un schéma. Soit $E \in D(\mathcal{O}_X)$ parfait
d'amplitude Tor contenue dans $[a, b]$ pour certains $a, b \in \mathbf{Z}$.
Soit $r \geq 0$.
Alors il existe un sous-schéma localement fermé $j : Z \to X$
caractérisé par les propriétés suivantes :
\begin{enumerate}
\item $H^a(Lj^*E)$ est un $\mathcal{O}_Z$-module localement libre de rang $r$ ;
\item un morphisme $f : Y \to X$ se factorise par $Z$
si et seulement si, pour tout morphisme $g : Y' \to Y$, le
$\mathcal{O}_{Y'}$-module $H^a(L(f \circ g)^*E)$ est localement libre
de rang $r$.
\end{enumerate}
De plus, $j : Z \to X$ est de présentation finie, et l'on a :
\begin{enumerate}
\item[(3)] si $f : Y \to X$ se factorise sous la forme $Y \xrightarrow{g} Z \to X$, alors
$H^a(Lf^*E) = g^*H^a(Lj^*E)$ ;
\item[(4)] si $\beta_a(x) \leq r$ pour tout $x \in X$, alors
$j$ est une immersion fermée et, pour tout $f : Y \to X$, les assertions suivantes
sont équivalentes :
\begin{enumerate}
\item $f : Y \to X$ se factorise par $Z$ ;
\item $H^a(Lf^*E)$ est un $\mathcal{O}_Y$-module localement libre de rang $r$ ;
\end{enumerate}
et, si $r = 1$, elles sont encore équivalentes à
\begin{enumerate}
\item[(c)] $\mathcal{O}_Y \to \SheafHom_{\mathcal{O}_Y}(H^a(Lf^*E), H^a(Lf^*E))$
est injectif.
\end{enumerate}
\end{enumerate}
\end{lemma}

\begin{proof}
Tout d'abord, soit $U \subset X$ le sous-schéma ouvert
localement constructible où la fonction
$\beta_a$ du lemme \ref{lemma-jump-loci} prend des valeurs $\leq r$.
Soit $f : Y \to X$ comme en (2). Alors, pour tout $y \in Y$,
on a $\beta_a(Lf^*E) = r$ ; le point $y$ est donc envoyé dans $U$ d'après le
lemme \ref{lemma-jump-loci}. Par conséquent, $f$ comme en (2) se factorise par $U$.
On peut ainsi remplacer $X$ par $U$ et supposer que
$\beta_a(x) \in \{0, 1, \ldots, r\}$ pour tout $x \in X$.
Nous allons montrer que, dans ce cas,
il existe un sous-schéma fermé $Z \subset X$ défini par un idéal
quasi-cohérent de type fini, caractérisé
par l'équivalence de (4)(a), (b), et de (4)(c) si $r = 1$,
et que (3) est satisfaite.
Cela achèvera la démonstration, car il en résultera a fortiori
que les morphismes comme en (2) se factorisent par $Z$.

\medskip\noindent
Si $x \in X$ et $\beta_a(x) < r$, il existe un
voisinage ouvert de $x$ sur lequel $\beta_a < r$
(lemme \ref{lemma-jump-loci}). On voit ainsi que,
au moins ensemblistement, $Z$ est un sous-ensemble fermé.

\medskip\noindent
Pour obtenir une structure schématique, considérons un point $x \in X$
tel que $\beta_a(x) = r$. Posons $\beta = \beta_{a + 1}(x)$.
D'après Compléments sur l'algèbre, lemme
\ref{more-algebra-lemma-lift-perfect-from-residue-field},
il existe un voisinage ouvert affine $U$ de $x$
tel que $K|_U$ soit représenté par un complexe
$$
\ldots \to 0 \to \mathcal{O}_U^{\oplus r}
\xrightarrow{(f_{ij})} \mathcal{O}_U^{\oplus \beta} \to
\ldots \to
\mathcal{O}_U^{\oplus \beta_{b - 1}(x)} \to
\mathcal{O}_U^{\oplus \beta_b(x)} \to 0
\to \ldots
$$
(On utilise également ici des résultats antérieurs pour ramener le problème à l'algèbre, par exemple les
lemmes \ref{lemma-affine-compare-bounded} et
\ref{lemma-perfect-affine}.) Or, si $g : Y \to U$ est un morphisme quelconque
de schémas tel que $g^\sharp(f_{ij})$ soit non nul pour un certain couple $i, j$, alors
$H^a(Lg^*E)$ n'est pas un $\mathcal{O}_Y$-module localement libre de rang $r$.
Voir Compléments sur l'algèbre, lemme \ref{more-algebra-lemma-coker-injective-free}.
Il est immédiat que $H^a(Lg^*E)$ est un $\mathcal{O}_Y$-module localement libre si
$g^\sharp(f_{ij}) = 0$ pour tous $i, j$. On voit donc que, sur $U$, le
sous-schéma fermé défini par tous les $f_{ij}$ satisfait
(3), et que l'on a l'équivalence de (4)(a) et (b).
La caractérisation de $Z$ montre que les morceaux construits localement
se recollent (détails omis). Enfin, si $r = 1$, alors
(4)(c) équivaut à (4)(b), car dans ce cas, localement,
$H^a(Lg^*E) \subset \mathcal{O}_Y$ est l'annulateur de l'idéal engendré
par les éléments $g^\sharp(f_{ij})$.
\end{proof}





\section{Applications}
\label{section-applications}

\noindent
Principalement des applications de la cohomologie et du changement de base. À l'avenir, nous pourrons
généraliser ces résultats à la situation examinée dans le
lemme \ref{lemma-base-change-tensor-perfect}.

\begin{lemma}
\label{lemma-jump-loci-geometric}
Soit $f : X \to S$ un morphisme propre de présentation finie.
Soit $\mathcal{F}$ un $\mathcal{O}_X$-module de présentation finie,
plat sur $S$. Pour $i \in \mathbf{Z}$ fixé, considérons la fonction
$$
\beta_i : S \to \{0, 1, 2, \ldots\},\quad
s \longmapsto \dim_{\kappa(s)} H^i(X_s, \mathcal{F}_s)
$$
Alors :
\begin{enumerate}
\item la formation de $\beta_i$ commute à tout changement de base ;
\item les fonctions $\beta_i$ sont semi-continues supérieurement ;
\item les ensembles de niveau de $\beta_i$ sont localement constructibles dans $S$.
\end{enumerate}
\end{lemma}

\begin{proof}
D'après la cohomologie et le changement de base (plus précisément d'après le
lemme \ref{lemma-flat-proper-perfect-direct-image-general}),
l'objet $K = Rf_*\mathcal{F}$ est un objet parfait de la catégorie dérivée
de $S$, dont la formation commute à tout changement de base.
En particulier, on a
$$
H^i(X_s, \mathcal{F}_s) = H^i(K \otimes_{\mathcal{O}_S}^\mathbf{L} \kappa(s))
$$
Le lemme résulte donc du
lemme \ref{lemma-jump-loci}.
\end{proof}

\begin{lemma}
\label{lemma-chi-locally-constant-geometric}
Soit $f : X \to S$ un morphisme propre de présentation finie.
Soit $\mathcal{F}$ un $\mathcal{O}_X$-module de présentation finie,
plat sur $S$. La fonction
$$
s \longmapsto \chi(X_s, \mathcal{F}_s)
$$
est localement constante sur $S$. La formation de cette fonction commute au
changement de base.
\end{lemma}

\begin{proof}
D'après la cohomologie et le changement de base (plus précisément d'après le
lemme \ref{lemma-flat-proper-perfect-direct-image-general}),
l'objet $K = Rf_*\mathcal{F}$ est un objet parfait de la catégorie dérivée
de $S$, dont la formation commute à tout changement de base.
Il faut donc montrer que l'application
$$
s \longmapsto \sum (-1)^i \dim_{\kappa(s)}
H^i(K \otimes^\mathbf{L}_{\mathcal{O}_S} \kappa(s))
$$
est localement constante sur $S$. C'est le lemme \ref{lemma-chi-locally-constant}.
\end{proof}

\begin{lemma}
\label{lemma-open-where-cohomology-in-degree-i-rank-r-geometric}
Soit $f : X \to S$ un morphisme propre de présentation finie.
Soit $\mathcal{F}$ un $\mathcal{O}_X$-module de présentation finie,
plat sur $S$. Fixons $i, r \in \mathbf{Z}$.
Alors il existe un sous-schéma ouvert
$U \subset S$ ayant la propriété suivante :
un morphisme $T \to S$ se factorise par $U$ si et seulement si
$Rf_{T, *}\mathcal{F}_T$ est isomorphe à un
module localement libre de type fini de rang $r$ placé en degré $i$.
\end{lemma}

\begin{proof}
D'après la cohomologie et le changement de base (plus précisément d'après le
lemme \ref{lemma-flat-proper-perfect-direct-image-general}),
l'objet $K = Rf_*\mathcal{F}$ est un objet parfait de la catégorie dérivée
de $S$, dont la formation commute à tout changement de base.
Ce lemme résulte donc immédiatement du
lemme \ref{lemma-open-where-cohomology-in-degree-i-rank-r}.
\end{proof}

\begin{lemma}
\label{lemma-vanishing-implies-locally-free}
Soit $f : X \to S$ un morphisme de présentation finie.
Soit $\mathcal{F}$ un $\mathcal{O}_X$-module de présentation finie,
plat sur $S$, dont le support est propre sur $S$. Si $R^if_*\mathcal{F} = 0$
pour $i > 0$, alors $f_*\mathcal{F}$ est localement libre et sa formation
commute à tout changement de base (voir la démonstration pour une explication).
\end{lemma}

\begin{proof}
D'après le lemme \ref{lemma-base-change-tensor-perfect},
l'objet $E = Rf_*\mathcal{F}$ de $D(\mathcal{O}_S)$
est parfait et sa formation commute à tout changement de base,
au sens où $Rf'_*(g')^*\mathcal{F} = Lg^*E$
pour tout diagramme cartésien
$$
\xymatrix{
X' \ar[r]_{g'} \ar[d]_{f'} &
X \ar[d]^f \\
S' \ar[r]^g &
S
}
$$
de schémas.
Puisqu'il n'y a jamais de cohomologie en degrés $< 0$, on voit que
$E$ a (localement) une amplitude de Tor contenue dans $[0, b]$ pour un certain $b$.
Si $H^i(E) = R^if_*\mathcal{F} = 0$ pour $i > 0$,
alors $E$ a une amplitude de Tor contenue dans $[0, 0]$. Par conséquent,
$E = H^0(E)[0]$. On en déduit que $H^0(E) = f_*\mathcal{F}$
est localement libre de type fini d'après
Compléments sur l'algèbre, lemme \ref{more-algebra-lemma-perfect}
(et la caractérisation des modules projectifs de type fini donnée dans
Algèbre, lemme \ref{algebra-lemma-finite-projective}).
La commutation au changement de base signifie que
$g^*f_*\mathcal{F} = f'_*(g')^*\mathcal{F}$ pour
un diagramme comme ci-dessus, et elle résulte de la commutation au changement
de base déjà établie pour $E$.
\end{proof}

\begin{lemma}
\label{lemma-proper-flat-h0}
Soit $f : X \to S$ un morphisme de schémas. Supposons que
\begin{enumerate}
\item $f$ soit propre, plat et de présentation finie ;
\item pour tout $s \in S$, on ait $\kappa(s) = H^0(X_s, \mathcal{O}_{X_s})$.
\end{enumerate}
Alors :
\begin{enumerate}
\item[(a)] $f_*\mathcal{O}_X = \mathcal{O}_S$, et
cela reste vrai après tout changement de base ;
\item[(b)] localement sur $S$, on a
$$
Rf_*\mathcal{O}_X = \mathcal{O}_S \oplus P
$$
dans $D(\mathcal{O}_S)$,
où $P$ est parfait d'amplitude de Tor contenue dans $[1, \infty)$.
\end{enumerate}
\end{lemma}

\begin{proof}
D'après la cohomologie et le changement de base
(lemme \ref{lemma-flat-proper-perfect-direct-image-general}),
le complexe $E = Rf_*\mathcal{O}_X$
est parfait et sa formation commute à tout changement de base.
Cela implique d'abord que $E$ a une amplitude de Tor contenue dans $[0, \infty)$.
Ensuite, cela implique que, pour $s \in S$, on a
$H^0(E \otimes^\mathbf{L} \kappa(s)) =
H^0(X_s, \mathcal{O}_{X_s}) = \kappa(s)$.
Il en résulte que l'application $\mathcal{O}_S \to Rf_*\mathcal{O}_X = E$
induit un isomorphisme
$\mathcal{O}_S \otimes \kappa(s) \to H^0(E \otimes^\mathbf{L} \kappa(s))$.
Ainsi, $H^0(E) \otimes \kappa(s) \to H^0(E \otimes^\mathbf{L} \kappa(s))$
est surjective, et l'on peut appliquer
Compléments sur l'algèbre, lemme \ref{more-algebra-lemma-better-cut-complex-in-two},
pour voir qu'après avoir remplacé $S$ par un voisinage ouvert affine de $s$,
on a une décomposition $E = H^0(E) \oplus \tau_{\geq 1}E$
où $\tau_{\geq 1}E$ est parfait d'amplitude de Tor contenue dans $[1, \infty)$.
Puisque $E$ a une amplitude de Tor contenue dans $[0, \infty)$, on trouve que
$H^0(E)$ est un $\mathcal{O}_S$-module plat.
Il en résulte que $H^0(E)$ est un $\mathcal{O}_S$-module plat et parfait,
donc localement libre de type fini ; voir
Compléments sur l'algèbre, lemme \ref{more-algebra-lemma-perfect}
(ainsi que le fait que les modules projectifs de type fini sont localement libres de type fini d'après
Algèbre, lemme \ref{algebra-lemma-finite-projective}).
Il en résulte que l'application $\mathcal{O}_S \to H^0(E)$ est
un isomorphisme, puisque cela se vérifie après tensorisation par les
corps résiduels (Algèbre, lemme \ref{algebra-lemma-cokernel-flat}).
\end{proof}

\begin{lemma}
\label{lemma-proper-flat-geom-red-connected}
Soit $f : X \to S$ un morphisme de schémas. Supposons que
\begin{enumerate}
\item $f$ soit propre, plat et de présentation finie ;
\item les fibres géométriques de $f$ soient réduites et connexes.
\end{enumerate}
Alors $f_*\mathcal{O}_X = \mathcal{O}_S$, et cela reste vrai
après tout changement de base.
\end{lemma}

\begin{proof}
D'après le lemme \ref{lemma-proper-flat-h0},
il suffit de montrer que $\kappa(s) = H^0(X_s, \mathcal{O}_{X_s})$
pour tout $s \in S$. Cela résulte de
Variétés, lemme
\ref{varieties-lemma-proper-geometrically-reduced-global-sections},
et du fait que $X_s$ est géométriquement connexe et géométriquement réduit.
\end{proof}

\begin{lemma}
\label{lemma-proper-idempotent-on-fibre}
Soit $f : X \to S$ un morphisme propre de schémas. Soit $s \in S$,
et soit $e \in H^0(X_s, \mathcal{O}_{X_s})$ un idempotent.
Alors $e$ appartient à l'image de l'application
$(f_*\mathcal{O}_X)_s \to H^0(X_s, \mathcal{O}_{X_s})$.
\end{lemma}

\begin{proof}
Soit $X_s = T_1 \amalg T_2$ la décomposition en union disjointe
où $T_1$ et $T_2$ sont non vides, ouverts et fermés dans $X_s$,
correspondant à $e$, c'est-à-dire telle que $e$ soit identiquement égal à $1$
sur $T_1$ et identiquement égal à $0$ sur $T_2$.

\medskip\noindent
Supposons $S$ noethérien. Nous utiliserons le théorème des fonctions formelles
sous la forme de Cohomologie des schémas, lemme
\ref{coherent-lemma-formal-functions-stalk}.
Il nous dit que
$$
(f_*\mathcal{O}_X)_s^\wedge = \lim_n H^0(X_n, \mathcal{O}_{X_n})
$$
où $X_n$ est le $n$-ième voisinage infinitésimal de $X_s$.
Puisque l'espace topologique sous-jacent à $X_n$ est égal à celui
de $X_s$, on obtient pour tout $n $ une décomposition de schémas en union disjointe
$X_n = T_{1, n} \amalg T_{2, n}$ où l'espace topologique sous-jacent
à $T_{i, n}$ est $T_i$ pour $i = 1, 2$. Cela signifie que
$H^0(X_n, \mathcal{O}_{X_n})$ contient un idempotent non trivial $e_n$,
à savoir la fonction identiquement égale à $1$ sur $T_{1, n}$ et
identiquement égale à $0$ sur $T_{2, n}$. Il est clair que $e_{n + 1}$
se restreint en $e_n$ sur $X_n$. Ainsi, $e_\infty = \lim e_n$
est un idempotent non trivial de la limite. Par conséquent, $e_\infty$
est un élément du complété de $(f_*\mathcal{O}_X)_s$
qui s'envoie sur $e$ dans $H^0(X_s, \mathcal{O}_{X_s})$.
Puisque l'application $(f_*\mathcal{O}_X)_s^\wedge \to H^0(X_s, \mathcal{O}_{X_s})$
se factorise par
$(f_*\mathcal{O}_X)^\wedge_s / \mathfrak m_s (f_*\mathcal{O}_X)_s^\wedge =
(f_*\mathcal{O}_X)_s / \mathfrak m_s (f_*\mathcal{O}_X)_s$
(Algèbre, lemme \ref{algebra-lemma-hathat-finitely-generated}),
on en déduit que $e$ appartient à l'image de
l'application $(f_*\mathcal{O}_X)_s \to H^0(X_s, \mathcal{O}_{X_s})$,
comme désiré.

\medskip\noindent
Cas général : on ramène le cas général au cas noethérien par des
arguments de limites. Nous invitons le lecteur à omettre la démonstration.
On peut remplacer $S$ par un voisinage ouvert affine de $s$.
On peut donc supposer, et l'on suppose, que $S$ est affine. D'après Limites, lemme
\ref{limits-lemma-proper-limit-of-proper-finite-presentation-noetherian},
on peut écrire $(f : X \to S) = \lim (f_i : X_i \to S_i)$, avec $f_i$
propre et $S_i$ noethérien. Notons $s_i \in S_i$ l'image de $s$.
Alors $s = \lim s_i$ ; voir
Limites, lemme \ref{limits-lemma-inverse-limit-irreducibles}.
Alors $X_s = X \times_S s = \lim X_i \times_{S_i} s_i = \lim X_{i, s_i}$,
car les limites commutent aux limites
(Catégories, lemme \ref{categories-lemma-colimits-commute}).
Ainsi, $e$ est l'image d'un certain idempotent
$e_i \in H^0(X_{i, s_i}, \mathcal{O}_{X_{i, s_i}})$
d'après Limites, lemme \ref{limits-lemma-descend-section}.
D'après le cas noethérien, il existe un élément $\tilde e_i$
de la fibre $(f_{i, *}\mathcal{O}_{X_i})_{s_i}$ qui s'envoie
sur $e_i$. En prenant l'image inverse de $\tilde e_i$, on obtient un
élément $\tilde e$ de $(f_*\mathcal{O}_X)_s$ qui s'envoie
sur $e$, ce qui achève la démonstration.
\end{proof}

\begin{lemma}
\label{lemma-proper-flat-geom-red}
Soit $f : X \to S$ un morphisme de schémas. Soit $s \in S$. Supposons que
\begin{enumerate}
\item $f$ soit propre, plat et de présentation finie ;
\item la fibre $X_s$ soit géométriquement réduite.
\end{enumerate}
Alors, après avoir remplacé $S$ par un voisinage ouvert de $s$, il
existe une décomposition en somme directe
$Rf_*\mathcal{O}_X = f_*\mathcal{O}_X \oplus P$
dans $D(\mathcal{O}_S)$, où $f_*\mathcal{O}_X$ est une algèbre finie \'etale
sur $\mathcal{O}_S$, et
$P$ est parfait d'amplitude de Tor contenue dans $[1, \infty)$.
\end{lemma}

\begin{proof}
La démonstration de ce lemme est analogue à celle du
lemme \ref{lemma-proper-flat-h0},
que nous suggérons au lecteur de lire d'abord.
D'après la cohomologie et le changement de base
(lemme \ref{lemma-flat-proper-perfect-direct-image-general}),
le complexe $E = Rf_*\mathcal{O}_X$
est parfait et sa formation commute à tout changement de base.
Cela implique d'abord que $E$ a une amplitude de Tor contenue dans $[0, \infty)$.

\medskip\noindent
Nous affirmons qu'après avoir remplacé $S$ par un voisinage ouvert de
$s$, on peut trouver une décomposition en somme directe
$E = H^0(E) \oplus \tau_{\geq 1}E$ dans $D(\mathcal{O}_S)$
où $\tau_{\geq 1}E$ a une amplitude de Tor contenue dans $[1, \infty)$.
Admettons cette assertion pour le moment et supposons le remplacement effectué,
de sorte que l'on dispose de la décomposition en somme directe.
Puisque $E$ a une amplitude de Tor contenue dans $[0, \infty)$, on trouve que
$H^0(E)$ est un $\mathcal{O}_S$-module plat.
Ainsi, $H^0(E)$ est un $\mathcal{O}_S$-module plat et parfait,
donc localement libre de type fini ; voir
Compléments sur l'algèbre, lemme \ref{more-algebra-lemma-perfect}
(ainsi que le fait que les modules projectifs de type fini sont localement libres de type fini d'après
Algèbre, lemme \ref{algebra-lemma-finite-projective}).
Bien entendu, $H^0(E) = f_*\mathcal{O}_X$ est une $\mathcal{O}_S$-algèbre.
D'après la cohomologie et le changement de base, on obtient
$H^0(E) \otimes \kappa(s) = H^0(X_s, \mathcal{O}_{X_s})$.
D'après Variétés, lemme
\ref{varieties-lemma-proper-geometrically-reduced-global-sections},
et l'hypothèse que $X_s$ est géométriquement réduit, on
voit que $\kappa(s) \to H^0(E) \otimes \kappa(s)$
est finie \'etale. D'après Morphismes, lemme
\ref{morphisms-lemma-set-points-where-fibres-etale},
appliqué au morphisme fini localement libre
$\underline{\Spec}_S(H^0(E)) \to S$,
on en déduit qu'après avoir rétréci $S$, la $\mathcal{O}_S$-algèbre
$H^0(E)$ est finie \'etale.

\medskip\noindent
Il reste à démontrer l'assertion. Pour cela, il suffit de montrer que l'application
$$
(f_*\mathcal{O}_X)_s
\longrightarrow
H^0(X_s, \mathcal{O}_{X_s}) = H^0(E \otimes^\mathbf{L} \kappa(s))
$$
est surjective ; voir
Compléments sur l'algèbre, lemme \ref{more-algebra-lemma-better-cut-complex-in-two}.
Choisissons un homomorphisme local plat d'anneaux $\mathcal{O}_{S, s} \to A$
tel que le corps résiduel $k$ de $A$ soit algébriquement clos ; voir
Algèbre, lemme \ref{algebra-lemma-flat-local-given-residue-field}.
D'après le changement de base plat (Cohomologie des schémas, lemme
\ref{coherent-lemma-flat-base-change-cohomology}),
on obtient $H^0(X_A, \mathcal{O}_{X_A}) =
(f_*\mathcal{O}_X)_s \otimes_{\mathcal{O}_{S, s}} A$
et $H^0(X_k, \mathcal{O}_{X_k}) =
H^0(X_s, \mathcal{O}_{X_s}) \otimes_{\kappa(s)} k$.
Il suffit donc de montrer que
$H^0(X_A, \mathcal{O}_{X_A}) \to H^0(X_k, \mathcal{O}_{X_k})$
est surjective. Puisque $X_k$ est un schéma propre réduit sur $k$
et que $k$ est algébriquement clos, on voit que
$H^0(X_k, \mathcal{O}_{X_k})$ est un produit fini de copies
de $k$ d'après Variétés, lemme déjà utilisé,
\ref{varieties-lemma-proper-geometrically-reduced-global-sections}.
Puisque, d'après le lemme \ref{lemma-proper-idempotent-on-fibre},
les idempotents de cette $k$-algèbre sont dans l'image
de $H^0(X_A, \mathcal{O}_{X_A}) \to H^0(X_k, \mathcal{O}_{X_k})$, on conclut.
\end{proof}





\section{Autres applications}
\label{section-other-applications}

\noindent
Dans cette section, nous énonçons et démontrons quelques résultats qui se déduisent
de la théorie développée ci-dessus.

\begin{lemma}
\label{lemma-cohomology-over-coherent-ring}
Soit $R$ un anneau cohérent. Soit $X$ un schéma de présentation finie sur $R$.
Soit $\mathcal{G}$ un $\mathcal{O}_X$-module de présentation finie,
plat sur $R$, dont le support est propre sur $R$. Alors
$H^i(X, \mathcal{G})$ est un $R$-module cohérent.
\end{lemma}

\begin{proof}
On combine le lemme \ref{lemma-base-change-tensor-perfect} avec les
lemmes de Compléments sur l'algèbre \ref{more-algebra-lemma-coherent-pseudo-coherent} et
\ref{more-algebra-lemma-perfect}.
\end{proof}

\begin{lemma}
\label{lemma-countable-cohomology}
Soit $X$ un schéma quasi-compact et quasi-séparé.
Soit $K$ un objet de $D_\QCoh(\mathcal{O}_X)$
tel que les faisceaux de cohomologie $H^i(K)$ aient des ensembles
dénombrables de sections sur les ouverts affines. Alors, pour tout ouvert quasi-compact
$U \subset X$ et tout objet parfait $E$ de $D(\mathcal{O}_X)$,
les ensembles
$$
H^i(U, K \otimes^\mathbf{L} E),\quad \Ext^i(E|_U, K|_U)
$$
sont dénombrables.
\end{lemma}

\begin{proof}
En utilisant Cohomologie, lemme \ref{cohomology-lemma-dual-perfect-complex},
on voit qu'il suffit de démontrer le résultat
pour les groupes $H^i(U, K \otimes^\mathbf{L} E)$.
Nous utiliserons le principe d'induction pour démontrer le lemme ; voir
Cohomologie des schémas, lemme \ref{coherent-lemma-induction-principle}.

\medskip\noindent
Montrons d'abord que l'assertion vaut lorsque $U = \Spec(A)$ est affine. En effet, on peut
représenter $K$ par un complexe de $A$-modules $K^\bullet$, et $E$ par un
complexe fini de $A$-modules projectifs de type fini $P^\bullet$.
Voir les lemmes \ref{lemma-affine-compare-bounded} et
\ref{lemma-perfect-affine},
ainsi que notre définition des complexes parfaits de $A$-modules
(Compléments sur l'algèbre, définition \ref{more-algebra-definition-perfect}).
Alors $(E \otimes^\mathbf{L} K)|_U$ est représenté par
le complexe total associé au complexe double
$P^\bullet \otimes_A K^\bullet$
(lemme \ref{lemma-quasi-coherence-tensor-product}).
En raisonnant par récurrence sur la longueur du complexe
$P^\bullet$ (ou en utilisant une suite spectrale appropriée),
on voit qu'il suffit de montrer que
$H^i(P^a \otimes_A K^\bullet)$ est dénombrable pour tout $a$.
Puisque $P^a$ est un facteur direct de $A^{\oplus n}$ pour
un certain $n$, cela résulte de l'hypothèse que
le groupe de cohomologie $H^i(K^\bullet)$ est dénombrable.

\medskip\noindent
Pour achever la démonstration, il suffit de montrer ceci : si $U = V \cup W$
et si le résultat vaut pour $V$, $W$ et $V \cap W$, alors
il vaut pour $U$. C'est une conséquence immédiate
de la suite de Mayer-Vietoris ; voir
Cohomologie, lemme \ref{cohomology-lemma-unbounded-mayer-vietoris}.
\end{proof}

\begin{lemma}
\label{lemma-countable}
Soit $X$ un schéma quasi-compact et quasi-séparé tel que
les ensembles de sections de $\mathcal{O}_X$ sur les ouverts affines soient dénombrables.
Soit $K$ un objet de $D_\QCoh(\mathcal{O}_X)$. Les
conditions suivantes sont équivalentes :
\begin{enumerate}
\item $K = \text{hocolim} E_n$, où $E_n$ est un objet parfait de
$D(\mathcal{O}_X)$ ;
\item les faisceaux de cohomologie $H^i(K)$ ont des ensembles
dénombrables de sections sur les ouverts affines.
\end{enumerate}
\end{lemma}

\begin{proof}
Si (1) est vraie, alors (2) l'est, car les colimites homotopiques commutent
à la prise des faisceaux de cohomologie
(d'après Catégories dérivées, lemme \ref{derived-lemma-cohomology-of-hocolim})
et parce qu'un complexe parfait est
localement isomorphe à un complexe fini de $\mathcal{O}_X$-modules libres de type fini
et satisfait donc (2) d'après l'hypothèse sur $X$.

\medskip\noindent
Supposons (2).
Choisissons un complexe K-injectif $\mathcal{K}^\bullet$ représentant $K$.
Choisissons un générateur parfait $E$ de $D_\QCoh(\mathcal{O}_X)$ et
représentons-le par un complexe K-injectif $\mathcal{I}^\bullet$.
D'après le théorème \ref{theorem-DQCoh-is-Ddga}
et sa démonstration, il existe une équivalence
de catégories triangulées $F : D_\QCoh(\mathcal{O}_X) \to D(A, \text{d})$,
où $(A, \text{d})$ est l'algèbre différentielle graduée
$$
(A, \text{d}) =
\Hom_{\text{Comp}^{dg}(\mathcal{O}_X)}
(\mathcal{I}^\bullet, \mathcal{I}^\bullet)
$$
qui envoie $K$ sur le module différentiel gradué
$$
M = \Hom_{\text{Comp}^{dg}(\mathcal{O}_X)}
(\mathcal{I}^\bullet, \mathcal{K}^\bullet)
$$
Notons que $H^i(A) = \Ext^i(E, E)$ et
$H^i(M) = \Ext^i(E, K)$.
De plus, puisque $F$ est une équivalence, elle et son quasi-inverse commutent
aux colimites homotopiques.
Il suffit donc d'écrire $M$ comme une colimite homotopique
d'objets compacts de $D(A, \text{d})$.
D'après Algèbre différentielle graduée, lemme \ref{dga-lemma-countable},
il suffit de montrer que $\Ext^i(E, E)$ et
$\Ext^i(E, K)$ sont dénombrables pour tout $i$.
Cela résulte du lemme \ref{lemma-countable-cohomology}.
\end{proof}

\begin{lemma}
\label{lemma-computing-sections-as-colim}
Soit $A$ un anneau. Soit $X$ un schéma de présentation finie sur $A$.
Soit $f : U \to X$ un morphisme plat de présentation finie. Alors
\begin{enumerate}
\item il existe un système projectif d'objets parfaits $L_n$ de
$D(\mathcal{O}_X)$ tel que
$$
R\Gamma(U, Lf^*K) = \text{hocolim}\ R\Hom_X(L_n, K)
$$
dans $D(A)$, fonctoriellement en $K$ dans $D_\QCoh(\mathcal{O}_X)$ ;
\item il existe un système d'objets parfaits $E_n$ de
$D(\mathcal{O}_X)$ tel que
$$
R\Gamma(U, Lf^*K) = \text{hocolim}\ R\Gamma(X, E_n \otimes^\mathbf{L} K)
$$
dans $D(A)$, fonctoriellement en $K$ dans $D_\QCoh(\mathcal{O}_X)$.
\end{enumerate}
\end{lemma}

\begin{proof}
D'après le lemme \ref{lemma-cohomology-base-change}, on a
$$
R\Gamma(U, Lf^*K) = R\Gamma(X, Rf_*\mathcal{O}_U \otimes^\mathbf{L} K)
$$
fonctoriellement en $K$. Observons que $R\Gamma(X, -)$ commute aux
colimites homotopiques, car il commute aux sommes directes d'après le
lemme \ref{lemma-quasi-coherence-pushforward-direct-sums}.
De même, $- \otimes^\mathbf{L} K$ commute aux colimites dérivées,
car $- \otimes^\mathbf{L} K$ commute aux sommes directes
(puisque les sommes directes dans $D(\mathcal{O}_X)$
sont données par les sommes directes des complexes qui les représentent).
Ainsi, pour démontrer (2), il suffit d'écrire
$Rf_*\mathcal{O}_U = \text{hocolim} E_n$ pour un système
d'objets parfaits $E_n$ de $D(\mathcal{O}_X)$. Cela fait,
on obtient (1) en posant $L_n = E_n^\vee$ ; voir
Cohomologie, lemme \ref{cohomology-lemma-dual-perfect-complex}.

\medskip\noindent
Écrivons $A = \colim A_i$, où $A_i$ est de type fini sur $\mathbf{Z}$.
D'après Limites, lemme \ref{limits-lemma-descend-finite-presentation},
on peut trouver un $i$ et des morphismes $U_i \to X_i \to \Spec(A_i)$
de présentation finie dont le changement de base à $\Spec(A)$ redonne
$U \to X \to \Spec(A)$.
Après avoir augmenté $i$, on peut supposer que $f_i : U_i \to X_i$ est
plat ; voir Limites, lemme \ref{limits-lemma-descend-flat-finite-presentation}.
D'après le lemme \ref{lemma-compare-base-change},
l'image inverse dérivée de $Rf_{i, *}\mathcal{O}_{U_i}$
par $g : X \to X_i$ est égale à $Rf_*\mathcal{O}_U$.
Puisque $Lg^*$ commute aux colimites dérivées, il suffit
de démontrer ce que l'on veut pour $f_i$. On peut donc supposer que
$U$ et $X$ sont de type fini sur $\mathbf{Z}$.

\medskip\noindent
Supposons que $f : U \to X$ soit un morphisme entre des schémas de type fini sur $\mathbf{Z}$.
Pour achever la démonstration, nous allons montrer que $Rf_*\mathcal{O}_U$ est une
colimite homotopique de complexes parfaits. Pour le voir, appliquons le lemme \ref{lemma-countable}.
Il suffit donc de montrer que $R^if_*\mathcal{O}_U$
a des ensembles dénombrables de sections sur les ouverts affines.
Cela résulte du lemme \ref{lemma-countable-cohomology}
appliqué au faisceau structural.
\end{proof}




\section{Caractérisation des complexes pseudo-cohérents, II}
\label{section-pseudo-coherent}

\noindent
Cette section prolonge la
section \ref{section-pseudo-coherent-hocolim}.
Nous y examinons des caractérisations des complexes pseudo-cohérents
en termes de cohomologie. On trouvera d'autres résultats de cette nature dans
Compléments sur les morphismes, section
\ref{more-morphisms-section-characterize-pseudo-coherent}.

\begin{lemma}
\label{lemma-pseudo-coherent-over-algebra}
Soit $A$ un anneau. Soit $R$ une $A$-algèbre (éventuellement non commutative)
qui est libre de type fini comme $A$-module. Alors tout objet $M$ de $D(R)$
qui est pseudo-cohérent dans $D(A)$ peut être représenté par un
complexe borné supérieurement de $R$-modules (à droite) libres de type fini.
\end{lemma}

\begin{proof}
Choisissons un complexe $M^\bullet$ de $R$-modules à droite représentant $M$.
Puisque $M$ est pseudo-cohérent, on a $H^i(M) = 0$ pour $i$ assez grand.
Soit $m$ le plus petit indice tel que $H^m(M)$ soit non nul.
Alors $H^m(M)$ est un $A$-module de type fini d'après
Compléments sur l'algèbre, lemme \ref{more-algebra-lemma-finite-cohomology}.
On peut donc choisir un $R$-module libre de type fini $F^m$ et une application
$F^m \to M^m$ tels que $F^m \to M^m \to M^{m + 1}$ soit nulle
et que $F^m \to H^m(M)$ soit surjective.
Diagramme :
$$
\xymatrix{
& F^m \ar[d]^\alpha \ar[r] & 0 \ar[d] \ar[r] & \ldots \\
M^{m - 1} \ar[r] & M^m \ar[r] & M^{m + 1} \ar[r] & \ldots
}
$$
Par récurrence descendante sur $n \leq m$, nous allons construire
des $R$-modules libres de type fini $F^i$ pour $i \geq n$, des différentielles
$d^i : F^i \to F^{i + 1}$ pour $i \geq n$, et des applications $\alpha : F^i \to K^i$
compatibles aux différentielles, telles que
(1) $H^i(\alpha)$ soit un isomorphisme pour $i > n$ et une surjection pour $i = n$, et
(2) $F^i = 0$ pour $i > m$. Diagramme :
$$
\xymatrix{
& F^n \ar[r] \ar[d]^\alpha & F^{n + 1} \ar[d]^\alpha \ar[r] & \ldots \ar[r]
& F^i \ar[d]^\alpha \ar[r] & 0 \ar[d] \ar[r] & \ldots \\
M^{n - 1} \ar[r] & M^n \ar[r] & M^{n + 1} \ar[r] & \ldots \ar[r] &
M^i \ar[r] & M^{i + 1} \ar[r] & \ldots
}
$$
Le cas initial $n = m$ a été traité ci-dessus.
Étape de récurrence. Soit $C^\bullet$ le cône de $\alpha$
(Catégories dérivées, définition \ref{derived-definition-cone}).
La suite exacte longue
de cohomologie montre que $H^i(C^\bullet) = 0$ pour $i \geq n$.
Observons que $F^\bullet$ est pseudo-cohérent comme complexe de $A$-modules,
car $R$ est libre de type fini comme $A$-module. Par conséquent,
d'après Compléments sur l'algèbre, lemme \ref{more-algebra-lemma-cone-pseudo-coherent},
on voit que $C^\bullet$ est $(n - 1)$-pseudo-cohérent comme complexe de
$A$-modules. D'après
Compléments sur l'algèbre, lemme \ref{more-algebra-lemma-finite-cohomology},
on voit que $H^{n - 1}(C^\bullet)$ est un $A$-module de type fini.
Choisissons un $R$-module libre de type fini $F^{n - 1}$ et une application
$\beta : F^{n - 1} \to C^{n - 1}$ tels que la composée
$F^{n - 1} \to C^{n - 1} \to C^n$ soit nulle et que $F^{n - 1}$
se surjecte sur $H^{n - 1}(C^\bullet)$. Puisque
$C^{n - 1} = M^{n - 1} \oplus F^n$,
on peut écrire $\beta = (\alpha^{n - 1}, -d^{n - 1})$. L'annulation de la
composée $F^{n - 1} \to C^{n - 1} \to C^n$ implique que
ces applications s'insèrent dans un morphisme de complexes
$$
\xymatrix{
& F^{n - 1} \ar[d]^{\alpha^{n - 1}} \ar[r]_{d^{n - 1}} &
F^n \ar[r] \ar[d]^\alpha &
F^{n + 1} \ar[d]^\alpha \ar[r] & \ldots \\
\ldots \ar[r] &
M^{n - 1} \ar[r] & M^n \ar[r] & M^{n + 1} \ar[r] & \ldots
}
$$
De plus, ces applications définissent un morphisme de triangles distingués
$$
\xymatrix{
(F^n \to \ldots) \ar[r] \ar[d] &
(F^{n - 1} \to \ldots) \ar[r] \ar[d] &
F^{n - 1} \ar[r] \ar[d]_\beta &
(F^n \to \ldots)[1] \ar[d] \\
(F^n \to \ldots) \ar[r] &
M^\bullet \ar[r] &
C^\bullet \ar[r] &
(F^n \to \ldots)[1]
}
$$
Notre choix de $\beta$ implique donc que le morphisme de complexes
$(F^{n - 1} \to \ldots) \to M^\bullet$ induit un isomorphisme en
cohomologie dans les degrés $\geq n$ et une surjection en degré $n - 1$.
Cela achève la démonstration du lemme.
\end{proof}

\begin{lemma}
\label{lemma-pseudo-coherent-on-projective-space}
Soit $A$ un anneau. Soit $n \geq 0$. Soit
$K \in D_\QCoh(\mathcal{O}_{\mathbf{P}^n_A})$.
Les conditions suivantes sont équivalentes :
\begin{enumerate}
\item $K$ est pseudo-cohérent ;
\item $R\Gamma(\mathbf{P}^n_A, E \otimes^\mathbf{L} K)$ est un objet pseudo-cohérent
de $D(A)$ pour tout objet pseudo-cohérent $E$ de
$D(\mathcal{O}_{\mathbf{P}^n_A})$ ;
\item $R\Gamma(\mathbf{P}^n_A, E \otimes^\mathbf{L} K)$ est un objet pseudo-cohérent
de $D(A)$ pour tout objet parfait $E$ de
$D(\mathcal{O}_{\mathbf{P}^n_A})$ ;
\item $R\Hom_{\mathbf{P}^n_A}(E, K)$ est un objet pseudo-cohérent
de $D(A)$ pour tout objet parfait $E$ de
$D(\mathcal{O}_{\mathbf{P}^n_A})$ ;
\item $R\Gamma(\mathbf{P}^n_A,
K \otimes^\mathbf{L} \mathcal{O}_{\mathbf{P}^n_A}(d))$ est un objet pseudo-cohérent
de $D(A)$ pour $d = 0, 1, \ldots, n$.
\end{enumerate}
\end{lemma}

\begin{proof}
Rappelons que
$$
R\Hom_{\mathbf{P}^n_A}(E, K) =
R\Gamma(\mathbf{P}^n_A, R\SheafHom_{\mathcal{O}_{\mathbf{P}^n_A}}(E, K))
$$
par définition ; voir Cohomologie, section \ref{cohomology-section-global-RHom}.
Ainsi, les assertions (4) et (3) sont équivalentes d'après
Cohomologie, lemme \ref{cohomology-lemma-dual-perfect-complex}.

\medskip\noindent
Puisque tout complexe parfait est pseudo-cohérent, il est clair que
(2) implique (3).

\medskip\noindent
Supposons (1). Alors $E \otimes^\mathbf{L} K$ est pseudo-cohérent
pour tout $E$ pseudo-cohérent ; voir
Cohomologie, lemme \ref{cohomology-lemma-tensor-pseudo-coherent}.
D'après le lemme \ref{lemma-flat-proper-pseudo-coherent-direct-image-general},
l'image directe d'un tel complexe pseudo-cohérent est pseudo-cohérente,
et l'on voit que (2) est vraie.

\medskip\noindent
L'assertion (3) implique (5), car on peut prendre $E = \mathcal{O}_{\mathbf{P}^n_A}(d)$
pour $d = 0, 1, \ldots, n$.

\medskip\noindent
Pour achever la démonstration, il faut montrer que (5) implique (1).
Soit $P$ comme dans (\ref{equation-generator-Pn}), et
$R$ comme dans (\ref{equation-algebra-for-Pn}).
D'après le lemme \ref{lemma-Pn-module-category}, on a une équivalence
$$
- \otimes^\mathbf{L}_R P :
D(R) \longrightarrow D_\QCoh(\mathcal{O}_{\mathbf{P}^n_A})
$$
Soit $M \in D(R)$ un objet tel que $M \otimes^\mathbf{L} P = K$.
D'après Algèbre différentielle graduée, lemme
\ref{dga-lemma-upgrade-tensor-with-complex-derived},
il existe un isomorphisme
$$
R\Hom(R, M) = R\Hom_{\mathbf{P}^n_A}(P, K)
$$
dans $D(A)$. En raisonnant comme ci-dessus, on obtient
$$
R\Hom_{\mathbf{P}^n_A}(P, K) = R\Gamma(\mathbf{P}^n_A,
R\SheafHom_{\mathcal{O}_{\mathbf{P}^n_A}}(E, K)) =
R\Gamma(\mathbf{P}^n_A,
P^\vee \otimes^\mathbf{L}_{\mathcal{O}_{\mathbf{P}^n_A}} K).
$$
Puisque $P^\vee$ est la somme directe des
$\mathcal{O}_{\mathbf{P}^n_A}(d)$ pour $d = 0, 1, \ldots, n$,
il résulte de (5) que $R\Hom(R, M)$ est pseudo-cohérent comme complexe
de $A$-modules. Bien entendu, $M = R\Hom(R, M)$ dans $D(A)$.
Ainsi, $M$ est pseudo-cohérent comme complexe de $A$-modules.
D'après le lemme \ref{lemma-pseudo-coherent-over-algebra},
on peut représenter $M$ par un complexe borné supérieurement $F^\bullet$
de $R$-modules libres de type fini. Alors
$F^\bullet = \bigcup_{p \geq 0} \sigma_{\geq p}F^\bullet$
est une filtration qui montre que $F^\bullet$ est un module différentiel
gradué sur $R$ ayant la propriété (P) ; voir
Algèbre différentielle graduée, section \ref{dga-section-P-resolutions}.
Par conséquent, $K = M \otimes^\mathbf{L}_R P$ est représenté
par $F^\bullet \otimes_R P$ (cela résulte de la construction du
foncteur de produit tensoriel dérivé ; voir, par exemple, la démonstration de
Algèbre différentielle graduée, lemme \ref{dga-lemma-tensor-with-complex-derived}).
Puisque $F^\bullet \otimes_R P$
est un complexe borné supérieurement dont les termes sont des sommes directes de
copies de $P$, on en déduit que le lemme est vrai.
\end{proof}

\begin{lemma}
\label{lemma-perfect-enough}
Soit $A$ un anneau. Soit $X$ un schéma quasi-compact
et quasi-séparé sur $A$. Soit $K \in D^-_\QCoh(\mathcal{O}_X)$.
Si $R\Gamma(X, E \otimes^\mathbf{L} K)$ est pseudo-cohérent
dans $D(A)$ pour tout $E$ parfait dans $D(\mathcal{O}_X)$,
alors $R\Gamma(X, E \otimes^\mathbf{L} K)$ est pseudo-cohérent
dans $D(A)$ pour tout $E$ pseudo-cohérent dans $D(\mathcal{O}_X)$.
\end{lemma}

\begin{proof}
Il existe un entier $N$ tel que
$R\Gamma(X, -) : D_\QCoh(\mathcal{O}_X) \to D(A)$
soit de dimension cohomologique $N$, comme expliqué dans le
lemme \ref{lemma-quasi-coherence-direct-image}.
Soit $b \in \mathbf{Z}$ tel que $H^i(K) = 0$ pour $i > b$.
Soit $E$ pseudo-cohérent sur $X$.
Il suffit de montrer que $R\Gamma(X, E \otimes^\mathbf{L} K)$
est $m$-pseudo-cohérent pour tout $m$.
Choisissons une approximation $P \to E$ par un complexe parfait $P$
de $(X, E, m - N - 1 - b)$. Cela est possible d'après le
théorème \ref{theorem-approximation}.
Choisissons un triangle distingué
$$
P \to E \to C \to P[1]
$$
dans $D_\QCoh(\mathcal{O}_X)$. Les faisceaux de cohomologie de $C$ sont nuls
dans les degrés $\geq m - N - 1 - b$. Par conséquent, les faisceaux de cohomologie
de $C \otimes^\mathbf{L} K$ sont nuls dans les degrés $\geq m - N - 1$.
Ainsi, les groupes de cohomologie de $R\Gamma(X, C \otimes^\mathbf{L} K)$
sont nuls dans les degrés $\geq m - 1$. Par conséquent,
$$
R\Gamma(X, P \otimes^\mathbf{L} K) \to R\Gamma(X, E \otimes^\mathbf{L} K)
$$
est un isomorphisme sur la cohomologie dans les degrés $\geq m$.
Par hypothèse, la source est pseudo-cohérente.
On en déduit que $R\Gamma(X, E \otimes^\mathbf{L} K)$
est $m$-pseudo-cohérent, comme désiré.
\end{proof}









\section{Objets relativement parfaits}
\label{section-relatively-perfect}

\noindent
Dans cette section, nous introduisons une notion tirée de
\cite{lieblich-complexes}.

\begin{definition}
\label{definition-relatively-perfect}
Soit $f : X \to S$ un morphisme de schémas plat et
localement de présentation finie. Un objet $E$ de $D(\mathcal{O}_X)$ est
{\it parfait relativement à $S$} ou
{\it $S$-parfait} si $E$ est pseudo-cohérent
(Cohomologie, définition \ref{cohomology-definition-pseudo-coherent}) et si
$E$ a localement une dimension de Tor finie comme objet de
$D(f^{-1}\mathcal{O}_S)$
(Cohomologie, définition \ref{cohomology-definition-tor-amplitude}).
\end{definition}

\noindent
Voir la remarque \ref{remark-discuss-rel-perfect} pour une discussion.

\begin{example}
\label{example-relatively-perfect-field}
Soit $k$ un corps. Soit $X$ un schéma de présentation finie sur $k$
(en particulier, $X$ est quasi-compact). Alors un objet $E$ de $D(\mathcal{O}_X)$
est $k$-parfait si et seulement s'il est borné et pseudo-cohérent
(par définition), c'est-à-dire si et seulement s'il appartient à $D^b_{\textit{Coh}}(X)$
(d'après le lemme \ref{lemma-identify-pseudo-coherent-noetherian}).
Ainsi, être relativement parfait ne signifie {\bf pas} « parfait sur les fibres ».
\end{example}

\noindent
La notion algébrique correspondante est étudiée dans
Compléments sur l'algèbre, section \ref{more-algebra-section-relatively-perfect}.
On peut relier la notion pour les schémas à la notion algébrique
comme suit.

\begin{lemma}
\label{lemma-affine-locally-rel-perfect}
Soit $f : X \to S$ un morphisme de schémas plat et
localement de présentation finie. Soit $E$ un objet de
$D_\QCoh(\mathcal{O}_X)$. Les conditions suivantes sont équivalentes :
\begin{enumerate}
\item $E$ est $S$-parfait ;
\item pour tout ouvert affine $U \subset X$ qui s'envoie dans un ouvert affine
$V \subset S$, le complexe $R\Gamma(U, E)$ est $\mathcal{O}_S(V)$-parfait ;
\item il existe un recouvrement ouvert affine $S = \bigcup V_i$
et, pour tout $i$, un recouvrement ouvert affine $f^{-1}(V_i) = \bigcup U_{ij}$
tels que le complexe $R\Gamma(U_{ij}, E)$ soit $\mathcal{O}_S(V_i)$-parfait.
\end{enumerate}
\end{lemma}

\begin{proof}
La pseudo-cohérence est une propriété locale, et
« avoir localement une dimension de Tor finie » est une propriété locale.
Le lemme se ramène donc immédiatement à l'énoncé suivant :
si $X$ et $S$ sont affines, alors $E$ est $S$-parfait
si et seulement si $K = R\Gamma(X, E)$ est $\mathcal{O}_S(S)$-parfait.
Écrivons $X = \Spec(A)$, $S = \Spec(R)$, et supposons que $E$ corresponde à
$K \in D(A)$, c'est-à-dire que $K = R\Gamma(X, E)$ ; voir le
lemme \ref{lemma-affine-compare-bounded}.

\medskip\noindent
Observons que $K$ est $R$-parfait si et seulement si $K$ est
pseudo-cohérent et a une dimension de Tor finie comme complexe
de $R$-modules (Compléments sur l'algèbre, définition
\ref{more-algebra-definition-relatively-perfect}).
D'après le lemme \ref{lemma-pseudo-coherent-affine},
on voit que $E$ est pseudo-cohérent si et seulement si
$K$ est pseudo-cohérent.
D'après le lemme \ref{lemma-tor-dimension-rel-affine}, on voit que
$E$ a une dimension de Tor finie sur $f^{-1}\mathcal{O}_S$
si et seulement si $K$ a une dimension de Tor finie comme complexe de
$R$-modules.
\end{proof}

\begin{lemma}
\label{lemma-triangulated}
Soit $f : X \to S$ un morphisme de schémas
plat et localement de présentation finie. La sous-catégorie pleine
de $D(\mathcal{O}_X)$ formée des objets $S$-parfaits est
une sous-catégorie triangulée saturée\footnote{Catégories dérivées, définition
\ref{derived-definition-saturated}.}.
\end{lemma}

\begin{proof}
Cela résulte de Cohomologie, lemmes
\ref{cohomology-lemma-cone-pseudo-coherent},
\ref{cohomology-lemma-summands-pseudo-coherent},
\ref{cohomology-lemma-cone-tor-amplitude} et
\ref{cohomology-lemma-summands-tor-amplitude}.
\end{proof}

\begin{lemma}
\label{lemma-perfect-relatively-perfect}
Soit $f : X \to S$ un morphisme de schémas plat et localement
de présentation finie. Tout objet parfait de $D(\mathcal{O}_X)$ est $S$-parfait.
Si $K, M \in D(\mathcal{O}_X)$, alors $K \otimes_{\mathcal{O}_X}^\mathbf{L} M$
est $S$-parfait si $K$ est parfait et $M$ est $S$-parfait.
\end{lemma}

\begin{proof}
Première démonstration : on se ramène au cas affine à l'aide du
lemme \ref{lemma-affine-locally-rel-perfect},
puis on applique Compléments sur l'algèbre, lemme
\ref{more-algebra-lemma-perfect-relatively-perfect}.
\end{proof}

\begin{lemma}
\label{lemma-base-change-relatively-perfect}
Soit $f : X \to S$ un morphisme de schémas plat et
localement de présentation finie.
Soit $g : S' \to S$ un morphisme de schémas. Posons $X' = S' \times_S X$
et notons $g' : X' \to X$ la projection.
Si $K \in D(\mathcal{O}_X)$ est $S$-parfait, alors $L(g')^*K$
est $S'$-parfait.
\end{lemma}

\begin{proof}
Première démonstration : on se ramène au cas affine à l'aide du
lemme \ref{lemma-affine-locally-rel-perfect},
puis on applique Compléments sur l'algèbre, lemme
\ref{more-algebra-lemma-base-change-relatively-perfect}.

\medskip\noindent
Deuxième démonstration : $L(g')^*K$ est pseudo-cohérent d'après
Cohomologie, lemme \ref{cohomology-lemma-pseudo-coherent-pullback},
et la propriété de dimension de Tor bornée résulte du
lemme \ref{lemma-tor-independence-and-tor-amplitude}.
\end{proof}

\begin{situation}
\label{situation-relative-descent}
Soit $S = \lim_{i \in I} S_i$ la limite d'un système filtrant de schémas
dont les morphismes de transition affines sont $g_{i'i} : S_{i'} \to S_i$.
On suppose que $S_i$ est quasi-compact et quasi-séparé pour tout $i \in I$.
On note $g_i : S \to S_i$ la projection. On fixe un élément $0 \in I$
et un morphisme plat de présentation finie $X_0 \to S_0$.
On pose $X_i = S_i \times_{S_0} X_0$ et $X = S \times_{S_0} X_0$,
et l'on note les morphismes de transition $f_{i'i} : X_{i'} \to X_i$
et les projections $f_i : X \to X_i$.
\end{situation}

\begin{lemma}
\label{lemma-relative-descend-homomorphisms}
Dans la situation \ref{situation-relative-descent},
soient $K_0$ et $L_0$ des objets de $D(\mathcal{O}_{X_0})$.
Posons $K_i = Lf_{i0}^*K_0$ et $L_i = Lf_{i0}^*L_0$ pour $i \geq 0$,
et posons $K = Lf_0^*K_0$ et $L = Lf_0^*L_0$. Alors l'application
$$
\colim_{i \geq 0} \Hom_{D(\mathcal{O}_{X_i})}(K_i, L_i)
\longrightarrow
\Hom_{D(\mathcal{O}_X)}(K, L)
$$
est un isomorphisme si $K_0$ est pseudo-cohérent et si
$L_0 \in D_\QCoh(\mathcal{O}_{X_0})$ a (localement)
une dimension de Tor finie comme objet de
$D((X_0 \to S_0)^{-1}\mathcal{O}_{S_0})$.
\end{lemma}

\begin{proof}
Pour tout ouvert quasi-compact $U_0 \subset X_0$, considérons la
condition $P$ selon laquelle
$$
\colim_{i \geq 0} \Hom_{D(\mathcal{O}_{U_i})}(K_i|_{U_i}, L_i|_{U_i})
\longrightarrow
\Hom_{D(\mathcal{O}_U)}(K|_U, L|_U)
$$
est un isomorphisme, où $U = f_0^{-1}(U_0)$ et $U_i = f_{i0}^{-1}(U_0)$.
Si $P$ vaut pour $U_0$, $V_0$ et $U_0 \cap V_0$, alors elle vaut
pour $U_0 \cup V_0$ par Mayer-Vietoris
pour les Hom dans la catégorie dérivée ; voir
Cohomologie, lemme \ref{cohomology-lemma-mayer-vietoris-hom}.

\medskip\noindent
Notons $\pi_0 : X_0 \to S_0$ le morphisme donné.
On peut d'abord considérer $U_0 = \pi_0^{-1}(W_0)$ avec
$W_0 \subset S_0$ ouvert quasi-compact. En appliquant le principe d'induction de
Cohomologie des schémas, lemme \ref{coherent-lemma-induction-principle},
aux ouverts quasi-compacts de $S_0$
et en utilisant la remarque ci-dessus, on constate qu'il suffit de démontrer
$P$ pour $U_0 = \pi_0^{-1}(W_0)$ avec $W_0$ affine.
Autrement dit, on s'est ramené au cas où $S_0$ est affine.
Ensuite, on applique de nouveau le principe d'induction, cette fois à tous les
ouverts quasi-compacts et quasi-séparés de $X_0$, pour se ramener également au
cas où $X_0$ est affine.

\medskip\noindent
Si $X_0$ et $S_0$ sont affines, le résultat découle de
Compléments sur l'algèbre, lemme \ref{more-algebra-lemma-colimit-relatively-perfect}.
En effet, d'après les lemmes \ref{lemma-pseudo-coherent} et
\ref{lemma-affine-compare-bounded},
l'énoncé se traduit en calculs de morphismes dans les
catégories dérivées de modules. Puis le
lemme \ref{lemma-pseudo-coherent-affine}
montre que le complexe de modules correspondant à $K_0$
est pseudo-cohérent. Et le
lemme \ref{lemma-tor-dimension-rel-affine}
montre que le complexe de modules correspondant à $L_0$
a une dimension de Tor finie sur $\mathcal{O}_{S_0}(S_0)$.
Les hypothèses de
Compléments sur l'algèbre, lemme \ref{more-algebra-lemma-colimit-relatively-perfect},
sont donc satisfaites, et l'on conclut.
\end{proof}

\begin{lemma}
\label{lemma-descend-relatively-perfect}
Dans la situation \ref{situation-relative-descent}, la catégorie des
objets $S$-parfaits de $D(\mathcal{O}_X)$ est la colimite des catégories
d'objets $S_i$-parfaits de $D(\mathcal{O}_{X_i})$.
\end{lemma}

\begin{proof}
Pour tout ouvert quasi-compact $U_0 \subset X_0$, considérons la condition $P$
selon laquelle le foncteur
$$
\colim_{i \geq 0} D_{S_i\text{-parfait}}(\mathcal{O}_{U_i})
\longrightarrow
D_{S\text{-parfait}}(\mathcal{O}_U)
$$
est une équivalence, où $U = f_0^{-1}(U_0)$ et $U_i = f_{i0}^{-1}(U_0)$.
Observons que nous savons déjà que ce foncteur est pleinement fidèle
d'après le lemme \ref{lemma-relative-descend-homomorphisms}. Il suffit donc de démontrer
la surjectivité essentielle.

\medskip\noindent
Supposons que $P$ vaille pour des ouverts quasi-compacts $U_0$, $V_0$ de $X_0$.
Nous affirmons que $P$ vaut pour $U_0 \cup V_0$. Nous utiliserons les notations
$U_i = f_{i0}^{-1}U_0$, $U = f_0^{-1}U_0$, $V_i = f_{i0}^{-1}V_0$,
et $V = f_0^{-1}V_0$, et nous emploierons abusivement le symbole
$f_i$ pour tous les morphismes $U \to U_i$, $V \to V_i$,
$U \cap V \to U_i \cap V_i$ et $U \cup V \to U_i \cup V_i$.
Soit $E$ un objet $S$-parfait de $D(\mathcal{O}_{U \cup V})$.
But : montrer que $E$ appartient à l'image essentielle du foncteur.
Par hypothèse,
on peut trouver $i \geq 0$, un objet $S_i$-parfait $E_{U, i}$ sur $U_i$,
un objet $S_i$-parfait $E_{V, i}$ sur $V_i$, et
des isomorphismes $Lf_i^*E_{U, i} \to E|_U$ et $Lf_i^*E_{V, i} \to E|_V$.
Soient
$$
a : E_{U, i} \to (Rf_{i, *}E)|_{U_i}
\quad\text{et}\quad
b : E_{V, i} \to (Rf_{i, *}E)|_{V_i}
$$
les morphismes adjoints aux isomorphismes $Lf_i^*E_{U, i} \to E|_U$
et $Lf_i^*E_{V, i} \to E|_V$.
Par pleine fidélité, après avoir augmenté $i$,
on peut trouver un isomorphisme
$c : E_{U, i}|_{U_i \cap V_i} \to E_{V, i}|_{U_i \cap V_i}$
dont l'image inverse donne les identifications
$$
Lf_i^*E_{U, i}|_{U \cap V} \to E|_{U \cap V} \to Lf_i^*E_{V, i}|_{U \cap V}.
$$
Appliquons Cohomologie, lemme \ref{cohomology-lemma-glue},
pour obtenir un objet $E_i$ sur $U_i \cup V_i$ et un morphisme $d : E_i \to Rf_{i, *}E$
qui se restreint aux morphismes $a$ et $b$ sur $U_i$ et $V_i$.
Il est alors clair que $E_i$ est $S_i$-parfait (car être
relativement parfait est une propriété locale) et que
$d$ est adjoint à un isomorphisme $Lf_i^*E_i \to E$.

\medskip\noindent
Par exactement le même argument que dans
la démonstration du lemme \ref{lemma-relative-descend-homomorphisms},
en utilisant le principe d'induction
(Cohomologie des schémas, lemme \ref{coherent-lemma-induction-principle}),
on se ramène au cas où $X_0$ et $S_0$ sont tous deux
affines. (On travaille d'abord avec les ouverts de $S_0$ pour se ramener à
$S_0$ affine, puis avec les ouverts de $X_0$ pour se ramener à
$X_0$ affine.) Dans le cas affine, le résultat découle de
Compléments sur l'algèbre, lemme \ref{more-algebra-lemma-colimit-relatively-perfect}.
La traduction en algèbre est donnée par le
lemme \ref{lemma-affine-locally-rel-perfect}.
\end{proof}

\begin{lemma}
\label{lemma-derived-pushforward-rel-perfect}
Soit $f : X \to S$ un morphisme de schémas plat, propre et
de présentation finie. Soit $E \in D(\mathcal{O}_X)$ un objet $S$-parfait.
Alors $Rf_*E$ est un objet parfait de $D(\mathcal{O}_S)$,
et sa formation commute à tout changement de base.
\end{lemma}

\begin{proof}
L'énoncé sur le changement de base est le lemme \ref{lemma-compare-base-change}.
Il suffit donc de montrer que $Rf_*E$ est un objet parfait. Nous allons nous ramener
au cas où $S$ est noethérien affine par un argument de limites.

\medskip\noindent
La question est locale sur $S$ ; on peut donc supposer $S$ affine.
Écrivons $S = \Spec(R)$. Écrivons $R = \colim R_i$ comme colimite filtrante
d'anneaux noethériens $R_i$. D'après Limites, lemme
\ref{limits-lemma-descend-finite-presentation},
il existe un $i$ et un schéma $X_i$ de présentation finie sur $R_i$
dont le changement de base à $R$ est $X$. D'après
Limites, lemmes \ref{limits-lemma-eventually-proper} et
\ref{limits-lemma-descend-flat-finite-presentation},
on peut supposer $X_i$ propre et plat sur $R_i$.
D'après le lemme \ref{lemma-descend-relatively-perfect},
on peut supposer qu'il existe un objet $R_i$-parfait $E_i$ de
$D(\mathcal{O}_{X_i})$ dont l'image inverse sur $X$ est $E$.
En appliquant le lemme \ref{lemma-perfect-direct-image}
à $X_i \to \Spec(R_i)$ et à $E_i$, puis en utilisant la
propriété de changement de base déjà démontrée, on obtient le résultat.
\end{proof}

\begin{lemma}
\label{lemma-compute-ext-rel-perfect}
Soit $f : X \to S$ un morphisme de schémas. Soient $E, K \in D(\mathcal{O}_X)$.
Supposons que
\begin{enumerate}
\item $S$ soit quasi-compact et quasi-séparé ;
\item $f$ soit propre, plat et de présentation finie ;
\item $E$ soit $S$-parfait ;
\item $K$ soit pseudo-cohérent.
\end{enumerate}
Alors il existe un $L \in D(\mathcal{O}_S)$ pseudo-cohérent tel que
$$
Rf_*R\SheafHom(K, E) = R\SheafHom(L, \mathcal{O}_S)
$$
et il en va de même après tout changement de base : étant donné
$$
\vcenter{
\xymatrix{
X' \ar[r]_{g'} \ar[d]_{f'} &
X \ar[d]^f \\
S' \ar[r]^g &
S
}
}
\quad\quad
\begin{matrix}
\text{cartésien, alors on a } \\
Rf'_*R\SheafHom(L(g')^*K, L(g')^*E) \\
= R\SheafHom(Lg^*L, \mathcal{O}_{S'})
\end{matrix}
$$
\end{lemma}

\begin{proof}
Puisque $S$ est quasi-compact et quasi-séparé, il en va de même de $X$.
D'après le lemme \ref{lemma-pseudo-coherent-hocolim}, on peut écrire
$K = \text{hocolim} K_n$, où $K_n$ est parfait et où $K_n \to K$ induit
un isomorphisme sur les troncatures $\tau_{\geq -n}$. Soit $K_n^\vee$
le complexe parfait dual
(Cohomologie, lemme \ref{cohomology-lemma-dual-perfect-complex}).
On obtient un système projectif $\ldots \to K_3^\vee \to K_2^\vee \to K_1^\vee$
d'objets parfaits. D'après le lemme \ref{lemma-perfect-relatively-perfect},
on voit que $K_n^\vee \otimes_{\mathcal{O}_X} E$ est $S$-parfait.
On peut donc appliquer le lemme \ref{lemma-derived-pushforward-rel-perfect}
à $K_n^\vee \otimes_{\mathcal{O}_X} E$, et l'on obtient un système projectif
$$
\ldots \to M_3 \to M_2 \to M_1
$$
de complexes parfaits sur $S$, avec
$$
M_n = Rf_*(K_n^\vee \otimes_{\mathcal{O}_X}^\mathbf{L} E) =
Rf_*R\SheafHom(K_n, E)
$$
De plus, la formation de ces complexes commute à tout
changement de base, à savoir $Lg^*M_n =
Rf'_*((L(g')^*K_n)^\vee \otimes_{\mathcal{O}_{X'}}^\mathbf{L} L(g')^*E) =
Rf'_*R\SheafHom(L(g')^*K_n, L(g')^*E)$.

\medskip\noindent
Puisque $K_n \to K$ induit un isomorphisme sur $\tau_{\geq -n}$, on voit que
$K_n \to K_{n + 1}$ induit un isomorphisme sur $\tau_{\geq -n}$.
Il en résulte que $K_{n + 1}^\vee \to K_n^\vee$
induit un isomorphisme sur $\tau_{\leq n}$, car
$K_n^\vee = R\SheafHom(K_n, \mathcal{O}_X)$.
Supposons que $E$ ait une amplitude de Tor contenue dans $[a, b]$ comme complexe
de $f^{-1}\mathcal{O}_Y$-modules. Il en va alors de même après
tout changement de base ; voir le lemme \ref{lemma-tor-independence-and-tor-amplitude}.
On constate que
$K_{n + 1}^\vee \otimes_{\mathcal{O}_X} E \to
K_n^\vee \otimes_{\mathcal{O}_X} E$
induit un isomorphisme sur $\tau_{\leq n + a}$,
et il en va de même après tout changement de base.
En appliquant le foncteur dérivé droit $Rf_*$,
on en déduit que les morphismes $M_{n + 1} \to M_n$
induisent des isomorphismes sur $\tau_{\leq n + a}$,
et il en va de même après tout changement de base.
Choisissons un triangle distingué
$$
M_{n + 1} \to M_n \to C_n \to M_{n + 1}[1]
$$
Prenons pour $S'$ le spectre du corps résiduel en un point
$s \in S$ et effectuons l'image inverse pour voir que
$C_n \otimes_{\mathcal{O}_S}^\mathbf{L} \kappa(s)$
n'a de cohomologie non nulle que dans les degrés $\geq n + a$. D'après
Compléments sur l'algèbre, lemme
\ref{more-algebra-lemma-lift-perfect-from-residue-field},
on voit que le complexe parfait $C_n$ a une amplitude de Tor contenue dans
$[n + a, m_n]$ pour un certain entier $m_n$.
En particulier, le complexe parfait dual $C_n^\vee$ a une amplitude de Tor contenue dans
$[-m_n, -n - a]$.

\medskip\noindent
Soit $L_n = M_n^\vee$ le complexe parfait dual. La
conclusion de la discussion du paragraphe précédent est que
$L_n \to L_{n + 1}$ induit des isomorphismes sur $\tau_{\geq -n - a}$.
Ainsi, $L = \text{hocolim} L_n$ est pseudo-cohérent ; voir le
lemme \ref{lemma-pseudo-coherent-hocolim}.
Puisque l'on a
$$
R\SheafHom(K, E) = R\SheafHom(\text{hocolim} K_n, E) =
R\lim R\SheafHom(K_n, E) = R\lim K_n^\vee \otimes_{\mathcal{O}_X} E
$$
(Cohomologie, lemme \ref{cohomology-lemma-colim-and-lim-of-duals})
et puisque $R\lim$ commute à $Rf_*$, on trouve que
$$
Rf_*R\SheafHom(K, E) = R\lim M_n = R\lim R\SheafHom(L_n, \mathcal{O}_S) =
R\SheafHom(L, \mathcal{O}_S)
$$
Cela démontre la formule sur $S$. Puisque la construction de $M_n$ est
compatible au changement de base, la formule reste vraie après
tout changement de base.
\end{proof}

\begin{remark}
\label{remark-compare-L}
Le lecteur aura peut-être remarqué la similitude entre le
lemme \ref{lemma-compute-ext-rel-perfect} et le
lemme \ref{lemma-compute-ext}.
En effet, le complexe pseudo-cohérent $L$ du
lemme \ref{lemma-compute-ext-rel-perfect}
peut être caractérisé comme l'unique complexe pseudo-cohérent
sur $S$ tel qu'il existe des isomorphismes fonctoriels
$$
\Ext^i_{\mathcal{O}_S}(L, \mathcal{F}) \longrightarrow
\Ext^i_{\mathcal{O}_X}(K,
E \otimes_{\mathcal{O}_X}^\mathbf{L} Lf^*\mathcal{F})
$$
compatibles aux morphismes de bord, lorsque $\mathcal{F}$ parcourt
$\QCoh(\mathcal{O}_S)$. Si cela nous est un jour nécessaire, nous
formulerons ici un résultat précis et en donnerons une démonstration détaillée.
\end{remark}

\begin{lemma}
\label{lemma-bounded-on-fibres}
Soit $f : X \to S$ un morphisme de schémas plat et
localement de présentation finie. Soit $E$ un objet pseudo-cohérent
de $D(\mathcal{O}_X)$. Les conditions suivantes sont équivalentes :
\begin{enumerate}
\item $E$ est $S$-parfait ;
\item $E$ est localement borné inférieurement et, pour tout point $s \in S$,
l'objet $L(X_s \to X)^*E$ de $D(\mathcal{O}_{X_s})$
est localement borné inférieurement.
\end{enumerate}
\end{lemma}

\begin{proof}
Puisque tout est local, on se ramène immédiatement au
cas où $X$ et $S$ sont affines ; voir le lemme
\ref{lemma-affine-locally-rel-perfect}.
Supposons que $X \to S$ corresponde à $\Spec(A) \to \Spec(R)$ et que
$E$ corresponde à $K$ dans $D(A)$. Si $s$ correspond à
l'idéal premier $\mathfrak p \subset R$, alors $L(X_s \to X)^*E$
correspond à $K \otimes_R^\mathbf{L} \kappa(\mathfrak p)$,
car $R \to A$ est plat ; voir, par exemple, le
lemme \ref{lemma-compare-base-change}.
On voit donc que notre lemme résulte du résultat algébrique
correspondant ; voir Compléments sur l'algèbre, lemme
\ref{more-algebra-lemma-bounded-on-fibres-relatively-perfect}.
\end{proof}

\begin{remark}
\label{remark-discuss-rel-perfect}
Notre définition \ref{definition-relatively-perfect} d'un
complexe relativement parfait est équivalente à celle donnée
dans \cite{lieblich-complexes} chaque fois que notre définition s'applique\footnote{Pour
le voir, utiliser le lemme \ref{lemma-affine-locally-rel-perfect} et
Compléments sur l'algèbre, lemme \ref{more-algebra-lemma-structure-relatively-perfect}.}.
Supposons maintenant que $f : X \to S$ soit seulement localement
de type fini (sans être nécessairement plat ni localement de
présentation finie). La définition de l'article cité ci-dessus est que
$E \in D(\mathcal{O}_X)$ est relativement parfait si
\begin{enumerate}
\item[(A)] localement sur $X$, l'objet $E$ est
quasi-isomorphe à un complexe fini de
modules plats sur $S$ et de présentation finie comme $\mathcal{O}_X$-modules.
\end{enumerate}
D'autre part, la généralisation naturelle de notre
définition \ref{definition-relatively-perfect} est la suivante :
\begin{enumerate}
\item[(B)] $E$ est pseudo-cohérent relativement à $S$
(Compléments sur les morphismes, définition
\ref{more-morphisms-definition-relative-pseudo-coherence})
et $E$ a localement une dimension de Tor finie comme objet de
$D(f^{-1}\mathcal{O}_S)$
(Cohomologie, définition \ref{cohomology-definition-tor-amplitude}).
\end{enumerate}
L'avantage de la condition (B) est qu'elle définit manifestement une sous-catégorie
triangulée de $D(\mathcal{O}_X)$, tandis que nous soupçonnons que ce n'est pas
le cas de la condition (A). L'avantage de la condition (A)
est qu'elle est plus facile à manier, en particulier en ce qui concerne les limites.
\end{remark}








\section{La propriété de résolution}
\label{section-resolution-property}

\noindent
Cette notion est étudiée dans l'article \cite{totaro_resolution} ;
la discussion se poursuit dans \cite{Gross-thesis},
\cite{Gross-surface} et \cite{Gross-stack}.
On ignore actuellement si tout schéma propre sur un corps
possède toujours la propriété de résolution ou si cela est faux.
Si vous connaissez la réponse à cette question, veuillez écrire à
\href{mailto:stacks.project@gmail.com}{stacks.project@gmail.com}.

\medskip\noindent
On peut donner la définition suivante, bien qu'il ne soit guère pertinent
de la considérer pour des schémas généraux.

\begin{definition}
\label{definition-resolution-property}
Soit $X$ un schéma. On dit que $X$ possède la {\it propriété de résolution}
si tout $\mathcal{O}_X$-module quasi-cohérent de type fini
est quotient d'un $\mathcal{O}_X$-module localement libre de type fini.
\end{definition}

\noindent
Si $X$ est un schéma quasi-compact et quasi-séparé, il suffit de vérifier que
tout $\mathcal{O}_X$-module de présentation finie (automatiquement
quasi-cohérent) est quotient d'un module localement libre de type fini
sur $\mathcal{O}_X$ ; voir Propriétés, lemme
\ref{properties-lemma-finite-directed-colimit-surjective-maps}.
Si $X$ est un schéma noethérien, les modules quasi-cohérents de type fini
sont exactement les $\mathcal{O}_X$-modules cohérents ; voir
Cohomologie des schémas, lemme \ref{coherent-lemma-coherent-Noetherian}.

\begin{lemma}
\label{lemma-resolution-property-ample}
Soit $X$ un schéma. Si $X$ possède un $\mathcal{O}_X$-module inversible ample,
alors $X$ possède la propriété de résolution.
\end{lemma}

\begin{proof}
C'est une conséquence immédiate de Propriétés, proposition
\ref{properties-proposition-characterize-ample}.
\end{proof}

\begin{lemma}
\label{lemma-resolution-property-ample-relative}
Soit $f : X \to Y$ un morphisme de schémas. Supposons que
\begin{enumerate}
\item $Y$ soit quasi-compact et quasi-séparé et possède la propriété de résolution ;
\item il existe un module inversible $f$-ample sur $X$.
\end{enumerate}
Alors $X$ possède la propriété de résolution.
\end{lemma}

\begin{proof}
Soit $\mathcal{F}$ un $\mathcal{O}_X$-module quasi-cohérent de type fini.
Soit $\mathcal{L}$ un module inversible $f$-ample.
Choisissons un recouvrement ouvert affine $Y = V_1 \cup \ldots \cup V_m$.
Posons $U_j = f^{-1}(V_j)$. D'après Propriétés, proposition
\ref{properties-proposition-characterize-ample},
pour tout $j$, on sait qu'il existe un nombre fini de morphismes
$s_{j, i} : \mathcal{L}^{\otimes n_{j, i}}|_{U_j} \to \mathcal{F}|_{U_j}$
qui sont conjointement surjectifs. Considérons les
$\mathcal{O}_Y$-modules quasi-cohérents
$$
\mathcal{H}_n =
f_*(\mathcal{F} \otimes_{\mathcal{O}_X} \mathcal{L}^{\otimes n})
$$
On peut considérer $s_{j, i}$ comme une section sur $V_j$ du faisceau
$\mathcal{H}_{-n_{j, i}}$. Supposons que l'on puisse trouver des
$\mathcal{O}_Y$-modules localement libres de type fini $\mathcal{E}_{i, j}$ et des morphismes
$\mathcal{E}_{i, j} \to \mathcal{H}_{-n_{j, i}}$ tels que
$s_{j, i}$ appartienne à leur image. Alors les morphismes correspondants
$$
f^*\mathcal{E}_{i, j}
\otimes_{\mathcal{O}_X}
\mathcal{L}^{\otimes n_{i, j}} \longrightarrow \mathcal{F}
$$
seront conjointement surjectifs, ce qui démontrera le lemme. D'après
Propriétés, lemme \ref{properties-lemma-quasi-coherent-colimit-finite-type},
pour tout $i, j$, on peut trouver un sous-module quasi-cohérent
de type fini
$\mathcal{H}'_{i, j} \subset  \mathcal{H}_{-n_{j, i}}$
qui contient la section $s_{i, j}$ sur $V_j$.
La propriété de résolution de $Y$ fournit donc des surjections
$\mathcal{E}_{i, j} \to \mathcal{H}'_{j, i}$, et l'on conclut.
\end{proof}

\begin{lemma}
\label{lemma-resolution-property-goes-up-affine}
Soit $f : X \to Y$ un morphisme affine ou quasi-affine de schémas, où
$Y$ est quasi-compact et quasi-séparé.
Si $Y$ possède la propriété de résolution, alors $X$ la possède aussi.
\end{lemma}

\begin{proof}
D'après Morphismes, lemme \ref{morphisms-lemma-quasi-affine-O-ample},
c'est un cas particulier du
lemme \ref{lemma-resolution-property-ample-relative}.
\end{proof}

\noindent
Voici un cas où l'on peut démontrer que la propriété de résolution descend.

\begin{lemma}
\label{lemma-resolution-property-goes-down-finite-flat}
Soit $f : X \to Y$ un morphisme de schémas surjectif et fini localement libre.
Si $X$ possède la propriété de résolution, alors $Y$ la possède aussi.
\end{lemma}

\begin{proof}
La condition signifie que $f$ est affine et que $f_*\mathcal{O}_X$
est un $\mathcal{O}_Y$-module localement libre de type fini et de rang positif.
Soit $\mathcal{G}$ un $\mathcal{O}_Y$-module quasi-cohérent
de type fini. Par hypothèse, il existe une surjection
$\mathcal{E} \to f^*\mathcal{G}$ pour un certain $\mathcal{O}_X$-module
$\mathcal{E}$ localement libre de type fini. Puisque $f_*$ est exact sur les
modules quasi-cohérents
(Cohomologie des schémas, lemme \ref{coherent-lemma-relative-affine-vanishing}),
on obtient une surjection
$$
f_*\mathcal{E}
\longrightarrow
f_*f^*\mathcal{G} = \mathcal{G} \otimes_{\mathcal{O}_Y} f_*\mathcal{O}_X
$$
En prenant les duaux, on obtient une surjection
$$
f_*\mathcal{E}
\otimes_{\mathcal{O}_Y}
\SheafHom_{\mathcal{O}_Y}(f_*\mathcal{O}_X, \mathcal{O}_Y)
\longrightarrow
\mathcal{G}
$$
Puisque $f_*\mathcal{E}$ est localement libre de type fini\footnote{En effet,
si $A \to B$ est un homomorphisme d'anneaux fini localement libre
et si $N$ est un $B$-module localement libre de type fini, alors $N$
est un $A$-module localement libre de type fini. Pour le voir, notons d'abord
que le fait que $N$ soit localement libre de type fini sur $B$ implique que $N$ est
plat et de présentation finie comme $B$-module ; voir
Algèbre, lemme \ref{algebra-lemma-finite-projective}. Ensuite,
$N$ est un $A$-module de présentation finie d'après
Algèbre, lemme \ref{algebra-lemma-finite-finitely-presented-extension},
et un $A$-module plat d'après Algèbre, lemme \ref{algebra-lemma-composition-flat}.
On conclut alors en appliquant
Algèbre, lemme \ref{algebra-lemma-finite-projective}, sur $A$.},
on conclut.
\end{proof}

\begin{lemma}
\label{lemma-resolution-property-finite-number}
Soit $X$ un schéma. Supposons donnés
\begin{enumerate}
\item un recouvrement ouvert affine fini $X = U_1 \cup \ldots \cup U_m$ ;
\item des idéaux quasi-cohérents de type fini $\mathcal{I}_j$
tels que $V(\mathcal{I}_j) = X \setminus U_j$.
\end{enumerate}
Alors $X$ possède la propriété de résolution si et seulement si $\mathcal{I}_j$
est quotient d'un $\mathcal{O}_X$-module localement libre de type fini
pour $j = 1, \ldots, m$.
\end{lemma}

\begin{proof}
Une implication du lemme est triviale. Pour l'autre,
supposons que $\mathcal{E}_j \to \mathcal{I}_j$ soit une surjection, avec
$\mathcal{E}_j$ localement libre de type fini. Dans le paragraphe suivant, nous
nous ramenons au cas noethérien ; nous suggérons au lecteur de l'omettre.

\medskip\noindent
La première observation est que $U_j \cap U_{j'}$ est quasi-compact,
car c'est le complémentaire du schéma des zéros de l'idéal quasi-cohérent de type fini
$\mathcal{I}_{j'}|{U_j}$ sur le schéma affine $U_j$ ; voir
Propriétés, lemme \ref{properties-lemma-quasi-coherent-finite-type-ideals}.
Ainsi, $X$ est quasi-compact et quasi-séparé ; voir
Schémas, lemme \ref{schemes-lemma-characterize-quasi-separated}.
D'après Limites, proposition \ref{limits-proposition-approximate},
on peut écrire $X = \lim X_i$ comme limite d'un système inductif de
schémas noethériens à morphismes de transition affines. Pour tout $j$,
on peut trouver un $i$ et un $\mathcal{O}_{X_i}$-module localement libre de type fini
$\mathcal{E}_{i, j}$ dont l'image inverse est $\mathcal{E}_j$ ; voir
Limites, lemme \ref{limits-lemma-descend-invertible-modules}.
Après avoir augmenté $i$, on peut supposer que la composée
$\mathcal{E}_j \to \mathcal{I}_j \to \mathcal{O}_X$ est l'image
inverse d'un morphisme $\mathcal{E}_{i, j} \to \mathcal{O}_{X_i}$ ; voir
Limites, lemme \ref{limits-lemma-descend-modules-finite-presentation}.
Notons $\mathcal{I}_{i, j} \subset \mathcal{O}_{X_i}$ l'image
de ce morphisme ; c'est un faisceau d'idéaux quasi-cohérent sur le schéma noethérien
$X_i$, dont l'image inverse sur $X$ est $\mathcal{I}_j$.
En notant $U_{i, j} \subset X_i$ les ouverts complémentaires, on
peut supposer qu'ils sont affines pour tous $i, j$ ; voir
Limites, lemme \ref{limits-lemma-limit-affine}.
Si l'on peut démontrer le lemme pour les ouverts $U_{i, j}$ et les
faisceaux d'idéaux $\mathcal{I}_{i, j}$ sur $X_i$, alors $X$, étant affine
sur $X_i$, possédera la propriété de résolution d'après le
lemme \ref{lemma-resolution-property-goes-up-affine}.
De cette manière, on se ramène au cas d'un schéma noethérien.

\medskip\noindent
Supposons $X$ noethérien. Pour tout module cohérent $\mathcal{F}$,
on peut choisir une liste finie de sections $s_{jk} \in \mathcal{F}(U_j)$,
$k = 1, \ldots, e_j$, qui engendrent la restriction de $\mathcal{F}$ à $U_j$.
D'après Cohomologie des schémas, lemme \ref{coherent-lemma-homs-over-open},
on peut prolonger $s_{jk}$ en un morphisme
$s'_{jk} : \mathcal{I}_i^{n_{jk}} \to \mathcal{F}$ pour un certain $n_{jk} \geq 1$.
On peut alors considérer les composées
$$
\mathcal{E}_j^{\otimes n_{jk}} \to \mathcal{I}_j^{n_{jk}} \to \mathcal{F}
$$
pour conclure.
\end{proof}

\begin{lemma}
\label{lemma-resolution-property-ample-family}
Soit $X$ un schéma. Si $X$ possède une famille ample de modules inversibles
(Morphismes, définition
\ref{morphisms-definition-family-ample-invertible-modules}),
alors $X$ possède la propriété de résolution.
\end{lemma}

\begin{proof}
Puisque $X$ est quasi-compact, il existe un $n$ et des couples
$(\mathcal{L}_i, s_i)$, $i = 1, \ldots, n$, où
$\mathcal{L}_i$ est un $\mathcal{O}_X$-module inversible et
$s_i \in \Gamma(X, \mathcal{L}_i)$ est une section
tels que l'ensemble des points $U_i \subset X$ où $s_i$ ne s'annule pas
soit affine et que $X = U_1 \cup \ldots \cup U_n$.
Soit $\mathcal{I}_i \subset \mathcal{O}_X$ l'image
de $s_i : \mathcal{L}_i^{\otimes -1} \to \mathcal{O}_X$.
En appliquant le lemme \ref{lemma-resolution-property-finite-number},
on constate que $X$ possède la propriété de résolution.
\end{proof}

\begin{lemma}
\label{lemma-regular-resolution-property}
Soit $X$ un schéma quasi-compact régulier à diagonale affine.
Alors $X$ possède la propriété de résolution.
\end{lemma}

\begin{proof}
On combine Diviseurs, lemme \ref{divisors-lemma-regular-ample-family}, et le
lemme \ref{lemma-resolution-property-ample-family} ci-dessus.
\end{proof}

\begin{lemma}
\label{lemma-resolution-property-descends}
Soit $X = \lim X_i$ la limite d'un système inductif de schémas quasi-compacts
et quasi-séparés à morphismes de transition affines.
Alors $X$ possède la propriété de résolution si et seulement si $X_i$ possède
la propriété de résolution pour un certain $i$.
\end{lemma}

\begin{proof}
Si $X_i$ possède la propriété de résolution, alors $X$ la possède d'après le
lemme \ref{lemma-resolution-property-goes-up-affine}.
Supposons que $X$ possède la propriété de résolution.
Choisissons $i \in I$. Choisissons un recouvrement ouvert affine fini
$X_i = U_{i, 1} \cup \ldots \cup U_{i, m}$.
Pour tout $j$, choisissons un faisceau quasi-cohérent de type fini
d'idéaux $\mathcal{I}_{i, j} \subset \mathcal{O}_{X_i}$
tel que $X_i \setminus V(\mathcal{I}_{i, j}) = U_{i, j}$ ; voir
Propriétés, lemme \ref{properties-lemma-quasi-coherent-finite-type-ideals}.
Notons $U_j \subset X$ l'image inverse de $U_{i, j}$,
et notons $\mathcal{I}_j \subset \mathcal{O}_X$
l'image inverse de $\mathcal{I}_{i, j}$.
Puisque $X$ possède la propriété de résolution, on peut
choisir des $\mathcal{O}_X$-modules localement libres de type fini $\mathcal{E}_j$
et des surjections $\mathcal{E}_j \to \mathcal{I}_j$.
D'après Limites, lemmes \ref{limits-lemma-descend-invertible-modules} et
\ref{limits-lemma-descend-modules-finite-presentation},
après avoir augmenté $i$, on peut trouver des
$\mathcal{O}_{X_i}$-modules localement libres de type fini $\mathcal{E}_{i, j}$
et des morphismes $\mathcal{E}_{i, j} \to \mathcal{O}_{X_i}$
dont les changements de base à $X$ redonnent les composées
$\mathcal{E}_j \to \mathcal{I}_j \to \mathcal{O}_X$, $j = 1, \ldots, m$.
Les images inverses des $\mathcal{O}_{X_i}$-modules de présentation finie
$\Coker(\mathcal{E}_{i, j} \to \mathcal{O}_{X_i})$
et $\mathcal{O}_{X_i}/\mathcal{I}_{i, j}$
sur $X$ coïncident comme quotients de $\mathcal{O}_X$.
Par conséquent, d'après Limites, lemme
\ref{limits-lemma-descend-modules-finite-presentation},
on peut supposer qu'ils coïncident, autrement dit que
l'image de $\mathcal{E}_{i, j} \to \mathcal{O}_{X_i}$
est égale à $\mathcal{I}_{i, j}$.
On en déduit que $X_i$ possède la propriété de résolution
d'après le lemme \ref{lemma-resolution-property-finite-number}.
\end{proof}

\begin{lemma}
\label{lemma-resolution-property-affine-diagonal}
\begin{reference}
Cas particulier de \cite[Proposition 1.3]{totaro_resolution}.
\end{reference}
Soit $X$ un schéma quasi-compact et quasi-séparé possédant
la propriété de résolution. Alors $X$ a une diagonale affine.
\end{lemma}

\begin{proof}
En combinant Limites, proposition \ref{limits-proposition-approximate},
et le lemme \ref{lemma-resolution-property-descends},
on se ramène au cas où $X$ est noethérien (un petit détail est omis).
Supposons $X$ noethérien.
Rappelons que $X \times X$ est recouvert par les ouverts affines
$U \times V$, où $U$, $V$ sont des ouverts affines de $X$ ; voir
Schémas, section \ref{schemes-section-fibre-products}.
Ainsi, pour montrer que la diagonale $\Delta : X \to X \times X$
est affine, il suffit de montrer que $U \cap V = \Delta^{-1}(U \times V)$
est affine pour tous ouverts affines $U$, $V$ de $X$ ; voir
Morphismes, lemme \ref{morphisms-lemma-characterize-affine}.
En particulier, il suffit de montrer que le morphisme d'inclusion
$j : U \to X$ est affine si $U$ est un ouvert affine de $X$.
D'après Cohomologie des schémas, lemme \ref{coherent-lemma-criterion-affine-morphism},
il suffit de montrer que $R^1j_*\mathcal{G} = 0$ pour tout
$\mathcal{O}_U$-module quasi-cohérent $\mathcal{G}$.
D'après la proposition \ref{proposition-Noetherian} (c'est ici que l'on utilise
la réduction au cas noethérien),
on peut représenter $Rj_*\mathcal{G}$ par un complexe
$\mathcal{H}^\bullet$ de $\mathcal{O}_X$-modules quasi-cohérents.
Supposons que
$$
H^1(\mathcal{H}^\bullet) =
\Ker(\mathcal{H}^1 \to \mathcal{H}^2)/\Im(\mathcal{H}^0 \to \mathcal{H}^1)
$$
soit non nul, afin d'obtenir une contradiction. On peut alors trouver un
$\mathcal{O}_X$-module cohérent $\mathcal{F}$ et un morphisme
$$
\mathcal{F} \longrightarrow
\Ker(\mathcal{H}^1 \to \mathcal{H}^2)
$$
tel que sa composée avec la projection sur $H^1(\mathcal{H}^\bullet)$
soit non nulle. En effet, on peut écrire $\Ker(\mathcal{H}^1 \to \mathcal{H}^2)$
comme réunion filtrante de ses sous-modules cohérents d'après
Propriétés, lemme \ref{properties-lemma-quasi-coherent-colimit-finite-type},
et l'un d'eux conviendra.
Ensuite, on choisit un $\mathcal{O}_X$-module localement libre de type fini
$\mathcal{E}$ et une surjection $\mathcal{E} \to \mathcal{F}$ en utilisant
la propriété de résolution de $X$.
Cela fournit un morphisme dans la catégorie dérivée
$$
\mathcal{E}[-1] \longrightarrow Rj_*\mathcal{G}
$$
qui est non nul sur les faisceaux de cohomologie et donc non nul dans $D(\mathcal{O}_X)$.
Par adjonction, cela revient à un morphisme
$$
j^*\mathcal{E}[-1] \to \mathcal{G}
$$
non nul dans $D(\mathcal{O}_U)$. Puisque $\mathcal{E}$ est localement libre
de type fini, cela revient à un élément non nul de
$$
H^1(U, j^*\mathcal{E}^\vee \otimes_{\mathcal{O}_U} \mathcal{G})
$$
où
$\mathcal{E}^\vee = \SheafHom_{\mathcal{O}_X}(\mathcal{E}, \mathcal{O}_X)$
est le module localement libre de type fini dual. Cependant, ce groupe est nul d'après
Cohomologie des schémas, lemme
\ref{coherent-lemma-quasi-coherent-affine-cohomology-zero},
ce qui est la contradiction recherchée.
(En cas de doute sur l'étape utilisant les duaux, voir le résultat plus général
Cohomologie, lemme \ref{cohomology-lemma-dual-perfect-complex}.)
\end{proof}





\section{La propriété de résolution et les complexes parfaits}
\label{section-resolution-property-perfect}

\noindent
Dans cette section, nous examinons la relation entre les complexes parfaits
et les complexes strictement parfaits sur les schémas qui possèdent la propriété de résolution.

\begin{lemma}
\label{lemma-construct-strictly-perfect}
Soit $X$ un schéma quasi-compact et quasi-séparé possédant la
propriété de résolution.
Soit $\mathcal{F}^\bullet$ un complexe borné inférieurement de
$\mathcal{O}_X$-modules quasi-cohérents représentant un objet parfait de
$D(\mathcal{O}_X)$. Alors il existe un complexe borné
$\mathcal{E}^\bullet$ de $\mathcal{O}_X$-modules localement libres de type fini
et un quasi-isomorphisme $\mathcal{E}^\bullet \to \mathcal{F}^\bullet$.
\end{lemma}

\begin{proof}
Soient $a, b \in \mathbf{Z}$ des entiers tels que $\mathcal{F}^\bullet$
ait une amplitude de Tor contenue dans $[a, b]$ et tels que $\mathcal{F}^n = 0$ pour
$n < a$. L'existence d'un tel couple d'entiers
résulte de Cohomologie, lemme \ref{cohomology-lemma-perfect},
et du fait que $X$ est quasi-compact.
Si $b < a$, alors $\mathcal{F}^\bullet$ est nul dans la
catégorie dérivée et le lemme en résulte.
Nous allons démontrer par récurrence sur
$b - a \geq 0$ qu'il existe un complexe
$\mathcal{E}^a \to \ldots \to \mathcal{E}^b$
où chaque $\mathcal{E}^i$ est localement libre de type fini,
et un quasi-isomorphisme $\mathcal{E}^\bullet \to \mathcal{F}^\bullet$.

\medskip\noindent
Le cas initial est celui où $b - a = 0$. Dans ce cas,
$H^b(\mathcal{F}^\bullet) = H^a(\mathcal{F}^\bullet) =
\Ker(\mathcal{F}^a \to \mathcal{F}^{a + 1})$
est localement libre de type fini. En effet, c'est un
$\mathcal{O}_X$-module de présentation finie, ayant une dimension de Tor égale à $0$, donc
localement libre de type fini. Voir Cohomologie, lemmes \ref{cohomology-lemma-perfect} et
\ref{cohomology-lemma-finite-cohomology}, ainsi que
Propriétés, lemme \ref{properties-lemma-finite-locally-free}.
Ainsi, on peut prendre pour $\mathcal{E}^\bullet$ le complexe
$H^b(\mathcal{F}^\bullet)$ placé en degré $b$.
Le reste de la démonstration est consacré à l'étape de récurrence.

\medskip\noindent
Supposons $b > a$. Observons que
$$
H^b(\mathcal{F}^\bullet) =
\Ker(\mathcal{F}^b \to \mathcal{F}^{b + 1})/
\Im(\mathcal{F}^{b - 1} \to \mathcal{F}^b)
$$
est un $\mathcal{O}_X$-module quasi-cohérent de type fini ; voir
Cohomologie, lemmes \ref{cohomology-lemma-perfect} et
\ref{cohomology-lemma-finite-cohomology}. On peut alors trouver un
$\mathcal{O}_X$-module quasi-cohérent de type fini $\mathcal{F}$ et un morphisme
$$
\mathcal{F} \longrightarrow
\Ker(\mathcal{F}^b \to \mathcal{F}^{b + 1})
$$
tel que sa composée avec la projection sur
$H^b(\mathcal{F}^\bullet)$ soit surjective.
En effet, on peut écrire $\Ker(\mathcal{F}^b \to \mathcal{F}^{b + 1})$
comme la réunion filtrante de ses sous-modules quasi-cohérents de type fini d'après
Propriétés, lemme \ref{properties-lemma-quasi-coherent-colimit-finite-type},
et l'un d'eux conviendra.
Ensuite, on choisit un $\mathcal{O}_X$-module localement libre de type fini
$\mathcal{E}^b$ et une surjection $\mathcal{E}^b \to \mathcal{F}$ en utilisant
la propriété de résolution de $X$. Considérons le morphisme de complexes
$$
\alpha : \mathcal{E}^b[-b] \to \mathcal{F}^\bullet
$$
et son cône $C(\alpha)^\bullet$ ; voir
Catégories dérivées, définition \ref{derived-definition-cone}.
Observons que $C(\alpha)^\bullet$ n'est non nul qu'en degrés
$\geq a$, a une amplitude de Tor contenue dans $[a, b]$ d'après
Cohomologie, lemme \ref{cohomology-lemma-cone-tor-amplitude},
et vérifie $H^b(C(\alpha)^\bullet) = 0$ par construction.
Ainsi, on trouve en fait que $C(\alpha)^\bullet$ a une amplitude de Tor
contenue dans $[a, b - 1]$. L'hypothèse de récurrence s'applique donc à
$C(\alpha)^\bullet$ et l'on obtient un morphisme de complexes
$$
(\mathcal{E}^a \to \ldots \to \mathcal{E}^{b - 1})
\longrightarrow
C(\alpha)^\bullet
$$
ayant les propriétés énoncées dans l'hypothèse de récurrence. En explicitant
la définition du cône, on obtient le diagramme commutatif
$$
\xymatrix{
\ldots \ar[r] &
\mathcal{E}^{b - 2} \ar[r] \ar[d] &
\mathcal{E}^{b - 1} \ar[r] \ar[d] &
0 \ar[r] \ar[d] &
\ldots \\
\ldots \ar[r] &
\mathcal{F}^{b - 2} \ar[r] &
\mathcal{F}^{b - 1} \oplus \mathcal{E}^b \ar[r] &
\mathcal{F}^b \ar[r] &
\ldots
}
$$
Il est clair que l'on obtient un morphisme de complexes
$(\mathcal{E}^a \to \ldots \to \mathcal{E}^b) \to \mathcal{F}^\bullet$.
Nous omettons de vérifier que ce morphisme est un quasi-isomorphisme.
\end{proof}

\begin{lemma}
\label{lemma-resolution-property-perfect-complex}
Soit $X$ un schéma quasi-compact et quasi-séparé possédant la
propriété de résolution. Alors tout objet parfait de $D(\mathcal{O}_X)$
peut être représenté par un complexe borné de
$\mathcal{O}_X$-modules localement libres de type fini.
\end{lemma}

\begin{proof}
Soit $E$ un objet parfait de $D(\mathcal{O}_X)$.
D'après le lemme \ref{lemma-resolution-property-affine-diagonal},
on voit que $X$ a une diagonale affine. Par conséquent, d'après la
proposition \ref{proposition-quasi-compact-affine-diagonal},
on peut représenter $E$ par un complexe $\mathcal{F}^\bullet$
de $\mathcal{O}_X$-modules quasi-cohérents.
Observons que $E$ appartient à $D^b(\mathcal{O}_X)$ parce que
$X$ est quasi-compact. Par conséquent, $\tau_{\geq n}\mathcal{F}^\bullet$
est un complexe borné inférieurement de $\mathcal{O}_X$-modules quasi-cohérents
qui représente $E$ si $n \ll 0$. On peut donc appliquer le
lemme \ref{lemma-construct-strictly-perfect} pour conclure.
\end{proof}

\begin{lemma}
\label{lemma-resolution-property-map-perfect-complex}
Soit $X$ un schéma quasi-compact et quasi-séparé possédant la
propriété de résolution. Soient $\mathcal{E}^\bullet$ et $\mathcal{F}^\bullet$
des complexes finis de $\mathcal{O}_X$-modules localement libres de type fini.
Alors tout
$\alpha \in \Hom_{D(\mathcal{O}_X)}(\mathcal{E}^\bullet, \mathcal{F}^\bullet)$
peut être représenté par un diagramme
$$
\mathcal{E}^\bullet \leftarrow \mathcal{G}^\bullet \to \mathcal{F}^\bullet
$$
où $\mathcal{G}^\bullet$ est un complexe borné dont chaque terme est localement libre de type fini en tant que
$\mathcal{O}_X$-module et où $\mathcal{G}^\bullet \to \mathcal{E}^\bullet$
est un quasi-isomorphisme.
\end{lemma}

\begin{proof}
D'après le lemme \ref{lemma-resolution-property-affine-diagonal},
on voit que $X$ a une diagonale affine. Par conséquent, d'après la
proposition \ref{proposition-quasi-compact-affine-diagonal},
on peut représenter $\alpha$ par un diagramme
$$
\mathcal{E}^\bullet \leftarrow \mathcal{H}^\bullet \to \mathcal{F}^\bullet
$$
où $\mathcal{H}^\bullet$ est un complexe de
$\mathcal{O}_X$-modules quasi-cohérents et où
$\mathcal{H}^\bullet \to \mathcal{E}^\bullet$
est un quasi-isomorphisme. Pour $n \ll 0$, les morphismes
$\mathcal{H}^\bullet \to \mathcal{E}^\bullet$ et
$\mathcal{H}^\bullet \to \mathcal{F}^\bullet$
se factorisent par le quasi-isomorphisme
$\mathcal{H}^\bullet \to \tau_{\geq n}\mathcal{H}^\bullet$
simplement parce que $\mathcal{E}^\bullet$ et $\mathcal{F}^\bullet$
sont des complexes bornés. On peut donc remplacer $\mathcal{H}^\bullet$
par $\tau_{\geq n}\mathcal{H}^\bullet$ et supposer que $\mathcal{H}^\bullet$
est borné inférieurement.
On peut alors appliquer le
lemme \ref{lemma-construct-strictly-perfect} pour conclure.
\end{proof}

\begin{lemma}
\label{lemma-resolution-property-homotopy-map-perfect-complex}
Soit $X$ un schéma quasi-compact et quasi-séparé possédant la
propriété de résolution. Soient $\mathcal{E}^\bullet$ et $\mathcal{F}^\bullet$
des complexes finis de $\mathcal{O}_X$-modules localement libres de type fini.
Soient $\alpha^\bullet, \beta^\bullet :\mathcal{E}^\bullet \to \mathcal{F}^\bullet$
deux morphismes de complexes définissant le même morphisme dans $D(\mathcal{O}_X)$.
Alors il existe un quasi-isomorphisme
$\gamma^\bullet : \mathcal{G}^\bullet \to \mathcal{E}^\bullet$
où $\mathcal{G}^\bullet$ est un complexe borné dont chaque terme est localement libre de type fini en tant que
$\mathcal{O}_X$-module,
tel que $\alpha^\bullet \circ \gamma^\bullet$ et
$\beta^\bullet \circ \gamma^\bullet$ soient des morphismes de complexes homotopes.
\end{lemma}

\begin{proof}
D'après le lemme \ref{lemma-resolution-property-affine-diagonal},
on voit que $X$ a une diagonale affine. Par conséquent, d'après la
proposition \ref{proposition-quasi-compact-affine-diagonal}
(et la définition de la catégorie dérivée),
il existe un quasi-isomorphisme
$\gamma^\bullet : \mathcal{G}^\bullet \to \mathcal{E}^\bullet$
où $\mathcal{G}^\bullet$ est un complexe de
$\mathcal{O}_X$-modules quasi-cohérents,
tel que $\alpha^\bullet \circ \gamma^\bullet$ et
$\beta^\bullet \circ \gamma^\bullet$ soient des morphismes de complexes homotopes.
Choisissons une homotopie $h^i : \mathcal{G}^i \to \mathcal{F}^{i - 1}$
attestant ce fait. Choisissons $n \ll 0$. Le morphisme
$\gamma^\bullet$ se factorise alors canoniquement par le morphisme quotient
$\mathcal{G}^\bullet \to \tau_{\geq n}\mathcal{G}^\bullet$
puisque $\mathcal{E}^\bullet$ est borné inférieurement. Pour exactement la même
raison, les morphismes $h^i$ se factorisent par les surjections
$\mathcal{G}^i \to (\tau_{\geq n}\mathcal{G})^i$.
On voit donc que l'on peut remplacer $\mathcal{G}^\bullet$
par $\tau_{\geq n}\mathcal{G}^\bullet$.
On peut alors appliquer le lemme \ref{lemma-construct-strictly-perfect} pour conclure.
\end{proof}

\begin{proposition}
\label{proposition-perfect-resolution-property}
Soit $X$ un schéma quasi-compact et quasi-séparé possédant la
propriété de résolution. Notons
\begin{enumerate}
\item $\mathcal{A}$ la catégorie additive des
$\mathcal{O}_X$-modules localement libres de type fini,
\item $K^b(\mathcal{A})$ la catégorie homotopique des complexes bornés
de $\mathcal{A}$ ; voir
Catégories dérivées, section \ref{derived-section-homotopy}, et
\item $D_{perf}(\mathcal{O}_X)$ la sous-catégorie strictement pleine, saturée
et triangulée de $D(\mathcal{O}_X)$ formée des
objets parfaits.
\end{enumerate}
Avec ces notations, le foncteur évident
$$
K^b(\mathcal{A}) \longrightarrow D_{perf}(\mathcal{O}_X)
$$
est un foncteur exact de catégories triangulées qui se factorise par une
équivalence $S^{-1}K^b(\mathcal{A}) \to D_{perf}(\mathcal{O}_X)$
de catégories triangulées,
où $S$ est le système multiplicatif saturé des quasi-isomorphismes
de $K^b(\mathcal{A})$.
\end{proposition}

\begin{proof}
Si vous parvenez à analyser l'énoncé de la proposition, sautez donc ce
premier paragraphe. Pour certaines des définitions employées, voir
Catégories dérivées, définition
\ref{derived-definition-triangulated-subcategory}
(sous-catégorie triangulée),
Catégories dérivées, définition \ref{derived-definition-saturated}
(sous-catégorie triangulée saturée),
Catégories dérivées, définition \ref{derived-definition-localization}
(système multiplicatif compatible avec la structure triangulée), et
Catégories,
définition \ref{categories-definition-saturated-multiplicative-system}
(système multiplicatif saturé).
Observons que $D_{perf}(\mathcal{O}_X)$ est une sous-catégorie triangulée saturée
de $D(\mathcal{O}_X)$ d'après
Cohomologie, lemmes \ref{cohomology-lemma-two-out-of-three-perfect} et
\ref{cohomology-lemma-summands-perfect}. Notons également que
$K^b(\mathcal{A})$ est une catégorie triangulée ; voir
Catégories dérivées, lemme
\ref{derived-lemma-bounded-triangulated-subcategories}.

\medskip\noindent
Il est clair que le foncteur transforme les triangles distingués en
triangles distingués, c'est-à-dire qu'il est exact. Alors $S$ est un système
multiplicatif saturé compatible avec la structure triangulée
de $K^b(\mathcal{A})$ d'après
Catégories dérivées, lemme \ref{derived-lemma-triangle-functor-localize}.
Par conséquent, la localisation $S^{-1}K^b(\mathcal{A})$ existe et est
une catégorie triangulée d'après
Catégories dérivées, proposition
\ref{derived-proposition-construct-localization}.
On obtient une factorisation par le foncteur exact
$S^{-1}K^b(\mathcal{A}) \to D_{perf}(\mathcal{O}_X)$ d'après
Catégories dérivées, lemme
\ref{derived-lemma-universal-property-localization}.
D'après les lemmes \ref{lemma-resolution-property-perfect-complex},
\ref{lemma-resolution-property-map-perfect-complex} et
\ref{lemma-resolution-property-homotopy-map-perfect-complex},
ce foncteur est une équivalence. Enfin, le foncteur
$S^{-1}K^b(\mathcal{A}) \to D_{perf}(\mathcal{O}_X)$
est une équivalence de catégories triangulées (au sens où
les triangles distingués se correspondent) d'après
Catégories dérivées, lemme \ref{derived-lemma-exact-equivalence}.
\end{proof}



\section{Groupes K}
\label{section-K-groups}

\noindent
Quelques mots au sujet de $K_0$ des diverses catégories associées aux schémas.
On trouvera des résultats antérieurs dans
Algèbre, section \ref{algebra-section-K-groups},
Homologie, section \ref{homology-section-K-groups},
Catégories dérivées, section \ref{derived-section-K-groups}, et
Compléments sur l'algèbre, lemme \ref{more-algebra-lemma-perfect-to-K-group-universal}.

\medskip\noindent
Par analogie avec Algèbre, section \ref{algebra-section-K-groups},
nous définirons deux groupes de $K$-théorie, $K'_0(X)$ et $K_0(X)$, pour tout schéma
noethérien $X$. Le premier fera intervenir les $\mathcal{O}_X$-modules cohérents
et le second les $\mathcal{O}_X$-modules localement libres de type fini.

\begin{lemma}
\label{lemma-Noetherian-Kprime}
Soit $X$ un schéma noethérien. Alors
$$
K_0(\textit{Coh}(\mathcal{O}_X)) =
K_0(D^b(\textit{Coh}(\mathcal{O}_X)) =
K_0(D^b_{\textit{Coh}}(\mathcal{O}_X))
$$
\end{lemma}

\begin{proof}
La première égalité est donnée par
Catégories dérivées, lemme \ref{derived-lemma-K-bounded-derived}.
On a $K_0(\textit{Coh}(\mathcal{O}_X)) =
K_0(D^b_{\textit{Coh}}(\mathcal{O}_X))$ d'après
Catégories dérivées, lemme \ref{derived-lemma-DBA-map-K}.
Ceci démontre le lemme.
(On peut aussi utiliser l'égalité
$D^b(\textit{Coh}(\mathcal{O}_X)) = D^b_{\textit{Coh}}(\mathcal{O}_X)$
de la proposition \ref{proposition-DCoh} pour obtenir la seconde
égalité.)
\end{proof}

\noindent
Voici la définition.

\begin{definition}
\label{definition-K-group}
Soit $X$ un schéma.
\begin{enumerate}
\item On note $K_0(X)$ le {\it groupe de Grothendieck de $X$}. C'est le
groupe de K-théorie d'indice zéro de la sous-catégorie triangulée
strictement pleine et saturée $D_{perf}(\mathcal{O}_X)$ de $D(\mathcal{O}_X)$
formée des objets parfaits. Autrement dit,
$$
K_0(X) = K_0(D_{perf}(\mathcal{O}_X))
$$
\item Si $X$ est localement noethérien, on note $K'_0(X)$ le
{\it groupe de Grothendieck des faisceaux cohérents sur $X$}. C'est le
groupe $K$ d'indice zéro de la catégorie abélienne
des $\mathcal{O}_X$-modules cohérents. Autrement dit,
$$
K'_0(X) = K_0(\textit{Coh}(\mathcal{O}_X))
$$
\end{enumerate}
\end{definition}

\noindent
Nous montrerons que notre définition de $K_0(X)$ coïncide avec la définition
couramment utilisée en termes de modules localement libres de type fini si $X$
possède la propriété de résolution (par exemple si $X$ possède un module
inversible ample). Voir le lemme \ref{lemma-K-is-old-K}.

\begin{lemma}
\label{lemma-K-agrees-affine}
Soit $X = \Spec(R)$ un schéma affine. Alors $K_0(X) = K_0(R)$
et, si $R$ est noethérien, $K'_0(X) = K'_0(R)$.
\end{lemma}

\begin{proof}
Rappelons que $K'_0(R)$ et $K_0(R)$ ont été définis dans
Algèbre, section \ref{algebra-section-K-groups}.

\medskip\noindent
D'après Compléments sur l'algèbre, lemme \ref{more-algebra-lemma-perfect-to-K-group-universal},
on a $K_0(R) = K_0(D_{perf}(R))$.
D'après les lemmes \ref{lemma-perfect-affine} et \ref{lemma-affine-compare-bounded},
on a $D_{perf}(R) = D_{perf}(\mathcal{O}_X)$.
Ceci démontre l'égalité $K_0(R) = K_0(X)$.

\medskip\noindent
L'égalité $K'_0(R) = K'_0(X)$ résulte du fait que
$\textit{Coh}(\mathcal{O}_X)$ est équivalente à la catégorie
des $R$-modules finis d'après Cohomologie des schémas, lemme
\ref{coherent-lemma-coherent-Noetherian}. De plus, les définitions
montrent immédiatement que $K'_0(R)$ est le groupe de K-théorie d'indice zéro de la catégorie
des $R$-modules finis.
\end{proof}

\noindent
Soit $X$ un schéma noethérien. Alors $K'_0(X)$ et $K_0(X)$
sont tous deux définis. Dans ce cas, il existe une application canonique
$$
K_0(X) = K_0(D_{perf}(\mathcal{O}_X))
\longrightarrow
K_0(D^b_{\textit{Coh}}(\mathcal{O}_X)) = K'_0(X)
$$
En effet, les complexes parfaits appartiennent à $D^b_{\textit{Coh}}(\mathcal{O}_X)$
(d'après le lemme \ref{lemma-identify-pseudo-coherent-noetherian}), le foncteur
d'inclusion
$D_{perf}(\mathcal{O}_X) \to D^b_{\textit{Coh}}(\mathcal{O}_X)$
induit un homomorphisme entre les groupes $K$ d'indice zéro
(Catégories dérivées, lemme \ref{derived-lemma-map-K}),
et l'égalité de droite est donnée par le
lemme \ref{lemma-Noetherian-Kprime}.

\begin{lemma}
\label{lemma-Kprime-K}
Soit $X$ un schéma noethérien régulier. Alors
l'application $K_0(X) \to K'_0(X)$ est un isomorphisme.
\end{lemma}

\begin{proof}
Cela résulte immédiatement du lemme \ref{lemma-perfect-on-regular}
et de la construction ci-dessus de l'application $K_0(X) \to K'_0(X)$.
\end{proof}

\noindent
Soit $X$ un schéma. Notons $\textit{Vect}(X)$ la catégorie
des $\mathcal{O}_X$-modules localement libres de type fini. Bien que
$\textit{Vect}(X)$ ne soit pas en général une catégorie abélienne, la notion
de suite exacte courte dans $\textit{Vect}(X)$ est claire.
Notons $K_0(\textit{Vect}(X))$ l'unique groupe abélien ayant les
propriétés suivantes\footnote{Le cadre général correct serait ici
de définir $K_0$ pour toute catégorie exacte ; voir
Injectifs, remarque \ref{injectives-remark-embed-exact-category}.} :
\begin{enumerate}
\item Pour tout $\mathcal{O}_X$-module localement libre de type fini $\mathcal{E}$,
un élément $[\mathcal{E}]$ de $K_0(\textit{Vect}(X))$ est donné,
\item pour toute suite exacte courte
$0 \to \mathcal{E}' \to \mathcal{E} \to \mathcal{E}'' \to 0$
de $\mathcal{O}_X$-modules localement libres de type fini, on a la relation
$[\mathcal{E}] = [\mathcal{E}'] + [\mathcal{E}'']$ dans $K_0(\textit{Vect}(X))$,
\item le groupe $K_0(\textit{Vect}(X))$ est engendré par les éléments
$[\mathcal{E}]$, et
\item toutes les relations dans $K_0(\textit{Vect}(X))$ entre les générateurs
$[\mathcal{E}]$ sont des combinaisons $\mathbf{Z}$-linéaires
des relations provenant des suites exactes comme ci-dessus.
\end{enumerate}
Nous omettons la construction détaillée de $K_0(\textit{Vect}(X))$.
Il existe une application naturelle
$$
K_0(\textit{Vect}(X)) \longrightarrow K_0(X)
$$
En effet, étant donné un $\mathcal{O}_X$-module localement libre de type fini $\mathcal{E}$,
notons $\mathcal{E}[0]$ le complexe parfait sur $X$ dans lequel
$\mathcal{E}$ est placé en degré $0$ et dont les autres termes sont nuls.
Étant donnée une suite exacte courte
$0 \to \mathcal{E} \to \mathcal{E}' \to \mathcal{E}'' \to 0$ de
$\mathcal{O}_X$-modules localement libres de type fini, on obtient un triangle distingué
$\mathcal{E}[0] \to \mathcal{E}'[0] \to \mathcal{E}''[0] \to \mathcal{E}[1]$ ;
voir Catégories dérivées, section \ref{derived-section-canonical-delta-functor}.
Cela montre que l'on obtient une application
$K_0(\textit{Vect}(X)) \to K_0(D_{perf}(\mathcal{O}_X)) = K_0(X)$
en envoyant $[\mathcal{E}]$ sur $[\mathcal{E}[0]]$,
en priant le lecteur d'excuser cette notation épouvantable.

\begin{lemma}
\label{lemma-K-is-old-K}
Soit $X$ un schéma quasi-compact et quasi-séparé possédant la
propriété de résolution. Alors l'application $K_0(\textit{Vect}(X)) \to K_0(X)$
est un isomorphisme.
\end{lemma}

\begin{proof}
Ce lemme résulte directement des
lemmes \ref{lemma-resolution-property-perfect-complex},
\ref{lemma-resolution-property-map-perfect-complex} et
\ref{lemma-resolution-property-homotopy-map-perfect-complex},
dont nous utiliserons les résultats sans autre mention.
Construisons une application réciproque
$$
c : K_0(X) = K_0(D_{perf}(\mathcal{O}_X)) \longrightarrow
K_0(\textit{Vect}(X))
$$
Tout objet de $D_{perf}(\mathcal{O}_X)$ peut être représenté
par un complexe borné $\mathcal{E}^\bullet$ de
$\mathcal{O}_X$-modules localement libres de type fini. On pose alors
$$
c([\mathcal{E}^\bullet]) = \sum (-1)^i[\mathcal{E}^i]
$$
Bien entendu, il faut montrer que cette définition est indépendante du choix.
Pour l'instant, nous considérons $c$ comme une application définie sur les complexes
bornés de $\mathcal{O}_X$-modules localement libres de type fini.

\medskip\noindent
Supposons que $\mathcal{E}^\bullet \to \mathcal{F}^\bullet$ soit une application surjective
de complexes bornés de $\mathcal{O}_X$-modules localement libres de type fini.
Soit $\mathcal{K}^\bullet$ son noyau. On obtient alors des suites exactes
courtes de $\mathcal{O}_X$-modules
$$
0 \to \mathcal{K}^n \to \mathcal{E}^n \to \mathcal{F}^n \to 0
$$
qui sont scindées localement parce que $\mathcal{F}^n$ est localement libre de type fini.
Ainsi, $\mathcal{K}^\bullet$ est lui aussi un complexe borné de
$\mathcal{O}_X$-modules localement libres de type fini, et l'on a
$c(\mathcal{E}^\bullet) = c(\mathcal{K}^\bullet) + c(\mathcal{F}^\bullet)$
dans $K_0(\textit{Vect}(X))$.

\medskip\noindent
Supposons donné un complexe borné $\mathcal{E}^\bullet$
de $\mathcal{O}_X$-modules localement libres de type fini qui soit acyclique.
Supposons que $\mathcal{E}^n = 0$ pour $n \not \in [a, b]$. On peut alors
décomposer $\mathcal{E}^\bullet$ en suites exactes courtes
$$
\begin{matrix}
0 \to \mathcal{E}^a \to \mathcal{E}^{a + 1} \to \mathcal{F}^{a + 1} \to 0,\\
0 \to \mathcal{F}^{a + 1} \to \mathcal{E}^{a + 2} \to
\mathcal{F}^{a + 3} \to 0, \\
\ldots \\
0 \to \mathcal{F}^{b - 3} \to \mathcal{E}^{b - 2} \to
\mathcal{F}^{b - 2} \to 0, \\
0 \to \mathcal{F}^{b - 2} \to \mathcal{E}^{b - 1} \to \mathcal{E}^b \to 0
\end{matrix}
$$
En raisonnant par récurrence descendante, on voit que
$\mathcal{F}^{b - 2}, \ldots, \mathcal{F}^{a + 1}$
sont des $\mathcal{O}_X$-modules localement libres de type fini et que
$$
c(\mathcal{E}^\bullet) = \sum (-1)[\mathcal{E}^n] =
\sum (-1)^n([\mathcal{F}^{n - 1}] + [\mathcal{F}^n]) = 0
$$
Ainsi, notre construction donne zéro sur les complexes acycliques.

\medskip\noindent
Il résulte des deux paragraphes précédents que $c$
est bien définie. En effet, supposons que les complexes bornés
$\mathcal{E}^\bullet$ et $\mathcal{F}^\bullet$ de $\mathcal{O}_X$-modules localement libres
de type fini représentent le même objet de $D(\mathcal{O}_X)$.
On peut alors trouver des quasi-isomorphismes
$a : \mathcal{G}^\bullet \to \mathcal{E}^\bullet$ et
$b : \mathcal{G}^\bullet \to \mathcal{F}^\bullet$,
où $\mathcal{G}^\bullet$ est un complexe borné de $\mathcal{O}_X$-modules localement libres
de type fini.
On obtient une suite exacte courte de complexes
$$
0 \to \mathcal{E}^\bullet \to C(a)^\bullet \to \mathcal{G}^\bullet[1] \to 0
$$
voir Catégories dérivées, définition \ref{derived-definition-cone}.
Puisque $a$ est un quasi-isomorphisme, le cône $C(a)^\bullet$ est
acyclique (cela résulte par exemple de la discussion dans
Catégories dérivées, section \ref{derived-section-canonical-delta-functor}).
Par conséquent,
$$
0 = c(C(f)^\bullet) = c(\mathcal{E}^\bullet) + c(\mathcal{G}^\bullet[1]) =
c(\mathcal{E}^\bullet) - c(\mathcal{G}^\bullet)
$$
comme voulu. Le même raisonnement appliqué à $b$ montre que
$0 = c(\mathcal{F}^\bullet) - c(\mathcal{G}^\bullet)$.
On obtient donc $c(\mathcal{E}^\bullet) = c(\mathcal{F}^\bullet)$,
et $c$ est bien définie.

\medskip\noindent
Un raisonnement analogue utilisant le cône d'une application
$\mathcal{E}^\bullet \to \mathcal{F}^\bullet$
de complexes bornés de $\mathcal{O}_X$-modules localement libres de type fini
montre que $c(Y) = c(X) + c(Z)$ si $X \to Y \to Z$ est un triangle distingué
dans $D_{perf}(\mathcal{O}_X)$. Les détails sont omis.
On obtient ainsi l'homomorphisme souhaité
de groupes abéliens $c : K_0(X) \to K_0(\textit{Vect}(X))$.

\medskip\noindent
Il est clair que la composée
$K_0(\textit{Vect}(X)) \to K_0(X) \to K_0(\textit{Vect}(X))$
est l'identité. D'autre part, soit $\mathcal{E}^\bullet$
un complexe borné de $\mathcal{O}_X$-modules localement libres de type fini.
L'existence des triangles distingués associés aux « troncatures bêtes »
(voir Homologie, section \ref{homology-section-truncations})
$$
\sigma_{\geq n}\mathcal{E}^\bullet \to
\sigma_{\geq n - 1}\mathcal{E}^\bullet \to
\mathcal{E}^{n - 1}[-n + 1] \to
(\sigma_{\geq n}\mathcal{E}^\bullet)[1]
$$
et une récurrence montrent que
$$
[\mathcal{E}^\bullet] = \sum (-1)^i[\mathcal{E}^i[0]]
$$
dans $K_0(X) = K_0(D_{perf}(\mathcal{O}_X))$, en priant le lecteur d'excuser la notation.
Ainsi, l'application $K_0(\textit{Vect}(X)) \to K_0(D_{perf}(\mathcal{O}_X)) = K_0(X)$
est surjective, ce qui achève la preuve.
\end{proof}

\begin{remark}
\label{remark-K-ring}
Soit $X$ un schéma. Le groupe de K-théorie $K_0(X)$ est canoniquement un anneau commutatif.
En effet, le produit tensoriel dérivé
$$
\otimes = \otimes^\mathbf{L}_{\mathcal{O}_X} :
D_{perf}(\mathcal{O}_X) \times D_{perf}(\mathcal{O}_X)
\longrightarrow
D_{perf}(\mathcal{O}_X)
$$
et Catégories dérivées, lemme \ref{derived-lemma-bilinear-map-K},
donnent une multiplication bilinéaire. Comme $K \otimes L \cong L \otimes K$,
ce produit est commutatif. Comme
$(K \otimes L) \otimes M = K \otimes (L \otimes M)$,
ce produit est associatif.
Enfin, l'unité de $K_0(X)$ est l'élément $1 = [\mathcal{O}_X]$.

\medskip\noindent
Si $\textit{Vect}(X)$ et $K_0(\textit{Vect}(X))$ sont comme ci-dessus, alors
il est clair que $K_0(\textit{Vect}(X))$ possède lui aussi une
structure d'anneau : si $\mathcal{E}$ et $\mathcal{F}$ sont des
$\mathcal{O}_X$-modules localement libres de type fini, on pose
$$
[\mathcal{E}] \cdot [\mathcal{F}] =
[\mathcal{E} \otimes_{\mathcal{O}_X} \mathcal{F}]
$$
Le lecteur vérifie aisément que cela définit bien un produit bilinéaire,
commutatif et associatif. Les détails sont omis. L'application
$$
K_0(\textit{Vect}(X)) \longrightarrow K_0(X)
$$
construite ci-dessus est un homomorphisme d'anneaux pour ces définitions.

\medskip\noindent
Supposons maintenant $X$ noethérien. Le produit tensoriel dérivé fournit aussi
une application
$$
\otimes = \otimes^\mathbf{L}_{\mathcal{O}_X} :
D_{perf}(\mathcal{O}_X) \times D^b_{\textit{Coh}}(\mathcal{O}_X)
\longrightarrow
D^b_{\textit{Coh}}(\mathcal{O}_X)
$$
En utilisant de nouveau Catégories dérivées, lemme \ref{derived-lemma-bilinear-map-K},
on obtient une multiplication bilinéaire $K_0(X) \times K'_0(X) \to K'_0(X)$,
car $K'_0(X) = K_0(D^b_{\textit{Coh}}(\mathcal{O}_X))$ d'après le
lemme \ref{lemma-Noetherian-Kprime}.
Le lecteur montre aisément que cela munit $K'_0(X)$ d'une structure
de module sur l'anneau $K_0(X)$.
\end{remark}

\begin{remark}
\label{remark-pushforward-K}
Soit $f : X \to Y$ un morphisme propre de schémas localement noethériens.
Il existe une application
$$
f_* : K'_0(X) \longrightarrow K'_0(Y)
$$
qui envoie $[\mathcal{F}]$ sur
$$
[\bigoplus\nolimits_{i \geq 0} R^{2i}f_*\mathcal{F}] -
[\bigoplus\nolimits_{i \geq 0} R^{2i + 1}f_*\mathcal{F}]
$$
Cette application est bien définie parce que les faisceaux $R^if_*\mathcal{F}$
sont cohérents (Cohomologie des schémas, proposition
\ref{coherent-proposition-proper-pushforward-coherent}), parce que localement
seul un nombre fini d'entre eux sont non nuls, et parce qu'une
suite exacte courte de faisceaux cohérents sur $X$ donne une suite exacte
longue de $R^if_*$ sur $Y$. Si $Y$ est quasi-compact (le seul
cas généralement utilisé en pratique), on peut réécrire ce qui précède sous la forme
$$
f_*[\mathcal{F}] = \sum (-1)^i[R^if_*\mathcal{F}] = [Rf_*\mathcal{F}]
$$
où l'on a utilisé l'égalité $K'_0(Y) = K_0(D^b_{\textit{Coh}}(Y))$ du
lemme \ref{lemma-Noetherian-Kprime}.
\end{remark}

\begin{lemma}
\label{lemma-projection-formula}
Soit $f : X \to Y$ un morphisme propre de schémas localement noethériens.
Alors $f_*(\alpha \cdot f^*\beta) = f_*\alpha \cdot \beta$
pour $\alpha \in K'_0(X)$ et $\beta \in K_0(Y)$.
\end{lemma}

\begin{proof}
Cela résulte du lemme \ref{lemma-cohomology-base-change}, de la discussion de la
remarque \ref{remark-pushforward-K} et de la définition du produit
$K'_0(X) \times K_0(X) \to K'_0(X)$ dans la remarque \ref{remark-K-ring}.
\end{proof}

\begin{remark}
\label{remark-perf-Z}
Soit $X$ un schéma. Soit $Z \subset X$ un sous-schéma fermé. Considérons la
sous-catégorie triangulée strictement pleine et saturée
$$
D_{Z, perf}(\mathcal{O}_X) \subset D(\mathcal{O}_X)
$$
formée des complexes parfaits de $\mathcal{O}_X$-modules
dont les faisceaux de cohomologie sont à support ensembliste dans $Z$.
Le groupe $K$ d'indice zéro $K_0(D_{Z, perf}(\mathcal{O}_X))$
de cette catégorie triangulée est parfois noté
$K_Z(X)$ ou $K_{0, Z}(X)$. En utilisant le produit tensoriel dérivé exactement
comme dans la remarque \ref{remark-K-ring}, on voit que $K_0(D_{Z, perf}(\mathcal{O}_X))$
possède une multiplication associative et commutative,
mais qu'en général $K_0(D_{Z, perf}(\mathcal{O}_X))$ n'a pas d'unité.
\end{remark}










\section{Déterminants des complexes}
\label{section-det-two-terms}

\noindent
Cette section fait suite à
Compléments d'algèbre, section \ref{more-algebra-section-determinants-complexes}.
Pour tout espace annelé $(X, \mathcal{O}_X)$, il existe un foncteur
$$
\begin{gathered}
\det :
\left\{
\begin{matrix}
\text{catégorie des complexes parfaits} \\
\text{les morphismes sont les isomorphismes}
\end{matrix}
\right\}
\\
\longrightarrow
\left\{
\begin{matrix}
\text{catégorie des modules inversibles} \\
\text{les morphismes sont les isomorphismes}
\end{matrix}
\right\}
\end{gathered}
$$
De plus, si $(L, F)$ est un objet de la catégorie dérivée filtrée
$DF(\mathcal{O}_X)$ dont la filtration est finie et dont les gradués
sont des complexes parfaits, il existe un isomorphisme canonique
$\det(\text{gr}L) \to \det(L)$. Voir \cite{determinant} pour l'exposé
original. Nous ajouterons ces développements ultérieurement (insérer une référence future).

\medskip\noindent
Pour le moment, nous présentons une construction ad hoc dans le cas
où $X$ est un schéma et où nous considérons des objets parfaits $L$ de
$D(\mathcal{O}_X)$ d'amplitude de Tor contenue dans $[-1, 0]$.

\begin{lemma}
\label{lemma-determinant-two-term-complexes}
Soit $X$ un schéma. Il existe un foncteur
$$
\begin{gathered}
\det :
\left\{
\begin{matrix}
\text{catégorie des complexes parfaits} \\
\text{d'amplitude de Tor contenue dans }[-1, 0] \\
\text{les morphismes sont les isomorphismes}
\end{matrix}
\right\}
\\
\longrightarrow
\left\{
\begin{matrix}
\text{catégorie des modules inversibles} \\
\text{les morphismes sont les isomorphismes}
\end{matrix}
\right\}
\end{gathered}
$$
De plus, si l'on se donne un objet parfait de rang $0$, noté $L$, de $D(\mathcal{O}_X)$,
d'amplitude de Tor contenue dans $[-1, 0]$, il existe un élément canonique
$\delta(L) \in \Gamma(X, \det(L))$ tel que, pour tout isomorphisme
$a : L \to K$ dans $D(\mathcal{O}_X)$, on ait $\det(a)(\delta(L)) = \delta(K)$.
De plus, la construction est donnée localement sur les ouverts affines par celle
de Compléments d'algèbre, section \ref{more-algebra-section-determinants-complexes}.
\end{lemma}

\begin{proof}
Soit $L$ un objet de la catégorie figurant à gauche. Si $\Spec(A) = U \subset X$
est un ouvert affine, alors $L|_U$ correspond à un complexe parfait $L^\bullet$
de $A$-modules d'amplitude de Tor contenue dans $[-1, 0]$, voir les
lemmes \ref{lemma-affine-compare-bounded},
\ref{lemma-tor-dimension-affine} et
\ref{lemma-perfect-affine}.
On peut alors considérer le $A$-module inversible $\det(L^\bullet)$ construit dans
Compléments d'algèbre, lemme \ref{more-algebra-lemma-determinant-two-term-complexes}.
Si $\Spec(B) = V \subset U$ est un autre ouvert affine contenu dans $U$,
alors $\det(L^\bullet) \otimes_A B = \det(L^\bullet \otimes_A B)$,
et cette construction est donc compatible avec les applications de restriction
(voir le lemme \ref{lemma-quasi-coherence-pullback} et noter que $A \to B$ est plat).
Ainsi, nous pouvons recoller ces modules inversibles pour obtenir un module inversible
$\det(L)$ sur $X$. La fonctorialité et les sections canoniques
se construisent exactement de la même manière. Détails omis.
\end{proof}

\begin{remark}
\label{remark-functorial-det}
La construction du lemme \ref{lemma-determinant-two-term-complexes}
est compatible avec les images inverses. Plus précisément, soient un morphisme
$f : X \to Y$ de schémas et un objet parfait $K$ de $D(\mathcal{O}_Y)$
d'amplitude de Tor contenue dans $[-1, 0]$. Alors $Lf^*K$ est un
objet parfait $K$ de $D(\mathcal{O}_X)$
d'amplitude de Tor contenue dans $[-1, 0]$ et nous avons une identification canonique
$$
f^*\det(K) \longrightarrow \det(Lf^*K)
$$
De plus, si $K$ est de rang $0$, alors l'image réciproque de $\delta(K)$ s'identifie à
$\delta(Lf^*K)$ par ce morphisme. Cela résulte immédiatement de la construction
du déterminant sur les ouverts affines.
\end{remark}




\section{Détection de la bornitude}
\label{section-detecting-boundedness}

\noindent
Dans cette section, nous montrons que les générateurs compacts de $D_\QCoh$ d'un
schéma quasi-compact et quasi-séparé, tels qu'ils sont construits dans la
section \ref{section-generating}, possèdent une propriété particulière.
Nous recommandons de lire d'abord cette dernière section, car elle est très semblable à celle-ci.

\begin{lemma}
\label{lemma-orthogonal-koszul-first-variant}
Dans la situation \ref{situation-complex}, notons $j : U \to X$ l'immersion
ouverte et soit $K$ l'objet parfait de $D(\mathcal{O}_X)$
correspondant au complexe de Koszul de $f_1, \ldots, f_r$ sur $A$.
Soient $E \in D_\QCoh(\mathcal{O}_X)$ et $a \in \mathbf{Z}$.
Considérons les conditions suivantes :
\begin{enumerate}
\item Le morphisme canonique $\tau_{\geq a}E \to \tau_{\geq a} Rj_*(E|_U)$
est un isomorphisme.
\item On a $\Hom_{D(\mathcal{O}_X)}(K[-n], E) = 0$ pour tout $n \geq a$.
\end{enumerate}
Alors (2) implique (1), et (1) implique (2) lorsque $a$ est remplacé par $a + 1$.
\end{lemma}

\begin{proof}
Choisissons un triangle distingué $N \to E \to Rj_*(E|_U) \to N[1]$.
Alors (1) implique $\tau_{\geq a + 1} N = 0$, tandis que
$\tau_{\geq a}N = 0$ implique (1). Observons que
$$
\Hom_{D(\mathcal{O}_X)}(K[-n], Rj_*(E|_U)) =
\Hom_{D(\mathcal{O}_U)}(K|_U[-n], E) = 0
$$
pour tout $n$, puisque $K|_U = 0$. Ainsi (2) équivaut à
$\Hom_{D(\mathcal{O}_X)}(K[-n], N) = 0$ pour tout $n \geq a$. 
Observons qu'il existe des triangles distingués
$$
\begin{aligned}
K^\bullet(f_1^{e_1}, \ldots, f_i^{e'_i}, \ldots, f_r^{e_r}) &\to
K^\bullet(f_1^{e_1}, \ldots, f_i^{e'_i + e''_i}, \ldots, f_r^{e_r}) \\
{}&\to
K^\bullet(f_1^{e_1}, \ldots, f_i^{e''_i}, \ldots, f_r^{e_r}) \to \ldots
\end{aligned}
$$
de complexes de Koszul, voir
Compléments sur l'algèbre, lemme \ref{more-algebra-lemma-koszul-mult}. Par conséquent,
$\Hom_{D(\mathcal{O}_X)}(K[-n], N) = 0$ pour tout $n \geq a$
équivaut à
$\Hom_{D(\mathcal{O}_X)}(K_e[-n], N) = 0$ pour tout $n \geq a$ et
tout $e \geq 1$, avec $K_e$ comme dans le
lemme \ref{lemma-represent-cohomology-class-on-closed}.
Puisque $N|_U = 0$, ce lemme implique que cette condition équivaut à son tour à
$H^n(X, N) = 0$ pour $n \geq a$. Nous en concluons que (2) équivaut
à $\tau_{\geq a}N = 0$, puisque $N$ est déterminé par le complexe de
$A$-modules $R\Gamma(X, N)$, voir le lemme \ref{lemma-affine-compare-bounded}.
Cela démontre le lemme.
\end{proof}

\begin{lemma}
\label{lemma-orthogonal-koszul-second-variant}
Dans la situation \ref{situation-complex}, notons $j : U \to X$ l'immersion
ouverte et soit $K$ l'objet parfait de $D(\mathcal{O}_X)$
correspondant au complexe de Koszul de $f_1, \ldots, f_r$ sur $A$.
Soient $E \in D_\QCoh(\mathcal{O}_X)$ et $a \in \mathbf{Z}$. Considérons
les conditions suivantes :
\begin{enumerate}
\item Le morphisme canonique $\tau_{\leq a}E \to \tau_{\leq a} Rj_*(E|_U)$
est un isomorphisme, et
\item $\Hom_{D(\mathcal{O}_X)}(K[-n], E) = 0$ pour tout $n \leq a$.
\end{enumerate}
Alors (2) implique (1), et (1) implique (2) lorsque $a$ est remplacé par $a - 1$.
\end{lemma}

\begin{proof}
Choisissons un triangle distingué $E \to Rj_*(E|_U) \to N \to E[1]$. Alors (1)
implique $\tau_{\leq a - 1}N = 0$, tandis que $\tau_{\leq a}N = 0$ implique (1).
Observons que
$$
\Hom_{D(\mathcal{O}_X)}(K[-n], Rj_*(E|_U)) =
\Hom_{D(\mathcal{O}_U)}(K|_U[-n], E) = 0
$$
pour tout $n$, puisque $K|_U = 0$. Ainsi (2) équivaut à
$\Hom_{D(\mathcal{O}_X)}(K[-n], N) = 0$ pour tout $n \leq a$. 
Observons qu'il existe des triangles distingués
$$
\begin{aligned}
K^\bullet(f_1^{e_1}, \ldots, f_i^{e'_i}, \ldots, f_r^{e_r}) &\to
K^\bullet(f_1^{e_1}, \ldots, f_i^{e'_i + e''_i}, \ldots, f_r^{e_r}) \\
{}&\to
K^\bullet(f_1^{e_1}, \ldots, f_i^{e''_i}, \ldots, f_r^{e_r}) \to \ldots
\end{aligned}
$$
de complexes de Koszul, voir
Compléments sur l'algèbre, lemme \ref{more-algebra-lemma-koszul-mult}. Par conséquent,
$\Hom_{D(\mathcal{O}_X)}(K[-n], N) = 0$ pour tout $n \leq a$
équivaut à
$\Hom_{D(\mathcal{O}_X)}(K_e[-n], N) = 0$ pour tout $n \leq a$ et tout $e \geq 1$,
avec $K_e$ comme dans le lemme \ref{lemma-represent-cohomology-class-on-closed}.
Puisque $N|_U = 0$, ce lemme implique que cette condition équivaut à son tour à
$H^n(X, N) = 0$ pour $n \leq a$. Nous en concluons que (2) équivaut à
$\tau_{\leq a}N = 0$, puisque $N$ est déterminé par le complexe de
$A$-modules $R\Gamma(X, N)$, voir le lemme \ref{lemma-affine-compare-bounded}.
Cela démontre le lemme.
\end{proof}

\begin{lemma}
\label{lemma-bounded-truncation}
Soit $X$ un schéma quasi-compact et quasi-séparé.
Soient $P \in D_{perf}(\mathcal{O}_X)$ et $E \in D_{\QCoh}(\mathcal{O}_X)$. 
Soit $a \in \mathbf{Z}$. Les conditions suivantes sont équivalentes :
\begin{enumerate}
\item $\Hom_{D(\mathcal{O}_X)}(P[-i], E) = 0$ pour $i \gg 0$, et
\item $\Hom_{D(\mathcal{O}_X)}(P[-i], \tau_{\geq a} E) = 0$ pour $i \gg 0$.
\end{enumerate}
\end{lemma}

\begin{proof}
En utilisant le triangle $ \tau_{< a} E \to E \to \tau_{\geq a} E \to$,
on voit qu'il suffit, pour obtenir l'équivalence, de montrer que
$$
\Hom_{D(\mathcal{O}_X)}(P[-i], \tau_{< a} E) =
\Hom_{D(\mathcal{O}_X)}(P, (\tau_{< a} E)[i]) = 0 
$$
pour $i \gg 0$. Comme $P$ est parfait, cela résulte du
lemme \ref{lemma-ext-from-perfect-into-bounded-QCoh}.
\end{proof}

\begin{lemma}
\label{lemma-bounded-below-truncation}
Soit $X$ un schéma quasi-compact et quasi-séparé. Soient
$P \in D_{perf}(\mathcal{O}_X)$ et $E \in D_{\QCoh}(\mathcal{O}_X)$.
Soit $a \in \mathbf{Z}$. Les conditions suivantes sont équivalentes :
\begin{enumerate}
\item $\Hom_{D(\mathcal{O}_X)}(P[-i], E) = 0$ pour $i \ll 0$, et
\item $\Hom_{D(\mathcal{O}_X)}(P[-i], \tau_{\leq a} E) = 0$ pour $i \ll 0$.
\end{enumerate}
\end{lemma}

\begin{proof}
En utilisant le triangle $ \tau_{\leq a} E \to E \to \tau_{> a} E \to$,
on voit qu'il suffit, pour obtenir l'équivalence, de montrer que
$$
\Hom_{D(\mathcal{O}_X)}(P[-i], \tau_{> a} E) =
\Hom_{D(\mathcal{O}_X)}(P, (\tau_{> a} E)[i]) = 0
$$
pour $i \ll 0$. Comme $P$ est parfait, cela résulte du
lemme \ref{lemma-ext-from-perfect-into-bounded-QCoh}.
\end{proof}

\begin{proposition}
\label{proposition-detecting-bounded-above}
Soit $X$ un schéma quasi-compact et quasi-séparé. Soit
$G \in D_{perf}(\mathcal{O}_X)$ un complexe parfait qui engendre 
$D_\QCoh (\mathcal{O}_X)$. Soit $E \in D_\QCoh (\mathcal{O}_X)$.
Les conditions suivantes sont équivalentes :
\begin{enumerate}
\item $E \in D^-_\QCoh (\mathcal{O}_X)$,
\item $\Hom_{D(\mathcal{O}_X)}(G[-i], E) = 0$ pour $i \gg 0$,
\item $\Ext^i_X(G, E) = 0$ pour $i \gg 0$,
\item $R\Hom_X(G, E)$ appartient à $D^-(\mathbf{Z})$,
\item $H^i(X, G^\vee \otimes_{\mathcal{O}_X}^\mathbf{L} E) = 0$
pour $i \gg 0$,
\item $R\Gamma(X, G^\vee \otimes_{\mathcal{O}_X}^\mathbf{L} E)$
appartient à $D^-(\mathbf{Z})$,
\item pour tout objet parfait $P$ de $D(\mathcal{O}_X)$,
\begin{enumerate}
\item les assertions (2), (3), (4) sont vraies lorsque $G$ est remplacé par $P$, et
\item $H^i(X, P \otimes_{\mathcal{O}_X}^\mathbf{L} E) = 0$ pour $i \gg 0$,
\item $R\Gamma(X, P \otimes_{\mathcal{O}_X}^\mathbf{L} E)$
appartient à $D^-(\mathbf{Z})$.
\end{enumerate}
\end{enumerate}
\end{proposition}

\begin{proof}
Supposons (1).\newline Puisque
$\Hom_{D(\mathcal{O}_X)}(G[-i], E) = \Hom_{D(\mathcal{O}_X)}(G, E[i])$,
ce groupe est nul pour $i \gg 0$ d'après le
lemme \ref{lemma-ext-from-perfect-into-bounded-QCoh}. Ceci démontre
que (1) implique (2).

\medskip\noindent
Les parties (2), (3), (4) sont équivalentes d'après la discussion de
Cohomologie, section \ref{cohomology-section-global-RHom}.
Les parties (5) et (6) sont équivalentes, puisque $H^i(X, -) = H^i(R\Gamma(X, -))$
par définition. Les conditions équivalentes (2), (3), (4) sont
équivalentes aux conditions équivalentes (5), (6) d'après
Cohomologie, lemme \ref{cohomology-lemma-dual-perfect-complex},
et le fait que $(G[-i])^\vee = G^\vee[i]$.

\medskip\noindent
Il est clair que (7) implique (2). Réciproquement, 
démontrons que les conditions équivalentes (2) -- (6) impliquent (7).
Rappelons que $G$ est un générateur classique de $D_{perf}(\mathcal{O}_X)$ d'après la
remarque \ref{remark-classical-generator}.
Pour $P \in D_{perf}(\mathcal{O}_X)$, notons $T(P)$ l'assertion selon laquelle
$R\Hom_X(P, E)$ appartient à $D^-(\mathbf{Z})$.
Il est clair que $T$ est stable par sommes directes,
vérifie la propriété des deux sur trois pour les triangles
distingués, et est stable par facteurs directs et par décalages.
Par conséquent, d'après Catégories dérivées, remarque \ref{derived-remark-check-on-generator},
on voit que (4) implique que $T$ vaut sur tout $D_{perf}(\mathcal{O}_X)$.
Le même argument vaut pour toutes les autres propriétés, si ce n'est que, pour les propriétés
(7)(b) et (7)(c), on utilise aussi que $P \mapsto P^\vee$ est une
autoéquivalence de $D_{perf}(\mathcal{O}_X)$. Un petit détail est omis.

\medskip\noindent
Nous allons démontrer que les conditions équivalentes (2) -- (7) impliquent (1)
en utilisant le principe de récurrence de
Cohomologie des schémas, lemme \ref{coherent-lemma-induction-principle}.

\medskip\noindent
D'abord, démontrons que (2) -- (7) $\Rightarrow$ (1) si $X$ est affine.
Posons $P = \mathcal{O}_X[0]$. La condition (7) donne $H^i (X, E) = 0$ pour $i \gg 0$.
Ainsi (1) en résulte, puisque $E$ est déterminé par $R\Gamma (X, E)$,
voir le lemme \ref{lemma-affine-compare-bounded}.

\medskip\noindent
Supposons maintenant $X = U \cup V$, où $U$ est un ouvert quasi-compact de $X$ et
$V$ un ouvert affine, et supposons que l'implication (2) -- (7) $\Rightarrow$ (1)
soit connue pour les schémas $U$, $V$ et $U \cap V$.
Supposons que $E \in D_\QCoh(\mathcal{O}_X)$ satisfasse aux conditions (2) -- (7).
D'après le lemme \ref{lemma-direct-summand-of-a-restriction} et le
théorème \ref{theorem-bondal-van-den-Bergh}, il existe un complexe parfait
$Q$ sur $X$ tel que $Q|_U$ engendre $D_\QCoh (\mathcal{O}_U)$. 
Soient $f_1, \dots , f_r \in \Gamma (V, \mathcal{O}_V)$
tels que $V \setminus U = V(f_1, \dots , f_r)$ comme sous-ensembles de
$V$. Soit $K \in D_{perf}(\mathcal{O}_V)$ l'objet
correspondant au complexe de Koszul de $f_1, \dots , f_r$.
Soit $K' \in D_{perf}(\mathcal{O}_X)$ défini par
\begin{equation}
\label{equation-detecting-bounded-above}
K' = R (V \to X)_* K = R (V \to X)_! K,
\end{equation}
voir Cohomologie, lemmes \ref{cohomology-lemma-pushforward-restriction} et
\ref{cohomology-lemma-pushforward-perfect}. C'est un complexe parfait
sur $X$ à support dans le fermé $X \setminus U \subset V$
et isomorphe à $K$ sur $V$. Par hypothèse, nous savons que
$R\Hom_{\mathcal{O}_X}(Q, E)$ et
$R\Hom_{\mathcal{O}_X}(K', E)$ sont bornés supérieurement.

\medskip\noindent
D'après la seconde description de $K'$ dans (\ref{equation-detecting-bounded-above}),
on a
$$
\Hom_{D(\mathcal{O}_V)}(K[-i], E|_V) = \Hom_{D(\mathcal{O}_X)}(K'[-i], E) = 0
$$
pour $i \gg 0$. Nous pouvons donc appliquer le
lemme \ref{lemma-orthogonal-koszul-first-variant} à $E|_V$ pour
obtenir un entier $a$ tel que
$\tau_{\geq a}(E|_V) = \tau_{\geq a} R (U \cap V \to V)_* (E|_{U \cap V})$.
Alors $\tau_{\geq a} E = \tau_{\geq a} R (U \to X)_* (E |_U)$
(vérifier que le morphisme canonique est un isomorphisme après restriction à
$U$ et à $V$). Par conséquent, en appliquant deux fois le lemme \ref{lemma-bounded-truncation},
on voit que
$$
\Hom_{D(\mathcal{O}_U)}(Q|_U [-i], E|_U) =
\Hom_{D(\mathcal{O}_X)}(Q[-i], R (U \to X)_* (E|_U)) = 0
$$
pour $i \gg 0$. Puisque la proposition est vraie pour $U$ et le générateur
$Q|_U$, on a $E|_U \in D^-_\QCoh(\mathcal{O}_U)$. Mais alors, puisque
le foncteur $R (U \to X)_*$ préserve $D^-_\QCoh$ 
(d'après le lemme \ref{lemma-quasi-coherence-direct-image}), on obtient
$\tau_{\geq a}E \in D^-_\QCoh(\mathcal{O}_X)$. Ainsi 
$E \in D^-_\QCoh (\mathcal{O}_X)$. 
\end{proof}

\begin{proposition}
\label{proposition-detecting-bounded-below}
Soit $X$ un schéma quasi-compact et quasi-séparé.
Soit $G \in D_{perf}(\mathcal{O}_X)$
un complexe parfait qui engendre $D_\QCoh (\mathcal{O}_X)$. Soit
$E \in D_\QCoh (\mathcal{O}_X)$. Les conditions suivantes sont équivalentes :
\begin{enumerate}
\item $E \in D^+_\QCoh (\mathcal{O}_X)$,
\item $\Hom_{D(\mathcal{O}_X)}(G[-i], E) = 0$ pour $i \ll 0$,
\item $\Ext^i_X(G, E) = 0$ pour $i \ll 0$,
\item $R\Hom_X(G, E)$ appartient à $D^+(\mathbf{Z})$,
\item $H^i(X, G^\vee \otimes_{\mathcal{O}_X}^\mathbf{L} E) = 0$
pour $i \ll 0$,
\item $R\Gamma(X, G^\vee \otimes_{\mathcal{O}_X}^\mathbf{L} E)$
appartient à $D^+(\mathbf{Z})$,
\item pour tout objet parfait $P$ de $D(\mathcal{O}_X)$,
\begin{enumerate}
\item les assertions (2), (3), (4) sont vraies lorsque $G$ est remplacé par $P$, et
\item $H^i(X, P \otimes_{\mathcal{O}_X}^\mathbf{L} E) = 0$ pour $i \ll 0$,
\item $R\Gamma(X, P \otimes_{\mathcal{O}_X}^\mathbf{L} E)$
appartient à $D^+(\mathbf{Z})$.
\end{enumerate}
\end{enumerate}
\end{proposition}

\begin{proof}
Supposons (1).\newline Puisque
$\Hom_{D(\mathcal{O}_X)}(G[-i], E) = \Hom_{D(\mathcal{O}_X)}(G, E[i])$,
ce groupe est nul pour $i \ll 0$ d'après le
lemme \ref{lemma-ext-from-perfect-into-bounded-QCoh}. Ceci démontre
que (1) implique (2).

\medskip\noindent
Les parties (2), (3), (4) sont équivalentes d'après la discussion de
Cohomologie, section \ref{cohomology-section-global-RHom}.
Les parties (5) et (6) sont équivalentes, puisque $H^i(X, -) = H^i(R\Gamma(X, -))$
par définition. Les conditions équivalentes (2), (3), (4) sont
équivalentes aux conditions équivalentes (5), (6) d'après
Cohomologie, lemme \ref{cohomology-lemma-dual-perfect-complex},
et le fait que $(G[-i])^\vee = G^\vee[i]$.

\medskip\noindent
Il est clair que (7) implique (2). Réciproquement, 
démontrons que les conditions équivalentes (2) -- (6) impliquent (7).
Rappelons que $G$ est un générateur classique de $D_{perf}(\mathcal{O}_X)$ d'après la
remarque \ref{remark-classical-generator}.
Pour $P \in D_{perf}(\mathcal{O}_X)$, notons $T(P)$ l'assertion selon laquelle
$R\Hom_X(P, E)$ appartient à $D^+(\mathbf{Z})$.
Il est clair que $T$ est stable par sommes directes,
vérifie la propriété des deux sur trois pour les triangles
distingués, et est stable par facteurs directs et par décalages.
Par conséquent, d'après Catégories dérivées, remarque \ref{derived-remark-check-on-generator},
on voit que (4) implique que $T$ vaut sur tout $D_{perf}(\mathcal{O}_X)$.
Le même argument vaut pour toutes les autres propriétés, si ce n'est que, pour les propriétés
(7)(b) et (7)(c), on utilise aussi que $P \mapsto P^\vee$ est une
autoéquivalence de $D_{perf}(\mathcal{O}_X)$. Un petit détail est omis.

\medskip\noindent
Nous allons démontrer que les conditions équivalentes (2) -- (7) impliquent (1)
en utilisant le principe de récurrence de
Cohomologie des schémas, lemme \ref{coherent-lemma-induction-principle}.

\medskip\noindent
D'abord, démontrons que (2) -- (7) $\Rightarrow$ (1) si $X$ est affine.
Posons $P = \mathcal{O}_X[0]$. La condition (7) donne $H^i (X, E) = 0$
pour $i \ll 0$. Ainsi (1) en résulte, puisque $E$ est
déterminé par $R\Gamma (X, E)$, voir le lemme \ref{lemma-affine-compare-bounded}.

\medskip\noindent
Supposons maintenant $X = U \cup V$, où $U$ est un ouvert quasi-compact de $X$ et
$V$ un ouvert affine, et supposons que l'implication (2) -- (7) $\Rightarrow$ (1)
soit connue pour les schémas $U$, $V$ et $U \cap V$.
Supposons que $E \in D_\QCoh(\mathcal{O}_X)$ satisfasse aux conditions (2) -- (7).
D'après le lemme \ref{lemma-direct-summand-of-a-restriction} et le
théorème \ref{theorem-bondal-van-den-Bergh}, il existe un complexe parfait
$Q$ sur $X$ tel que $Q|_U$ engendre $D_\QCoh (\mathcal{O}_U)$. 
Soient $f_1, \dots , f_r \in \Gamma (V, \mathcal{O}_V)$
tels que $V \setminus U = V(f_1, \dots , f_r)$ comme sous-ensembles de
$V$. Soit $K \in D_{perf}(\mathcal{O}_V)$ l'objet
correspondant au complexe de Koszul de $f_1, \dots , f_r$.
Soit $K' \in D_{perf}(\mathcal{O}_X)$ défini par
\begin{equation}
\label{equation-detecting-bounded-below}
K' = R (V \to X)_* K = R (V \to X)_! K,
\end{equation}
voir Cohomologie, lemmes \ref{cohomology-lemma-pushforward-restriction} et
\ref{cohomology-lemma-pushforward-perfect}. C'est un complexe parfait
sur $X$ à support dans le fermé $X \setminus U \subset V$
et isomorphe à $K$ sur $V$. Par hypothèse, nous savons que
$R\Hom_{\mathcal{O}_X}(Q, E)$ et
$R\Hom_{\mathcal{O}_X}(K', E)$ sont bornés inférieurement.

\medskip\noindent
D'après la seconde description de $K'$ dans (\ref{equation-detecting-bounded-below}),
on a
$$
\Hom_{D(\mathcal{O}_V)}(K[-i], E|_V) = \Hom_{D(\mathcal{O}_X)}(K'[-i], E) = 0
$$
pour $i \ll 0$. Nous pouvons donc appliquer le
lemme \ref{lemma-orthogonal-koszul-second-variant} à $E|_V$ pour
obtenir un entier $a$ tel que
$\tau_{\leq a}(E|_V) = \tau_{\leq a} R(U \cap V \to V)_*(E|_{U \cap V})$.
Alors $\tau_{\leq a} E = \tau_{\leq a} R(U \to X)_*(E|_U)$
(vérifier que le morphisme canonique est un isomorphisme après restriction à
$U$ et à $V$). Par conséquent, en appliquant deux fois le lemme \ref{lemma-bounded-below-truncation},
on voit que
$$
\Hom_{D(\mathcal{O}_U)}(Q|_U [-i], E|_U) =
\Hom_{D(\mathcal{O}_X)}(Q[-i], R (U \to X)_* (E|_U)) = 0
$$
pour $i \ll 0$. Puisque la proposition est vraie pour $U$ et le générateur
$Q|_U$, on a $E|_U \in D^+_\QCoh (\mathcal{O}_U)$. Mais alors, puisque
le foncteur $R(U \to X)_*$ préserve les objets bornés inférieurement
(voir Cohomologie, section \ref{cohomology-section-derived-functors}), on obtient
$\tau_{\leq a} E \in D^+_\QCoh(\mathcal{O}_X)$. Ainsi 
$E \in D^+_\QCoh (\mathcal{O}_X)$.
\end{proof}







\section{Objets quasi-cohérents dans la catégorie dérivée}
\label{section-QC}

\noindent
Soit $X$ un schéma. Rappelons que $X_{affine, Zar}$
désigne la catégorie des ouverts affines de $X$ munie de la topologie donnée
par les recouvrements de Zariski standards ; voir
Topologies, définition \ref{topologies-definition-big-small-Zariski}.
Rappelons au lecteur que le topos de $X_{affine, Zar}$
est le petit topos de Zariski de $X$ ; voir Topologies, lemme
\ref{topologies-lemma-alternative-zariski}. Le site $X_{affine, Zar}$
est muni d'un faisceau structural $\mathcal{O}$ et il existe une
équivalence de topos annelés
$$
(\Sh(X_{affine, Zar}), \mathcal{O})
\longrightarrow
(\Sh(X_{Zar}), \mathcal{O})
$$
Voir Descente, équation (\ref{descent-equation-alternative-small-ringed}),
ainsi que la discussion qui l'entoure dans
Descente, section \ref{descent-section-alternative-quasi-coherent},
où une notation légèrement différente est employée.

\medskip\noindent
Dans cette section, nous notons $X_{affine}$ la catégorie sous-jacente à
$X_{affine, Zar}$, munie de la topologie chaotique, de sorte que les faisceaux
coïncident avec les préfaisceaux. En particulier, le faisceau structural $\mathcal{O}$
devient également un faisceau sur $X_{affine}$.
Nous obtenons un morphisme de sites annelés
$$
\epsilon :
(X_{affine, Zar}, \mathcal{O})
\longrightarrow
(X_{affine}, \mathcal{O})
$$
comme dans Cohomologie sur les sites, section \ref{sites-cohomology-section-compare}.
Dans cette section, nous identifierons $D_\QCoh(\mathcal{O}_X)$ à
la catégorie $\mathit{QC}(X_{affine}, \mathcal{O})$ introduite dans
Cohomologie sur les sites, section \ref{sites-cohomology-section-modules-cohomology}.

\begin{lemma}
\label{lemma-DQCoh-alternative-small}
Dans la situation ci-dessus, il existe des équivalences exactes canoniques entre
les catégories triangulées suivantes :
\begin{enumerate}
\item $D_\QCoh(\mathcal{O}_X)$,
\item $D_\QCoh(X_{Zar}, \mathcal{O})$,
\item $D_\QCoh(X_{affine, Zar}, \mathcal{O})$,
\item $D_\QCoh(X_{affine}, \mathcal{O}_X)$, et
\item $\mathit{QC}(X_{affine}, \mathcal{O})$.
\end{enumerate}
\end{lemma}

\begin{proof}
Si $U \subset V \subset X$ sont des ouverts affines, alors le morphisme d'anneaux
$\mathcal{O}(V) \to \mathcal{O}(U)$ est plat. Par conséquent,
l'équivalence entre (4) et (5) est un cas particulier de
Cohomologie sur les sites, lemme
\ref{sites-cohomology-lemma-cartesian-flat-quasi-coherent}
(la démonstration éclaire également l'énoncé).

\medskip\noindent
Le site annelé $(X_{Zar}, \mathcal{O})$ et l'espace annelé
$(X, \mathcal{O}_X)$ ont la même catégorie de modules d'après
Descente, remarque \ref{descent-remark-Zariski-site-space}.
Par cette équivalence, les modules quasi-cohérents se correspondent d'après
Descente, proposition \ref{descent-proposition-equivalence-quasi-coherent}.
On obtient donc une équivalence exacte canonique entre les catégories triangulées
de (1) et (2).

\medskip\noindent
La discussion précédant le lemme montre que nous avons une équivalence de
topos annelés
$(\Sh(X_{affine, Zar}), \mathcal{O}) \to (\Sh(X_{Zar}), \mathcal{O})$
et donc une équivalence entre les catégories abéliennes de modules.
Comme la notion de module quasi-cohérent est intrinsèque
(Modules sur les sites, lemme \ref{sites-modules-lemma-special-locally-free}),
nous voyons que cette équivalence préserve les sous-catégories
de modules quasi-cohérents. Ainsi, nous obtenons une équivalence exacte canonique
entre les catégories triangulées de (2) et (3).

\medskip\noindent
Pour obtenir une équivalence exacte entre les catégories triangulées
de (3) et (4), nous allons appliquer Cohomologie sur les sites, lemme
\ref{sites-cohomology-lemma-compare-topologies-derived-adequate-modules},
au morphisme $\epsilon : (X_{affine, Zar}, \mathcal{O}) \to
(X_{affine}, \mathcal{O})$ ci-dessus. Nous prenons
$\mathcal{B} = \Ob(X_{affine})$ et nous prenons pour
$\mathcal{A} \subset \textit{PMod}(X_{affine}, \mathcal{O})$
la sous-catégorie pleine des préfaisceaux $\mathcal{F}$
tels que $\mathcal{F}(V) \otimes_{\mathcal{O}(V)} \mathcal{O}(U)
\to \mathcal{F}(U)$ soit un isomorphisme. Observons que, d'après
Descente, lemme \ref{descent-lemma-quasi-coherent-alternative-small},
les objets de $\mathcal{A}$ sont exactement les faisceaux pour la topologie de Zariski
qui sont des modules quasi-cohérents sur $(X_{affine, Zar}, \mathcal{O})$.
D'autre part, d'après Modules sur les sites, lemme
\ref{sites-modules-lemma-chaotic-quasi-coherent},
les objets de $\mathcal{A}$ sont exactement les modules quasi-cohérents
sur $(X_{affine}, \mathcal{O})$, c'est-à-dire pour la topologie chaotique.
Ainsi, si nous montrons que Cohomologie sur les sites, lemme
\ref{sites-cohomology-lemma-compare-topologies-derived-adequate-modules},
s'applique, alors nous obtenons bien l'équivalence canonique entre les
catégories de (3) et (4) au moyen de $\epsilon^*$ et $R\epsilon_*$.

\medskip\noindent
Nous devons vérifier quatre conditions :
\begin{enumerate}
\item Tout objet de $\mathcal{A}$ est un faisceau pour la topologie de Zariski.
Nous l'avons vu ci-dessus.
\item $\mathcal{A}$ est une sous-catégorie de Serre faible de
$\textit{Mod}(X_{affine, Zar}, \mathcal{O})$. Nous avons vu
ci-dessus que $\mathcal{A} = \QCoh(X_{affine, Zar}, \mathcal{O})$ et
nous avons également vu que ses objets,
par l'équivalence $\textit{Mod}(X_{affine, Zar}, \mathcal{O}) =
\textit{Mod}(X, \mathcal{O}_X)$, correspondent aux modules quasi-cohérents
sur $X$. Le résultat découle donc de la discussion de
Schémas, section \ref{schemes-section-quasi-coherent}.
\item Tout objet de $X_{affine}$ possède un recouvrement pour la topologie
chaotique dont les membres sont des éléments de $\mathcal{B}$. C'est le cas
car $\mathcal{B}$ contient tous les objets.
\item Pour tout objet $U$ de $X_{affine}$ et tout $\mathcal{F}$ de $\mathcal{A}$,
nous avons $H^p_{Zar}(U, \mathcal{F}) = 0$ pour $p > 0$.
Cela résulte de l'annulation de la cohomologie des modules quasi-cohérents
sur les schémas affines ; voir
Cohomologie des schémas, lemme
\ref{coherent-lemma-quasi-coherent-affine-cohomology-zero}.
\end{enumerate}
Ceci achève la démonstration.
\end{proof}

\begin{remark}
\label{remark-QC-compare}
Soit $S$ un schéma.\newline Nous montrerons plus loin que
$\mathit{QC}((\textit{Aff}/S), \mathcal{O})$
est également canoniquement équivalente à $D_\QCoh(\mathcal{O}_S)$.
Voir Faisceaux sur les champs, proposition \ref{stacks-sheaves-proposition-QC-compare}.
\end{remark}








\begin{multicols}{2}[\section{Autres chapitres}]
\noindent
Préliminaires
\begin{enumerate}
\item \hyperref[introduction-section-phantom]{Introduction}
\item \hyperref[conventions-section-phantom]{Conventions}
\item \hyperref[sets-section-phantom]{Théorie des ensembles}
\item \hyperref[categories-section-phantom]{Catégories}
\item \hyperref[topology-section-phantom]{Topologie}
\item \hyperref[sheaves-section-phantom]{Faisceaux sur les espaces}
\item \hyperref[sites-section-phantom]{Sites et faisceaux}
\item \hyperref[stacks-section-phantom]{Champs}
\item \hyperref[fields-section-phantom]{Corps}
\item \hyperref[algebra-section-phantom]{Algèbre commutative}
\item \hyperref[brauer-section-phantom]{Groupes de Brauer}
\item \hyperref[homology-section-phantom]{Algèbre homologique}
\item \hyperref[derived-section-phantom]{Catégories dérivées}
\item \hyperref[simplicial-section-phantom]{Méthodes simpliciales}
\item \hyperref[more-algebra-section-phantom]{Compléments d'algèbre}
\item \hyperref[smoothing-section-phantom]{Lissage des morphismes d'anneaux}
\item \hyperref[modules-section-phantom]{Faisceaux de modules}
\item \hyperref[sites-modules-section-phantom]{Modules sur les sites}
\item \hyperref[injectives-section-phantom]{Objets injectifs}
\item \hyperref[cohomology-section-phantom]{Cohomologie des faisceaux}
\item \hyperref[sites-cohomology-section-phantom]{Cohomologie sur les sites}
\item \hyperref[dga-section-phantom]{Algèbre différentielle graduée}
\item \hyperref[dpa-section-phantom]{Algèbre à puissances divisées}
\item \hyperref[sdga-section-phantom]{Faisceaux différentiels gradués}
\item \hyperref[hypercovering-section-phantom]{Hyperrecouvrements}
\end{enumerate}
Schémas
\begin{enumerate}
\setcounter{enumi}{25}
\item \hyperref[schemes-section-phantom]{Schémas}
\item \hyperref[constructions-section-phantom]{Constructions de schémas}
\item \hyperref[properties-section-phantom]{Propriétés des schémas}
\item \hyperref[morphisms-section-phantom]{Morphismes de schémas}
\item \hyperref[coherent-section-phantom]{Cohomologie des schémas}
\item \hyperref[divisors-section-phantom]{Diviseurs}
\item \hyperref[limits-section-phantom]{Limites de schémas}
\item \hyperref[varieties-section-phantom]{Variétés}
\item \hyperref[topologies-section-phantom]{Topologies sur les schémas}
\item \hyperref[descent-section-phantom]{Descente}
\item \hyperref[perfect-section-phantom]{Catégories dérivées de schémas}
\item \hyperref[more-morphisms-section-phantom]{Compléments sur les morphismes}
\item \hyperref[flat-section-phantom]{Compléments sur la platitude}
\item \hyperref[groupoids-section-phantom]{Schémas en groupoïdes}
\item \hyperref[more-groupoids-section-phantom]{Compléments sur les schémas en groupoïdes}
\item \hyperref[etale-section-phantom]{Morphismes étales de schémas}
\end{enumerate}
Thèmes de théorie des schémas
\begin{enumerate}
\setcounter{enumi}{41}
\item \hyperref[chow-section-phantom]{Homologie de Chow}
\item \hyperref[intersection-section-phantom]{Théorie de l'intersection}
\item \hyperref[pic-section-phantom]{Schémas de Picard des courbes}
\item \hyperref[weil-section-phantom]{Théories cohomologiques de Weil}
\item \hyperref[adequate-section-phantom]{Modules adéquats}
\item \hyperref[dualizing-section-phantom]{Complexes dualisants}
\item \hyperref[duality-section-phantom]{Dualité pour les schémas}
\item \hyperref[discriminant-section-phantom]{Discriminants et différentes}
\item \hyperref[derham-section-phantom]{Cohomologie de de Rham}
\item \hyperref[local-cohomology-section-phantom]{Cohomologie locale}
\item \hyperref[algebraization-section-phantom]{Géométrie algébrique et formelle}
\item \hyperref[curves-section-phantom]{Courbes algébriques}
\item \hyperref[resolve-section-phantom]{Résolution des surfaces}
\item \hyperref[models-section-phantom]{Réduction semi-stable}
\item \hyperref[functors-section-phantom]{Foncteurs et morphismes}
\item \hyperref[equiv-section-phantom]{Catégories dérivées de variétés}
\item \hyperref[pione-section-phantom]{Groupes fondamentaux de schémas}
\item \hyperref[etale-cohomology-section-phantom]{Cohomologie étale}
\item \hyperref[crystalline-section-phantom]{Cohomologie cristalline}
\item \hyperref[proetale-section-phantom]{Cohomologie pro-étale}
\item \hyperref[relative-cycles-section-phantom]{Cycles relatifs}
\item \hyperref[more-etale-section-phantom]{Compléments de cohomologie étale}
\item \hyperref[trace-section-phantom]{La formule des traces}
\end{enumerate}
Espaces algébriques
\begin{enumerate}
\setcounter{enumi}{64}
\item \hyperref[spaces-section-phantom]{Espaces algébriques}
\item \hyperref[spaces-properties-section-phantom]{Propriétés des espaces algébriques}
\item \hyperref[spaces-morphisms-section-phantom]{Morphismes d'espaces algébriques}
\item \hyperref[decent-spaces-section-phantom]{Espaces algébriques décents}
\item \hyperref[spaces-cohomology-section-phantom]{Cohomologie des espaces algébriques}
\item \hyperref[spaces-limits-section-phantom]{Limites d'espaces algébriques}
\item \hyperref[spaces-divisors-section-phantom]{Diviseurs sur les espaces algébriques}
\item \hyperref[spaces-over-fields-section-phantom]{Espaces algébriques sur des corps}
\item \hyperref[spaces-topologies-section-phantom]{Topologies sur les espaces algébriques}
\item \hyperref[spaces-descent-section-phantom]{Descente et espaces algébriques}
\item \hyperref[spaces-perfect-section-phantom]{Catégories dérivées d'espaces}
\item \hyperref[spaces-more-morphisms-section-phantom]{Compléments sur les morphismes d'espaces}
\item \hyperref[spaces-flat-section-phantom]{Platitude sur les espaces algébriques}
\item \hyperref[spaces-groupoids-section-phantom]{Groupoïdes dans les espaces algébriques}
\item \hyperref[spaces-more-groupoids-section-phantom]{Compléments sur les groupoïdes dans les espaces}
\item \hyperref[bootstrap-section-phantom]{Amorçage}
\item \hyperref[spaces-pushouts-section-phantom]{Sommes amalgamées d'espaces algébriques}
\end{enumerate}
Thèmes de géométrie
\begin{enumerate}
\setcounter{enumi}{81}
\item \hyperref[spaces-chow-section-phantom]{Groupes de Chow des espaces}
\item \hyperref[groupoids-quotients-section-phantom]{Quotients de groupoïdes}
\item \hyperref[spaces-more-cohomology-section-phantom]{Compléments sur la cohomologie des espaces}
\item \hyperref[spaces-simplicial-section-phantom]{Espaces simpliciaux}
\item \hyperref[spaces-duality-section-phantom]{Dualité pour les espaces}
\item \hyperref[formal-spaces-section-phantom]{Espaces algébriques formels}
\item \hyperref[restricted-section-phantom]{Algébrisation des espaces formels}
\item \hyperref[spaces-resolve-section-phantom]{Résolution des surfaces revisitée}
\end{enumerate}
Théorie des déformations
\begin{enumerate}
\setcounter{enumi}{89}
\item \hyperref[formal-defos-section-phantom]{Théorie formelle des déformations}
\item \hyperref[defos-section-phantom]{Théorie des déformations}
\item \hyperref[cotangent-section-phantom]{Le complexe cotangent}
\item \hyperref[examples-defos-section-phantom]{Problèmes de déformation}
\end{enumerate}
Champs algébriques
\begin{enumerate}
\setcounter{enumi}{93}
\item \hyperref[algebraic-section-phantom]{Champs algébriques}
\item \hyperref[examples-stacks-section-phantom]{Exemples de champs}
\item \hyperref[stacks-sheaves-section-phantom]{Faisceaux sur les champs algébriques}
\item \hyperref[criteria-section-phantom]{Critères de représentabilité}
\item \hyperref[artin-section-phantom]{Axiomes d'Artin}
\item \hyperref[quot-section-phantom]{Espaces de Quot et de Hilbert}
\item \hyperref[stacks-properties-section-phantom]{Propriétés des champs algébriques}
\item \hyperref[stacks-morphisms-section-phantom]{Morphismes de champs algébriques}
\item \hyperref[stacks-limits-section-phantom]{Limites de champs algébriques}
\item \hyperref[stacks-cohomology-section-phantom]{Cohomologie des champs algébriques}
\item \hyperref[stacks-perfect-section-phantom]{Catégories dérivées de champs}
\item \hyperref[stacks-introduction-section-phantom]{Introduction aux champs algébriques}
\item \hyperref[stacks-more-morphisms-section-phantom]{Compléments sur les morphismes de champs}
\item \hyperref[stacks-geometry-section-phantom]{La géométrie des champs}
\end{enumerate}
Thèmes de théorie des espaces de modules
\begin{enumerate}
\setcounter{enumi}{107}
\item \hyperref[moduli-section-phantom]{Champs de modules}
\item \hyperref[moduli-curves-section-phantom]{Espaces de modules de courbes}
\end{enumerate}
Divers
\begin{enumerate}
\setcounter{enumi}{109}
\item \hyperref[examples-section-phantom]{Exemples}
\item \hyperref[exercises-section-phantom]{Exercices}
\item \hyperref[guide-section-phantom]{Guide bibliographique}
\item \hyperref[desirables-section-phantom]{Desiderata}
\item \hyperref[coding-section-phantom]{Style de programmation}
\item \hyperref[obsolete-section-phantom]{Obsolète}
\item \hyperref[fdl-section-phantom]{Licence de documentation libre GNU}
\item \hyperref[index-section-phantom]{Index généré automatiquement}
\end{enumerate}
\end{multicols}


\bibliography{my}
\bibliographystyle{amsalpha}

\end{document}
